
%\documentclass[a4paper,11pt,oneside,headinclude,titlepage,bibliography=totoc,listof=totoc]{scrbook}
\documentclass[a4paper,11pt,oneside,headinclude,titlepage,bibliography=totoc,listof=totoc,cleardoublepage=current]{scrbook}
%,headings = normal, headings = onelinechapter, headings = onelineappendix, cleardoublepage=plain

%\documentclass[a4paper,11pt,titlepage,bibliography=totoc,listof=totoc]{scrreprt}
%\usepackage[DIV=15]{typearea}
%\usepackage[inner=1.5cm,outer=1.5cm,top=1.5cm,bottom=1.5cm,includeheadfoot]{geometry}
%\setlength{\topskip}{\ht\strutbox} % behebt Warnung von geometry
%\setlength{\headheight}{1.1\baselineskip}   % Headheight hochsetzten
\usepackage[a4paper]{geometry}
\geometry{paper=a4paper,left=30mm,right=30mm,top=35mm,bottom=35mm}
%\renewcommand{\chapterpagestyle}{scrheadings}
%\renewcommand*\chapterheadstartvskip{\vspace*{-0.9cm}}
\renewcommand*\chapterheadstartvskip{\vspace*{-0.5cm}}

\usepackage{float}
\usepackage{subfig}

\usepackage[latin1]{inputenc}
\usepackage[T1]{fontenc}			% erm�glicht die Silbentrennung von W�rtern mit Umlauten
\usepackage[ngerman]{babel} 	% erm�glicht deutsche Silbentrennung und direkte Eingabe von Umlauten, ...

\usepackage[pdftex]{hyperref} % PDF wird mit Lesezeichen (verlinktes Inhaltsverzeichnis) versehen (bei Betrachtung mit Acrobat Reader sichtbar)
\hypersetup{% �=\304; �=\326; �=\334; �=\344; �=\366; �=\374; �=\377
  pdftitle={Masterarbeit},
  pdfauthor={Maximilian H\366nig},
  pdfsubject={Prozedurale Spielweltgenerierung: Dungeons},
  pdfkeywords={Computergrafik, Prozedurale Geometrie},
  colorlinks=true, 		% aktiviert farbige Referenzen
  linkcolor=black,
  citecolor=black,
  urlcolor=black,
  %plainpages=false, % weniger Warnungen durch eindeutigere Anker
  %pdfpagelabels			% das gleiche
}

\usepackage[numbers]{natbib} 		% URLs etc. mit Bibtex, numbers = Zitierart (natbib -> f�r Naturwissenschaften)
\bibliographystyle{alphadin} 		% Bibtex style, natdin

\usepackage{graphicx}
%\usepackage{epstopdf} 					% eps-Dateien mit pdftex nutzen

\usepackage{amsmath} 
\usepackage{amssymb}
\usepackage{floatflt}

\usepackage{wasysym}						% Winkel Zeichen
\usepackage{dsfont}							% Zeichen f�r ganze Zahlen, nat�rlich Zahlen etc. mit \mathds{Z} etc.

\usepackage{esvect}							% bessere Vektor-Darstellung mit \vv{Vektor}

\usepackage{algorithm}					% Algortihmen per \begin{algorithm}
\floatname{algorithm}{Algorithmus} % Caption-Rename
\usepackage{algpseudocode} 			% Pseudcode per \begin{algorithmic}
\renewcommand{\algorithmicrequire}{\textbf{Eingabe: }} 
\renewcommand{\algorithmicensure}{\textbf{Ausgabe: }} 

\usepackage{listings}						% f�r Quellcode
\usepackage{color}							% farbiges Syntaxhighlighting

\usepackage{paralist}						% kompakte Listen mit \begin{compactitem}, \begin{compactenum}

% Tabellen farbig setzen
\usepackage{color}
\definecolor{FarbeZeit}{rgb}{0.8,1.0,0.8}
\usepackage{colortbl}

%\usepackage{tabularx}					% Tabellen mit konstanter Gesamtbreite

%\usepackage{algorithmicx}

%\usepackage{fullpage}

%\usepackage{subfigure}

\interfootnotelinepenalty=10000 % Strafwert f�r Verteilen von Fu�noten auf 2 Seiten

%///////////////////////////////////////////////////

\begin{document}

\pagenumbering{roman}
\title{Masterarbeit}
\author{Maximilian H�nig}

\begin{titlepage}
\vspace*{1.5cm}
%\vspace*{0.5cm}
\begin{center}

\large
Masterarbeit\\

\vspace{1cm}
\Huge
Prozedurale Spielweltgenerierung: Dungeons\\

\Large
\vspace{1cm}
Maximilian H�nig\\

\large
\vspace{1cm}
{4. Juli 2012} \\

\vspace{1.5cm}
betreut durch:\\
Dr. rer. nat. habil. Peter Schenzel\\

\vspace{1cm}
Arbeitsgruppe Computergrafik \\ Naturwissenschaftliche Fakult�t III \\ Institut f�r Informatik \\

\vspace{0.5cm}
Martin-Luther-Universit�t Halle-Wittenberg\\ Halle (Saale) \\ Sachsen-Anhalt, Deutschland

\vspace{0.5cm}
\includegraphics[width=0.5\textwidth]{Bilder/Uni_Logo}

\end{center}
\end{titlepage}
%\newpage
%\section*{}
%\thispagestyle{empty}

\chapter*{ }
%\addcontentsline{toc}{chapter}{Erkl�rung}

Dies ist eine erg�nzte und korrigierte Version der urspr�nglichen Arbeit.
Diese Version stammt vom 27. August 2012.%\today.

\chapter*{Eidesstattliche Erkl�rung}
%\addcontentsline{toc}{chapter}{Erkl�rung}

Ich erkl�re hiermit, dass ich diese Arbeit selbst�ndig verfasst,
noch nicht anderweitig f�r andere Pr�fungszwecke vorgelegt,
keine anderen als die angegebenen Quellen und Hilfsmittel benutzt
und die den benutzten Werken w�rtlich oder sinngem�� entnommenen Stellen als
solche gekennzeichnet habe.

\vspace{3cm}

%Halle (Saale), \today


\begin{center}
\line(1,0){150} \hspace*{4cm}
\line(1,0){150} \\

%\ \\
\textit{Ort, Datum} \hspace*{7.3cm} \textit{Unterschrift}
\end{center}

\newpage
\section*{}
%\thispagestyle{empty}

%\begin{abstract}
%Hier steht sp�ter der Abstract.
%\end{abstract}
\tableofcontents

%\pagenumbering{arabic}
%\pagestyle{scrheadings}
%\KOMAoptions{DIV=12}
%\KOMAoptions{headinclude}
%\KOMAoptions{headlines=10}
\newtheorem{TheoremDefinition}{Definition}

\pagenumbering{arabic}

\chapter{Einleitung}

%1. Bedeutung Computerspiele, 2. prozedurale Generierung, 3. Dungeons

Diese Arbeit besch�ftigt sich mit der prozeduralen Generierung von Dungeons, die als Spielwelten in Computerspielen verwendet werden k�nnen.
Computerspiele stellen ein interaktives Unterhaltungsmedium dar, welches in den letzten Jahren massiv an Bedeutung gewonnen hat.
Die entwickelten Spiele sind mitunter Gro�produktionen, die in ihrem Budget und Aufwand mit Hollywood-Blockbustern vergleichbar sind.
%Zahlen nennen mit !!!QUELLE
Herausforderungen bei der Entwicklung von Computerspielen sind unter anderem: die Einhaltung des Budgets,
die Erstellung von vielen Spielinhalten und das Einhalten von Speicherplatzbeschr�nkungen.
Diese Punkte stehen allerdings im Widerspruch zueinander.

Dieser Widerspruch kann durch den Einsatz prozeduralen Generierungsmethoden entspannt werden.
Mittels entsprechender Algorithmen lassen sich hiermit aus wenigen Eingabeparametern komplexe Inhalte, wie beispielsweise Landschaften, erzeugen.
Somit k�nnen Spielinhalte vergleichsweise schnell und einfach erstellt werden.
Die Algorithmen werden mitunter direkt in Spiele integriert.
Durch die Verwendung zuf�lliger Werte f�r ausgew�hlte Generierungsparameter kann bei jedem Start des Spiels eine neue Spielwelt erzeugt werden.
Hierdurch wird f�r eine hohe Wiederspielbarkeit gesorgt.
Methoden der prozeduralen Generierung erlauben es ebenfalls, gro�e Mengen an Inhalten mit wenig Speicherverbrauch abzulegen,
da nur die Generierungsparameter gespeichert werden m�ssen und nicht das Gesamtresultat.

Dungeons, w�rtlich �bersetzt Verliese oder Kerker, stellen ein wichtiges Element in Computerspielen dar.
Sie bestehen aus untereinander verbundenen H�hlen, G�ngen und R�umen.
Insbesondere in der sehr weit verbreiteten Gruppe der Rollenspiele bilden Dungeons oft den Hauptteil der Spielwelt.
Bis vor einigen Jahren waren diese prim�r zweidimensional aufgebaut,
jedoch gewinnen zunehmend dreidimensionale Strukturen an Bedeutung.
Im Gegensatz zu zweidimensionalen Dungeons, f�r die es bereits viele Generierungswerkzeuge gibt,
ist die Erstellung dreidimensionaler Dungeons deutlich weniger gut erforscht.

%Ein wichtige Untergruppe der Computerspiele stellen Rollenspiele dar.
%%meist in einem Fantasy-Setting angesiedelt
%Die Erforschung und Bew�ltigung von Dungeons stellt hier oft den Hauptanteil des Spiels dar.
%Dungeons, w�rtlich �bersetzt Verliese oder Kerker, bestehen aus einem H�hlensystem sowie damit verbundenen G�ngen und R�umen.

%Bis vor einigen Jahren waren diese prim�r zweidimensional aufgebaut.
%Die prozedurale Generierung solcher 2D-Dungeons ist gut erforscht.
%Zunehmend gewinnen jedoch dreidimensionale Strukturen an Bedeutung.
%Dies erfordert Verfahren, um auch 3D-Dungeons prozedural generieren zu k�nnen.

%Sie sind schon seit langem in Computerspielen zu finden. %!!!QUELLERogue
%Im Gegensatz zu zweidimensionale Dungeons, f�r die es viele Generierungswerkzeuge gibt,
%ist die Erstellung dreidimensionaler Dungeons deutlich weniger gut erforscht.
%Es besteht also ein Bedarf an der prozeduralen Generierung dreidimensionaler Dungeons.

%Dungeons, auch H�hlenverliese geannnt
%
%
%Erkl�rung prozedurale Generierung (kurz),
%Erkl�rung prozedurale Spielweltgenerierung,
%Erkl�rung: Dungeon (kurz),
%%incl. Bilder echter H�hlen, G�nge, R�ume
%Bedeutung f�r Computerspiele,
%%
%Warum diese Arbeit? 
%- Prozeduralit�t spart Arbeit, gr��ere Vielfalt m�glich
%- prozedurale 2D-Dungeons gut erforscht
%-> prozedurale 3D-Dungeons
%%
%%\section{Motivation}
%1)
%- Spiele entwickeln teuer
%- Verschiedene Bereiche: u.a. Spielinhalt - eng verkn�pft mit Leveldesign, welches die Spielwelt aufbaut
%- Kosten k�nne durch automatische Erstellung von Spielinhalten reduziert
%%
%2)
%- Wiederspielbarkeit
%- automatische Erzeugung von Spielinhalten mit Zufallsfaktoren
%- dadurch hohe Anzahl verschiedener Level m�glich
%%
%- Bedeutung von Dungeons in Computerspielen
%- besonders wichtig in RPG's

\section{Zielsetzung}

Ziel der Arbeit ist es, ein Generierungsverfahren f�r realistische 3D-Dungeons zu entwickeln. 
Zu diesen Dungeons geh�ren zerkl�ftete H�hlensysteme, gewundene G�nge und passend ausgestattete R�ume.
Schwerpunkt der Arbeit liegt auf der Generierung der Geometrien dieser Objekte sowie der passenden Platzierung und nahtlosen Verbindung aller Bestandteile.
Die entwickelten Algorithmen sollen effizient bez�glich ihrer Laufzeit und des Speicherverbrauchs sein und konkret in einem Programm implementiert werden.
%
%
%- LOD
%- effiziente Algorithmen
%- konkrete Imlempentierung
%- Export
%
%-----
%
%Aufgabenstellung, Prozedurale Generierung von Dungeons (Dungeons (H�hlenverliese) erkl�ren [QUELLE])
%Schwerpunkt: prozedurale Geometrie (keine prozeduralen Texturen etc.)
%
%- 3D-Dungeons
%- realistische, zerkl�ftete 3D-H�hlen-Systeme
%- System von G�ngen und R�umen, mit H�hle verbunden
%- incl. M�glichkeit zum Platzieren von Detailobjekten
%- exportierbar zur Nutzung in anderen Programmen (vorrangig Spiele, aber auch g�ngige 3D-Modelling-Programme)
%- effiziente Algorithmen: Speicherbedarf nicht zu hoch, schnell -> praktikabel
%
%- Implementierung dieser Algorithmen -> Dungeongenerator-Programm
%- um zu aufzuzeigen, ob diese auch in der Praxis funktionieren

\section{Aufbau der Arbeit}

Der Aufbau der Arbeit folgt, nach Betrachtung der Grundbegriffe \emph{H�hle}, \emph{Dungeon} und \emph{prozedurale Generierung} und Beispielen f�r diese,
dem eigentlichen Generierungsprozesses des Dungeons.
Dabei werden die jeweiligen Grundlagen in den Abschnitten erl�utert, in denen diese Verwendung finden.
Dies hat den Vorteil, dass gleich darauf Bezug genommen werden kann.

\ \\
\emph{Kapitel drei} beschreibt einige selbst entwickelte Herangehensweisen f�r die Dungeongenerierung, inklusive des finalen Ansatzes.
Weiterhin werden wichtige Bemerkungen zur verwendeten Engine, dem Koordinatensystem, den Abstandsma�en und dem Zufallsgenerator getroffen.
Kapitel vier bis acht bilden den Kern der Arbeit und beschreiben die eigentlichen Schritte des Generierungsprozesses.

\ \\
\emph{Kapitel vier} erkl�rt die Erstellung von Voxelh�hlen auf der Basis von L-Systemen.
Grundlagen hierf�r sind die Konzepte Fraktal, L-System, Voxel und 3D-Bresenham-Linien-Algo"-rithmus.
Basierend darauf wird ein Verfahren entwickelt, um Turtle-Grafiken im Voxelraum zu zeichnen.
Weiterhin werden zwei Algorithmen zur Nachbearbeitung der Voxelgrafiken ausgearbeitet.

\ \\
\emph{Kapitel f�nf} beschreibt die Umwandlung von Voxelgrafiken in Dreiecksnetze und deren Darstellung.
Grundlagen hierf�r sind die Konzepte Sweepverfahren, Octree, Normalenberechnung, 3D-Texturen
und das Lambertsche Beleuchtungsmodell.
Es werden Verfahren entwickelt, um diese Umwandlung durchzuf�hren, passende Koordinaten der Eckpunkte zu berechnen,
Normalen von �quivalenten Eckpunkten an Meshgrenzen zusammenzufassen und das Dreiecksnetz zu beleuchten.

\ \\
\emph{Kapitel sechs} beschreibt die Erstellung von G�ngen und R�umen.
Grundlagen hierf�r sind die Konzepte kubische Hermite-Splines, Bogenl�nge und das Bisektionsverfahren.
Es werden Algorithmen zur Erstellung polygonaler Schl�uche und polygonaler Adapter hergeleitet,
das Platzieren von Objekten entlang Kurven erl�utert und das Konzept der Subszene definiert.

\ \\
\emph{Kapitel sieben} beschreibt die Platzierung und das Verbinden von G�ngen und R�umen mit und um die Voxelh�hle herum.
Grundlagen hierf�r sind die Konzepte B�zierkurven und konvexe H�lle.
Es werden Algorithmen entwickelt, um diese Platzierung vorzunehmen und die Objekte korrekt zu verbinden.

\ \\
\emph{Kapitel acht} beschreibt Techniken zur Reduzierung des Renderaufwandes.
Grundlagen hierf�r sind die Konzepte Culling, Level of Detail und die Reduktion von Dreiecksnetzen.
Es wird ein Verfahren entwickelt, um einen Sichtbarkeitsgraphen aufzubauen und die Sichtbarkeitseigenschaften der generierten Objekte zu bestimmen.
Weiterhin wird ein eigener Reduktionsalgorithmus f�r m�glichst regelm��ige Meshes aufgezeigt.

\ \\
\emph{Kapitel neun} beschreibt den strukturellen Aufbau des entwickelten Generatorprogramms.
In \emph{Kapitel zehn} werden Tests bez�glich verschiedener Generierungsparameter durchgef�hrt,
eine Zusammenfassung gegeben und m�gliche Punkte f�r weiterf�hrende Arbeiten aufgef�hrt.


%Erst Algorithmen incl. Zielsetzung und Beschreibung
%sp�ter (in Auswertung): weitere Analyse des Einflusses verschiedener Parameter auf konkrete Resultate
%
%- Warum ist der Aufbau so?
%  - der Pipeline folgend: erst Grundlagen, dann Weiterentwicklung und Anwendung
%%- Wo sind theoretische Konzepte untergebracht?
%%- Was war bisher?
%%- Wo sind eigene Ans�tze erarbeitet? -> eventuell in Kapitel "`Wertung"' oder "`Kurz�berblick"' als zus�tzliches Kapitel der Einleitung
%%- Was wurde weiterentwickelt?------
%
%%- Einleitung im ersten Kapitel, Begriffskl�rung
%- im zweiten Kapitel Begriff prozeduraler Generierung erkl�rt
%und Beispiele aus Computerspielen, sowie anderen Programmen genannt mit Schwerpunkt Dungeons und Landschaften erl�utert
%
%- im dritten Kapitel verschiedenen eigenentwickelten Ans�tze der Dungeongenerierung erl�utert, sowie die finale Pipeline erkl�rt
%weiterhin Anmerkungen �ber die verwendete Irrlichengine und weitere Vorbemerkungen
%
%- danach folgen in Kapitel vier bis acht die Erl�uterung der dem Dungeongenerator zugrundeliegenden Generierungsmethoden
%- die Kapitel sind dabei thematisch am Ablauf der Pipeline orientiert
%- verwendete bestehende Konzepte werden im jeweiligen Kapitel besprochen
%- auf diesen Konzepten aufbauend werden die f�r diese Arbeit eigenen entwickelten Ans�tze und Algorithmen beschrieben
%
%- Kapitel neun beschreibt den prinzipiellen Aufbau des entwickelten Generierungsprogramms unter Zuhilfenahme von UML
%
%- die Analyse des Einflusses verschiedener Parameter auf die Dungeongestaltung und die Performanz erfolgt in Kapitel zehn
%ebenfalls wird dort eine Auswertung getroffen und ein Ausblick �ber m�gliche Weiterentwicklungen gegeben

\section{Voraussetzungen}

Voraussetzung f�r das Verst�ndnis der Arbeit
sind fundierte Kenntnisse der Grundkonzepte der Informatik, wie dem bin�re Zahlensystem und der Bedeutung der Landauschen Symbole.
Ebenfalls Voraussetzung ist das Verst�ndnis der Grundkonzepte der Computergrafik, wie Vektorrechnung und Dreiecksnetze,
sowie gute Kenntnisse der Mathematik.
Hilfreich sind weiterhin Kenntnisse in algorithmischer Geometrie, als auch vertiefte Kenntnisse der Computergrafik.


\section{Ergebnisse in Kurzform}

Es wurde eine Methode entwickelt, um dreidimensionale Dungeons generieren zu k�nnen.
Es sind Generierungsverfahren f�r zerkl�ftete H�hlensysteme und gewundenen G�nge geschaffen worden.
F�r die R�ume wurde ein Konzept �berlegt, wie sich vordefinierte Szenen als Vorlagen f�r diese einbinden lassen.
Der Realismusgrad der erzeugten H�hlen lie�e sich durch passende prozedurale 3D-Texturen weiter steigern.

Beim Entwurf der Algorithmen wurde darauf geachtet, dass diese m�glichst auch f�r andere Zwecke eingesetzt werden k�nnen.
Es ist m�glich, L-Systeme und Turtle-Grafiken im Voxelraum darzustellen und daraus ein Dreiecksnetz inklusive Reduktionsstufen zu berechnen.
Die Methode zur Erstellung polygonaler Schl�uche entlang von Kurven sowie das Konzept der Adapter zum Anbinden an andere Geometrien erlauben ebenfalls weitere Einsatzm�glichkeiten.
Die entwickelten Verfahren sind effizient, zumeist handelt es sich um Linearzeitalgorithmen.

Sie wurden sowohl in einem Programm implementiert als auch getestet.
Die Ergebnisse sind visuell ansprechend und die Generierungsmethoden auch in der Praxis effektiv.
% und unter Nutzung passender Texturierung f�r die Verwendung in Computerspielen geeignet.
F�r die Erstellung eines kompletten Dungeons wird auf handels�blichen PCs weniger als eine Minute Rechenzeit ben�tigt.							% Einleitung
%\include{K_Begriffe}								% Begriffe
\chapter{Vorbetrachtung}

In diesem Kapitel erfolgt zuerst die Erkl�rung der Begriffe H�hle und Dungeon
sowie das Auff�hren einiger Beispiele hierf�r.
Danach wird der Begriff prozedurale Generierung n�her betrachtet
und es werden Beispiele daf�r aus dem Bereich der Spielweltgenerierung aufgezeigt.
% (mit Schwerpunkt auf Dungeons, Landschaften und Computerspielen)

%=========================================================================
\section{H�hlen und Dungeons}

\emph{H�hlen} seien im Folgenden analog zu den Ausf�hrungen von \cite[S.6 ff.]{B_Hoehlen} beschrieben.
Hier wird erl�utert, dass die Auffassungen dar�ber,
was unter dem Begriff H�hle zu verstehen ist, relativ subjektiv sind und weit auseinander gehen.
\cite{B_Hoehlen} beschreibt H�hlen, als "`vom Menschen begehbare und durch Naturprozesse gebildete Hohlr�ume,
die ganz oder teilweise von festen Gestein umgeben sind"'.

Es wird zwischen prim�ren und sekund�ren H�hlen unterschieden.
Prim�re H�hlen entstehen schon beim Prozess der Gesteinsbildung.
Ein Beispiel sind Lavah�hlen, die entstehen, wenn erkaltende Lava oberfl�chlich erstarrt
und sich tiefere Teile des Lavastroms weiterbewegen.
Sekund�re H�hlen entstehen erst nach dem Prozess der Gesteinsbildung.
Ein Beispiel hierf�r sind Wasserh�hlen, deren Entstehung auf das
erosive Wirken von Wasser zur�ckzuf�hren ist.

In \cite{W_HoehleDef} wird eine H�hle als
"`nat�rlicher unterirdischer Hohlraum, der gro� genug ist von Menschen betreten zu werden"'
beschrieben.
Hier werden Gletscherh�hlen explizit als H�hlen aufgef�hrt.
Beispiele f�r H�hlen sind in Abbildung \ref{B_Hoehlenbsp} dargestellt.

\begin{figure}
  \centering % trim: links unten rechts oben
 
  \hfill
  \subfloat{\includegraphics[width=7.45cm]{Bilder/Hoehle_Gletscher2}} %[Eis-H�hle][Gletscherh�hle \cite{W_HoehleGletscher}]
  \hfill
  \subfloat{\includegraphics[width=7.45cm]{Bilder/Hoehle_Lava2}} 			%[Lava-H�hle]  [Lavah�hle \cite{W_HoehleLava}]
  \hfill \\

  \hfill
  \subfloat{\includegraphics[height=5.0cm]{Bilder/Hoehle_Mammut}}%[Die Mammuth�hle im  Dachstein \cite{W_HoehleMammut}]
  \hfill
  \subfloat{\includegraphics[height=5.0cm]{Bilder/Hoehle_Wueste}}    %[Kreideh�hle in W�ste][Kreideh�hle in der Weissen W�ste \cite{W_HoehleWueste}]
  \hfill  \\
     
  \hfill
  \subfloat{\includegraphics[width=15cm]{Bilder/Hoehle_Map_CrystalCave_geschnitten_neu_sharpen}} %[Grundriss der Crystal Cave in Wisconsin \cite{W_HoehleCrystalCaveGrundriss}]
  \hfill
  
	\caption[Diverse in der Natur vorkommende H�hlen]{\emph{Diverse in der Natur vorkommende H�hlen}: v.l.n.r \& v.o.n.u. \\
	(a) Gletscherh�hle \cite{W_HoehleGletscher}, (b) Lavah�hle \cite{W_HoehleLava}, (c) die Mammuth�hle im Dachstein \cite{W_HoehleMammut},
	(d) Kreideh�hle in der Wei�en W�ste \cite{W_HoehleWueste}, \\ (e) Grundriss der Crystal Cave in Wisconsin \cite{W_HoehleCrystalCaveGrundriss}}
	\label{B_Hoehlenbsp}
\end{figure} 


\ \\
\emph{Dungeons} bestehen meist aus einer oder mehreren H�hlen, 
R�umen und G�ngen zur Verbindung dieser Bestandteile.
Die R�ume und G�nge sind dabei zumeist, wie auch die H�hle, unterirdisch angelegt.
Es m�ssen nicht alle dieser Bestandteile vorhanden sein.
Es gibt auch Dungeons, die nur aus R�umen und G�ngen bestehen,
als auch solche, die nur ein H�hlensystem bilden.
Diese Ausarbeitung besch�ftigt sich mit der Generierung von Dungeons,
die aus einer H�hle sowie mehreren G�ngen und R�umen aufgebaut sind.
Die generierten H�hlen besitzen keine direkte Verbindung zur Au�enwelt,
sind also geschlossen und nur durch G�nge betretbar.

Dungeons finden in Spielen vielfach Verwendung,
insbesondere in Rollenspielen mit Fantasyhintergrund.
Sie werden von gef�hrlichen Monstern bev�lkert,
sind gespickt mit t�dlichen Fallen
und reich an wertvollen Sch�tzen.
Einige Beispiele f�r Dungeons sind in Abbildung \ref{B_Dungeonbsp} aufgef�hrt.
Dungeons k�nnen zu Fu� erkundet werden,
aber auch fliegend und kletternd.

\begin{figure}
  \centering % trim: links unten rechts oben
 
  \hfill
  \subfloat{\includegraphics[width=7.4cm]{Bilder/Dungeon_Uldaman}} 
  \hfill
  \subfloat{\includegraphics[width=7.4cm]{Bilder/Dungeon_Schwarzfelstiefen}}
  \hfill \\
  
  \hfill
  \subfloat{\includegraphics[height=5.0cm]{Bilder/Dungeon_XulJarak}} 
  \hfill
  \subfloat{\includegraphics[height=5.0cm]{Bilder/Dungeon_BlackViper}}
  \hfill 
  \hfill\\
  
  \hfill
  \subfloat{\includegraphics[height=7.5cm]{Bilder/Dungeon_Durimek}} 
  \hfill
  \subfloat{\includegraphics[height=7.5cm]{Bilder/Dungeon_Mega}}
  \hfill
  \hfill
  
  
	\caption[Diverse Dungeons]{\emph{Diverse Dungeons}: v.l.n.r \& v.o.n.u. \\
	(a)+(b) aus dem Computerspiel World of Warcraft \cite{W_WorldOfWarcraft}, \\
	(c)\cite[S.22]{B_DnDGruumsch}+(d)\cite[S.169]{B_DnDSerp}+(e)\cite[S.3]{B_DnDRaiders} aus dem Pen-and-Paper-Rollenspiel Dungeons \& Dragons, \\
	(f) Beispiel eines Megadungeons \cite{W_FantasticWorlds}	
	}
	\label{B_Dungeonbsp}
	
\end{figure} 



% Gletscher1:
% https://naturfotografen-forum.de/o247753-Blick+aus+einer+Gletscherh%C3%B6hle
% https://naturfotografen-forum.de/data/media/19/IMG_2301080825%20b::Dr.Karl-Heinz_Limmer_hoehle_gletscher.jpg

% W�ste1:
% https://naturfotografen-forum.de/o43935-My+Cave
% https://naturfotografen-forum.de/data/media/24/eg_bcave_::Dionys_Moser_wueste.jpg

% Lava1:
% https://naturfotografen-forum.de/o520782-Red+stone+cave
% https://naturfotografen-forum.de/data/o/104/520782/image::Philip_Klinger_lavahoehle_lava_langzeitbelichtung_island_hoehle.jpg

% Hoehle_Mammut:
% http://www.foto360.at/vr_panorama_fotos_dyn/fullscreen/vollbild_qtvr_panorama_photo.php?id=616

% Crystal Cave Map:
% http://www.uwec.edu/jolhm/Cave2007/TeamC/TeamC/CrystalCave/Crystal%20Cave%20Map.jpg

% Grand Canyon Caverns :
% http://www.flickr.com/photos/lauri_vain/5025847717/in/photostream/

% Hoehle_Lava2:  http://www.planet-wissen.de/natur_technik/hoehlen/hoehlenforschung/wie_hoehlen_entstehen.jsp

% Hoehle_Gletscher2: http://www.fotocommunity.de/pc/pc/display/19475695

% W�ste2:
% http://www.fotocommunity.de/pc/pc/display/5288904

%~~~~~~~~~~~~~~~~~~~~~~~~~~~~~~~~~~~~~~~~~~~~~~~~~~~~~~~~~~~~~~~~

%Dungeon aus World of Warcraft: Uldaman
%http://wow.4fansites.de/dungeons_uldaman.php

%Dungeon aus World of Warcraft: Teil der Schwarzfelstiefen
%http://wow.4fansites.de/dungeons_schwarzfelstiefen.php

% Dungeons and Dragons:
% Black Viper Altar, Serpent Kingdoms.pdf, S.169 ***
% The Well of Darkness, Silver Marches.pdf, S.153
% X�l-Jarak, Ebene 4, Sons Of Gruumsch.pdf S.22 ***
% Bau eines Bronzedrachens, Draconomicon.pdf S.221
% Durimek, Races of Stone.pdf, Dungeons & Dragons: Raiders of the High Citadel - Web Enhancement for Races of Stone S.3 ***
% The Devil's Cellar, Ghostwalk.pdf S.199 **

% Megadungeon
% http://cimorra.blogspot.de/2010/12/philosophie-megadungeons.html

%=========================================================================
%=========================================================================

\section{Prozedurale Generierung}

Unter \emph{prozeduraler Generierung}, auch als \emph{prozedurale Synthese} bezeichnet,
versteht man die automatische Generierung von Daten nach einer festgelegten Generierungsprozedur.
Die Daten werden also in ihrer endg�ltigen Form nicht manuell erzeugt,
sondern unter Vorgabe entsprechender Parameter mittels eines Algorithmus.
Solche Daten sind oft medialer Natur, beispielsweise Texturen, Geometrien oder Kl�nge.
%Charakteristisch ist eine gr��ere Menge von Ausgabedaten als Eingabedaten
Der Prozess ist meist deterministisch, allerdings k�nnen auch Zufallsparameter einbezogen werden.

%Professionale Anwendung findet solche beispielsweise in den Landschaftsgeneratoren "`World Machine"' \cite{W_WorldMachine}
%und "`Terragen"' \cite{W_Terragen}, wobei letzterer in vielen Hollywood-Filmproduktionen engesetzt wird.

%- zwei M�glichkeiten: in Echtzeit (dynamisch) oder vorberechnet (in externem Programm oder intern)
%%- grundlegende prozedurale Techniken
%Definition von \cite{W_PCGW_Main} ?
%
%Beispiel:

%=========================================================================

\subsection{Prozedurale Inhalte in Computerspielen}

%alternativ:
%\subsection{Unterteilung: Echtzeit oder vorberechnet}
%\subsection{Prozedurale Dungeons}

Prozedurale Inhalte sind in Computerspielen weit verbreitet.
Es gibt zwei grundlegende M�glichkeiten zur Anwendung:
die Daten werden bei Bedarf in Echtzeit berechnet oder sie sind vorberechnet.
Die Vorberechnung kann in einem externen Generatorprogramm geschehen
oder intern im Spiel selbst, z.B. beim Start des Programms.

Durch prozedurale Generierung kann Speicherplatz eingespart werden,
da die Generierungsparameter i.A. deutlich weniger Speicher ben�tigen als
die generierten Daten.
Weiterhin ist eine hohe Vielfalt bez�glich dieser Daten erzielbar,
die bei manueller Erstellung nicht m�glich w�re.

%
%Unterscheidung:
%- Prozedurale Inhalte in Echtzeit:
%  - besonders wichtig bei sehr gro�en Welten oder starken Speicherplatzeinschr�nkungen\\
%  - dynamische on-the-fly Berechnung der Inhalte, da nicht alles in Speicher passen w�rde\\
%  - Berechnungsverfahren i.d.R. deterministisch, also immer gleiche Resultate
%- Vorberechnete Prozedurale Inhalte

%\subsubsection{Elite}
%
%Elite \cite{W_ElitePage} ist eine von 1984 stammende Weltraumsimulation.
%Die Spielwelt besteht aus 8 Galaxien mit je 256 Planeten, die viele spezifische Eigenschaften wie Handelsg�ter oder Regierungsformen besitzen.
%Das Spiel ist lauff�hig auf 8-Bit-Rechnern mit sehr begrenztem Speicherplatz.
%Die Daten werden prozedural bei Bedarf berechnet.
%
%Als Methode zur Generierung werden laut \cite{A_EliteGen} und \cite{W_EliteGen}
%%- Vorgabe eines Randomseeds -> fest Eingestellt damit Universum �berall identisch\\
%Pseudozufallszahlen mittels einer generalisierten Fibonacci-Sequenz erzeugt.
%Durch numerische Beschr�nkungen werden die Zahlen modulo 65536 genommen, welches einer generalisierten Pisano-Sequenz entspricht.
%Unterschiedlich interpretiert ergeben diese Zufallszahlen Sternensysteme, Planetennamen, Handelsg�ter, Alienrassen und weitere Daten.

\subsubsection{Minecraft}

Minecraft \cite{W_MinecraftPage} ist eine Adventure- und Aufbausimulation.
Das Gel�nde ist dreidimensional aus Bl�cken aufgebaut.
Jeder Block besteht aus einem bestimmtem Material wie Fels, Erde, Sand, Wasser oder Luft.
Das Gel�nde wird on-the-fly prozedural generiert, die theoretische Spielweltgr��e ist riesig.
Es wird immer nur das Gel�nde in der N�he der Spielfigur erzeugt.
Bei �nderungen am Gel�nde, beispielsweise durch Aktionen des Spielers, werden diese �nderungen gespeichert.

Der Basisansatz des Generierungsverfahrens gestaltet sich nach \cite{W_MinecraftGen} folgenderma�en:
Mittels eines 3D-Perlin-Noise wird die Dichte jedes Blocks bestimmt,
wobei zus�tzlich die H�he des Wasserpegels einbezogen wird.
Ist die Dichte kleiner null, handelt es sich bei dem Block um Luft,
ist sie gr��er oder gleich null, handelt es sich um festes Material.

\subsubsection{.kkrieger}

.kkrieger ist ein 3D-Ego-Shooter von der Codergruppe Farbrausch/.theprodukkt \cite{W_Farbrausch}.
Die Besonderheit ist, dass das Programm inklusive aller Daten nur 96 KiB gro� ist.
Die ben�tigte Daten werden direkt nach Programmstart prozedural generiert.
Dazu geh�ren die 3D-Modelle, Texturen, Musik, Soundeffekte und der Aufbau der Spielwelt selbst.

Die Generierung funktioniert nach dem Prinzip modularer Synthesizer \cite{W_kkriegerPostmortem}.
Hierbei sind einzelne Generierungsbausteine, wie Perlin Noise oder Bump Mapping,
�ber einen kreisfreien Graphen miteinander verkn�pft.
Die Ausgabe eines Generatorbausteins ist dabei die Eingabe des n�chsten.
Gespeichert werden nur diese Operatoren und ihre Verschaltungen.
Das Spiel kann mit diesen Informationen alle Daten prozedural generieren.

%=========================================================================
\subsection{Roguelikes und prozedurale Dungeons}

Roguelikes \cite{W_RoguelikeDefWiki} sind Computerrollenspiele, die auf dem Spielprinzip des Spiels "`Rogue"' basieren.
Hierzu geh�rt vor allem eine zuf�llig generierte Spielwelt, prim�r aufgebaut aus zweidimensionalen Dungeons,
und das Hindurchk�mpfen durch die diese Dungeons bewohnenden Monsterhorden.
Die 2D-Dungeons in Roguelikes bestehen aus einzelnen Feldern, die in einem quadratischen Gitter angeordnet sind.
Diese Felder sind als bestimmte Typen charakterisiert, beispielsweise als R�ume, G�nge oder W�nde.

\ \\
Viele Computerrollenspiele basieren heute noch auf diesem Prinzip.
Beispiele hierf�r sind Diablo II und Diablo III,
f�r die Screenshots in Abbildung \ref{B_ShotsDiablo} gezeigt sind.
Es gibt eine Vielzahl verschiedener Verfahren, um solche Dungeons zu erstellen.

\begin{figure}[htbp]
  \centering  
  \hfill
  \subfloat{\includegraphics[height=4.75cm]{Bilder/Spielscreenshot_Diablo2_heller}}
  \hfill
  \subfloat{\includegraphics[height=4.75cm]{Bilder/Spielscreenshot_Diablo3_heller}}  
  \hfill
	\caption[Bilder von Roguelikes]{\emph{Bilder von Roguelikes}: v.l.n.r. Screenshot von (a) Diablo II \cite{W_Diablo2}, (b) Diablo III \cite{W_Diablo3}}
	\label{B_ShotsDiablo}
\end{figure}

%\ \\
\noindent Der folgende Algorithmus f�r die Generierung eines Dungeons zur Verwendung in Roguelikes
wird in \cite{W_RoguelikeGen} beschrieben: \\

\begin{compactenum}
	\item Die Karte besteht aus quadratischen Feldern, initial alles W�nde bzw. Felsen, also unpassierbares Gel�nde.
	\item Teile die Karte durch ein Gitternetz auf, jedes Gitterst�ck beinhaltet einen Raum.
	\item Ordne jedem Raum zu: Flag "`verbunden"' (initial sind alle R�ume nicht-verbunden), Array zum Speichern von Verbindungen zu Nachbarr�umen.
	
	\item Bestimme einen zuf�lligen Raum, markiere diesen als "`verbunden"' und mache ihn zum aktuellen Raum.
	\item Solange es nicht-verbundene Nachbarr�ume gibt:
		\begin{compactitem}
			\item	Stelle eine Verbindung zu einem Nachbarraum her.
			\item Markiere diesen als "`verbunden"' und mache ihn zum aktuellen Raum.
		\end{compactitem}		
	\item Solange es nicht-verbundene R�ume gibt:
		\begin{compactitem}
			\item	W�hle einen nicht-verbundenen Raum.
			\item	Versuche zu einem zuf�lligen Nachbar zu verbinden, der schon "`verbunden"' ist.
			\item	In diesen Fall markiere den aktuellen Raum als "`verbunden"', ansonsten w�hle n�chsten nicht-verbunden Raum aus.
		\end{compactitem}		
	\item Alle R�ume sind nun verbunden.
	\item F�ge weitere, zuf�llige Verbindungen zwischen benachbarten R�umen hinzu.
	
	\item Felder setzen: Freir�ume
		\begin{compactitem}
			\item	Zeichne R�ume auf der Karte ein.
			\item	Zeichne Korridore zwischen miteinander verbundenen benachbarten R�umen.
		\end{compactitem}	
	\item Felder setzen: T�ren
		\begin{compactitem}
			\item	Suche nach Feldern mit 2 angrenzenden Wandfeldern, 1-2 angrenzenden Raumfeldern und 0-1 angrenzenden Korridorfeldern.
			\item	Wandle diese in T�ren um.
		\end{compactitem}	
	\item Felder setzen: Treppen
		\begin{compactitem}
			\item	Setze Treppe nach oben im ersten als "`verbunden"' markierten Raum.
			\item	Setze Treppe nach unten im letzten als "`verbunden"' markierten Raum.
		\end{compactitem}	
		Durch die gew�hlte Methode des Treppensetzens wird i.d.R. ein langer Laufweg erzielt, der durch viele R�ume hindurchf�hrt.
\end{compactenum}

\subsubsection{H�hlen per zellul�rem Automat}

In \cite{W_DungCellGen} wird eine Methode f�r die Generierung von 2D-H�hlen beschrieben.
Das Spielfeld besteht aus einzelnen quadratischen Feldern.
Die Methode basiert auf zellul�ren Automaten: \\

\begin{compactenum}
	\item Erstelle initiale Karte durch zuf�lliges Setzen der Felder nach folgender Verteilung:
		\begin{compactitem}
			\item 60\% der Felder sind Gestein.
			\item 40\% der Felder sind Freiraum.
		\end{compactitem}	 
	\item Wende die 4-5 Regel f�r zellul�re Automaten auf jedes Feld der Karte an (eventuell mehrfache Durchl�ufe):
		\begin{compactitem}
			\item Bei weniger als 4 angrenzenden Felsfeldern (in 8er Nachbarschaft): Feld wird zu Freiraum.
			\item Bei mehr als 5 angrenzenden Felsfeldern: Feld wird zu Fels.
			\item Ansonsten: Feld bleibt so belassen.
		\end{compactitem} 	 
	\item Verbinde die einzelnen Kavernen zu einer Gesamth�hle:
		\begin{compactitem}
			\item Bestimme die einzelnen Kavernen.
			\item Zeichne von einem Freiraumfeld jeder Kaverne eine m�andrische Linie von Freiraumfeldern in Richtung Mitte, bis eine andere Kaverne getroffen wird.
		\end{compactitem} 
\end{compactenum}

\subsubsection{DungeonMaker}

Der DungeonMaker \cite{W_DungMakMan} ist eigenst�ndiger Dungeongenerator zum Erstellen gitterbasierten 2D-Dungeons.
Er ist beispielsweise f�r Pen-and-Paper-Rollenspiele oder zur Einbindung in Computerspiele nutzbar.

Der Generierungsprozess verwendet Mechaniken von virtuellem Leben, um Dungeons zu erstellen.
Zuerst platziert der Nutzer einige Designelemente auf der Karte, z.B. R�ume und Mauern.
Nun generiert der DungeonMaker um diese vorplatzierten Elemente den Dungeon.
Dies geschieht mit den Buildern, kleinen Agenten, die mittels KI gesteuert werden.
Subklassen der Builder sind WallCrawler (bauen W�nde), Tunneler (heben Tunnel aus) und Roomies (bauen R�ume).

Ein Vorschlag f�r das Aufbrechen der zugrundeliegenden achsenparallelen Gitterstruktur ist in \cite{W_DungMakSug} aufgef�hrt:
Die entstehende Karte wird in 3D gerendert und dabei verzerrt.


%\subsubsection{Dwarf Fortress}
%- Dwarf Fortress \cite{W_Dwarffortress} simuliert komplette Fantasywelt \\
%- extrem hohe Komplexit�t der Welt \\
%- Welt wird prozedural generiert: Hintergrundgeschichte , Geographie, St�dte, H�hlen, Tierwelt \\

% http://www.mxac.com.au/drt/OblivionProcedural.htm
% http://properundead.com/2009/03/cave-generator.html

%=========================================================================
%\subsection{Prozedurale Landschaften}
%Lanschaftsgeneratoren:
%z.B. World Machine -> kann auch f�r vorberechnete Inhalte in Computerspielen verwendet werden (QUELLE: siehe Making Games Magazin 06/2011, Seite 10)
%Terragen
%
%
%%\section{Grundlegende Prozedurale Techniken}
%%- CG2-�bung

						% Prozedurale Generierung erkl�rt, andere Computerspiele

\chapter{Aufbau der Generatorpipeline}

Wie im vorherigen Kapitel dargelegt wurde,
ist die Generierung von auf einem Gitterlayout basierenden 2D-Dungeons gut erforscht. % (Roguelikes)
Ziel ist es, realistische 3D-Dungeons ohne achsenparallele Gitterstruktur zu erstellen.
Da Spiele i.A. mit Dreiecksnetzen als Basis f�r Geometrien arbeiten,
sollen als Output eben solche plus eventuelle Zusatzdaten erzeugt werden.
Das bedeutet, dass eine Konvertierung aller durch andere Modellierungsmethoden erzeugten Geometrien in Dreiecksnetze erforderlich ist.

\section{Erste Ans�tze}

\subsection{Hierarchische bausteinartige Generierung}
\label{Pipeline-Ansatz1}
%!!!n�her erl�utern: (siehe gro�es Notizbuch):

Der erste Ansatz verwendet eine hierarchische Generierung.
Ein Generator ist f�r das Levellayout zust�ndig.
Er legt Lage, Eigenschaften und Verbindungen einzelner Levelbausteine fest.
Die einzelne Levelbausteine werden durch weitere Generatoren erstellt.
Die Generierung der Bausteine erfolgt nach den bereits festgelegten Eigenschaften.
Dazu geh�ren u.a. �ffnungen bzw. Verbindungen an definierten Stellen
sowie die Vorgabe der Gr��e und des Typs des Bausteins, z.B. Raum, Gangst�ck oder H�hlenteil.
Weitere Generatoren generieren die Feingeometrien, also Teile der Levelbausteine,
wie beispielsweise die Wand eines Raums, Tische oder Torb�gen.

Zu den m�glichen Probleme z�hlen die relativ regelm��igen erzeugten Strukturen,
die �hnlich zu den Roguelikes eine Art Kachel-Layout darstellen w�rden.
Problematisch ist ebenfalls, die H�hlenstruktur realistisch und exakt in einen Bereich einzupassen,
auch in Bezug auf �ffnungen und Verbindungen der H�hle zum Rest des Dungeons.
Durch diese Problematik ist der Realismusgrad eventuell nicht hoch genug.

\subsection{�bereinanderlegung und Vereinigung polygonaler Netze}
\label{Pipeline-Ansatz2}
%!!!siehe Masterarbeit-Konzept:

Der n�chste Ansatz orientiert sich bei der Generierung mehr an der Realit�t und damit an der Frage: Wie w�rde ein Dungeon entstehen?
Zuerst erfolgt die Generierung einer zerkl�fteten, gewundenen H�hle.
Nach der H�hlenerstellung geschieht das Hereinbauen von R�umen und G�ngen um und in die fertige H�hlenstruktur.
Insgesamt sind folgende Schritte angedacht:

\begin{figure}[htbp]
  \centering  
  \hfill
  \subfloat{\includegraphics[width=7.4cm]{Bilder/Scan_LSys2_Skelett}}
  \hfill
  \subfloat{\includegraphics[width=7.4cm]{Bilder/Scan_LSys2_Radius14}} 
  \hfill
	\caption[L-Systeme f�r H�hlenstrukturen]{\emph{L-Systeme f�r H�hlenstrukturen}: 
	v.l.n.r. (a) dargestellt ist die sechste Iteration eines zuf�llig erzeugten L-Systems $G$  mit\\
	$G=\langle \{F,X,Y,+,-\},FF,\{F\rightarrow FXF+YX-F-, X\rightarrow YFY+,$ \\
	$ Y\rightarrow -X-YXFF\} \rangle$ und $\alpha_{Gier} = 276^\circ$		\\
	 (b) dieselbe Iteration desselben L-Systems, gezeichnet mit r�umlicher Ausdehnung der Striche}
	\label{B_LSys_Hoehlenbasis}
\end{figure}

\ \\
\emph{1. Schritt}: Generierung der H�hle.
Ein Beispiel f�r die zugrundeliegende Idee ist in Abbildung \ref{B_LSys_Hoehlenbasis} gezeigt:
Eine fraktal�hnliche Struktur auf der Basis eines Lindenmayer-Systems (s. Kapitel  \ref{K_GenLsys})
ergibt durch r�umliche Expansion ein zerkl�ftetes H�hlensystem.
Hierzu erfolgt zuerst die Erstellung eines fraktalartigen H�hlenskeletts mittels eines passenden L-Systems.
Danach wird dieses Skelett zu einem Dreiecksnetz aufgeblasen und entstehende �berlappungen werden entfernt.

Problematisch ist das Entfernen der Selbst�berschneidung der H�hlenstruktur.
Hierzu m�ssen alle Schnitte zwischen Dreiecken gefunden,
die schneidenden Dreiecke zerlegt
und danach s�mtliche inneren Dreiecke entfernt werden.
Dieser Prozess ist sehr aufw�ndig.
%!!!Quellen ?[proc08.pdf/rosen][V1Nos1to4\_54.pdf] falls O(nlogn) -> besonders bei vielen Dreiecken lange Laufzeit

\ \\
\emph{2. Schritt}: Platzierung von G�ngen und R�umen an passenden Stellen.
Zuerst werden die passenden Stellen gesucht.
Dann erfolgt die Auswahl oder auch Generierung der G�nge und R�ume.
Abschlie�end werden diese platziert, wobei sich die Objekte an den Verbindungsstellen �berlappen.

\ \\
\emph{3. Schritt}: Verschmelzung.
Die H�hle, G�nge und R�ume werden zu einem Dreiecksnetz verschmolzen.
Dies entspricht der Booleschen Vereinigung des Inneren der Ausgangsnetze.

Es ergibt sich ein �hnliches Problem wie in Schritt 1.
Zus�tzlich herausfordernd ist die Beibehaltung der passenden Texturierung beim Zerlegen von Raum- oder Gangdreiecken.
Die H�hle kann per prozeduraler 3D-Textur texturiert werden, wof�r keine expliziten Texturkoordinaten berechnet werden m�ssen.
Bei G�ngen und R�umen sind 2D-Texturen besser geeignet.

%Siehe Konzept: (das folgende �berarbeiten)
%Schritte:
%Generierung der H�hle
%---------------------
%//- Erstellung eines 3D-Fraktalskeletts f�r die grundlegende Struktur der H�hle
%//- Aufblasen des Skelett zu einem 3D-Gebilde im Voxelraum
%//- Umwandlung der Voxelh�hle in Dreiecksnetz
%//-[BILD]
%//- passende Texturierung der H�hle
%
%- Erstellung von Fraktal-Skelett (Lindenmayer-System) f�r H�hle
%- Aufblasung des Skeletts zu Dreiecksnetz
%- Entfernung von �berlappungen
%
%Hineinschneiden der R�ume
%--------------------------
%- R�ume, G�nge und andere Gebilde aus Dreicksnetzen werden gemeinsam mit dem Dreiecksnetz der H�hle im 3D-Raum platziert
%- Verschmelzung aller Dreicksnetze
%  - ein zusammenh�ngendes Dreicksnetz entsteht (das des Dungeons)
%	- Beibehaltung der Texturierung der Teilnetze (Teilnetz: zusammenh�ngend, zu einem Ausgangobjekt geh�rig)
%- Herausforderung: �berg�nge zwischen Teilnetzen, Entfernung �bersch�ssiger Polygone
%-[BILD]
%
%Platzierung zus�tzlicher Detail- und Interaktionsobjekte
%--------------------------------------------------------
%- Partikeleffekte
%- Fallen, Monster
%- 3D-Objekte ohne Integration ins Dreiecksnetz der H�hle: z.B. Truhen, Pflanzen
%	- sinnvoll f�r Objekte die mehrfach auftauchen: m�ssen nur einmal gespeichert werden
%		(bei Intergration ins Dreicksnetz w�rde mehrfach gespeichert werden)
%		
%Zerlegung des Dungeons 
%----------------------
%- Zerschneiden des Dungeons in mehrere St�cke (Sektoren)
%- Berechnung mittels Dijkstra: welchen Abstand haben je 2 Sektoren
%	(gemeint ist die Anzahl von Sektoren die man im Dungeon durchqueren m�sste, um von einem Sektor zum anderen zu gelangen)
%	- sinnvoll f�r Level-of-Detail Algorithmen
%
%Generierung von Dreiecksnetzen verschiedene Aufl�sung f�r die Dungeonteile
%--------------------------------------------------------------------------
%- Nutzung eines Reduktionsalgorithmus, um  Sektoren in mehren Detailstufen zu speichern
%	- Herausforderung: benachbarte Sektoren passen an den R�ndern zusammen, egal welche Detailstufe
%	
%Export:
%-------
%- Speichern der resultierenden Dreiecksnetze	
%- Speichern der Zusatzinformationen
%	- Position und Ausrichtung von Detail- und Interaktionsobjekten
%	- Abstandsmatrix der Sektoren
%	
%	
%Probleme: 
%- Verschmelzungsoperation (�berlappende H�hlenteile, R�ume, G�nge) aufw�ndig
%  - Laufzeit nicht linear:
%   
%  	O(n*log n) (ERKL�RUNG mit Quelle www.cs.swarthmore.edu/~adanner/cs97/s08/papers/rosen.pdf)   	
%  	
%	- Datenstrukturen daf�r ben�tigen viel Speicher (Doubly Connected Edge List o.�.)
%- Raumstruktur wird durch die Verschmelzung ver�ndert, was nicht immer gew�nscht ist
%- G�nge sind nicht prozedural -> wenig Flexibilit�t -> Dungeonstruktur nicht vielf�ltig genug

\section{Iterative Erstellung mittels Subgeneratoren f�r einzelne Dungeonbestandteile}

Das umgesetzte Verfahren enth�lt Elemente aus beiden Ans�tzen.
Die Generierungsmethode ist iterativ und der nat�rlichen Entstehung eines Dungeons folgend.
Es werden Subgeneratoren f�r die Erstellung der einzelnen Geometrien verwendet,
wobei die Teile des Dungeons so konstruiert werden, dass sie nahtlos aneinander passen.
Das Verfahren besteht aus den folgenden Schritten:

\ \\
\emph{1. Schritt}: Es werden Turtle-Grafik-Zeichenanweisungen mittels eines Lindenmayer-Systems erstellt (s. Kapitel \ref{K_GenLsys}).
									Diese beschreiben die grundlegende H�hlenstruktur.									

\ \\
\emph{2. Schritt}: Die Turtle-Grafik wird im Voxelraum gezeichnet (s. Kapitel \ref{K_TurtleVoxel}).
	Voxel haben gegen�ber der Verwendung eines Dreicksnetzes (wie in Abschnitt \ref{Pipeline-Ansatz2}) den Vorteil,
	dass die �berlagerungen der einzelnen Strukturen keine Konflikte verursachen.
	%- Vorteile: CSG-Operationen/�berlagerung sehr einfach umsetzbar
	%- hohe Flexibilit�t bei der Gestaltung	Das entstehende Gebilde stellt die H�hle dar.

\ \\
\emph{3. Schritt}: Es erfolgt die Erosion der H�hle im Voxelraum, um die nat�rliche Zerkl�ftung zu verst�rken (s. Kapitel \ref{K_Erosion}).

\ \\
\emph{4. Schritt}: Die Filterung schwebender Fragmente im Voxelraum wird durchgef�hrt (s. Kapitel \ref{K_Filterung}).
									Hierbei werden solche Fragmente entfernt,	da sie der nat�rlichen Gravitation widersprechen.

\ \\
\emph{5. Schritt}: R�ume und G�nge werden in den Freir�umen um die H�hle hinzugef�gt (s. Kapitel \ref{K_GaengeRaeumeKonzept} und \ref{K_Architekt}).
						Die R�ume werden dabei �ber die G�nge untereinander und mit der H�hle verbunden.
						Dieser Schritt beinhaltet ebenfalls die Berechnung verschiedener Detailstufen
						und Sichtbarkeitsinformationen f�r die G�nge (s. Kapitel \ref{K_LOD}).
	
\ \\
\emph{6. Schritt}: Eine zweite Filterung schwebender Fragmente im Voxelraum wird durchgef�hrt,
						um die eventuell durch das Ausfr�sen von Gang�ffnungen entstandenen Fragmente zu entfernen.

\ \\
\emph{7. Schritt}: Die Umwandlung der Voxelh�hle in ein Dreiecksnetz bzw. mehrere Subnetze findet statt (s. Kapitel \ref{KK_Umwandlung}).
						Dieser Schritt beinhaltet ebenfalls die Berechnung verschiedener Detailstufen f�r die entstehenden Netze
						und der Sichtbarkeitsberechnungen f�r die H�hle (s. Kapitel \ref{K_LOD}).

\ \\
Diese Umwandlung der H�hle in ein Dreiecksnetz erfolgt erst im letzten Schritt,
da in Schritt 5 und 6 �nderungen an der Voxelh�hle vorgenommen werden.
In Schritt 5 werden �ffnungen f�r G�nge hinzugef�gt und in Schritt 6 dadurch entstandene schwebende Fragmente entfernt.

Auf die Erosion und das Filtern schwebender Fragmente kann bei der Erstellung eines Dungeons verzichtet werden.
Wenn Schritt 6 durchgef�hrt werden soll, muss auch Schritt 4 durchgef�hrt werden,
da sonst R�ume in Fragmenten platziert werden k�nnen, die sp�ter entfernt werden.

%==================================================================
%==================================================================

\section{Allgemeine Anmerkungen}


\subsection{Die Irrlicht-Engine}

Das entwickelte Dungeongenerator-Programm nutzt die Irrlicht-Engine \cite{W_Irrlicht} f�r die
Darstellung, f�r die GUI, als Exporter und Importer, zur Speicherung von Dreiecksnetzen und f�r mathematische Hilfsroutinen.
Die Vorteile der Irrlicht-Engine im Vergleich zu anderen Engines werden in \cite{T_EngineVergleich} dargelegt.
Zu nennen ist hier beispielsweise die gro�e Anzahl an unterst�tzten Mesh- und Texturformaten.

Der wichtigste Grund f�r die Verwendung ist das Vorhandensein eines eigenen 3D-Szenenformats, dem .irr-Format,
welches auf XML aufbaut und den Export sowie Import kompletter Szenen erm�glicht.
Damit k�nnen erstellte Dungeons exportiert werden.
Eine solche Szene kann dann als Spielwelt in einem Computerspiel importiert werden,
wie es unter anderem in \cite{T_BachelorHoenig} getan wird.

%Die Irrlicht-Engine l�sst sich auch prakt f�r das

\ \\
Die Szenen in Irrlicht werden als Szenegraph verwaltet.
Der Aufbau dieses Graphen entspricht einer Baumstruktur.
Jeder Knoten enth�lt eine Transformation und optional ein Szenenobjekt,
wie z.B. ein zu zeichnendes 3D-Objekt.
Die Berechnung der resultierenden Transformation f�r jedes Objekt geschieht durch Aufmultiplizierung aller Transformationsmatrizen
von der Wurzel des Szenegraphen bis zum Objekt selbst.

\ \\
Die Transformationen sind entweder direkt als $4\times4$ - Matrix $M$ angebbar
oder einzeln als 3D-Vektoren f�r Skalierung, Rotation und Translation.
Die Rotationswinkel $\left(\alpha_x,\alpha_y,\alpha_z \right)$ sind im Gradma� f�r die jeweilige Achse im mathematisch positiven Sinn angeben.
%also im UZS von Oben auf Ebene schauend da linksh�ndiges System: siehe http://de.wikipedia.org/wiki/Drehrichtung

%Falls die Transformation durch Skalierung, Rotation und Translation angegeben ist,
%berechnet Irrlicht daraus die entsprechende Transformationsmatrix $M$.

Der Aufbau von $M$ entspricht dem aus Formel \ref{F_IrrMatrix}.
Die Matrizen der Skalierung $S$, der Translation $T$
und der Rotationen gem�� Eulerwinkeln $R_x$ (Rotation um X-Achse), $R_y$ (Rotation um Y-Achse) und $R_z$ (Rotation um Z-Achse)
seien dabei aus den entsprechenden 3D-Vektoren gebildet.

\begin{equation}
 \begin{aligned}
 \label{F_IrrMatrix}
 		M = T \cdot R_z \cdot R_y \cdot R_x \cdot S
	\end{aligned}
\end{equation}


\subsubsection{Dreiecksnetze}
\label{KK_Dreiecksnetze}

Die Ablage der Dreiecksnetze, die auch als Dreiecksmeshes oder Meshes bezeichnet werden,
geschieht mittels einer Menge von Vertices und einer Menge von Dreiecken.
Die \emph{Vertices} werden in einem Array gespeichert.
Jeder Vertex hat die Eigenschaften Position, Normale und Texturkoordinaten.
Die \emph{Dreiecke} sind durch ihre Eckpunkte definiert.
In einem Array werden die Indices der Vertices abgespeichert, aus denen die Dreiecke bestehen.
Jedes Dreieck wird aus drei aufeinanderfolgenden Eintr�gen gebildet.

Durch die Verwendung eines unsigned 16 Bit-Datentypes f�r die Vertexindices der Dreiecke
ist die maximale Anzahl der Vertices je Netz auf $65536$ begrenzt.
Meshobjekte bestehen in Irrlicht aus einem oder mehreren Meshbuffern, die einzelne Dreicksnetze darstellen.
Die Begrenzung der Vertices gilt je Meshbuffer.

%- !!!Erkl�rung: Dreiecksnetz
%- Dreiecksnetze in vielen Bereichen prim�res Ablageformat f�r Geometrien:
%  - Computerspiele, 3D-Modelling-Programme
%-> erzeugte Geometrien des Dungeongenerators sollen zum Schluss als Dreiecksnetze verf�gbar sein
% SCHON AM ANFANG DES KAPITELS ERL�UTERT
  
\subsubsection{Koordinatensysteme}

In vielen 3D-Modelling-Programmen und der meisten Literatur wird ein rechtsh�ndiges Koordinatensystem mit Z als H�henachse verwendet.
Irrlicht verwendet ein linksh�ndiges Koordinatensystem mit Y als H�henachse.\footnote{Die
Idee hinter einem solchen Koordinatensystem ist es, das zweidimensionale XY-Koordinatensystem des Bildschirms um die Tiefenachse Z zu erweitert.}
Diese Ausarbeitung und das Dungeongeneratorprogramm verwenden ebenfalls ein derartiges linksh�ndiges System,
wodurch Umrechnungen innerhalb des Generatorprogramms vermieden werden k�nnen.

Falls eine Konvertierung erforderlich ist (beispielsweise beim Export in ausgew�hlte 3D-Formate),
kann zur Umrechnung beider Koordinatensystem die Y- und Z-Achse vertauscht werden.
In homogenen Koordinaten ist dies durch die Multiplikation mit Matrix $M_k$ m�glich:

\begin{equation}
M_k =  \begin{pmatrix}
    1 & 0 & 0 & 0 \\
    0 & 0 & 1 & 0 \\
    0 & 1 & 0 & 0 \\
    0 & 0 & 0 & 1 \\
  \end{pmatrix}
\end{equation}

\ \\
Unter Umst�nden muss nach dieser Umrechnung die Umlaufrichtung der Dreiecke ge�ndert werden,
damit die Fl�chennormalen in die korrekte Richtung zeigen.

\subsection{Abstandsma�e}

F�r die weitere Betrachtung ist eine Erkl�rung der verwendeten Abstandsma�e relevant.
Solange nichts anderes erw�hnt wird, soll mit "`Abstand"' der euklidische Abstand bezeichnet werden.
Der euklidische Abstand zwischen zwei Punkten $\vv{x}=(x_1,x_2,...,x_n)$ und $\vv{y}=(y_1,y_2,...,y_n)$ im n-dimensionalen Raum ist wie folgt definiert:

\begin{equation}
 \label{F_AbstandEuklidisch}
 		D_{e}(\vv{x},\vv{y})= \sqrt{\sum\limits_{i=1}^{n} \left(x_{i}-y_{i}\right)^2}
\end{equation}

\ \\
Der Manhattenabstand zwischen zwei solchen Punkten im n-dimensionalen Raum ist folgenderma�en definiert:

\begin{equation}
 \label{F_AbstandManhatten}
 			D_{m}(\vv{x},\vv{y})= \sum\limits_{i=1}^{n} \left|x_{i}-y_{i}\right|
\end{equation}

\subsection{Zufallsgenerator}

Damit Teile des Dungeons zuf�llig erzeugt werden k�nnen und
somit eine gro�e Vielfalt von Form und Aufbau leicht zug�nglich ist,
m�ssen im Generatorprogramm Zufallszahlen erzeugt werden.
Wichtig f�r diese ist eine optionale deterministische Generierungsmethode,
damit jeder Dungeon aus wenigen Daten komplett erstellt werden kann.

%\ \\
Als Methode wird der \emph{lineare Kongruenzgenerator} verwendet, der 1949 von D.H. Lehmer eingef�hrt wurde
und unter anderem in \cite[S.10 ff.]{B_Knuth} beschrieben ist.
Die Methode verwendet vier Parameter: den Modulus $m$, den Multiplikator $a$, den Inkrementor $c$ und den
Startwert $x_0$, welcher auch als Seed bezeichnet wird.
Die Sequenz der Zufallszahlen $x_i$ berechnet sich daraus wie in Formel \ref{F_Zufallszahlen} angegeben,
wobei mit $x_0$ begonnen wird.
Durch feste Werte f�r $a$, $c$, $m$ und $x_0$ ist die Sequenz genau definiert, also deterministisch.
Durch zuf�llige Festlegung der Werte, i.A. des Startwerts $x_0$, lassen sich variable Sequenzen erzeugen.

\begin{equation}
		x_{n+1}= \left(a\cdot x_n+c\right)\ \text{mod}\ m
		\label{F_Zufallszahlen}
\end{equation}

\ \\
Sinnvoll ist es, die Parameter so zu w�hlen, dass die Periode der erzeugten Sequenz m�glichst gro� ist.
Diese gibt an, nach wie vielen Zahlen sich die gesamte Sequenz wiederholt.
Nach dem Satz von Knuth \cite[S.17]{B_Knuth} wird die maximale Periodenl�nge $m$ erreicht,
wenn folgende drei Bedingungen zutreffen: \\

\begin{compactenum}
 \item $c$ ist teilerfremd zu $m$
 \item jeder Primfaktor vom $m$ ist Teiler von $a-1$ 
 \item wenn $m$ ein Vielfaches von $4$ ist, dann auch $a-1$
\end{compactenum}

\ \\
Beispielwerte hierf�r sind die Parameter $a=1103515245$, $c=12345$ und $m=32768$,
die somit eine Periodenl�nge von $32768$ erzeugen.

\ \\
Anzumerken ist, dass es sich bei den generierten Zahlen nicht um echte Zufallszahlen handelt, sondern um Pseudozufallszahlen,
da sie nach festen Vorschriften berechnet werden.
Echte Zufallszahlen sind mit heutiger Rechentechnik nicht erzeugbar.

								% Pipeline
\chapter{Erstellung der H�hle im Voxelraum}

Das Grundgedanke der H�hlenerstellung beruht darauf, dass eine (zuf�llige) fraktal�hnliche Grundstruktur,
chaotisch, aber trotzdem mit einer inh�renten Ordnung, durch das Zeichnen von Strichen mit einem bestimmten Radius so �berlagert wird,
dass eine H�hlenstruktur entsteht.
Zuerst erfolgt hierzu die Erstellung einer H�hlenstruktur-Beschreibung mittels eines L-Systems.
Diese H�hle wird dann im Voxelraum gezeichnet und danach nachbearbeitet, um sie realistischer zu gestalten.

%================================================================================================
\section{Generierung von Lindenmayer-System-Fraktalen}
\label{K_GenLsys}

Der Begriff \emph{Fraktal} wurde von Beno�t B. Mandelbrot entscheidend gepr�gt.
Besonders hervorzuheben ist hier sein Werk "`Die fraktale Geometrie der Natur"' \cite{B_Mandelbrot},
auf welchem die folgende Erkl�rung des Fraktalbegriffes basiert.

Um Fraktale zu definieren, ben�tigt man den Begriff der \emph{Dimension}.
Die \emph{topologische Dimension} stellt stets eine ganze Zahl dar.
F�r Kurven ist sie gleich $1$, f�r Fl�chen gleich $2$ und f�r K�rper gleich $3$.
Es gibt verschiedene allgemeinere Charakterisierungen der topologischen Dimension,
eine gebr�uchliche davon ist die Lebesgue'sche �berdeckungsdimension.

% !!!QUELLE f�r Def. der Lebesgue'sche �berdeckungsdimension
%bedeckt man ein Objekt mit Kugel des Radius $\epsilon$, wobei $\epsilon$ beliebig klein, aber nicht $0$ sein darf,
%so liegt ein beliebiger Punkt des Objektes mindestens zwischen $N$ Kugeln
%http://analysis.math.uni-kiel.de/lasse/DSA/Vortraege/node3.html

Nach \cite[S.49]{B_Mandelbrot} kann f�r selbst�hnliche Figuren die \emph{Hausdorff-Besicovitch-Dimension}
mittels der \emph{�hnlichkeitsdimension} abgesch�tzt werden.
Diese sind f�r solche F�lle oftmals identisch, wobei die �hnlichkeitsdimension einfacher zu ermitteln ist.
Die Hausdorff-Besicovitch-Dimension sowie die �hnlichkeitsdimension sind grunds�tzlich mindestens so gro� wie die topologische Dimension.

Kann eine Figur in $N$ Teile aufgespalten werden und diese Teile mittels einer �hnlichkeitsabbildung mit dem Quotienten $r$
wieder dem Original entsprechen, so hat diese Figur die �hnlichkeitsdimension $D$, wobei gilt:

\begin{equation}
 \begin{aligned}
 \label{F_DimAehnlichkeit}
		D = -\frac{\log N}{\log r}
	\end{aligned}
\end{equation}

\ \\
Ein Beispiel f�r eine solche Figur ist die Koch-Kurve, deren Konstruktion in Abbildung \ref{B_Kochkurve} dargestellt wird.
Die Erstellung beginnt mit dem \emph{Initiator}.
Nun wird in jeder Iteration jede Strecke der vorhergehenden Iteration durch den \emph{Generator} ersetzt.
Der Generator besteht im Fall der Koch-Kurve aus vier Strecken, die gegen�ber der zu ersetzenden Strecke um ein Drittel verkleinert wurden,
%Diese Strecken bilden einen Streckenzug mit den Winkeln $60^\circ, -120^\circ, 60^\circ$ in den Verbindungspunkten.
und hat die Form der ersten Ite\-ra\-ti\-on aus Abbildung \ref{B_Kochkurve}.
Die eigentliche Koch-Kurve entsteht nach unendlich vielen solchen Iterationen.
Als �hnlichkeitsdimension der Kurve ergibt sich $D = -\log 4 / \log \frac{1}{3} \approx 1,26$.
Es gibt eine Vielzahl weiterer solcher mittels Generatoren konstruierten selbst�hnlichen Kurven.

\begin{figure}[htbp]
  \centering  % trim: links unten rechts oben
	\includegraphics[width=14.5cm]{Bilder/koch}
	\caption[Konstruktion der Koch-Kurve]{Konstruktion der Koch-Kurve \cite{W_Quadsoft}}
	\label{B_Kochkurve}
\end{figure}

%\ \\
Mit den Dimensionsbegriffen kann nun der Begriff des Fraktals definiert werden.
Die exakte Definition nach \cite[S.27]{B_Mandelbrot} ist in Definition \ref{D_Fraktal} wiedergegeben.
Die Koch-Kurve ist somit ein Fraktal.

\begin{TheoremDefinition}[Fraktal]
    Ein Fraktal ist eine Menge, deren Hausdorff-Besicovitch-Dimen"-sion echt die topologische Dimension �bersteigt.
  \label{D_Fraktal}
\end{TheoremDefinition}

%!!!BILDER: Mandelbrotmenge, 3D-Fraktal (Julia Menge?)
%================================================================================================
\subsection{Lindenmayer-Systeme}
\label{KK_Lindenmayer-Systeme}

Lindenmayer-Systeme, oder auch \emph{L-Systeme} genannt, wurden vom Biologen Aristid Lindenmayer entwickelt.
Die Erl�uterung der L-Systeme folgt in dieser Ausarbeitung im Wesentlichen den Ausf�hrungen von \cite{B_AlgBoP},
wo diese zur Generierung von Pflanzenstrukturen angewendet werden.

Ein L-System stellt eine Grammatik dar, bei der im Gegensatz zu formalen Chomsky-Gram"-matiken,
die Ersetzungsregeln (Produktionen) gleichzeitig auf ein ganzes Wort angewendet werden, statt sequentiell.
Ein kontextfreies L-System (OL-System) wird nach \cite[S.4]{B_AlgBoP} durch die Definitionen \ref{D_L-System} und \ref{D_L-SystemAbl} beschrieben.
In dieser Arbeit finden ausschlie�lich deterministische kontextfreie L-Systeme Anwendung.
Die Ableitung der L�nge $n$ von $\omega$ sei auch als $n$-te Iteration des L-Systems bezeichnet.

\begin{TheoremDefinition}[Kontextfreies L-System]
    Sei $V$ ein Alphabet, $V^*$ die Menge aller W�rter �ber $V$ und $V^+$ die Menge aller nichtleeren W�rter �ber $V$.
    Ein OL-System ist ein geordnetes Tripel $G=\left\langle V,\omega,P\right\rangle$,
    wobei $V$ das Alphabet des Systems sei, $\omega \in V^+$ ein nichtleeres Wort, welches Axiom genannt wird,
    und $P\subset V\times V^*$ eine endliche Menge an Produktionen.
    Eine Produktion $(a,x) \in P$ wird geschrieben als $a \rightarrow x$.
    Das Zeichen $a$ hei�t Predecessor und das Wort $x$ Successor dieser Produktion.
    F�r jedes Zeichnen $a \in V$ soll mindestens ein Wort $x \in V^*$ existieren, mit $a \rightarrow x$.
    Wenn f�r einen Predessecor $a \in V$ keine Produktion angegeben ist, so wird angenommen, dass die Produktion $a \rightarrow a$ existiert.
    Ein OL-System ist deterministisch gdw. f�r jedes $a \in V$ genau ein Wort $x \in V^*$ mit $a \rightarrow x$ existiert.    
  \label{D_L-System}
\end{TheoremDefinition}

\begin{TheoremDefinition}[Ableitungen von $\omega$]
	Sei $\mu = a_1 \cdots a_m$ ein beliebiges Wort �ber $V$.
	Das Wort $\nu = x_1 \cdots x_m \in V^*$ ist direkt abgeleitet von $\mu$, notiert als $\mu \Rightarrow \nu$,
	gdw. $\forall i = 1,...,m: \exists\ a_i \rightarrow x_i$.
	Ein Wort $\nu$ wird durch $G$ in einer Ableitung der L�nge $n$ generiert, wenn es eine
	Entwicklungssequenz von W�rtern $\mu_0,\mu_1,\cdots,\mu_n$ gibt, so dass
	$\mu_0 = \omega$, $\mu_n = \nu$ und $\mu_0 \Rightarrow \mu_1 \Rightarrow \cdots \Rightarrow \mu_n$.	
  \label{D_L-SystemAbl}
\end{TheoremDefinition}

%W�rter mit einer Ableitung der L�nge $n$ seien auch als \emph{$n$-te Iteration} oder \emph{Iterationsstufe $n$} des L-Systems bezeichnet.
%und das Axiom als \emph{Startwort} bezeichnet. !!!PR�FEN ob Bezeichner verwendet werden

Die Symbole der L-Systeme k�nnen nach \cite[S.6]{B_AlgBoP} als Turtle-Grafik-Zeichen"-anweisungen interpretiert werden.
Dabei stellen die Zeichen des Alphabets die einzelne Anweisungen f�r die "`Turtle"' dar.
Eine Folge dieser Anweisungen, notiert durch eine Ableitung von $\omega$, wird nacheinander Zeichen f�r Zeichen abgearbeitet, wodurch die eigentliche Grafik entsteht.
Die gebr�uchlichsten Symbole und ihre Interpretationen sind beispielsweise in \cite{B_AlgBoP}, \cite{B_LindSys} und \cite{B_GraDV2} beschrieben.
Tabelle \ref{T_Zeichenanweiungen} zeigt die in dieser Arbeit verwendeten Befehle auf,
wobei jeweils das Symbol nach \cite[S.209]{B_AlgBoP} angeben ist und das Symbol, welches im Generatorprogramm genutzt wird.
Die Ab�nderung erfolgt zur Vermeidung von XML-Sonderzeichen, da die Daten u.a. als XML-Dateien gespeichert werden.
Symbol '$\$$' wurde in seiner Interpretation modifiziert, um eine M�glichkeit zu haben, die Turtle horizontal auszurichten
und so einfacher ebene, begehbare Segmente innerhalb von H�hlen schaffen zu k�nnen.
F�r die Tabelle seien folgende Festlegungen getroffen:
die "`Turtle"' habe die Eigenschaften: lokales Koordinatensystem $\zeta_{lok}$, mit den Achsen $\vv{Vorn}$, $\vv{Oben}$, $\vv{Links}$ und der Position $\vv{P}$,
sowie den Radius $r$.
Die Winkel $\alpha_{Gier}$, $\alpha_{Nick}$ und $\alpha_{Roll}$ seien als zus�tzliche Parameter gegeben.

Weitere Zeichen mit anderen Interpretationen sind prinzipiell beliebig m�glich.
Ein vielfach erw�hntes, aber nicht verwendetes Symbol ist f�r das Vorw�rtsbewegen ohne zu Zeichnen zust�ndig.
Hierauf wurde verzichtet, damit die entstehenden Strukturen in jedem Fall zusammenh�ngend sind.
Alle nicht interpretierbaren Symbole werden beim Zeichnen grunds�tzlich ignoriert,
k�nnen aber zur Schaffung von Ersetzungsregeln verwendet werden.

\ \\
Das generelle Konzept der L-Systeme ist das fortlaufende iterative Ersetzen von Objektteilen mittels einer Menge von Ersetzungsregeln, ausgehend von einem Initialobjekt.
Ein Beispiel daf�r ist die oben beschriebene Koch-Kurve, die ebenfalls durch solche Ersetzungen generiert wird.
Diese l�sst sich �ber das L-System \\ $G=\left\langle \{F,+,-\},F,\{F\rightarrow F+F--F+F\} \right\rangle$ mit $\alpha_{Gier} = 60^\circ$ beschreiben.
In jeder neuen Ableitung wird die Strichl�nge auf $\frac{1}{3}$ der Strichl�nge der vorherigen Ableitung gesetzt.

Bei �ber Generatoren erstellten Fraktalen k�nnen die Ersetzungen mittels entsprechender Produktionsregeln dargestellt werden.
Mittels entsprechender L-Systeme lassen sich also Fraktale beschreiben, wobei aber nicht jedes L-System zwangsl�ufig ein solches definiert.

% Kochsche Schneeflocke ist nicht selbst�hnlich, aber trotzdem Fraktal!
%-> selbst�hnliche mittels L-Systemen
%!!!Anpassung in Zeichenbereich sonst kein Fraktal!!!?

\begin{table}
\centering
\begin{tabular}[hc]{|c|c|c|c|}
  \hline
  Symbol & Verwendet & Interpretation\\
  \hline\hline
  $F$ & $F$ & Bewege vorw�rts und zeichne Strich mit Radius $r$\\

  $+$ & $+$ &Drehe um $\vv{Oben}$ um $\alpha_{Gier}$ (nach links)\\
  $-$ & $-$ &Drehe um $\vv{Oben}$ um $-\alpha_{Gier}$ (nach rechts)\\
  \^{ } & $o$ &Drehe um $\vv{Links}$ um $-\alpha_{Nick}$ (nach oben)\\
  $\&$ & $u$ &Drehe um $\vv{Links}$ um $\alpha_{Nick}$ (nach unten)\\
  $\backslash$& $g$  &Drehe um $\vv{Vorn}$ um $-\alpha_{Roll}$ (gegen UZS, nach vorn blickend)\\
  $/$ & $z$ &Drehe um $\vv{Vorn}$ um $\alpha_{Roll}$ (im UZS, nach vorn blickend)\\
  $|$ & $|$ &Drehe um $\vv{Oben}$ um $180^\circ$\\
  $\$$ & $\$$ &Drehe die Turtle so, dass $\vv{Oben}$ gleich $(0,1,0)$ ist \\
      &      &(In \cite{B_AlgBoP}: so rotieren, dass $\vv{Links}$ senkrecht zu $(0,1,0)$)\\
  $[$ & $[$ &Speichere $\zeta_{lok}$, $\vv{P}$ und $r$ auf dem Stack \\
  $]$ & $]$ &Lade $\zeta_{lok}$, $\vv{P}$ und $r$ vom Stack \\
  $!$ & $!$ &Reduziere Radius $r$ \\
  \hline
\end{tabular}
\caption{Symbole f�r Turtle-Grafik-Zeichenanweisungen}
 {v.l.n.r: Symbol nach \cite{B_AlgBoP}, Symbol in dieser Arbeit, Interpretation des Symbols}
 \label{T_Zeichenanweiungen}
\end{table}


%- !!!Erkl�rung: L-Systeme
%- Erkl�rung der verwendeten Grammatik
%- Erkl�rung des Prinzips der Generatoren im Kontext der L-Systeme
%- Fortf�hrung des Beispiels der Kochkurve, nun als L-System
%- m�gliche Befehle
%- verwendete Befehle (Begr�ndung)
%- Art der verwendeten L-Systems (kontextfrei,...)
%================================================================================================
\subsection{Fraktalgenerator}

Der Fraktalgenerator ist ein Generierungsbaustein des Dungeongenerator-Programms, welcher
dazu dient, zuf�llige L-Systeme zu erstellen sowie Turtle-Grafik-Zeichenanweisungen
aus gegebenen L-Systemen zu generieren.
Zum einen lassen sich Parameter, wie die Winkel $\alpha_{Gier}$, $\alpha_{Nick}$ ,$\alpha_{Roll}$,
zuf�llig bestimmen, zum anderen k�nnen ein zuf�lliges Axiom sowie zuf�llige Produktionen f�r gegebene Symbole
erzeugt werden.

F�r die Generierung einer Folge von Zeichenanweisungen werden ein L-System und die gew�nschte L�nge der Ableitung $n$ ben�tigt.
Ausgehend vom Axiom als Startstring wird nun iteriert.
In jeder Iterationsstufe wird ein neuer String erzeugt, in dem �ber den alten String gelaufen wird 
und die Ersetzungen mittels der Produktionen vorgenommen werden.
Diese Iterationen werden insgesamt $n$ mal durchgef�hrt.
Die resultierende Folge von Zeichenanweisungen ist der zuletzt erzeugte String.

Anmerkung:
Diese Zeichenanweisungen erzeugen auch im Fall von fraktalbeschreibenden L-Systemen keine echten Fraktale,
da die Ableitungsfolge nicht unendlich lang ist.\footnote{Solche unendlich langen Ableitungsfolgen lassen sich verst�ndlicherweise nicht generieren.}
Die resultierenden Strukturen sind aber i.d.R. fraktal�hnlich, �hneln also in ihrer Form echten Fraktalen.
%Diese Fraktal�hnlichkeit sei hierbei durch Formen definiert, die scheinen als ob sie eine h�here Dimension haben als tats�chlich vorliegt.
%Dies ist durch an den R�ndern s

%================================================================================================
\section{Turtle-Grafik im Voxelraum}
\label{K_TurtleVoxel}

Der Fraktalgenerator liefert Turtle-Grafik-Zeichenanweisungen, die nun im 3D-Raum gezeichnet werden m�ssen.
Als Grundlage f�r den 3D-Raum werden Voxel verwendet.
\emph{Voxel} ("`volume element"') bilden das dreidimensionale �quivalent zu den zweidimensionalen Pixeln ("`picture element"').
Der Voxelraum ist nach \cite[S.549]{B_Foley} per gleichm��igem Gitter in kleine identische Zellen aufgeteilt, die meistens kubischer Form sind.
Jede dieser Zellen stellt einen Voxel dar.
Abbildung \ref{B_Voxeltorus} zeigt einen aus Voxeln aufgebauten Torus.

Die in dieser Arbeit verwendeten Voxel sollen ausschlie�lich kubischer Natur sein.
Die Schreibweise $Voxel(x,y,z)$ bezeichne die Belegung des Voxels an Position $(x,y,z)$ mit $x,y,z \in \mathds{N}_0$.
Der Voxelraum soll den Bereich aller Voxel in $[0,X_{vmax}] \times [0,Y_{vmax}]\times [0,Z_{vmax}]$ umfassen.

%Def. Voxel: "`volumetric pixel"' Datenpunkt einer dreidimensionalen Rastergrafik
%- W�rfel im Raum, im folgenden sei die Kantenl�nge eines Voxels mit 1 festgesetzt

�ber den Manhattenabstand (siehe Formel \ref{F_AbstandManhatten}) lassen sich verschiedene Nachbarschaftsbeziehungen zwischen Voxeln definieren.
Diese sind dabei an die gebr�uchlichen Begriffe 4er-Nachbarschaft und 8er-Nachbarschaft aus der Bildverarbeitung angelehnt (vgl. \cite[S.173-174]{B_DigBildv}),
die sich auf den zweidimensionalen Pixelraum beziehen.
%!!!TODO: Vierer, Sechser etc. (ausschreiben, Abk�zungen angeben)

\begin{TheoremDefinition}[6er-Nachbarschaft]
    Bez�glich eines Voxel an Position $\vv{a}$ liegt ein Voxel an Position $\vv{b}$ in 6er-Nachbarschaft, wenn gilt:
    $D_{m}(\vv{a},\vv{b}) =1$ und $D_{e}(\vv{a},\vv{b}) < 2$ .
  \label{VoxelNachbarschaft6er}
\end{TheoremDefinition}
\begin{TheoremDefinition}[18er-Nachbarschaft]
    Bez�glich eines Voxel an Position $\vv{a}$ liegt ein Voxel an Position $\vv{b}$ in 18er-Nachbarschaft, wenn gilt:
    $1\leq D_{m}(\vv{a},\vv{b}) \leq 2$ und $D_{e}(\vv{a},\vv{b}) < 2$ .
  \label{VoxelNachbarschaft18er}
\end{TheoremDefinition}
\begin{TheoremDefinition}[26er-Nachbarschaft]
    Bez�glich eines Voxel an Position $\vv{a}$ liegt ein Voxel an Position $\vv{b}$ in 26er-Nachbarschaft, wenn gilt:
     $1\leq D_{m}(\vv{a},\vv{b}) \leq 3$ und $D_{e}(\vv{a},\vv{b}) < 2$ .
  \label{VoxelNachbarschaft26er}
\end{TheoremDefinition}

F�r das Zeichnen der H�hle soll ein nicht gesetzter Voxel ($0$) �quivalent zu Felsgestein sein, welches gleichzeitig der Initialwert f�r alle Voxel ist.
Ein gesetzter Voxel ($1$) bezeichnet den Freiraum der H�hle.
Tempor�r sind andere Markierungen f�r bestimmte Algorithmen m�glich, die entsprechend erw�hnt werden.
Bei der Implementierung wird die Gr��e des Voxelraums auf $512\times512\times512$ Einheiten festgelegt, welches $134.217.728$ Voxeln entspricht.
Zur Speicherung wird ein Byte je Voxel verwendet, wodurch hierf�r insgesamt $128$ MiB n�tig sind.
Da heutzutage oftmals mehrere GiB Arbeitspeicher vorhanden sind, ist dies problemlos umsetzbar.
Der gesamte Voxelraum wird als 3D-Array abgelegt.

Alternativ l�sst sich ein Bit je Voxel mit acht Voxeln in einem Byte verwenden, wodurch f�r gleiche Voxelraum-Gr��e nur noch $16$ MiB Speicherplatz ben�tigt werden w�rden.
Diese Speicherungsvariante kann f�r gr��ere Voxelr�ume relevant sein, da bei Verdopplung der Kantenl�nge des Voxelraums der Speicherverbrauch um das Achtfache steigt.
Allerdings erfolgt der Zugriff auf einzelne Voxel hier per Bin�roperatoren (bin�res Und, bin�res Oder), was unperformanter als der direkte Zugriff ist.
Da f�r bestimmte verwendete Algorithmen andere Markierungen au�er $0$ und $1$ n�tig sind, ist hier ebenfalls weiterer Speicher n�tig.

\begin{figure}[htbp]
  \centering  
	\includegraphics[width=12cm]{Bilder/voxel-model}
	\caption[Aus Voxeln aufgebauter Torus]{Aus Voxeln aufgebauter Torus \cite[S.549]{B_Foley}}
	\label{B_Voxeltorus}
\end{figure}

%\ \\
\noindent Die Turtle habe die in Kapitel \ref{KK_Lindenmayer-Systeme} genannten Eigenschaften.
F�r das Zeichnen der Turtle-Grafik seien weiterhin als Parameter gegeben:
die Winkel $\alpha_{Gier}$, $\alpha_{Nick}$ ,$\alpha_{Roll}$, Radiusfaktor $r_f$ und Radiusdekrementor $r_d$.
Der Stack f�r das Speichern der Turtle-Zust�nde sei initial leer.

\subsection{Drehungen und Radius�nderungen}

F�r die Umsetzung der Rotationen (Symbole '$+$','$-$','$o$','$u$','$g$','$z$') sei gefordert,
dass die Vektoren des lokalen Koordinatensystems immer senkrecht aufeinander stehen und immer die L�nge $1$ besitzen.
Dies erfordert auch bei vielen hintereinander ausgef�hrten Rotationen numerische Stabilit�t.
Dazu erfolgt bei einer Drehung zuerst die Errechnung \emph{einer} neuen Achse des lokalen Koordinatensystems.
Die Konstruktion der zweiten sich ver�ndernden Achse erfolgt per Kreuzprodukt.
In jedem Schritt wird der entsprechende Vektor normalisiert.\footnote{Bei absolut exakter Berechnung k�nnte das Normalisieren weggelassen werden.}
Formel \ref{F_TurtleDrehung} gibt das Beispiel f�r die Drehung um Achse $\vv{Oben}$ mit Winkel $\alpha_{Gier}$ an.
Die Achse $\vv{Oben}$ bleibt bei dieser Drehung unver�ndert, die beiden anderen Achsen werden wie angegeben modifiziert.

\begin{equation}
 \begin{aligned}
 \label{F_TurtleDrehung}
		\vv{Vorn_{neu}} =& \frac{\cos \alpha_{Gier} \cdot \vv{Vorn}+ \sin \alpha_{Gier} \cdot \vv{Links}}
														{\left|\cos \alpha_{Gier} \cdot \vv{Vorn}+ \sin \alpha_{Gier} \cdot \vv{Links}\right|} \\
		\vv{Links_{neu}} =& \frac{\vv{Vorn_{neu}} \times \vv{Oben}}{\left| \vv{Vorn_{neu}} \times \vv{Oben}\right|}
 \end{aligned}
\end{equation}

\ \\
Bei der Ausrichtung von $\vv{Oben}$ (Symbol '$\$$') nach $(0,1,0)$ erfolgt zuerst die Neubelegung von $\vv{Oben}$.
Um die nachfolgende Berechnung numerisch stabil zu gestalten, wird nun bestimmt, welcher der beiden Vektoren $\vv{Vorn}$ oder $\vv{Links}$
den h�chsten Betrag des Skalarprodukts mit $\vv{Oben}$ aufweist, also m�glichst weit senkrecht auf $\vv{Oben}$ steht.
Der entsprechende Vektor sei Vektor $\vv{a}$, der andere Vektor $\vv{b}$.
Vektor $\vv{a}$ wird nun f�r die Berechnung des neuen Werts f�r $\vv{b}$ verwendet, die durch die Bildung des Kreuzprodukts mit $\vv{Oben}$ erfolgt.
Danach wird der neue Wert f�r $\vv{a}$ ebenfalls per Kreuzprodukt ermittelt.

F�r die Drehung um $180^\circ$ (Symbol '$|$') brauchen nur $\vv{Vorn}$ und $\vv{Links}$ mit $-1$ multipliziert zu werden.

\ \\
Die Reduzierung des Radius (Symbol '$!$') erfolgt nach Formel \ref{F_FrakRadius}.
Ziel ist es eine pro\-zen\-tu\-ale (per $r_f$), als auch eine absolute (per $r_d$) Radiusreduzierung zu erm�glichen.
Der Radius kann per Definition nicht kleiner als $0$ werden. In der Implementierung des Generator-Programms ist ein minimaler Radius von $1$ vorgesehen.
Der Wert des Radius wird als Gleitkommazahl gespeichert und beim Zeichnen von Strichen auf eine ganze Zahl abgerundet.

\begin{equation}
 \begin{aligned}
  \label{F_FrakRadius}
		r_{neu} = r_f \cdot r - r_d
 \end{aligned}
\end{equation}

%================================================================================================
\subsection{Zeichnen von Strichen mit gegebenem Radius}

Das Zeichen der Striche erfordert die Aktualisierung der Position. 
Die neue Position berechnet sich wie in Formel \ref{F_FrakPos} angegeben.
Die Strichl�nge sei $L$.

\begin{equation}
 \begin{aligned}
 \label{F_FrakPos}
		\vv{P_{neu}} = \vv{P} + L \cdot \vv{Vorn}
 \end{aligned}
\end{equation}

\ \\
Weiterhin m�ssen die durch den Strich abgedeckten Voxel auf $1$ gesetzt werden.
Der erste Ansatz bestand darin, die gesamte Turtle-Grafik zuerst mittels des 3D-Bresenham-Algorithmus
nur mit Strichdicke $1$ zu zeichnen und dieses Skelett dann zur finalen Ausdehnung zu expandieren.
Dazu wurde eine Priorit�tswarteschlange als Datenstruktur f�r die expandierenden Voxel verwendet.
Die Abarbeitungspriorit�t der Voxel entsprach \emph{(gew�nschter Radius - aktueller Abstand vom Skelett)}.
Problematisch war die Laufzeitkomplexit�t von $\mathcal O(n \log n)$, mit $n$ als Anzahl zu zeichnender Voxel, und dem hohen Speicherverbrauch pro Voxel.
Diese Methode hatte bereits f�r einfache Turtle-Grafiken mit wenigen Strichen eine Laufzeit von mehreren Minuten.
Im Vergleich dazu schaffte es der finale Bresenham-Zylinder-Kugel-Algorithmus diese Grafiken in weniger als einer Sekunde zu zeichnen.

%- 2 Methoden werden betrachtet
%- beide nutzen den f�r drei Dimensionen erweiterten Bresenham-Algorithmus
%- in der erste Methode wird zuerst die Turtle-Grafik mit Strichdicke 1 gezeichnet und danach zu den passenden Strichdicken erweitert
%  - hier wird kein Voxel mehrfach gezeichnet
%- bei der zweiten und finalen Methode wird jeder Strich sofort mit der passenden Dicke gezeichnet
%  - hier kommt es mit unter zu �berdeckungen von Strichen, d.h. Voxel werden mehrfach gezeichnet

\subsubsection{Exkurs: Bresenham-Linien-Algorithmus in 3D}

%!!!vollst�ndiger: Bresenham-Linien-Algortihmus
Der \emph{Bresenham-Linien-Algorithmus}, auch kurz Bresenham-Algorithmus oder Bresenham genannt, ist ein Algorithmus zum Zeichnen von Linien in zweidimensionalen Rastergrafiken,
welcher 1965 von J. E. Bresenham in \cite{bresenham1965algorithm} vorgestellt wurde.
Der Algorithmus ist sehr schnell, da er prinzipiell nur mit Additionen und Entscheidungen auskommt und ausschlie�lich ganze Zahlen verwendet.
Damit stellt er quasi den Standardalgorithmus zum Zeichnen von Linien dar.
%- Rundungsfehler die durch Verwendung kontinuierlicher Koordianten entstehen w�rden werden eliminiert

Die Rasterung einer Linie von $(x_1,y_1)$ nach $(x_2,y_2)$, mit $dx\geq dy$ (wobei $dx=x_2-x_1$, $dy=y_2-y_1$), nach dem Bresenham-Algorithmus wird in Abbildung \ref{B_Bresenham} dargestellt.
Zus�tzliche Voraussetzungen sind $x_1\leq x_2$ und $y_1\leq y_2$. F�r alle anderen F�lle wird der Algorithmus entsprechend angepasst.
Begonnen wird in $(x_1,y_1)$. In jedem Zeichenschritt wird die Koordinate $x$ um $1$ erh�ht.
F�r die Koordinate $y$ gibt es zwei M�glichkeiten: entweder wird sie so belassen oder es findet ebenfalls eine Erh�hung um $1$ statt.
Die Idee besteht darin, dies anhand des Fehlers zu entscheiden, der entstehen w�rde,
und dabei die Variante zu w�hlen, bei der der gezeichnete Pixel n�her an der echten Linie liegt.
Der Bresenham-Algorithmus f�hrt hierzu eine Fehlervariable ein, anhand der entschieden wird.
Diese Variable wird zus�tzlich zur Positions�nderung je nach Entscheidung unterschiedlich aktualisiert.

Diese Fehlervariable sei $e_y$. Falls $e_y\geq0$ wird die Y-Koordinate um $1$ erh�ht, ansonsten so gelassen.
In \cite{bresenham1965algorithm} wird gezeigt, dass wenn anfangs $e_y = 2\cdot dy-dx$ gesetzt wird,
die nachfolgende Aktualisierung dem Prinzip aus Formel \ref{F_Bresen2D} gen�gen muss.

\begin{equation}
 \begin{aligned}
  \label{F_Bresen2D}
		e_y^{neu} = \begin{cases} 
									e_y + 2\cdot dy- 2 \cdot dx & falls\ e_y\geq0 \\ 
									e_y + 2\cdot dy & sonst 
								\end{cases} 
 \end{aligned}
\end{equation}

\begin{figure}[htbp]
  \centering  
	\includegraphics[width=14.5cm]{Bilder/BresenhamLine}
	\caption[Rasterung einer Linie nach dem Bresenham-Verfahren]{Rasterung einer Linie nach dem Bresenham-Verfahren \cite{W_WikiBresen}}
	\label{B_Bresenham}
\end{figure}

\ \\
Der Bresenham-Algorithmus l�sst sich auf dreidimensionale Rastergrafiken, also Voxelgrafiken, erweitern.
In \cite{A_Bresen3D} wird diese 3D-Variante beschrieben.
Algorithmus \ref{Bresenham3D} stellt eine leicht modifizierte Version dar, mit beliebiger Reihenfolge der Start- und Endpunkte.
Dennoch muss gelten, dass $|dx|\geq |dy|$ und $|dx|\geq |dz|$, also ist $x$ die Prim�rkoordinate.
F�r die Anwendung auf alle F�lle wird eine Funktion vorgeschaltet,
die den passenden Fall mit $x$, $y$ oder $z$ als Prim�rkoordinate bestimmt und daraufhin den entsprechend angepassten Algorithmus aufruft.
Die Anpassung erfolgt durch den Rollentausch der Koordinaten $x$, $y$ und $z$.

Der Unterschied zum originalen Bresenham-Algorithmus besteht darin, dass in jedem Zeichenschritt zus�tzlich die Z-Koordinate modifiziert werden kann.
Hierzu gibt es die Fehlervariable $e_z$, anhand der entschieden wird. Deren Aktualisierung findet �quivalent zu der von $e_y$ statt.

\begin{algorithm}
\caption{Bresenham-Linien-Algorithmus in 3D (C++ Listing)}
\label{Bresenham3D}
\algorithmicrequire{\textit{/*Endpunkte (mit X als l�ngster Achse)*/}\\ int xs,int ys,int zs,int xe,int ye,int ze} \\
\algorithmicensure{gezeichnete Linie}
\begin{algorithmic}[1]	
	\State \textit{/* Richtungen festlegen */}
	\State \textbf{int} xinc = 1; \textbf{if}(xe<xs) xinc = -1;
	\State \textbf{int} yinc = 1; \textbf{if}(ye<ys) yinc = -1;
	\State \textbf{int} zinc = 1; \textbf{if}(ze<zs) zinc = -1;
	\State \textit{/* Berechne die delta-Werte */}
	\State \textbf{int} dx=\textbf{abs}(xe-xs); dy=\textbf{abs}(ye-ys); dz=\textbf{abs}(ze-zs);
	\State \textit{/* Berechne die Initialwerte der Entscheidungsparameter */}
	\State \textbf{int} ey=2*dy-dx; ez=2*dz-dx;
	\State \textit{/* Berechne Konstanten */}
	\State \textbf{int} twoDy=2*dy; twoDyDx=2*(dy-dx);
	\State \textbf{int} twoDz=2*dz; twoDzDx=2*(dz-dx);
	\State \textit{/* Setze Startposition */}
	\State \textbf{int} x=xs; y=ys; z=zs;
	\State \textit{/* Annahme: dx>=dy, dx>=dz */}
	\State \textbf{for} (i=0; i<=dx; i++)
	\State \{
	\State \hspace{6mm}\textbf{output}(x,y,z); \textit{/* Zeichne die aktuelle Position*/}
	\State \hspace{6mm}x+=xinc; \textit{/* Start der Berechnung der n�chsten Koordinaten */}
	\State \hspace{6mm}\textbf{if} (ey>=0) \{y+=yinc; ey+=twoDyDx;\}
	\State \hspace{12mm}\textbf{else} ey+=twoDy;
	\State \hspace{6mm}\textbf{if} (ez>=0) \{z+=zinc; ez+=twoDzDx;\}
	\State \hspace{12mm}\textbf{else} ez+=twoDz; \textit{/* Ende der Berechnung der n�chsten Koordinaten */}
	\State \}	
\end{algorithmic}
\end{algorithm}

\subsubsection{Kugel und Kreiszylinder}

Die finale Methode zum Zeichnen von Strichen mit gew�nschtem Radius $r$ im Voxelraum basiert auf dem Zeichnen von zwei grafischen Primitiven.
Zuerst erfolgt das Zeichnen eines geraden Kreiszylinders mit Radius $r$.
Danach wird eine abschlie�ende Kugel mit gleichem Radius am Zylinderende gezeichnet, damit bei Drehungen nachfolgender Striche keine Spalten entstehen
und eine weichere �berblendung bei Radiusreduzierungen erreicht wird.

Der Algorithmus f�r die Kugel ist vergleichsweise einfach umzusetzen.
Es erfolgt das Ablaufen des W�rfels mit dem Mittelpunkt gleich dem Kugelmittelpunkt und der Seitenl�nge gleich dem Kugeldurchmesser
mittels einer dreifach geschachtelte Schleife, wobei jede Schachtelung f�r eine Dimension des W�rfels zust�ndig ist.
Im Inneren der Schleife wird f�r jeden Voxel ein Test durchgef�hrt, ob der Abstand vom Mittelpunkt der Kugel h�chstens so gro� wie der Radius der Kugel ist.
Wenn dieser Test positiv verl�uft, wird der Voxel gezeichnet, also auf $1$ gesetzt.
Zur Optimierung werden bei allen Berechnungen und Tests quadratische Abst�nde verwendet
und Rechenoperationen m�glichst weit aus den Schleifenschachtelungen herausgezogen.

\ \\
Das lokale Koordinatensystem ist durch die Turtle bekannt, die Vektoren sind normiert.
Die Idee des Kreiszylinder-Algorithmus ist es, diese Vektoren zu nutzen, um den Zylinder aufzuspannen.
$\vv{Oben}$ und $\vv{Links}$ multipliziert mit dem Zylinderradius spannen die Grundfl�che auf
und die Zylinderachse verl�uft von der Startposition zur Endposition des Striches,
die beide ebenfalls bekannt sind.

Das Aufspannen der Vektoren geschieht mittels des 3D-Bresenham-Algorithmus.
F�r jede der drei Achsen wird ein Bresenham verwendet, der in einen Buffer zeichnet, statt direkt in den Voxelraum.
Danach wird eine dreifach geschachtelte Schleife durchlaufen, die jeweils �ber alle Buffer-Elemente l�uft, die Koordinaten zusammenaddiert
und ggf. den entsprechenden Voxel zeichnet, falls er wirklich im Zylinder liegt.
Das Verfahren ist in Algorithmus \ref{VoxelZylinder} dargestellt.
  
Durch die Spiegelung eines Viertels der Zylindergrundfl�che l�sst sich die Zahl n�tiger Tests reduzieren.
Allerdings wird die ZeichneVoxel()-Routine an Positionen mit $\vv{P_{1}}$ oder $\vv{P_{2}}$ gleich dem Nullvektor doppelt
und mit $\vv{P_{1}}$ und $\vv{P_{2}}$ gleich dem Nullvektor sogar vierfach ausgef�hrt.
Dies ist aber bei gr��eren Zylinderradien vernachl�ssigbar,
da der Anteil dieser F�lle im Verh�ltnis zur Gesamtvoxelzahl abnimmt.
Ein Abfangen der F�lle vorm Zeichnen der Voxel ist i.d.R. nicht empfehlenswert,
da bedingte Anweisungen vergleichsweise teure Operationen darstellen.

\begin{algorithm}
\caption{Gerader Kreiszylinder im Voxelraum}
\label{VoxelZylinder}
\algorithmicrequire{Radius $r$, Startpunkt $\vv{P_s}$, Endpunkt $\vv{P_e}$, Lokales Koordinatensystem}  \\
\algorithmicensure{gezeichneter Zylinder}
\begin{algorithmic}[1]

	%\State \Comment {Deckflaeche vertikal}
	\State $Buffer1 \gets Bresenham3D(\vv{0}, r\cdot\vv{Oben})$ \Comment {$\vv{0}$ bezeichnet Nullvektor}
	\State $Buffer2 \gets Bresenham3D(\vv{0}, r\cdot\vv{Links})$ 
	\State $Buffer3 \gets Bresenham3D(\vv{P_s}, \vv{P_e})$

	\ForAll{$\vv{P_{1}} \in Buffer1$}
		\ForAll{$\vv{P_{2}} \in Buffer2$}
			\If{$(\vv{P_{1}}+\vv{P_{2}})^2 \leq r^2$}
				\ForAll{$\vv{P_{3}} \in Buffer3$}				
					\State $ZeichneVoxel(\vv{P_{3}}+\vv{P_{1}}+\vv{P_{2}})$
					\State $ZeichneVoxel(\vv{P_{3}}-\vv{P_{1}}-\vv{P_{2}})$ \Comment {Spiegelung}
					\State $ZeichneVoxel(\vv{P_{3}}+\vv{P_{1}}-\vv{P_{2}})$ \Comment {Spiegelung}
					\State $ZeichneVoxel(\vv{P_{3}}-\vv{P_{1}}+\vv{P_{2}})$ \Comment {Spiegelung}
				\EndFor				
			\EndIf	
		\EndFor
	\EndFor

\end{algorithmic}
\end{algorithm}

%\ \\
\noindent Problematisch ist, dass es beim Anwenden von Algorithmus \ref{VoxelZylinder} unter der Verwendung einer ZeichneVoxel()-Routine,
die genau einen Voxel an gegebener Position zeichnet, zur L�ckenbildung im Zylinder kommen kann.
Die Ursache liegt darin, dass der 3D-Bresenham-Algorithmus beim Zeichnen von Linien die 26er-Nachbarschaftsbeziehung verwendet,
�quivalent zum 2D-Bresenham, der die 8er-Nachbarschaft nutzt.
Abbildung \ref{B_Voxelluecken} zeigt die Entstehung von L�cken beim Zeichnen einer Fl�che.
�quivalent kommt es bei der Kombination der 3D-Bresenhams zu einem volumetrischen Gebilde potentiell zu L�cken.
Jede dieser L�cken besitzt aber grunds�tzlich mindestens einen gesetzten Voxel in 6er-Nachbarschaft.

\begin{figure}[hbtp]
  \centering  
	\includegraphics[width=15cm]{Bilder/Voxelluecken}
	\caption[Entstehung von L�cken beim Zeichnen des Zylinders]{\emph{Entstehung von L�cken beim Zeichnen des Zylinders}:
	 betrachtet wird 2D-Fall: Fig. (a) und (b) zeigen Teilauschnitte einer Linie. Werden deren Punkte als Startpunkte neuer Linien verwendet,
	 so k�nnen je nach Linienrichtung L�cken entstehen.}
	 \label{B_Voxelluecken}
\end{figure}

Dadurch ergibt sich als m�gliche L�sung, dass Voxelbl�cke statt einzelner Voxel gezeichnet werden.
Diese Bl�cke sollten ann�hernd kugelf�rmig bez�glich eines Mittelpunktsvoxels sein, damit die Symmetrie des Zylinders gewahrt bleibt.
Wenn man einen 6er-Nachbarschafts-Voxelblock zeichnet (Bezugsvoxel und alle Voxel in 6er-Nachbarschaft),
so werden die L�cken innerhalb des Zylinders geschlossen.
Allerdings k�nnen, wie in Abbildung \ref{B_Voxelluecken_Rand} zu sehen ist, noch Scharten am Rand des Zylinders auftreten.
Diese k�nnen durch Verwendung von 18er-Nachbarschafts-Voxelbl�cken eliminiert werden.
Durch das Verwenden dieser Voxelbl�cke muss der Parameter $r$ beim Aufruf von Algorithmus \ref{VoxelZylinder}
um $1$ verringert werden, um den korrekten resultierenden Radius zu erhalten.

\begin{figure}[hbtp]
  \centering  
	\includegraphics[width=14cm]{Bilder/Voxelluecken_Rand}
	\caption[Scharten am Rand des Zylinders]{\emph{Scharten am Rand des Zylinders}:
	 betrachtet wird Zylinderrand, die Zahlenwerte geben die relative H�he an:
	 Sind initial L�cken am Rand (Fig. (a), rote Quadrate), so werden diese nicht per 6er-Nachbarschafts-Voxelblock geschlossen (Fig. (b)).}
	 \label{B_Voxelluecken_Rand}
\end{figure}

Eine Verbesserungsidee stellt das Schlie�en der L�cken schon in der Grundfl�che dar, also vor dem eigentlichen Zeichnen des Zylinders.
Dadurch k�nnen einzelne Voxel statt Bl�cken gezeichnet werden, der Zylinder ist runder und damit exakter.
Diese Methode ist je nach Aufwand der Modifikation m�glicherweise langsamer,
besonders bei kurzen Zylindern, bei denen die Berechnung der Grundfl�che einen relativ hohen Anteil am Gesamtaufwand hat.
Zudem w�ren mehr Elemente in den Buffern, die durchlaufen werden m�ssten.

\ \\
Anmerkung: Algorithmus \ref{VoxelZylinder} ist nicht auf gerade Kreiszylinder beschr�nkt.
Wenn $\left(\vv{P_e}-\vv{P_s}\right)$ nicht senkrecht auf der Ebene, die durch $\vv{Oben}$ und $\vv{Links}$ aufgespannt wird, steht, so wird ein schiefer Zylinder gezeichnet.
Durch Verwendung von Vektoren $\vv{Oben}$ und $\vv{Links}$ mit einer L�nge ungleich $1$, lassen sich Zylinder mit elliptischer Grundfl�che zeichnen.
Der Test in Zeile 6 des Algorithmus muss dabei wie in Formel \ref{ElliptischerZylinder} angepasst werden,
wobei $r$ nicht mehr f�r den Radius, sondern f�r einen Skalierungsfaktor steht.
\begin{equation}
 \label{ElliptischerZylinder}
  \left( \frac{\vv{P_{1}}}{\left|\vv{Oben}\right|} + \frac{\vv{P_{2}}} {\left|\vv{Links}\right|} \right)^2 \leq r^2
\end{equation}

%\ \\
%Laufzeitbetrachtung:
%- beide Algorithmen: linear in Anzahl bearbeiteter Voxel, deren Anzahl von Radius abh�ngig
%- Zylinder: O(r*r*l) l...L�nge, r...Radius
%- Kugel: O(r*r*r)

%================================================================================================
\subsection{Abarbeitung der Zeichenanweisungen}

Das Zeichen der Gesamtstruktur erfolgt in zwei Schritten.
In beiden Schritten sind die Vektoren des lokalen Koordinatensystems der Turtle initial $\vv{Vorn}=(1,0,0)$, $\vv{Oben}=(0,1,0)$ und $\vv{Links}=(0,0,1)$.

Im ersten Schritt werden die Anweisungen zwar vollst�ndig abgearbeitet,
aber statt Striche wirklich zu zeichnen, wird nur die Ausdehnung der Gesamtstrukur bestimmt,
um sie danach im vordefinierten Zeichenbereich, also dem Voxelraum, darstellen zu k�nnen.

Hierzu werden vor Abarbeitung der Anweisungen die Strichl�nge $L=1$ und die Startposition $\vv{P_0}=(0,0,0)$ gesetzt.
Eine Bounding Box wird erstellt, die anfangs nur die Startposition beinhaltet.
Nach jeder Positions�nderung wird die Bounding Box aktualisiert, so dass diese die Gesamtstruktur umschlie�t.
Am Ende von Schritt 1 erfolgt die Berechnung der neuen Daten.
Die Startposition wird so gesetzt, dass der Mittelpunkt der Struktur im Mittelpunkt des Voxelraums liegt.
Die Strichl�nge wird so angepasst, dass die gesamte Struktur in den vordefinierten Voxelraum passt.
Dies geschieht unter Ber�cksichtigung des maximalen Strichradius, der dem Startradius entspricht.
Der Startradius selbst wird nicht modifiziert.

\ \\
Im zweiten Schritt wird tats�chlich gezeichnet.
Zuerst wird eine Kugel mit dem Startradius an der Startposition $\vv{P_0}$ zum Abschlie�en des ersten Striches gesetzt.
Dann erfolgt die schrittweise Abarbeitung der Zeichenanweisungen unter Verwendung der im ersten Schritt bestimmten Werte f�r $L$ und $\vv{P_0}$.

Zur Optimierung werden direkt aufeinander folgende Striche zu einem Strich zusammengefasst.
Dazu werden Folgen von '$F$' und '$[$' kombiniert.
Bei Symbol '$F$' wird die aktuelle Position errechnet, bei Symbol '$[$' wird der aktuelle Zustand gespeichert.
Der Strich wird erst gezeichnet, wenn er die komplette L�nge erreicht hat.
Dadurch werden Verbindungskugeln zwischen den einzelnen Strichen und das Berechnen der Grundfl�che f�r weitere Zylinder eingespart.

F�r die Position der Turtle werden Gleitkommazahlen benutzt.
Bei gr��eren Strukturen kann die Strichl�nge kleine Werte annehmen,
wodurch sich die Position nicht zwangsl�ufig �ndert, wenn eine Zeichenanweisung '$F$' vorliegt.
Deswegen werden Striche nur dann gezeichnet, wenn, in Voxelkoordinaten umgerechnet, die neue Position ungleich der alten Position ist.
Falls sich die neue Position von der alten Position in jeder Dimension um h�chstens einen Voxel unterscheidet,
reicht es, nur eine Kugel zu zeichnen, statt Zylinder und Kugel.

\ \\
Den Hauptanteil an der Laufzeitkomplexit�t des Verfahrens tr�gt im Normalfall der Zeichenalgorithmus f�r Striche,
welcher auf der Laufzeitkomplexit�t f�r das Zeichnen von Kugel und Zylinder beruht.
Er ist linear in der Anzahl der gezeichneten Voxel, wobei sich aber Striche untereinander �berlagern k�nnen
und so Voxel mehrfach gezeichnet werden.
Mit einer Strichl�nge von $L$ und einem Startradius von $r$ ergibt sich eine Laufzeitkomplexit�t pro Strich von
$\mathcal O(L\cdot r^2)$ f�r den Zylinder plus $\mathcal O(r^3)$ f�r die Kugel, also $\mathcal O(L\cdot r^2 + r^3)$ insgesamt.

Mit $n$ Strichen und $m$ weiteren Anweisungen ergibt sich als Gesamtlaufzeitkomplexit�t
$\mathcal O(m + n \cdot L\cdot r^2 + n \cdot r^3)$.

%Auf Grund von �berlagerungen werden allerdings Voxel mehrfach gezeichnet.
%bezogen auf den Radius $r$ ergibt sich somit $\mathcal O(r^2L)$ f�r den Zylinder plus $\mathcal O(r^3)$ f�r die Kugel, also $\mathcal O(r^3)$ insgesamt.
%Da die Strichl�nge $L$ von Struktur zu Struktur stark variiert, ist die Einbeziehung dieser schwer umsetzbat.
%Stattdessen l�sst sich der Radius als Bezugs-Parameter f�r die Laufzeitkomplexit�t sehr gut verwenden,
%da eine Ver�nderung des Radius die Grundstruktur einer Zeichnung �hnlich l�sst.Als Laufzeitkomplexit�t bezogen auf den Radius $r$ ergibt sich somit $\mathcal O(r^2)$ f�r den Zylinder plus $\mathcal O(r^3)$ f�r die Kugel,
%also $\mathcal O(r^3)$ insgesamt.

%================================================================================================
\section{Nachbearbeitung}

Die nun gezeichnete Voxelgrafik wird nachbearbeitet, um die Struktur h�hlenartiger und realistischer zu gestalten.
Zu den Nachbearbeitungsschritten geh�rt zum einen die Erosion und zum anderen die Filterung schwebender Fragmente.
F�r diese Schritte sind die Begriffe \emph{Randvoxel} und \emph{Randvoxel im erweiterten Sinn} relevant.
Ein \emph{Randvoxel} ist ein Voxel mit Belegung $0$ mit mindestens einem Nachbarvoxel in 6er-Nachbarschaft, der die Belegung $1$ besitzt.
Ein \emph{Randvoxel im erweiterten Sinn} ist ein Voxel mit Belegung $0$ mit mindestens einem Nachbarvoxel in 18er-Nachbarschaft, der die Belegung $1$ besitzt.

\subsection{Erosion}
\label{K_Erosion}
\begin{figure}[htbp]
  \centering  
  \hfill
  \subfloat{\includegraphics[width=7.4cm]{Bilder/Screenshot_Erosion_ohne}}
  \hfill
  \subfloat{\includegraphics[width=7.4cm]{Bilder/Screenshot_Erosion_mit}}  
  \hfill
	\caption[Erosion im Voxelraum]{\emph{Erosion im Voxelraum}: v.l.n.r. (a) ohne Erosion, (b) mit Erosion (ein Erosionschritt mit $p_e=0,5$)}
	\label{B_Erosion}
\end{figure}
\noindent Die Erosion bewirkt Zerkl�ftung an den R�ndern der Grafik und ist mehrfach anwendbar.
Die Auswirkungen sind in Abbildung \ref{B_Erosion} zu sehen.
Erodiert werden grunds�tzlich nur \emph{Randvoxel} der Struktur.

Bei der Durchf�hrung des Verfahrens wird jeder Randvoxel mit einer Wahrscheinlichkeit, die Erosionswahrscheinlichkeit $p_e$ genannt wird, auf $1$ gesetzt.
Im ersten Schritt wird mittels einer Schleife �ber den kompletten Voxelraum gelaufen und jeder Randvoxel mit der Wahrscheinlichkeit $p_e$ mit $2$ markiert.
Im zweiten Schritt werden alle Voxel des Voxelraums mit Markierung $2$ auf den Wert $1$ gesetzt, also erodiert.
Durch Aufteilung in zwei Schritte werden Konflikte vermieden, da bei nur einem Schritt die Nachbarn neu erodierter Voxel sofort wieder erodiert werden k�nnten.
Dadurch k�nnte die Erosion bei schon einer Anwendung des Erosionsverfahrens prinzipiell beliebig weit voranschreiten, statt nur einen Voxel tief.

Der Erosionsprozess hat eine Laufzeitkomplexit�t von $\mathcal O(n)$, mit $n$ als Anzahl aller Voxel im Voxelraum.

\subsection{Filterung schwebender Fragmente}
\label{K_Filterung}
\begin{figure}[htbp]
  \centering  
  \hfill
  \subfloat{\includegraphics[width=7.4cm]{Bilder/Screenshot_schwebendes_Fragment}}
  \hfill
  \subfloat{\includegraphics[width=7.4cm]{Bilder/Screenshot_schwebendes_Fragment2}}  
  \hfill
	\caption[Schwebende Fragmente im Voxelraum]{Schwebende Fragmente im Voxelraum}
	\label{B_SchwebendesFragment}
\end{figure}
\noindent Schwebende Fragmente entstehen bei der Umschlie�ung von Gebieten durch gezeichnete Striche,
bei der Erosion oder durch das Hineinschneiden von Gang�ffnungen in die Voxelstruktur (siehe Kapitel \ref{KK_AndockenHoehle}).
Abbildung \ref{B_SchwebendesFragment} zeigt solche Fragmente.
Sie sind i.d.R. f�r H�hlen unrealistisch, da sie dem Wirken der Gravitation widersprechen.
Durch die Filterung werden diese Fragmente entfernt.

Es sei im Folgenden angenommen, dass jeder Voxel am Rand des Voxelraums die Belegung $0$ habe.
Um ein schwebendes Fragment zu charakterisieren, wird Definition \ref{Zusammenhangskomponente} ben�tigt.
Bei schwebenden Fragmenten handelt es sich um Zusammenhangskomponenten, die nicht der �u�eren Zusammenhangskomponente entsprechen.

\begin{TheoremDefinition}[Zusammenhangskomponente]
		Eine Zusammenhangskomponente \\ besteht ausschlie�lich aus Voxeln mit Belegung $0$ und
		jeder Voxel mit Belegung $0$ ist Teil einer Zusammenhangskomponente.
		Ein Voxel geh�rt zu einer Zusammenhangskomponente M, wenn er einen Nachbarvoxel in 6er-Nachbarschaft besitzt, der zu M geh�rt.
		Ein Voxel kann nicht mehreren Zusammenhangskomponenten angeh�ren.
		
		Die �u�ere Zusammenhangskomponente ist die Zusammenhangskomponente, der der Voxel mit den Koordinaten $(0,0,0)$ angeh�rt.
    \label{Zusammenhangskomponente}
\end{TheoremDefinition}

Der erste Ansatz des Filterungsalgorithmus verwendet eine Tiefensuche oder wahlweise Breitensuche �ber alle $0$-Voxel von $(0,0,0)$ aus beginnend.
Dabei wird der Voxelraum als Graph aufgefasst, mit den Voxeln als Knoten und Kanten zwischen Voxeln gleicher Belegung in 6er-Nachbarschaft.
Alle gefundenen Voxel werden mit $2$ markiert, wodurch genau und ausschlie�lich die �u�ere Zusammenhangskomponente komplett mit $2$ markiert wird.
%
Im zweiten Schritt werden alle $0$-Voxel zu $1$-Voxeln abge�ndert, da diese zu schwebenden Fragmenten geh�ren.
Danach werden alle $2$-Voxel wieder zu $0$-Voxeln umgelabelt.
	
Problematisch ist, dass die �u�ere Zusammenhangskomponente eine meist hohe Voxelanzahl besitzt.
Dies bedeutet einen hohen Speicherverbrauch, da im Stack (f�r die Tiefensuche) bzw. in der Queue (f�r die Breitensuche)
die Position der Voxel gespeichert werden muss.
Diese umfasst je nach Aufl�sung des Voxelraums mehrere Byte.\footnote{Bei der verwendeten Aufl�sung von $512$ Voxeln Seitenl�nge mindestens $27$ Bit.}

\ \\
Die zweite und implementierte Variante wird in Algorithmus \ref{VoxelFilter} aufgezeigt.%\footnote{Die
%tats�chlichen Schleifenschranken lassen sich hierbei noch etwas sch�rfer setzen.}
Gegeben sei der Voxelraum mit seiner Ausdehnung von $[0,X_{vmax}] \times [0,Y_{vmax}]\times [0,Z_{vmax}]$. \\
\emph{Schritt 1}:
Beginnend vom Rand des Voxelraums wird ein Randvoxel der �u�eren Zusammenhangskomponente gesucht.
F�r einen schnellen Test wird nur gepr�ft, ob der aktuelle Voxel$(x,y,z)=0$ und Voxel$(x+1,y,z)=1$ gilt, anstatt alle Nachbarn zu betrachten.
Die Suche wird mit $x=0$ begonnen und mit aufsteigenden $x$-Werten fortgesetzt.
Der erste so gefundene Voxel muss zur �u�eren Zusammenhangskomponente geh�ren. \\
\emph{Schritt 2}:
Von diesem Voxel aus wird eine Breitensuche �ber alle Randvoxel im erweiterten Sinn durchgef�hrt und alle gefunden Voxel auf $2$ gesetzt. 
Damit wird der komplette Rand der �u�eren Zusammenhangskomponente, also die \emph{�u�ere H�lle}, mit $2$ markiert. \\
\emph{Schritt 3}:
Nun wird der restliche Voxelraum gepr�ft und erneut nach Randvoxeln gesucht.
Dadurch, dass schwebende Fragmente im Inneren von 1er-Voxel-Gebieten liegen, kann der zu durchsuchende Bereich
in jeder Dimension um einige Voxel eingeschr�nkt werden.
Wenn ein Randvoxel gefunden wird, muss dieser zu einem schwebenden Fragment geh�ren.
Nun wird eine Breitensuche �ber die komplette Zusammenhangskomponente durchgef�hrt und alle gefundenen Voxel werden auf $1$ gesetzt,
wodurch das schwebende Fragment entfernt wird. \\
\emph{Schritt 4}:
Zum Schluss werden alle Voxel mit Belegung $2$ wieder auf $0$ abge�ndert.

\ \\
Bei dieser Variante werden nur die Voxel des Rands der �u�eren Zusammenhangskomponente
sowie die Voxel der schwebenden Fragmente in die Queue geladen und dies einzeln Komponente f�r Komponente.
Obwohl beide Varianten einen Komplexit�t bzgl. des Speicherverbrauchs von $\mathcal O(n)$ haben,
wird der letztere Algorithmus in der Praxis i.A. weniger Speicher ben�tigen als der erstere.
Beide Verfahren haben eine Laufzeitkomplexit�t von $\mathcal O(n)$ .

%- �u�ere H�lle: Randvoxel im erweiterten Sinn der �u�eren Zusammenhangskomponente

%\begin{figure}[hbtp]
%  \centering  
%	\includegraphics[width=10cm]{Bilder/Filterung}
%	\caption[Filterung schwebender Fragmente - Gebietsgrenzen]{\emph{Filterung schwebender Fragmente - Gebietsgrenzen}: 
%	Voxel sind wei� markiert, Voxel mit Belegung 1 grau, Voxel der �u�eren H�lle rot. \\
%	Die Grenzen der Gebiete sind eingezeichnet: 
%	Der Startvoxel der �u�eren H�lle wird im Bereich  $[0,X_{vmax}-2] \times [1,Y_{vmax}-1]\times [1,Z_{vmax}-1]$
%	Die �u�ere H�lle erstreckt sich im Bereich $[E,X_{vmax}] \times [0,Y_{vmax}]\times [0,Z_{vmax}]$.
%	Die schwebenden Fragmente liegen im Bereich $[E+2,X_{vmax}-2] \times [2,Y_{vmax}-2]\times [2,Z_{vmax}-2]$.
%	}
%	\label{B_Filterung}
%\end{figure}

\begin{algorithm}
\caption{Filterung schwebender Fragmente im Voxelraum}
\label{VoxelFilter}
\algorithmicrequire{ungefilterter Voxelraum, mind. 1 Voxel breiter Rand des Voxelraums mit $0$ belegt} \\
\algorithmicensure{gefilterter Voxelraum}
\begin{algorithmic}[1]

	\State $\triangleright$ Schritt 1: Finde Startvoxel f�r �u�ere H�lle
	\For{$x = 0 \to X_{vmax}-2$} \Comment{Suche mit aufsteigender X-Koordinate}
		\ForAll{$(y,z) \in [1,Y_{vmax}-1]\times [1,Z_{vmax}-1]$}
			\If{$Voxel(x,y,z)=0 \wedge Voxel(x+1,y,z)=1$}
				\State \textbf{put} $(x,y,z)$ \textbf{in} Queue	
				\State $E \gets x$		\Comment{niedrigste X-Koordinate der �u�eren H�lle}
				\State \textbf{break}	\Comment{Abbrechen der Suche}
			\EndIf
		\EndFor
	\EndFor
	
	%\State
	\State $\triangleright$ Schritt 2: Markierung der �u�eren H�lle
	\While{$Queue\neq leer$}
		\State \textbf{take} $\vv{a}$ \textbf{from} Queue	
		\ForAll{$\vv{b} \in 6er\ Nachbarschaft$ \textbf{of} $\vv{a}$}		
			\If{$\vv{b} \equiv Randvoxel\ im\ erweiterten\ Sinn$}
			%\If{$Voxel(\vv{b}) = 0 \wedge \exists  \left(Voxel(\vv{c})=1\right) \in 18er\ Nachbarschaft$ \textbf{of} $\vv{b}$}
				\State \textbf{put} $\vv{b}$ \textbf{in} Queue	
				\State $Voxel(\vv{b}) \gets 2$	\Comment{Markieren des Voxels}
			\EndIf		
		\EndFor
	\EndWhile
	
	%\State
	\State $\triangleright$ Schritt 3: L�schung freischwebender Komponenten
	\ForAll{$(x,y,z) \in [E+2,X_{vmax}-2] \times [2,Y_{vmax}-2]\times [2,Z_{vmax}-2]$}
		\If{$Voxel(x,y,z)=0 \wedge Voxel(x+1,y,z)=1$}
			\State \textbf{put} $(x,y,z)$ \textbf{in} Queue
			\State $Voxel(x,y,z) \gets 1$	\Comment{L�schen des Voxels}
		
			\While{$Queue\neq leer$}
				\State \textbf{take} $\vv{a}$ \textbf{from} Queue	
				\ForAll{$\vv{b} \in 6er\ Nachbarschaft$ \textbf{of} $\vv{a}$}		
					\If{$Voxel(\vv{b}) = 0$}
						\State \textbf{put} $\vv{b}$ \textbf{in} Queue	
						\State $Voxel(\vv{b}) \gets 1$	\Comment{L�schen des Voxels}
					\EndIf		
				\EndFor
			\EndWhile		
		
		\EndIf
	\EndFor	
	
	%\State
	\State $\triangleright$ Schritt 4: Originalmarkierung der �u�eren H�lle wiederherstellen
	\ForAll{$(x,y,z) \in [E,X_{vmax}] \times [0,Y_{vmax}]\times [0,Z_{vmax}]$}
	 	\If{$Voxel(x,y,z) = 2$}
	 		\State $Voxel(x,y,z) \gets 0$
	 	\EndIf		
	\EndFor
	
\end{algorithmic}
\end{algorithm}


\subsection{Anmerkung zur Optimierung}
\label{K_Voxeloptimierung}

Algorithmus \ref{VoxelFilter} ben�tigt zur Markierung der �u�eren H�lle mindestens einen 1 Voxel breiten Rand aus Voxeln mit Belegung 0.
Zur Vermeidung von Sonderf�llen an R�ndern ist es sinnvoll, diesen Rand auf drei Voxel zu erweitern.
Ein Voxel wird dabei f�r die Markierung der �u�eren H�lle ben�tigt plus einen Voxel f�r die Nachbarschaft
plus einen Voxel f�r die Pr�fung, ob dieser Nachbarvoxel ebenfalls ein Randvoxel im erweiterten Sinn ist.

Dieser Rand ist ebenfalls ausreichend, damit die Erosion, die im n�chsten Kapitel folgende Umwandlung
sowie das Markieren von Andockstellen (Kapitel \ref{K_Architekt}) ohne Randpr�fungs-Sonderf�lle arbeiten k�nnen.
In anf�nglichen Implementierungen wurde hierdurch ein Speedup von bis zu Faktor $5$ erzielt. %$3$ bis

%
%Bei der Implementierung des Zeichenalgorithmus als auch der Erosion ist es sinnvoll, einen minimalen Rand des Voxelgebildes zu den Voxelraumgrenzen zu wahren,
%welcher immer eine Belegung ungleich $1$ besitzt.
%Dies erm�glicht es den Algorithmen der Erosion, des Filterns, der im n�chsten Kapitel folgenden Umwandlung sowie dem Markieren von Andockstellen (Kapitel \ref{K_Architekt})
%keine Sonderf�lle f�r die R�nder pr�fen zu m�ssen.
%Dies brachte in anf�nglichen Implementierungen einen durchschnittlichen Speedup von Faktor $3$ bis $5$.
%
%Damit komplett ohne Sonderf�lle gearbeitet werden kann, muss der Voxelrand mindestens drei Voxel betragen.
%Dies ist auf den Filterungsalgorithmus zur�ckzuf�hren:
%Dieser Voxelrand ist ebenfalls f�r Erosion, Umwandlung und dem Markieren der Andockstellen ausreichend.
%=====================================
%Implementierung
%=====================================
%Code: Grundlegende Hierarchie
%
%Fraktalh�hle erzeugen:
%----------------------
%
%VoxelZeichnen() .. Zeichnet Voxel und schiebt ihn in die Queue
%
%LinieZeichnen() ... zeichnet 3D-Linie mit 3D-Bresenham in den Voxelspace, nutzt zum Zeichnen VoxelZeichnen()
%
%PrimitivZeichnen() ...  zeichnet grafisches Primitiv , nutzt zum Zeichnen LinieZeichnen()
%
%GrafikZeichnen() ... zeichnet eien Turtle-Grafik (hier: Fraktal) indem String ausgelesen wird, nutzt zum Zeichnen PrimitivZeichnen()
			% Erstellung der H�hle
\chapter{Umrechnung der Voxelh�hle in Dreiecksnetz}
\label{KK_Umwandlung}

Die H�hlenstruktur liegt nun als Voxelzeichnung vor.
Diese soll in ein entsprechendes Dreiecksnetz umgewandelt werden.

Jeder Voxel soll einem W�rfel der Kantenl�nge $1$ entsprechen,
wobei ein Voxel mit Position $(x,y,z)$ den W�rfel $[x,x+1]\times[y,y+1]\times[z,z+1]$ widerspiegelt.
F�r das Dreiecksnetz sind nur die $0\Leftrightarrow 1$-�berg�nge, also die Grenzen des H�hlenfreiraums zu festem Gestein, relevant.
Jeder dieser �berg�nge entspricht einer quadratischen Fl�che, die durch zwei Dreiecke dargestellt werden kann.
Jedes dieser Dreieckspaare ben�tigt vier Vertices, die allerdings mit anderen Dreieckspaaren geteilt werden k�nnen.
Die Beziehung zwischen Voxeln und Vertices ist in Abbildung \ref{B_VoxelZuVertex}(a) aufgezeigt.

Gesucht ist ein Algorithmus mit der Aufgabe, aus allen �berg�ngen das resultierende Dreicksnetz zu bilden.
Dabei ist eine Aufteilung der H�hle in mehrere Subnetze notwendig,
da die Irrlicht Engine auf maximal $65536$ Vertices pro Meshbuffer begrenzt ist (siehe Kapitel \ref{KK_Dreiecksnetze}).
Die Aufteilung der H�hle in mehrere Subnetze ist ebenso sinnvoll,
da sich in Abh�ngigkeit von Position und Ausrichtung des Betrachters so
Teile der H�hle leichter ausblenden oder in niedrigeren Detailstufen darstellen lassen.

\begin{figure}[hbtp]
  \centering  
	\includegraphics[width=15cm]{Bilder/VoxelZuVertex}
	\caption[Voxelkoordinaten zu Vertexkoordinaten]{\emph{Voxelkoordinaten zu Vertexkoordinaten}: v.l.n.r. \\
	(a) einzelner Voxel mit Voxelkoordinaten $(x,y,z)$ mit seinen acht begrenzenden Vertices (gr�n) und ihren Vertexkoordinaten \\
	(b) Umwandlung eines Voxelgebildes per Sweepverfahren
	}
	\label{B_VoxelZuVertex}
\end{figure}

%================================================================================================
%================================================================================================
\section{Umwandlungsalgorithmus}
 
\emph{Sweep-Verfahren} funktionieren nach dem Prinzip der Dimensionsreduzierung \cite[S.51 ff.]{B_AlgGeo}.
Ein statisches $d$-dimensionales Problem wird in ein dynamisches $(d-1)$-dimen"-sionales Problem umgewandelt.
Eine r�umliche Dimension wird so zu einer zeitlichen Dimension.

Im Fall des 3D-Raums ergibt sich eine \emph{Sweepebene}. Die zeitliche Dimension ist das Sweep\-en entlang einer Achse.
Sinn und Zweck von Sweep-Verfahren ist das Reduzieren von Laufzeiten (da weniger Elementbeziehungen gleichzeitig betrachtet werden m�ssen)
oder das Einsparen von Speicherplatz.
Im Umwandlungsalgorithmus wird es zur Reduzierung des ben�tigten Speicherplatzes verwendet.

Ein \emph{Octree} ist eine Struktur, die ein gegebenes quaderf�rmiges Raumgebiet in acht kleinere gleichgro�e Raumgebiete aufteilt,
die Oktanten genannt werden \cite[S.237-238]{B_3DGameProg}.
Die Aufteilung des Raumgebietes erfolgt jeweils mittig des Quaders durch drei senkrecht aufeinanderliegende Schnittebenen.
Die Oktanten werden bei Bedarf auf die gleiche Weise weiter unterteilt.
Die entstehende Baumstruktur bildet den Octree.
Die Raumgebiete stellen die Knoten des Baums dar, die Kanten verlaufen vom jeweiligen Raumgebiet zu den zugeh�rigen Oktanten.

%- Voxel: M�glichkeiten zur Umwandlung
%	-	W�rfel (max. 12 Dreiecke)
%	- Tetraeder (max. 8 Dreiecke)
%	- ein Vertex pro Voxel:
%		- Vorteil: weniger Vertexe (nur auf den ersten Blick: W�rfelmethode ebenso sparsam!)
%		- aber: Vermeidung zweidimensionaler Geometrie beachten
%	[BILD: Vergleich der Methoden in 2D]
%	-> Vergleichstabelle: Anzahl der m�glichen Vertexe, Genauigkeit, Winkelvielfalt
%		 W�rfel: Anzahl der Voxel, perfekt, wenig (nur 90�)
%		 Tetraeder: 3 * Anzahl der Voxel  , schlecht, viel (45�,90�)
%		 Vertex: Anzahl der Voxel, schlecht, viel (45�,90�)
%-> Nutzung der W�rfelabtastung (Grund: Exaktheit)

%- nun Umrechnung der gesetzten Voxel in Dreiecksnetz
%- nur Randvoxel (1-Voxel mit 0-Voxel-Nachbarn) brauchen betrachtet werden
\ \\
Der Octree wird zur Aufteilung des Netzes verwendet, so dass in jedem Teilnetz maximal $V_{max}$ viele Vertices vorhanden sind.
In der Implementierung wird $V_{max} = 65000$ gew�hlt, welches unter dem Maximum der Irrlicht-Engine je Meshbuffer liegt.
Die Tiefe des erzeugten Octrees liegt in $\mathcal O(\log n)$ mit $n$ als Anzahl aller Voxel.

Die klassische Methode zum Aufbau des Octrees ist die \emph{Top-Down}-Variante.
Begonnen wird mit dem gesamten Voxelraum als Startknoten.
F�r jeden Knoten erfolgt das Z�hlen der ben�tigten Vertices $V$ mittels Algorithmus \ref{VoxelVerticesZaehlen}.
Wenn $V > V_{max}$ ist, dann wird der entsprechende Knoten in acht gleichgro�e Kindknoten aufgeteilt,
mit denen �quivalent verfahren wird.
Die Laufzeitkomplexit�t der Methode ist $\mathcal O(n \log n)$, da in jeder Ebene des Octrees schlimmstenfalls alle Voxel gepr�ft werden m�ssen.

Die schnellere Variante ist die \emph{Bottom-Up} Methode.
Hier wird zuerst der Octree soweit aufgebaut, bis jeder Knoten maximal so gro� ist, dass seine \emph{theoretische} maximale Vertexanzahl h�chsten $V_{max}$ betr�gt.
F�r $V_{max} = 65000$ w�re dies bei einer Ausgangsgr��e des Voxelraums von $512^3$ eine Knotengr��e von $32^3$, welches maximal $33^3=35937$ Vertices entsprechen w�rde.
Eine Z�hlung mittels Algorithmus \ref{VoxelVerticesZaehlen} findet nur f�r die unterste Ebene des Octrees, also die Bl�tter, statt.
Nun werden von den Kindknoten aus bis zur Wurzel hin die Werte $V$ der Knoten zusammenaddiert.
Dabei berechnet sich $V$ eines Elternknotens aus der Summe der $V$ der Kindknoten:

\begin{equation}
 \label{OctreeBottomUpSumme}
 			V_{Eltern} = \sum\limits_{i=0}^{7} \left(V_{Kind\ i}\right) 
\end{equation}

\ \\
Durch dieses Verfahren kann $V$ eines Nichtblattknotens h�her liegen als die tats�chlich Anzahl an Vertices, die f�r diesen Knoten ben�tigt werden w�rden.
Keinesfalls liegt sie aber darunter. Die Werte $V$ stellen somit obere Schranken dar.
Dieses Verfahren hat eine Laufzeitkomplexit�t von $\mathcal O(n)$ und liegt damit im Punkt Geschwindigkeit vor der Top-Down-Variante.\footnote{Ausnahme:
Octree hat nur eine Hierarchieebene. Hier liegt die Top-Down-Variante minimal vor der Bottom-Up-Variante.}
In der Implementierung des Programms wird die Bottom-Up-Methode verwendet.

\begin{algorithm}
\caption{Umwandlung Voxelgebiet in Dreiecksnetz: Z�hlen der Vertices}
\label{VoxelVerticesZaehlen}
\algorithmicrequire{Voxelraum, Grenzen des Voxelgebietes: $X_{min}$, $X_{max}$, $Y_{min}$, $Y_{max}$, $Z_{min}$, $Z_{max}$. $S[a]$ und $S[b]$ sind mit $0$ initialisiert.}\\
\algorithmicensure{Anzahl ben�tiger Vertices: $V$}
\begin{algorithmic}[1]
	\State $V \gets 0$

	\For{$x = X_{min} \to X_{max}$}	\Comment{Sweepen entlang der X-Achse}
		\ForAll{$(y,z) \in [Y_{min},Y_{max}]\times [Z_{min},Z_{max}]$}
			  \If{$Voxel(x,y,z) = 1$}
			  	\If{$Voxel(x-1,y,z) = 0$}
			  		\State $S[a](y,z) \gets 1\ ,\ S[a](y+1,z) \gets 1$
			  		\State $S[a](y,z+1) \gets 1\ ,\ S[a](y+1,z+1) \gets 1$
			  	\EndIf			  	
			  	\If{$Voxel(x+1,y,z) = 0$}
			  	  \State $S[b](y,z) \gets 1\ ,\ S[b](y+1,z) \gets 1$
			  		\State $S[b](y,z+1) \gets 1\ ,\ S[b](y+1,z+1) \gets 1$
			  	\EndIf			  	 
			  	\If{$Voxel(x,y-1,z) = 0$}
			  	  \State $S[a](y,z) \gets 1\ ,\ S[a](y,z+1) \gets 1$
			  	  \State $S[b](y,z) \gets 1\ ,\ S[b](y,z+1) \gets 1$
			  	\EndIf			  	
			  	\If{$Voxel(x,y+1,z) = 0$}
			  	 	\State $S[a](y+1,z) \gets 1\ ,\ S[a](y+1,z+1) \gets 1$
			  	  \State $S[b](y+1,z) \gets 1\ ,\ S[b](y+1,z+1) \gets 1$
			  	\EndIf			  				  	
			  	\If{$Voxel(x,y,z-1) = 0$}
			  	  \State $S[a](y,z) \gets 1\ ,\ S[a](y+1,z) \gets 1$
			  	  \State $S[b](y,z) \gets 1\ ,\ S[b](y+1,z) \gets 1$
			  	\EndIf			  				  	
			  	\If{$Voxel(x,y,z+1) = 0$}
			  	  \State $S[a](y,z+1) \gets 1\ ,\ S[a](y+1,z+1) \gets 1$
			  	  \State $S[b](y,z+1) \gets 1\ ,\ S[b](y+1,z+1) \gets 1$
			  	\EndIf			  
			 	\EndIf
		\EndFor
		
		\ForAll{$(y,z) \in [Y_{min},Y_{max}+1]\times [Z_{min},Z_{max}+1]$}
			\State $V \gets V+S[a](y,z)$ \Comment{Vertices z�hlen}
			\State $S[a](y,z) \gets 0 $	\Comment{Sweepebene a l�schen}
		\EndFor		
		\State $a\leftrightarrow b$	\Comment{Sweepebenen vertauschen}		
	\EndFor \Comment{Sweep abgeschlossen}
	
	\ForAll{$(y,z) \in [Y_{min},Y_{max}+1]\times [Z_{min},Z_{max}+1]$}
		\State $V \gets V+S[a](y,z)$			\Comment{verbleibende Vertices z�hlen}
	\EndFor		
\end{algorithmic}
\end{algorithm}

\ \\
Algorithmus \ref{VoxelVerticesZaehlen} nutzt f�r das Z�hlen der ben�tigten Vertices ein Sweepverfahren, um Speicherplatz zu sparen.
Statt $\mathcal O\left(\left(X_{max}-X_{min}\right)\cdot \left(Y_{max}-Y_{min}\right) \cdot \left(Z_{max}-Z_{min}\right)\right)$ zus�tzlichen Speicherplatz zu ben�tigen,
wird nur $\mathcal O\left(\left(Y_{max}-Y_{min}\right) \cdot \left(Z_{max}-Z_{min}\right)\right)$ ben�tigt.
Der Algorithmus durchl�uft das entsprechende Voxelgebiet und testet auf $0\Leftrightarrow 1$ �berg�nge zwischen benachbarten Voxeln.
Jeder dieser �berg�nge ben�tigt vier Vertices an entsprechender Position, um durch zwei Dreiecke dargestellt werden zu k�nnen.
Diese Vertices werden in den beiden Sweepebenen markiert. Bei jedem neuen Sweepschritt entlang der X-Achse wird die Anzahl der ben�tigten Vertices aus der �lteren
der beiden Sweepebenen ausgelesen und zur Gesamtanzahl dazuaddiert.
Das Sweepverfahren ist in Abbildung \ref{B_VoxelZuVertex}(b) skizziert.

\ \\
Nach dem Aufbau des Octrees findet der Umwandlungsschritt statt. Der Octree wird nun Top-Down durchlaufen.
F�r jeden Knoten wird getestet, ob $V \leq V_{max}$, wobei die vorher ausgerechneten Werte f�r $V$ verwendet werden.
Falls der Test positiv verl�uft, wird das zum Knoten geh�rige Gebiet umgewandelt.
Bei negativem Test wird stattdessen mit den Kindknoten weiter verfahren.

Zur Umwandlung eines Voxelgebietes in ein Dreiecksnetz wird Algorithmus \ref{VoxelVerticesZaehlen} entsprechend angepasst.
Statt den Eintr�gen der Sweepebenen eine $1$ zuzuweisen, wird zuerst �berpr�ft, ob ein Eintrag vorliegt.
Wenn nicht, wird ein neuer Vertex mit den entsprechenden Koordinaten (Sweepebene $a$: $x_{res} = x$, Sweepebene $b$: $x_{res} = x+1$, $y$ und $z$ werden �bernommen)
erstellt und ein Verweis auf diesen Vertex in der zugeh�rigen Sweepebene gespeichert.
Die Berechnung der finalen Koordinaten des Vertex findet wie in Kapitel \ref{K_VertexkoordinatenBerechnen} beschrieben statt.

Nach dem Testen der vier Sweepebeneneintr�ge werden zwei Dreiecke erstellt, die die jeweilige Voxelgrenzfl�che darstellen.
F�r die Eckpunkte der zwei Dreiecke werden die vier Vertices entsprechend der Eintr�ge in den beiden Sweepebenen verwendet.
Danach wird der n�chste der sechs F�lle abgearbeitet.

Das Z�hlen der Vertices entf�llt in diesem Algorithmus.
Der Umwandlungsprozess hat eine Laufzeitkomplexit�t von $\mathcal O(n)$ mit $n$ als Anzahl aller Voxel.

%================================================================================================
%================================================================================================

\section{Berechnung der Vertexkoordinaten}
\label{K_VertexkoordinatenBerechnen}

Die Berechnung der Vertexkoordinaten wird �ber eine Funktion realisiert, die aus Eingangskoordinaten Ausgangskoordinaten berechnet.
Dabei sind die Eingangskoordinaten ganzzahlig, da es sich um Voxeleckpunkte handelt, und die Ausgangskoordinaten Gleitkommazahlen.
Das Prinzip Eingangskoordinaten in Ausgangskoordinaten umzuwandeln
entspricht dem von Raumverzerrungsmethoden.
Andere oder auch zus�tzliche Raumverzerrungsverfahren zur hier verwendeten Funktion sind denkbar.

Wichtig dabei ist, dass die Funktion deterministisch arbeitet, d.h. bei gleichen Eingangskoordinaten die gleichen Ausgangskoordinaten ausgibt.
Dies ist relevant, damit die verschiedenen Subnetze der H�hle an den R�ndern zusammenpassen,
d.h. die korrespondierenden Randvertices m�ssen identische Koordinaten besitzen.
Weiterhin k�nnen so G�nge exakt an die H�hle angedockt werden (siehe Kapitel \ref{KK_AndockenHoehle}).

%, z.B. periodisches W�lben der XZ-Ebene des Raums mittels Sinus und Cosinusfunktionen
%- eine andere Verwacklungsfunktion k�nnte auf einem Perlin-Noise-Rauschen (!!!Quelle) basieren  
\ \\
Bei der Basisfunktion entsprechen die Ausgangskoordinaten genau den Eingangskoordinaten,
wodurch die W�rfelstruktur des Voxelraums erhalten bleibt.
Eine weitere Anpassung der Koordinaten f�r zus�tzlichen Realismus ist sinnvoll.
Statt die Voxelstruktur nur durch die W�rfelform wiederzugeben,
ist es besser, diese Struktur zu verzerren, so dass sie nat�rlicher wirkt.
Als grundlegende Strategie werden die Vertices \emph{verwackelt},
um dem Dreiecksnetz der H�hle eine zerkl�ftetere Struktur zu geben.

Beim Verwackeln wird der Eingangsvertex entlang aller drei Achsen per deterministischer Zufallsfunktion verschoben.
Die Zufallsfunktion wird mit einem einzigartigen Seed f�r jede m�gliche Eingangsposition initialisiert.
Jede Eingangsposition erh�lt so einzigartige Zufallswerte, die Funktion ist deterministisch.

Dann werden die Abweichungen von der Ausgabeposition zur Eingangsposition wie in Formel \ref{VerwacklungDeltaX} berechnet.
Die Formeln f�r $\Delta_{y}$ und $\Delta_{z}$ sind �quivalent.
Hierbei stellt $K$ die Verwacklungsst�rke dar, die zur Erhaltung der Symmetrie f�r alle drei Achsen gleich ist.
Die Funktion $Zufallswert[0,1]$ liefert eine zuf�llige Gleitkommazahl im Intervall $[0,1]$.
Resultat f�r $\Delta_{x}$, $\Delta_{y}$ und $\Delta_{z}$ sind Werte zwischen $-K$ und $K$.
\begin{equation}
 \label{VerwacklungDeltaX}
 			\Delta_{x} = -K + 2 \cdot K \cdot Zufallswert[0,1]
\end{equation}

\ \\
Es muss $K < 0,5$ gelten, damit keine Vertices aufeinanderliegen k�nnen, siehe hierzu Abbildung \ref{B_Clamping}(a).
Dies ist aber nicht ausreichend, da trotzdem Konflikte auftreten k�nnen, wie in (b) zu sehen ist.
Hier w�rden sich beide Dreiecke der Fl�che (grau) ineinander schieben.
Dem Problem kann vorgebeugt werden, indem der Vertex nur so verschoben wird, dass der Manhattenabstand zur Ausgangsposition
in jeder Ausgangsebene der Dreiecksfl�chen (XY, XZ, YZ) weniger als $0,5$ betr�gt, wie in (c) zu sehen ist.
Es muss also gelten:

\begin{equation}
 \label{VerwacklungManhattenXY}
 			\left|\Delta_{x}\right|+\left|\Delta_{y}\right| < 0,5
 			\ \wedge \ \left|\Delta_{x}\right|+\left|\Delta_{z}\right| < 0,5
 			\ \wedge \ \left|\Delta_{y}\right|+\left|\Delta_{z}\right| < 0,5
\end{equation}

\ \\
Der erste L�sungsansatz besteht in Begrenzung der Verwacklungsst�rke auf $< 0,25$.
Nachteil ist, dass Potential verschenkt wird, da Koordinaten so nur recht wenig abge�ndert werden.
Der neue Ansatz erlaubt weiterhin Verwacklungsst�rken $< 0,5$.
Hier wird nach der Berechnung von $\Delta_{x}$, $\Delta_{y}$ und $\Delta_{z}$ getestet, ob ein Problemfall vorliegt,
also der zul�ssige Manhattenabstand �berschritten wurde.
Wenn ja, dann wird dieser Konflikt aufgel�st.

\begin{figure}[hbtp]
  \centering  
	\includegraphics[width=15cm]{Bilder/Clamping}
	\caption[Verwacklung der Vertexkoordinaten]{\emph{Verwacklung der Vertexkoordinaten}:
	v.l.n.r. \\
	(a) �nderung einzelner Koordinaten $< 0,5$ (Grenzen: Quadrate,\\ rot gestrichelt), euklidischer Abstand $< 0,5$ (Kreise, rot) \\
	(b) beides kann zu Konflikten f�hren \\
	(c) keine Konflikte bei Manhattenabstand $< 0,5$ (Quadrate, rot)}
	\label{B_Clamping}
\end{figure}

\ \\
Dies geschieht mittels \emph{Clamping}.\footnote{In
der Computergrafik bezeichnet der Begriff "`Clamping"' das Begrenzen von Werten auf einen bestimmten Bereich (vgl. Color Clamping in OpenGL).}
Wenn der maximal zul�ssige Manhattenabstand $M$ �berschritten wird,
dann beschneide die Koordinatenabweichung so, dass keine �berschreitung von $M$ mehr vorliegt.
In der Implementierung wurde der Wert $M = 0,49$ gew�hlt. F�r $K$ muss gelten $0\leq K\leq M$.

Es k�nnen grunds�tzlich drei verschiedene Problemf�lle auftreten, die unterschiedliche Clampingstrategien sinnvoll machen.
Wichtig f�r die entwickelten Clampingstrategien ist es, die Ausgangswerte m�glichst wenig und m�glichst gleichm��ig zu ver�ndern.

\ \\
\emph{Fall 1}: Es liegt eine Verletzung der Bedingung in allen drei Ebenen vor.
L�sungsstrategie: Normalisiere alle Abweichungswerte ($\Delta_{x},\Delta_{y},\Delta_{z}$).
Nutze dazu folgendes Verfahren:
\begin{equation}
 \label{VerwacklungClamp3}
 \begin{aligned}
 			n &= \text{max}\left( \left|\Delta_{x}\right|+\left|\Delta_{y}\right|,\left|\Delta_{x}\right|+\left|\Delta_{z}\right|,\left|\Delta_{y}\right|+\left|\Delta_{z}\right| \right)\\
			{\Delta}^{neu}_{x} &= \Delta_{x} \cdot M / n \\
 			{\Delta}^{neu}_{y} &= \Delta_{y} \cdot M / n \\
 			{\Delta}^{neu}_{z} &= \Delta_{z} \cdot M / n \\ 	
 \end{aligned}
\end{equation}

\ \\
\emph{Fall 2}: Es liegt eine Verletzung der Bedingung in genau zwei Ebenen vor.
L�sungsstrategie: Clampe die gemeinsame Koordinate.
Nutze dazu folgendes Verfahren (Beispiel f�r XY und XZ $\rightarrow$ Clamping von $\Delta_{x}$):
\begin{equation}
 \label{VerwacklungClamp2}
 \begin{aligned}
 			n &= \text{min}\left( M-\left|\Delta_{y}\right|,M-\left|\Delta_{z}\right| \right) \\
			{\Delta}^{neu}_{x} &= \begin{cases} 
									n & falls\ \Delta_{x}>0 \\ 
									-n & sonst 
									\end{cases} 
 \end{aligned}
\end{equation}

\ \\
\emph{Fall 3}: Es liegt eine Verletzung der Bedingung in exakt einer Ebene vor.
L�sungsstrategie: Normalisiere beide zur Ebene geh�rigen Abweichungswerte.
Nutze dazu folgendes Verfahren (Beispiel f�r XY-Ebene):
\begin{equation}
 \label{VerwacklungClamp1}
 \begin{aligned}
 			n &= \left|\Delta_{x}\right|+\left|\Delta_{y}\right| \\
			{\Delta}^{neu}_{x} &= \Delta_{x} \cdot M / n \\
 			{\Delta}^{neu}_{y} &= \Delta_{y} \cdot M / n \\
 \end{aligned}
\end{equation}

\ \\
Nach dem Clamping werden die resultierenden Ausgangskoordinaten $\left(x_{aus},y_{aus},z_{aus}\right)$ aus den Eingangskoordinaten
$\left(x_{ein},y_{ein},z_{ein}\right)$ berechnet:

\begin{equation}
 \label{VerwacklungKoordRes}
 		\left(x_{aus},y_{aus},z_{aus}\right) = \left(x_{ein}+\Delta_{x},y_{ein}+\Delta_{y},z_{ein}+\Delta_{z}\right)
\end{equation}

%================================================================================================
\subsection{Gl�ttung des Dreiecksnetzes}

\begin{figure}[htbp]
  \centering  
  \subfloat{\includegraphics[width=14cm]{Bilder/Screenshot_Verwackeln_kein_Glaetten}}\\    
  \subfloat{\includegraphics[width=14cm]{Bilder/Screenshot_Verwackeln_Glaetten}}  
	\caption[Modifizierung der Vertexkoordinaten]{\emph{Modifizierung der Vertexkoordinaten}:
	 $K=0,35$ , v.o.n.u. (a) Verwackeln ohne Gl�tten, (b) Verwackeln mit Gl�tten}
	 \label{B_VerwackelnGlaettenShots}
\end{figure} 

Falls die H�hlenw�nde eine glattere, aber dennoch unregelm��ige Struktur besitzen sollen,
w�re eine optionale Gl�ttung des verwackelten Dreiecksnetzes von Vorteil.
Die �berlegung besteht darin, die Verwacklungsrichtungen so zu gestalten, dass durch den Verwacklungsprozess diese Gl�ttung erreicht wird.

%Jeder Vertex an Ausgangsposition $(x,y,z)$ geh�rt potentiell zu acht Voxeln.
%Diese haben die Voxelkoordinaten im Bereich $[x-1,x]\times[y-1,y]\times[z-1,z]$
\ \\
Dazu wird die Nachbarschaft des jeweiligen Vertex betrachtet.
Zu einem gegebenen Vertex sei ein anderer Vertex benachbart, wenn er mit diesem �ber eine Voxelkante verbunden ist.
In Abbildung \ref{B_Kantenwichtung} sind solche Nachbarschaften aufgezeigt.
Ein Vertex kann �ber maximal sechs Nachbarvertices verf�gen, in jeder Achsenrichtung einen.
Folgende Beobachtung l�sst sich treffen: f�r glattere Struktur muss ein Vertex i.d.R. in Richtung seiner Nachbarvertices bewegt werden.
In entgegengesetzter Richtung w�rden Zacken entstehen.

Die Idee ist es, diese Kanten zu Nachbarvertices zu z�hlen und gegeneinander aufzuwiegen,
um die finale Verwacklungsrichtung zu bestimmen.
Durch Betrachtung der m�glichen F�lle (Auswahl in Abbildung \ref{B_Kantenwichtung})
wird ersichtlich, dass die Kanten dabei unterschiedlich gewichtet werden m�ssen:
Liegt die Kante zu einem Nachbarvertex inmitten einer Fl�che, so z�hlt dies zweifach,
liegt sie dagegen auf einer Ecke, z�hlt sie nur einfach.
Wenn keine Kante existiert, so z�hlt dies nullfach.

Dadurch ergeben sich sechs Werte als \emph{Kantengewichte}, f�r jede Achse jeweils in positive und negative Richtung:
$W_{x+}$,$W_{x-}$,$W_{y+}$,$W_{y-}$,$W_{z+}$,$W_{z-}$

\begin{figure}[hbtp]
  \centering  
	\includegraphics[width=15cm]{Bilder/Kantenwichtung}
	\caption[Wichtung der Kanten]{\emph{Wichtung der Kanten}: Vertices gr�n, Kanten zu Nachbarvertices rot, Zahlen geben Kantengewichte an}
	\label{B_Kantenwichtung}
\end{figure}

%\ \\
\noindent F�r jede Achse werden die Kantengewichte f�r positive und negative Richtung gegeneinander aufgewogen,
um die erlaubten Verwacklungsrichtungen zu bestimmen.
Die Richtung mit dem h�herem Gewicht ist erlaubt, bei gleichem Gewicht sind beide Richtungen erlaubt.
Damit berechnen sich die \emph{Verwacklungsrichtungsparameter} $R_{x}$, $R_{y}$ und $R_{z}$ wie in Formel \ref{GlaettenRichtungX} angeben (Beispiel f�r $x$, $y$ und $z$ �quivalent):

\begin{equation}
 \label{GlaettenRichtungX}
 \begin{aligned}
		R_{x} = \begin{cases} 
						0 & falls\ W_{x+} = W_{x-} \\ 
						1 & falls\ W_{x+} > W_{x-} \\ 
						-1 &	sonst
						\end{cases} 
 \end{aligned}
\end{equation}

\ \\
Die Verwacklungsrichtungsparameter geben an, in welche Richtung die Koordinaten verwackelt werden d�rfen.
Der Verwacklungsschritt aus Formel \ref{VerwacklungDeltaX} wird nun wie in Formel \ref{VerwacklungGlaettenX} angepasst (Beispiel f�r $x$, $y$ und $z$ �quivalent):

\begin{equation}
 \label{VerwacklungGlaettenX}
 \begin{aligned}
			\Delta_{x} = \begin{cases} 
									-K + 2 \cdot K \cdot Zufallswert[0,1] & falls\ R_{x} = 0 \\ 
									K \cdot Zufallswert[0,1] & falls\ R_{x} > 0 \\ 
									-K \cdot Zufallswert[0,1] & sonst
									\end{cases} 
 \end{aligned}
\end{equation}

\ \\
Zur Optimierung wird eine Look-Up-Tabelle zu Speicherung von $R_x$, $R_y$ und $R_z$ verwendet.
F�r die Berechnung des Tabellenindex wird die Belegung $Voxel_i$ ($0$ oder $1$) aller acht den Vertex umgebenden Voxel verwendet.
Jeder Voxel repr�sentiert dabei ein Bit des Index.
\begin{equation}
 \label{GlaettenIndex}
 \begin{aligned}
		Index = \sum\limits_{i=0}^{7} \left(2^i \cdot Voxel_i\right) 
 \end{aligned}
\end{equation}

\ \\
Der Aufbau der Tabelle erfolgt bei Initialisierung des Programms.
Alle m�glichen Indices von  $0$ bis $255$ werden durchlaufen und dabei aus jedem Index die zugeh�rige Belegung der Voxel bestimmt.
F�r diese Voxelkonfiguration werden die Verwacklungsrichtungsparameter berechnet und in der Tabelle gespeichert.

Bei der Berechnung der Vertexkoordinaten werden die den Vertex umgebenden Voxel betrachtet und daraus der Index bestimmt.
Der entsprechende Eintrag der Look-Up-Tabelle wird ausgelesen und die Werte werden angewendet.

\ \\
Ein Vergleich der Resultate zwischen der Umwandlung mit Verwackeln ohne Gl�tten
und der mit Verwackeln mit Gl�tten ist in Abbildung \ref{B_VerwackelnGlaettenShots} zu sehen.

%================================================================================================
%================================================================================================
\section{Berechnung der Normalen}
\label{KK_VoxelNormalen}

F�r die Beleuchtung der H�hle ist es erforderlich, die Normale in jedem Vertex zu kennen.
Dabei wird jede Vertexnormale aus den Normalen der Dreiecke, zu denen dieser Vertex zugeh�rig ist, bestimmt.

Das Gesamtverfahren sieht wie folgt aus: Anfangs ist jede Vertexnormale mit dem Nullvektor initialisiert.
Nun werden alle Dreiecke abgelaufen und jeweils die unnormalisierte Dreiecksnormale $\vv{N_d}$ nach Formel \ref{F_NormaleDreieck} berechnet.
Die Eckpunkte des entsprechenden Dreiecks seien $\vv{P_{1}},\vv{P_{2}},\vv{P_{3}}$.

\begin{equation}
 \label{F_NormaleDreieck}
 		\vv{N_d} = \left(\vv{P_{2}}-\vv{P_{1}}\right)\times \left(\vv{P_{3}}-\vv{P_{1}}\right)
\end{equation}

\ \\
Es gibt nun unterschiedliche Arten, diese Dreiecksnormalen f�r die Berechnung der Vertexnormalen einzubeziehen.
Drei dieser Methoden werden in Tabelle \ref{TabelleNormalenWichtung} verglichen.
Je nach Verfahren wird die Dreiecksnormale normalisiert.
Daraufhin folgt die Wichtung der Normale f�r jeden der drei Eckpunkte und die Aufaddierung dieser gewichteten Normale zu der Normalen jedes Eckpunktes.
Nach dem Verfahren besitzt jeder Vertex seine \emph{Rohnormale}, die aus der Summe der gewichteten Dreiecksnormalen besteht.

Diese Rohnormalen werden am Ende des gesamten Berechnungsprozesses normalisiert.

\begin{table}
\centering
\begin{tabular}[c]{|l|c|c|c|c|}
  \hline
  Verfahren & Normalisieren & Wichtung f�r $\vv{P_1}$ & W. f�r $\vv{P_2}$ & W. f�r $\vv{P_3}$\\
  \hline\hline  
  Uniform& ja & 1 & 1 & 1\\  
  per Winkel & ja & $\varangle \left( \vv{P_{2}}-\vv{P_{1}} , \vv{P_{3}}-\vv{P_{1}} \right)$ &
  $\varangle$ in $\vv{P_2}$ &
  $\varangle$ in $\vv{P_3}$ \\ 
  per Fl�cheninhalt & nein & 1 & 1 & 1\\  
  \hline
\end{tabular}
\caption{Verschiedene Wichtungsmethoden von Normalen bei der Berechnung von Vertexnormalen aus Dreiecksnormalen}
 \label{TabelleNormalenWichtung}
\end{table}

\ \\
Die \emph{uniforme Wichtung} stellt die grundlegendste Methode dar.
Die Normale eines Vertex berechnet sich aus dem Durchschnitt der Normalen der zugeh�rigen Dreiecke.
Die normalisierte Normale jedes Dreiecks wird dabei mit Faktor $1$ gewichtet.

Die \emph{Wichtung per Winkel} wird beispielsweise von Irrlicht \cite{W_IrrlichtDokuMeshMani} verwendet.
Die Normale eines Vertex berechnet sich aus den normalisierten Normalen der zugeh�rigen Dreiecke,
die mit dem Winkel des jeweiligen Dreicks in diesem Vertex gewichtet werden.
Die Winkel k�nnen per Kosinussatz berechnet werden.

Die \emph{Wichtung per Fl�cheninhalt} wichtet die Normale jedes Dreiecks mit dem Fl�cheninhalt des Dreiecks.
Die L�nge der unnormalisierten Normale entspricht dem Fl�cheninhalt des Parallelogramms,
welches von $\left(\vv{P_{2}}-\vv{P_{1}}\right)$ und $\left(\vv{P_{3}}-\vv{P_{1}}\right)$ aufgespannt wird.
Dieser Fl�cheninhalt entspricht dem doppelten des Dreicksfl�cheninhaltes. 
Die Normale braucht somit nicht normalisiert zu werden, da sie bereits korrekt gewichtet ist.
Der Faktor $2$ kann beibehalten werden, da er einen konstanten Faktor f�r jede Normale darstellt.

\ \\
Durch den Wegfall der Normalisierung ist die Wichtung per Fl�cheninhalt die schnellste aller drei Varianten.
Die Vorteile gegen�ber den beiden anderen Methoden liegen dabei nicht nur in der schnelleren Berechnung:
Degenerierte Dreiecke, wie Linien oder Punkte, ohne bestimmbare Normalen werden hierbei von vornherein nicht in die Vertexnormalen einbezogen,
da der Fl�cheninhalt f�r diese Dreiecke null ist.
F�r die beiden anderen Varianten sei hier die Festlegung getroffen, dass der normalisierte Nullvektor wieder gleich dem Nullvektor sei.\footnote{Diese
Variante ist auch in Irrlicht implementiert.}
Da die unnormalisierte Dreiecksnormale jedes degenerierten Dreiecks dem Nullvektor entspricht,
werden diese so auch nicht in die Vertexnormalen einbezogen.

Die Methode Wichtung per Winkel erscheint intuitiv, da Ecken mit gr��erem Winkel einen gr��eren Einfluss auf ein Dreieck haben.
Bei der Wichtung per Fl�cheninhalt ist der Beitrag zur Normalen umso gr��er, je gr��er das Dreieck ist, was ebenfalls intuitiv erscheint.

Im Vergleich der Ergebnisse (siehe Abbildung \ref{B_Shot_Normalen})
liefert die Wichtung per Fl�cheninhalt die gleichm��igsten Normalen und die Wichtung per Winkel die weichsten �berg�nge an Kanten.
Insgesamt �hneln sich die Ergebnisse dennoch.
Dies ist nicht verwunderlich: Die Voxelmeshes haben eine vergleichsweise regelm��ige Struktur.
F�r einen Mesh mit ausschlie�lich gleichgro�en gleichseitigen Dreiecken w�rden alle drei Methoden das gleiche Resultat liefern.
Andere Methoden sind denkbar. \cite{W_VertexNormal} beschreibt unter anderem die "`Wichtung per inversen Fl�cheninhalt"'
(Idee: kleine Dreiecke "`wissen"' mehr �ber die lokale Oberfl�che, als gro�e Dreiecke mit weit auseinanderliegenden Eckpunkten)
und eine "`Methode der kleinsten Quadrate"'
(nimm eine Ebene, um alle Vertices in der Nachbarschaft zu approximieren, und nutze deren Normale).

\begin{figure}[htbp]
  \centering  % trim: links unten rechts oben
  %trim = 0mm 60mm 0mm 45mm
  \subfloat{\includegraphics[width=14cm,trim = 0mm 110mm 30mm 10mm,clip]{Bilder/Screenshot_Normalen_uniform}} \\
  \subfloat{\includegraphics[width=14cm,trim = 0mm 110mm 30mm 10mm,clip]{Bilder/Screenshot_Normalen_per_Winkel}} \\
  \subfloat{\includegraphics[width=14cm,trim = 0mm 110mm 30mm 10mm,clip]{Bilder/Screenshot_Normalen_per_Flaecheninhalt}} 
	\caption[Vergleich von Normalenberechnungsmethoden]{\emph{Vergleich von Normalenberechnungsmethoden}: v.o.n.u.
	 (a) uniforme Wichtung, (b) Wichtung per Winkel, (c) Wichtung per Fl�cheninhalt}
	\label{B_Shot_Normalen}
\end{figure}

\subsection{Normalen an Meshgrenzen}

Die einzelnen Dreiecksnetze der H�hle teilen sich u.U. Vertices an ihren Grenzen.
Diese Grenzvertices sind dann in mehreren Dreiecksnetzen vorhanden, besitzen aber identische Koordinaten.
Es ist erforderlich, dass diese Vertices ebenfalls identische Normalen besitzen, da sie ein und den selben Vertex repr�sentieren.
Die entsprechenden Normalen m�ssen also zu einer einzigen Normale zusammengefasst werden.

\ \\
Die Methode, um die Vertices zusammenzufassen, funktioniert wie folgt:
Beim Umwandeln der Voxel in ein Dreiecksnetz werden bereits alle Grenzvertices markiert.
Dazu wird getestet, ob sich der Vertex am Rand des umgewandelten Gebietes befindet und ob das benachbarte Netz
an dieser Position ebenfalls einen Vertex generieren w�rde (dies kann durch Betrachtung der acht den Vertex umgebenden Voxel bestimmt werden).
Wenn ja, wird die X-Texturkoordinate $\vv{T}.x$ auf $1,0$ gesetzt, ansonsten auf $-1,0$.\footnote{Da
die H�hlenmeshes keine 2D-Texturen verwenden, k�nnen deren Texturkoordinaten f�r andere Zwecke genutzt werden.}

Nach dem Umwandeln und der Berechnung der \emph{Rohnormalen} f�r jeden Vertex (s.o.) werden
alle Rohnormalen der Vertices mit $\vv{T}.x = 1,0$ in eine Tabelle geschrieben.
Die Koordinaten $(x,z)$ des jeweiligen Vertex definieren die Adresse in der Tabelle.
Nach dem Eintragen werden alle Tabelleneintr�ge nach ihrer Y-Koordinate sortiert.

Nun wird jeder $(x,z)$-Eintrag der Tabelle durchlaufen und die Normalen von aufeinander folgenden Vertices gleicher Y-Koordinate zusammengefasst.
Diese repr�sentieren den gleichen Vertex, da sie identische $(x,y,z)$-Koordinaten haben.
Beim Zusammenfassen werden alle Normalen zusammenaddiert und diese Summe dann den zugeh�rigen Vertices als neue Normale zugewiesen.
Die Normalisierung aller Normalen findet zum Schluss des Prozesses statt.

\ \\
Dadurch ergibt sicht folgendes Gesamtverfahren:
\begin{compactenum}
 \item Berechne alle Rohnormalem
 \item Fasse die Rohnormalen aller Grenzvertices zusammen
 \item Normalisiere alle Normalen
\end{compactenum}

\ \\
Die Laufzeitkomplexit�t der Normalenberechnung betr�gt somit $\mathcal O(n \log y)$,
mit $n$ als Anzahl der Vertices und $y$ als Anzahl der maximalen Anzahl der Tabelleneintr�ge.
Das Sortierverfahren kostet hierbei die meiste Zeit, ohne Sortieren w�re die Laufzeitkomplexit�t linear in $n$.
Als Verbesserungsm�glichkeit lie�e sich statt einer $(x,z)$-Tabelle auch eine Hashtabelle verwendet.
N�here Ausf�hrungen zur Wahl von Hashfunktionen sind in \cite[S.398 ff.]{B_AlgDaten} und \cite[S.264 ff.]{B_AlgoEinf} gelistet.
Die Herausforderung liegt hier darin, die Anzahl der Eintr�ge der Tabelle nicht zu hoch zu w�hlen (Speicherplatz),
aber auch keine zu hohe Anzahl an Vertices pro Tabelleneintrag zu generieren (Laufzeit).

%\ \\
%Der Schritt des Normalenzusammenfassens findet nach der Aufsummierung der gewichteten Dreiecksnormalen,
%aber vor Normalisierung aller Vertexnormalen statt.
%Dadurch ist gew�hrleistet, dass die gew�hlte Wichtungsmethode auch f�r die Grenzvertices angewendet wird.
%Dadurch, dass der Normalenzusammenfassungsprozess die Normalen der prinzipiell identischen Vertices zusammenaddiert,
%werden deren bereits aufsummierten gewichteten Normalen nocheinmal mesh�bergreifend aufsummiert.
%!!!verst�ndlicher erkl�ren, als 3 Schritt Prozess?



%================================================================================================
%================================================================================================
\section{Texturierung und Beleuchtung der H�hle}

Die H�hle mittels einer 2D-Textur zu texturieren, gestaltet sich problematisch.
Zum einen sind passende Texturkoordinaten schwierig zu berechnen,
zum anderen w�rden viele Wiederholungen der verwendeten Textur auftreten
oder aber sehr gro�e Texturen n�tig sein, also sehr speicheraufw�ndige Megatexturen.

Viel sinnvoller ist die Verwendung prozeduraler 3D-Texturen.
Eine 3D-Textur besitzt drei Texturkoordinaten statt wie im 2D-Fall nur zwei.
Als Texturkoordinaten k�nnen im Fall der H�hle direkt die Positionswerte $(x,y,z)$ der einzelnen Punkte verwendet werden.
Eine Texturfunktion weist einem solchen Texturkoordinaten-Tripel $(x,y,z)$ einen konkreten Farbwert zu.
Dieser Farbwert wird entsprechend der Beleuchtung modifiziert und dann dargestellt.
Je nach Wahl der Funktion werden verschiedene Strukturen wie z.B. Marmor oder Holz erzeugt,
Beispiele hierf�r werden in \cite[S.213 ff.]{B_GraDV2} beschrieben.

\ \\
Kristalline Strukturen lassen sich in Verbindung mit der zerkl�fteten unregelm��igen Geometrie der H�hle durch passende Beleuchtung vergleichsweise einfach darstellen.
Ausgegangen wird von der Grundformel \ref{F_3DTextureLambert} f�r diffuse Beleuchtung nach dem Lambertschen Gesetz \cite[S.723-724]{B_Foley},
mit $\vv{L}$ als Vektor in Richtung des einfallenden Lichtstrahls und $\vv{N}$ als Normale der Oberfl�che. 
$\vv{L}$ und $\vv{N}$ seien normiert.
$C_{res}$ ist die resultierende Farbe, $C_{basis}$ die gew�nschte Basisfarbe mit $0\cdot C_{basis}$ als schwarzer Farbe (unbeleuchtet).

\begin{equation}
 \label{F_3DTextureLambert}
 		C_{res} = \begin{cases} 
						\left(\vv{L} \cdot \vv{N}\right) \cdot C_{basis} & falls\ \left(\vv{L} \cdot \vv{N}\right) \geq 0 \\ 
						0\cdot C_{basis} & sonst
						\end{cases}
\end{equation}

\ \\
Der effektiv zul�ssige Winkel zwischen $\vv{L}$ und $\vv{N}$ f�r Formel \ref{F_3DTextureLambert} betr�gt $0$ bis $90$ Grad.
Durch die zerkl�ftete Struktur der H�hle ist dieser Winkel unter Verwendung von Vertexnormalen auch bei sichtbaren Oberfl�chenteilen oftmals gr��er.
Mit Beleuchtung nach Formel \ref{F_3DTextureLambert} w�rde hier sehr fr�h zu schwarz �berblendet werden.
Eine feinere Abstufung der Beleuchtung w�rde entfallen.
Statt zwischen $\vv{L}$ und $\vv{N}$ nur einen Winkel $0$ bis $90$ Grad zuzulassen, seien nun beliebige Winkel (also $0$ bis $180$ Grad) zul�ssig.

Als Lichtquelle wird eine Punktlichtquelle verwendet, die die Position der Kamera innehat.
Somit zeigt $\vv{L}$ vom entsprechendem Punkt der Oberfl�che in Richtung Kamera.
Die Normale wird invertiert, wodurch ein Verschlucken des Lichtes bei frontaler Betrachtung der jeweiligen Struktur erreicht wird.
An den �u�eren R�ndern tritt jedoch das Licht hervor, welches schimmernde kristalline Strukturen, wie z.B. Eis, darstellt.
Diese Modifikationen f�hren zu Formel \ref{F_3DTextureKristall}.
Ein Vergleich der Methoden ist in Abbildung \ref{B_BeleuchtungVergleich} zu sehen.

\begin{equation}
 \label{F_3DTextureKristall}
 		C_{res} = \left(1 -\left(\vv{L} \cdot \vv{N}\right)\right) \cdot \frac{C_{basis}}{2}
\end{equation}

\ \\
F�r das Darstellen der H�hle im Dungeongenerator-Programm wird als Texturfunktion eine einheitliche Basisfarbe f�r alle $(x,y,z)$ verwendet.
Die Beleuchtung wird mittels der Methode aus Formel \ref{F_3DTextureKristall} durchgef�hrt.
Diese Kombination erlaubt ein performantes Rendering der H�hlenstrukturen, da nur wenige Berechnungen notwendig sind,
und ist einfach mittels Shadern implementierbar.

\begin{figure}[htbp]
  \centering  % trim: links unten rechts oben
  \hfill
  \subfloat{\includegraphics[width=49mm,trim = 200mm 60mm 0mm 45mm,clip]{Bilder/Screenshot_Beleuchtung_Lambert}}
  \hfill
  \subfloat{\includegraphics[width=49mm,trim = 200mm 60mm 0mm 45mm,clip]{Bilder/Screenshot_Beleuchtung_Kristall_braun}}
  \hfill
  \subfloat{\includegraphics[width=49mm,trim = 200mm 60mm 0mm 45mm,clip]{Bilder/Screenshot_Beleuchtung_Kristall_blau}} 
  \hfill
	\caption[Vergleich von Beleuchtungsmethoden]{\emph{Vergleich von Beleuchtungsmethoden}: v.l.n.r.
	 (a) lambertsche Beleuchtung, (b) Kristall braun (Bernstein), (c) Kristall blau (Eis)}
	 \label{B_BeleuchtungVergleich}
\end{figure}

	% Erstellung der H�hle
\chapter{Konzept der G�nge und R�ume}
\label{K_GaengeRaeumeKonzept}

In den beiden letzten Kapiteln wurde beschrieben, wie die H�hle f�r den Dungeon erstellt wird.
Dieses Kapitel soll beschreiben, wie die G�nge und R�ume als weitere Bestandteile des Dungeons konzipiert sind.

%Bedeutung von R�umen und G�nge in Dungeons
%!!!Bilder von R�umen und G�ngen

\section{G�nge als polygonale Schl�uche}

G�nge dienen dazu, R�ume untereinander und mit der H�hle zu verbinden.
Ein Gang soll durch einen polygonalen Schlauch (aufgebaut aus Dreiecken) mit definierbarem Querschnitt wiedergegeben werden.
Dieser Schlauch soll in seinem Verlauf einer fest vorgegebenen Kurve $\vv{P(t)}$, mit $t \in [0,1]$, folgen.

Zur Konstruktion des Gangs seien $n$ zweidimensionale Punkte $\vv{Q_0},...,\vv{Q_{n-1}}$ gegeben, die den Querschnitt des Gangs repr�sentieren.
%Diese Punkte seien nach ihrer korrekten Reihenfolge im Querschnittsverlauf angeordnet. %zirkul�r aufsteigend im Uhrzeigersinn
Diese Punkte seien um den Punkt $(0,0)$ im Uhrzeigersinn herum angeordnet.
%- die Koordinaten jedes Punktes liegen in $[-1,1]\times[-1,1]$
Die Gangkurve $\vv{P(t)}$ habe die erste Ableitung $\vv{P'(t)}$.
F�r eine gleichm��ige Struktur des Gangs und eine gute Kontrolle �ber die Anzahl erzeugter Dreiecke
ist es sinnvoll, den Abstand zwischen zwei aufeinanderfolgenden Gangsegmenten gleich gro� zu gestalten.
Dieser Abstand sei $w$.  
% gleichm��ige Unterteilung f�r Kollisionsabfragen sinnvoll, f�r Detailobjekte, evtl. f�r Beleuchtung(?)

\ \\
Algorithmus \ref{A_Gang} beschreibt die Konstruktion des Gangs.
F�r ein aktuelles $t$ wird zun�chst das lokale Koordinatensystem bestimmt, �hnlich dem Frenetschen Koordinatensystem.
Der Unterschied ist die Einbeziehung des H�henvektors $(0,1,0)$, der die Richtung "`Oben"' f�r den Gang definiert.
Das Frenetsche Koordinatensystem verwendet an dieser Stelle die zweite Ableitung der Kurve.
%
Nun wird ein Gangsegment aus 3D-Punkten erzeugt und mittels Dreiecken mit dem vorherigen Gangsegment verbunden. %, soweit dieses vorhanden ist
Dabei werden immer zwei benachbarte Punkte des aktuellen Gangsegments mit den �quivalenten Punkten des letzten Gangsegments durch
eine viereckige Fl�che, bestehend aus zwei Dreiecken, verbunden (Zeilen 16-19 des Algorithmus \ref{A_Gang}).
Danach erfolgt die Berechnung des n�chsten $t$.
Dieses soll so berechnet werden, dass der euklidische Abstand zwischen $\vv{P(t_{neu})}$ und $\vv{P(t_{alt})}$ genau $w$ betr�gt.
So wird der Gang st�ckweise aufgebaut.
Der Algorithmus terminiert bei $t=1$.
Die Bedingung des Abstands $w$ muss f�r das letzte Segment nicht unbedingt eingehalten werden,
da dieser hier typischerweise kleiner ist als f�r die anderen Segmente.

Von einer konstanten Anzahl von Gangsegmentpunkten und einer konstanten Laufzeit f�r das Finden des n�chsten $t$ ausgehend,
besitzt Algorithmus \ref{A_Gang} eine Laufzeitkomplexit�t von $\mathcal O(n)$, mit $n$ als Anzahl der erstellten Segmente.
Diese Anzahl ist sowohl vom Kurvenverlauf als auch von $w$ abh�ngig.

\ \\
Potentiell problematisch ist das �berschreiten des Limits f�r die Anzahl von Vertices je Meshbuffer (siehe Abschnitt \ref{KK_Dreiecksnetze}).
Um dieses Problem zu l�sen, wird ein weiterer Test eingef�gt.
Nach der Vertex- und Dreiecksgenerierung wird gepr�ft, ob f�r das n�chste Gangsegment noch genug Reserve vorhanden ist.
Wenn nicht, wird ein neuer Meshbuffer angelegt und das aktuelle Gangsegment f�r den neuen Buffer ein zweites Mal erstellt.
Danach wird wie gehabt mit dem Algorithmus fortgefahren.

\ \\
Zur Texturierung des Gangs wird f�r jeden Punkt des Querschnitts zus�tzlich eine X-Texturkoordinate angegeben.
Die Y-Texturkoordinate wird entlang des Gangs entsprechend des Abstands zwischen $\vv{P(t_{neu})}$ und $\vv{P(t_{alt})}$ weitergez�hlt.
Dabei ist die Erh�hung der Y-Texturkoordinate pro Abstand $1$ angebbar.

%- Texturkoordinaten f�r Gangquerschnitt bei jedem Punkt angebbar
%  - so gute Kontrolle �ber Texturierung
%  --> Textur-Koordinate X

%	- Dimension des Querschnitts in Voxeln: dim*dim mit dim <= 2*minStrichradius
%- erste Idee: man verwendet abwechselnd 0 und 1 f�r jeden Vertex als Texturkoordinaten
%- sowohl in X-Richtung: um den Querschnitt, als auch in Y-Richtung: entlang des Querschnitts
%- Problem: ungleichm��ige Abst�nde der Vertices sorgen f�r ungleichm��ige Texturierung
%- Problem bei ungerader Anzahl der Vertices des Querschnitts (zwei 1en oder 0en nebeneinander)

\begin{algorithm}
\caption{Erstellung eines polygonalen Schlauchs entlang Kurve $P(t)$, $t \in [0,1]$}
\label{A_Gang}
\algorithmicrequire{Kurve $P(t)$, Querschnitt des Schlauchs, Gangsegmentabstand $w$}\\
\algorithmicensure{Dreiecksnetz des polygonalen Schlauchs}
\begin{algorithmic}[1]

	\State $t \gets 0$
	\State $\vv{Vorn} \gets \left\|\vv{P'(t)}\right\|$	\Comment{lokales Koordsys. berechnen}
	\State $\vv{Links} \gets \left\|\vv{Vorn} \times (0,1,0)\right\|$ \Comment{$\left\|\vv{v}\right\|$ ... $\vv{v}$ normalisiert}
	\State $\vv{Oben} \gets \vv{Links} \times \vv{Vorn}$
	
	\State
	\ForAll{$\vv{Q_i} \in Querschnitt$}	\Comment{Vertices erzeugen}
			\State $\textbf{create}\ Vertex\ V_{i}$
			\State $\vv{V_i.Pos} \gets \vv{P(t)} - \vv{Q_i}.x \cdot \vv{Links} + \vv{Q_i}.y \cdot \vv{Oben}$ \Comment{Position berechnen}
			\State $\vv{V_i.N} \gets \left\|\vv{P(t)}-\vv{V_{i}.Pos}\right\|$ \Comment{Normale berechnen}
	\EndFor
	
	\State
	\While{$t < 1$}
		\State $t \gets t_{neu}\ mit\ D_{e}\left(\vv{P(t)},\vv{P(t_{neu})}\right)=w$ \Comment{n�chstes t bestimmen}
		
		\State $\textbf{calculate}\ lokales\ Koordinatensystem$ \Comment{s.o.}
		\State $\textbf{create}\ Vertices$				\Comment{s.o.}
		
		\For{$i = 1 \to n-1$}
				\State $\textbf{create}\ ViereckAusZweiDreiecken(V_{i-1},V_{i},V^{letztes}_{i-1},V^{letztes}_{i})$	
		\EndFor		
		\State $\textbf{create}\ ViereckAusZweiDreiecken(V_{n-1},V_{0},V^{letztes}_{n-1},V^{letztes}_{0})$	
	
	\EndWhile
\end{algorithmic}
\end{algorithm}


%================================================================================================

\subsection{Verlauf des Gangs per kubischem Hermite-Spline}

F�r die Form des Gangs w�re ein geschwungener Verlauf w�nschenswert.
Als Vorgabe sollen die Endpunkte und die Richtung des Gangs in diesen Punkten angegeben werden.
Hierf�r eignet sich ein dreidimensionaler kubischer Hermite-Spline als Kurve $\vv{P(t)}$.
Dieser stellt ein Polynom dritter Ordnung dar und erm�glicht somit eine gewundene und abgerundete Struktur.
Er ist durch die Vorgaben der Positionen und ersten Ableitungen der Endpunkte konstruierbar.
Die ersten Ableitungen beschreiben die Richtung der Kurve im jeweiligen Punkt.
Abbildung \ref{B_HermiteSpline} zeigt einen solchen Hermite-Spline,
Abbildung \ref{B_HermiteSplineGang} zeigt einen per Hermite-Spline konstruierten Gang.

%!!!Warum nicht h�herdimensional? weitere Kontrollpunkte n�tig -> wie finden?
\begin{figure}[hbtp]
  \centering  
	\includegraphics[width=13cm]{Bilder/HermiteSpline}
	\caption[Kubischer Hermite-Spline]{\emph{Kubischer Hermite-Spline}: mit Anfangs- und Endpunkt sowie den ersten Ableitungen in diesen Punkten (gestrichelt)}
	\label{B_HermiteSpline}
\end{figure}

\ \\
Die folgende Herleitung der Formeln f�r $\vv{P(t)}$ und $\vv{P'(t)}$ des kubischen Hermite-Splines ist im Wesentlichen an \cite[S.111 ff.]{salomon2006curves} angelehnt.
Die Formel f�r ein Polynom dritter Ordnung lautet:
\begin{equation}
		\vv{P(t)} = \vv{a}\cdot t^3 + \vv{b}\cdot t^2 + \vv{c}\cdot t + \vv{d}
\end{equation}

\ \\
Daraus ergibt sich als erste Ableitung:
\begin{equation}
		\vv{P'(t)} = 3\cdot\vv{a}\cdot t^2 + 2\cdot\vv{b}\cdot t + \vv{c}
\end{equation}

\ \\
Gegeben seien die zwei Punkte $\vv{P_1} = \vv{P(0)}$, $\vv{P_2} = \vv{P(1)}$ und die ersten Ableitungen in diesen Punkten $\vv{P'_1} = \vv{P'(0)}$, $\vv{P'_2} = \vv{P'(1)}$.
Es gilt somit:

\begin{equation}
 \begin{aligned}
		\vv{P_1} &= \vv{a}\cdot 0^3 + \vv{b}\cdot 0^2 + \vv{c}\cdot 0 + \vv{d}  \\
		\vv{P_2} &= \vv{a}\cdot 1^3 + \vv{b}\cdot 1^2 + \vv{c}\cdot 1 + \vv{d}  \\
		\vv{P'_1} &= 3\cdot\vv{a}\cdot 0^2 + 2\cdot\vv{b}\cdot 0 + \vv{c} \\
		\vv{P'_2} &= 3\cdot\vv{a}\cdot 1^2 + 2\cdot\vv{b}\cdot 1 + \vv{c} 
	\end{aligned}
\end{equation}

\ \\
Als L�sungen dieses Gleichungssystems ergeben sich:
\begin{equation}
 \begin{aligned}
		\vv{a} &= 2\cdot \vv{P_1} - 2\cdot \vv{P_2} + \vv{P'_1} + \vv{P'_2} \\
		\vv{b} &= -3\cdot \vv{P_1} + 3\cdot \vv{P_2} - 2\cdot \vv{P'_1} - \vv{P'_2} \\
		\vv{c} &= \vv{P'_1} \\
		\vv{d} &= \vv{P_1} 
	\end{aligned}
\end{equation}

\ \\
Man erh�lt eingesetzt und in Matrizenschreibweise umgeformt:

\begin{equation}
 \label{F_HermitePunkt}
\vv{P(t)} = \left(t^3,t^2,t,1\right)
  \begin{pmatrix}
    2 & -2 & 1 & 1 \\
    -3 & 3 & -2 & -1 \\
    0 & 0 & 1 & 0 \\
    1 & 0 & 0 & 0 \\
  \end{pmatrix}
    \begin{pmatrix}
    \vv{P_1}\\
    \vv{P_2} \\
    \vv{P'_1} \\
    \vv{P'_2} \\
  \end{pmatrix}
\end{equation}

\begin{equation}
\vv{P'(t)} = \left(t^3,t^2,t,1\right)
  \begin{pmatrix}
    0 & 0 & 0 & 0 \\
    6 & -6 & 3 & 3 \\
    -6 & 6 & -4 & -2 \\
    0 & 0 & 1 & 0 \\
  \end{pmatrix}
    \begin{pmatrix}
    \vv{P_1}\\
    \vv{P_2} \\
    \vv{P'_1} \\
    \vv{P'_2} \\
 \end{pmatrix}
\end{equation}

\ \\
Damit lassen sich f�r jedes $t$ die zugeh�rige Position $\vv{P(t)}$ und die erste Ableitung $\vv{P'(t)}$ bestimmen.
Bei der Konstruktion der Parameter werden im Programm die Positionen der Endpunkte $\vv{P_1}$, $\vv{P_2}$
und die Richtungen des Gangs vom Endpunkt zum Ganginneren $\vv{P'_1}$, $-\vv{P'_2}$ angegeben.
Der Betrag der ersten Ableitung wirkt als Gewichtung f�r die Richtung des Gangs.
Je h�her diese ist, desto h�her ist auch der Einfluss der Richtungsvorgabe.

\begin{figure}[htbp]
  \centering  
  \hfill
  \subfloat{\includegraphics[width=7.4cm]{Bilder/Screenshot_Gang_Innen}}
  \hfill
  \subfloat{\includegraphics[width=7.4cm]{Bilder/Screenshot_Gang_Aussen}}  
  \hfill
	\caption[G�nge per kubischem Hermite-Spline]{\emph{G�nge per kubischem Hermite-Spline}: v.l.n.r. (a) Gang von innen, (b) Gang von au�en. Verwendet wird eine Testtextur.}
	\label{B_HermiteSplineGang}
\end{figure}

\subsubsection{�quidistante Unterteilung}

Algorithmus \ref{A_Gang} verwendet �quidistante Abst�nde $w$ zwischen aufeinanderfolgenden Punkten $\vv{P(t)}$.
Es stellt sich die Frage, wie diese Unterteilung zu bestimmen ist.
Eine �quidistante Unterteilung von $t$ wird i.A. keine solche Abst�nde liefern.
Eine Approximation mittels der Bogenl�nge $s$ als euklidischem Abstand w�re denkbar.
Die Kurve m�sste also in gleichgro�e St�cke bez�glich ihrer Bogenl�nge aufgeteilt werden.
Hierf�r ist eine Reparametrisierung von $\vv{P(t)}$ nach $\vv{P(s)}$ n�tig.
Das Ablaufen der Kurve w�rde mittels gleichm��iger Erh�hung des Parameters $s$ erfolgen.
F�r die Reparametrisierung muss die Funktion $t(s)$ bekannt sein.

\ \\
Die Bogenl�nge einer Kurve $\vv{P(u)}$, mit $\vv{P'(u)}$ als erster Ableitung, zwischen den Punkten $\vv{P(a)}$ und $\vv{P(b)}$
berechnet sich wie folgt:
\begin{equation}
		s(a,b) = \int\limits_{a}^b \sqrt{\vv{P'(u)}^2}\ du
\end{equation}

\ \\
Verwendet man f�r $\vv{P(u)}$ ein Polynom dritter Ordnung, wie es beim kubischen Hermite-Spline der Fall ist, so erh�lt man f�r $\vv{P'(u)}$ eines zweiter Ordnung.
Dieses wird durch die Quadrierung zu einem Polynom vierter Ordnung unter der Wurzel, mit $k_i$ als den entsprechenden Koeffizienten:
\begin{equation}
		s(a,b) = \int\limits_{a}^b \sqrt{k_1u^4 + k_2u^3 + k_3u^2 + k_4u + k_5}\ du
\end{equation}

\ \\
F�r $s(t)$ als Bogenl�nge der Kurve zwischen $\vv{P(0)}$ und $\vv{P(t)}$ resultiert daraus:
\begin{equation}
		s(t) = \int\limits_{0}^t \sqrt{k_1u^4 + k_2u^3 + k_3u^2 + k_4u + k_5}\ du
\end{equation}

\ \\
Nun ergibt sich folgendes Problem:
L�st man die Formel f�r $s(t)$ und die Umkehrfunktion $t(s)$ (beispielsweise mit einem Computeralgebrasystem),
dann erh�lt man Formeln, die mehrere Seiten lang und �u�erst komplex sind.
Dieser Ansatz ist somit i.A. f�r kubische Polynome nicht realisierbar.

%- wichtig: �quidistante Unterteilung des Splines -> damit Gang �quidistant unterteilt
%(->WARUM?): gleichm��ige Unterteilung bez�glich Parameter t in AUsgangsformel liefert i.A.
%keine sinnvolle Unterteilung des Splines (!Beispiel zeigen).
%�quidistante Unterteilung des Splines erm�glicht mehr Kontrolle �ber Qualit�t und Polygonanzahl
%Auch m�glich: an Stellen starker �nderung (Betrag 2. Ableitung sehr hoch -> starke Kurven vom Spline bzw. Gang) feinere Unterteilung
%- Herleitung der Formel:
%Integrate[Sqrt[a*u\^4 + b*u\^3 + c*u\^2 + d*u + e], u]
%Bogenl�nge konkret in Abh�ngigkeit von t:
%s(t) == Integrate[Sqrt[a*u\^4 + b*u\^3 + c*u\^2 + d*u + e], {u, 0, t}]
%-->
%t(s)=
%Solve [s == Integrate[Sqrt[a*u\^4 + b*u\^3 + c*u\^2 + d*u + e], {u, 0, t}], s] 
\ \\
Die Alternative, welche im Generatorprogramm Anwendung findet, ist die Verwendung einer numerischen L�sung.
Genutzt wird das Bisektionsverfahren, welches u.a. in \cite[S.49-50]{B_NumAnalysis} beschrieben wird.
Das Verfahren halbiert ein gegebenes Intervall fortlaufend und verwendet dabei in der n�chsten Iteration
jeweils die H�lfte weiter, in der die L�sung liegen muss.
Dies wird so lange fortgef�hrt, bis die gew�nschte Genauigkeit f�r die L�sung erreicht ist.
In der Praxis erweist sich diese Methode als sehr schnell und robust.

Gegeben sind $t_{akt}$, der gew�nschte euklidischer Abstand $w_{ziel}$ zwischen $\vv{P(t_{akt})}$ und $\vv{P(t_{neu})}$ sowie die Genauigkeit $\epsilon$.
Gesucht ist $t_{neu}$.

Zuerst wird getestet, ob $D_e\left(\vv{P(t_{akt})},\vv{P(1)}\right)\leq w_{ziel}$ gilt.
Wenn dem so ist, dann ist $t_{neu}=1$ das R�ckgabeergebnis.
%Dieser Wert kann dann nicht weiter verbessert werden.
Der Rest des Verfahrens verl�uft folgenderma�en: \\

\begin{compactenum}
 \item Starte mit Intervall $[m,n] = [t_{akt},1]$
 
 \item Setze $t_{neu} =  m+\frac{(n-m)}{2}$
 
 \item Berechne aktuellen Abstand $w_{akt} = D_e\left( \vv{P(t_{akt})},\vv{P(t_{neu})}\right)$
 
 \item Teste, ob $w_{akt} \in [w_{ziel}-\epsilon,w_{ziel}+\epsilon]$\\
       ja: Gew�nschte Genauigkeit erreicht. Abbruch mit Ergebnis $t_{neu}$.\\           
       nein: Fahre fort.
         
 \item Wenn $w_{akt} > w_{ziel}$, dann setze $n = t_{neu}$ (Wert zu gro� $\rightarrow$ Weiterverwendung der unteren Intervallh�lfte), \\
 			 sonst setze $m = t_{neu}$ (Wert zu klein $\rightarrow$ Weiterverwendung der oberen Intervallh�lfte).
 
 \item Gehe zu 2.

\end{compactenum}

\ \\
%Zus�tzlich mit Ergebnis $t_{neu}$ abgebrochen, wenn $n-m < \phi$ (mit $\phi$ als Abbruchtoleranz), um zu verhindern dass der richtige Wert von $t$ �bersprungen wurde.
Als Verbesserung lie�e sich eine \emph{adaptive Unterteilung} verwenden.
Grundidee ist eine feinere Unterteilung an den Stellen, an denen sich die Kurvenrichtung stark �ndert,
um diese Abschnitte exakter wiedergeben zu k�nnen.
Diese Stellen sind durch Betrachtung der zweiten Ableitung der Kurve bestimmbar.
Charakteristisch ist ein hoher Betrag der zweiten Ableitung und eine Richtung dieser, die sich stark von der Richtung der ersten Ableitung unterscheidet.

%================================================================================================
    
\subsection{Adapter}
\label{KK_Adapter}
Der Querschnitt des Gangs hat i.A. eine andere Form als die �ffnungen von R�umen oder der H�hle,
wodurch eine Anpassung der Geometrien notwendig ist.
Die Geometrien der �ffnungen in R�umen und der H�hle seien als \emph{Andockstellen} bezeichnet.
Diese Andockstellen bilden exakt die Form der �ffnung ab.
An den Enden des Gangs muss der Querschnitt so angepasst werden, dass er zur Form der Andockstelle �berblendet.
Die dazu verwendete Geometrie sei \emph{Adapter} genannt.

\ \\
Das \emph{Gangsegment} in $\vv{P(0)}$ bzw. $\vv{P(1)}$ sei gegeben durch eine Folge von Vertices und ein lokales Koordinatensystem.
F�r die Vertices sind Positionen, Normalen und Texturkoordinaten bekannt.
\\
%
Die \emph{Andockstelle} sei gegeben durch eine Folge von Vertices, einen Mittelpunkt $\vv{M_a}$ und eine normalisierte Normale $\vv{N_a}$.
Die Vektoren $\vv{Links_a}$ und $\vv{Oben_a}$ des lokalen Koordinatensystems der Andockstelle werden berechnet.
Die Normale der Andockstelle �bernimmt dabei die Rolle von $\vv{Vorn}$.
F�r die Andockstellenvertices sind nur die Positionen, aber keine Texturkoordinaten oder Normalen bekannt.

\begin{figure}[hbtp]
  \centering  
	\includegraphics[width=11.5cm]{Bilder/Adapter}
	\caption[Prinzip des Adapters]{\emph{Prinzip des Adapters}: Verbindung zwischen ersten bzw. letztem Gangsegment und einer Andockstellengeometrie.
	 Triangulierung (blau) erfolgt �ber Zuordnung der Vertices (gr�n) per Winkel�hnlichkeit.}
	 \label{B_AdapterPrinzip}
\end{figure}

\noindent Das Prinzip der Adaptererstellung ist in Abbildung \ref{B_AdapterPrinzip} dargestellt.
Die Triangulation zwischen Gangsegment und Andockstelle erfolgt auf Basis der Winkel�hnlichkeit ihrer Vertices.
Zur sp�teren Zuordnung der Vertices wird hierf�r jedem Vertex ein Winkel $\alpha$ zugewiesen.
Die Winkelmessung findet mit den auf die jeweilige Ebene\footnote{Gangsegmentebene:
$\vv{P(0)}$ bzw. $\vv{P(1)}$ als Mittelpunkt und $\vv{P'(0)}$ bzw. $\vv{P'(1)}$ normalisiert als Normale. \\
Andockstellenebene: $\vv{M_a}$ als Mittelpunkt und $\vv{N_a}$ als Normale.}
projizierten Vektoren zwischen Vertex und Ebenenmittelpunkt statt.
Der entsprechende Vektor $\vv{Links}$ steht f�r $0^\circ$, gemessen wird im Uhrzeigersinn.
Die Winkel f�r die Punkte des Gangsegments werden wie in Formel \ref{F_WinkelGangsegment} berechnet.
F�r die Berechnung k�nnen die 2D-Koordinaten der Querschnittspunkte $\vv{Q_i}$ genutzt werden,
aus denen die 3D-Punkte des Segments gebildet wurden.
\begin{equation}
 \begin{aligned}
		\alpha = \text{atan2} \left(\vv{Q_{i}}.y, -\vv{Q_{i}}.x\right) \\
	\end{aligned}
	\label{F_WinkelGangsegment}
\end{equation}

\ \\
Die Winkel f�r die Punkte der Andockstelle werden wie folgt berechnet, mit $\vv{V_i.Pos}$ als 3D-Koordinaten der Andockstellenvertices:
\begin{equation}
 \begin{aligned}
		\alpha = \text{atan2} \left(  \left(\vv{V_i.Pos}-\vv{M_a}\right)\cdot \vv{Oben_a}, \left(\vv{V_i.Pos}-\vv{M_a}\right)\cdot \vv{Links_a} \right) \\
	\end{aligned}
\end{equation}

\ \\
F�r die Korrektheit der nachfolgenden Algorithmen sei gefordert, dass der Gang sowie die Andockstelle so konstruiert sind,
dass sowohl die Gangsegmentvertices als auch die Andockstellenvertices zirkul�r aufsteigende Winkel besitzen.
Weiterhin ist gefordert, dass der auf die jeweilige Ebene
projizierte Kantenzug, der alle Vertices ihrer Reihenfolge nach miteinander verbindet,
den Mittelpunkt der Ebene umschlie�t.\footnote{Dieser Kantenzug ist geschlossen, d.h. erster und letzter Vertex werden ebenfalls verbunden.}
%dass der resultierende Winkelunterschied zwischen zwei aufeinanderfolgenden Vertices weniger als $180^\circ$ betragen soll.
Zur einfacheren Berechnung werden die Vertices zun�chst so umsortiert, dass der kleinste Index auch den kleinsten Winkel enth�lt.
Dies wird �ber eine Indexverschiebung realisiert, d.h. alle Elemente werden um eine bestimmte Anzahl von Positionen verschoben.

\ \\
Der erste Ansatz verwendet nur die bestehenden Vertices zur Triangulation des Adapters.
Zun�chst werden die Normalen der Andockstellenvertices $V_{i}$ berechnet.
Diese entsprechen dem normalisierten, auf die Andockstellenebene projizierten Vektor zwischen Mittelpunkt und Vertex:
\begin{equation}
 \begin{aligned}
		\vv{V_i.N} =  \frac{\left(\vv{M_a} - \vv{V_i.Pos}\right)	- \left(\vv{N_a} \cdot \left(\vv{M_a} - \vv{V_i.Pos}\right)\right) \cdot \vv{N_a}}
												 {\left| \left(\vv{M_a} - \vv{V_i.Pos}\right)	- \left(\vv{N_a} \cdot \left(\vv{M_a} - \vv{V_i.Pos}\right)\right) \cdot \vv{N_a} \right|}
	\end{aligned}
	\label{F_NormalenADS}
\end{equation}

\ \\
Danach erfolgt die Berechnung der Texturkoordinaten.
F�r die Y-Texturkoordinate w�re ein Weiterz�hlen dieser Koordinate entsprechend des Abstands zum letzten Gangsegment denkbar.
Problematisch ist die X-Texturkoordinate.
Abbildung \ref{B_AdapterStuetzvertices} zeigt m�gliche Probleme.
Ein einfaches �bernehmen der X-Texturkoordinate per Winkel�hnlichkeit (Fig. (a)) kann zu Verdrehungen der Textur f�hren.
Interpoliert man die X-Texturkoordinaten (Fig. (b)), kann die Textur �berdehnen.

Die L�sung besteht in dem Einf�gen zus�tzlicher St�tzvertices mit den entsprechenden Winkeln und X-Texturkoordinaten (Fig. (c)).
Die Vertices zwischen den St�tzvertices interpolieren diese Koordinaten.

\begin{figure}[hbtp]
  \centering  
	\includegraphics[width=15cm]{Bilder/Adapter_Stuetzvertices}
	\caption[Notwendigkeit von St�tzvertices f�r Andockstellen]{\emph{Notwendigkeit von St�tzvertices f�r Andockstellen}:
	betrachtet wird Gangsegment (schwarz) und Andockstelle (blau):	v.l.n.r. \\
	(a) Zuordnung der Texturkoordinaten per Winkel�hnlichkeit f�hrt zu Verdrehung,	(b) per Interpolation zu �berdehnung \\
	(c) durch St�tzvertices k�nnen die Texturkoordinaten korrekt interpoliert werden}
	\label{B_AdapterStuetzvertices}
\end{figure}

%\ \\
\subsubsection{Finales Verfahren}
Dies f�hrt zu folgendem Verfahren: Erzeuge aus der Vertexmenge des Gangsegments und der Andockstelle die finalen Andockstellenvertices $V'_{i}$.
Diese werden statt der urspr�nglichen Vertices der Andockstelle f�r die Triangulation des Adapters verwendet.
Zur Erzeugung der $V'_{i}$ wird �ber die komplette Vertexmenge (Gangsegment und Andockstelle) per aufsteigendem Winkel gelaufen
und es werden daraus die Vertices mit entsprechenden Parametern generiert.

\ \\
F�r die nachfolgenden Formeln werden folgende Notationen verwendet:
\begin{tabbing}
\hspace*{3.5cm}\=\kill
$V_{av}$ \> bezeichnet den Vertex der Andockstelle mit n�chstkleinerem \\ \>  Winkel bez�glich des Winkels des aktuellen Vertex.\\
$V_{an}$ \> den mit n�chstgr��erem Winkel.\\
$V_{gv}$ \> bezeichnet den Vertex des Gangsegments mit n�chstkleinerem \\ \> Winkel bez�glich des Winkels des aktuellen Vertex.\\
$V_{gn}$ \> den mit n�chstgr��erem Winkel.\\
$TY_{Gangende}$ \> bezeichnet die Y-Texturkoordinate der Gangsegmentpunkte.\footnotemark \\
$TY_{IncPro1}$  \> gibt die Erh�hung der Y-Texturkoordinate pro Abstand $1$ an. \\
$\Delta_{\varangle}\left(\gamma_1,\gamma_2\right)$  \> bestimmt den resultierenden Winkelunterschied zwischen \\ \> zwei Winkeln $\gamma_1$ und $\gamma_2$.
Dieser Winkelunterschied \\ \> wird als Winkel zwischen $0^\circ$ und $180^\circ$ angegeben. \\ \\
$\vv{V'_i.Pos}$  \> gibt die Position des Vertex an. \\
$\vv{V'_i.T}$  \> gibt die Texturkoordinaten des Vertex an. \\
$V'_i.Markierung$  \> dient zur Markierung, ob es sich um einen Textur-St�tzvertex \\ \>  handelt ($Markierung = wahr$). \\
$V'_i.Korrespondenz$  \> gibt den winkel�hnlichsten Vertex des Gangsegments an, \\ \> welches der sp�teren Triangulierung dient.
\end{tabbing}
\footnotetext{Diese ist f�r alle Punkte eines Gangsegments identisch.}
Zur Berechnung der Y-Texturkoordinate $\vv{V'_i.T}.y$ wird entweder eine Addition oder Subtraktion verwendet, je nach Ende des Gangs.
Sie berechnet sich �ber den Abstand des Andockstellenvertex zur Position auf dem Gangsegment mit identischem Winkel.

Der Interpolationsfaktor $f$ sei im Folgenden stets so gew�hlt, dass die interpolierte Position den gleichen Winkel besitzt wie $V_{akt}$.        
Faktor $f$ kann durch trigonometrische Betrachtungen bestimmt werden.
Bei der Betrachtung der Vertices gibt es drei m�gliche F�lle:

\ \\
\emph{Fall 1}: Der aktuelle Vertex $V_{akt}$ ist ein Gangsegmentvertex.
Vorgehen: Erzeuge einen neuen St�tzvertex mit gleicher X-Texturkoordinate und gleichem Winkel.

\begin{equation}
 \begin{aligned}
		\vv{V'_i.Pos} &= (1-f) \cdot \vv{V_{av}.Pos} + f \cdot \vv{V_{an}.Pos} \\
		\vv{V'_i.T}.x &=  \vv{V_{akt}.T}.x \\
		\vv{V'_i.T}.y &= TY_{Gangende} \pm \left|\vv{V'_i.Pos} - \vv{V_{akt}.Pos}\right| \cdot TY_{IncPro1} \\
		V'_i.Markierung &= wahr \\
		V'_i.Korrespondenz &= V_{akt}
	\end{aligned}
\end{equation}

\ \\
\emph{Fall 2}: Der aktuelle Vertex $V_{akt}$ ist ein Andockstellenvertex.
Vorgehen: Der Vertex wird weiterverwendet.

\begin{equation}
 \begin{aligned}
		\vv{V'_i.Pos} &= \vv{V_{akt}.Pos} \\
		\vv{temp} &= (1-f) \cdot \vv{V_{gv}.Pos} + f \cdot \vv{V_{gn}.Pos} \\
		\vv{V'_i.T}.y &= TY_{Gangende} \pm \left|\vv{V'_i.Pos} - \vv{temp}\right| \cdot TY_{IncPro1} \\
		V'_i.Markierung &= falsch \\
		V'_i.Korrespondenz &= \begin{cases} 
													%V_{gv} & falls\ \left| V_{akt}.\alpha - V_{gv}.\alpha\right| < \left| V_{akt}.\alpha - V_{gn}.\alpha\right| \\ 
													V_{gv} & falls\ \Delta_{\varangle}\left( V_{akt}.\alpha , V_{gv}.\alpha\right) \leq \Delta_{\varangle}\left( V_{akt}.\alpha,V_{gn}.\alpha\right) \\ 
													V_{gn} & sonst
													\end{cases}
	\end{aligned}
\end{equation}

\ \\
\emph{Fall 3}: Der aktuelle Vertex ist Andockstellenvertex  $V_{akt}^a$ und Gangsegmentvertex $V_{akt}^g$ (Winkelgleichheit).
Vorgehen: Der Andockstellenvertex �bernimmt die X-Texturkoordinate des Gangsegmentvertex und wird als St�tzvertex weiterverwendet.

\begin{equation}
 \begin{aligned}
		\vv{V'_i.Pos} &= \vv{V_{akt}^a.Pos} \\
		\vv{V'_i.T}.x &=  \vv{V_{akt}^g.T}.x \\
		\vv{V'_i.T}.y &= TY_{Gangende} \pm \left|\vv{V'_i.Pos} - \vv{V_{akt}^g.Pos}\right| \cdot TY_{IncPro1} \\
		V'_i.Markierung &= wahr \\
		V'_i.Korrespondenz &= V_{akt}^g
	\end{aligned}
\end{equation}

\ \\
Nun erfolgt die Berechnung der X-Texturkoordinaten der Nicht-St�tzvertices der Andockstelle.
Dazu wird die Texturkoordinaten eines solchen Vertex $V'_k$ aus den beiden n�chstliegenden St�tzvertices $V'_m$ und $V'_n$ interpoliert.
$V'_m$ sei dabei der St�tzvertex mit n�chstkleinerem Winkel und $V'_n$ der St�tzvertex mit n�chstgr��erem Winkel bez�glich des Winkel von $V'_k$.
Die Winkel werden dabei im Uhrzeigersinn zirkul�r betrachtet.
Der Interpolationsfaktor bestimmt sich aus den L�ngenverh�ltnissen der auf die Andockstellenebene projizierten Kantenz�ge, die zum einen 
$V'_m$ mit $V'_k$, als auch $V'_m$ mit $V'_n$ verbinden.

\begin{equation}
 \begin{aligned}
 		\vv{V'_k.T}.x = (1-t) \cdot \vv{V'_m.T}.x + t \cdot \vv{V'_{n}.T}.x\\
 		mit\ t =  \frac{
							  \sum\limits_{i=m+1}^{k} \left|
								\left(\vv{V'_i.Pos}-\vv{V'_{i-1}.Pos}\right)
							- \left(\vv{N_a} \cdot \left(\vv{V'_i.Pos}-\vv{V'_{i-1}.Pos}\right)\right) \cdot \vv{N_a}
							\right| 
							}{
							\sum\limits_{i=m+1}^{n} \left|
								\left(\vv{V'_i.Pos}-\vv{V'_{i-1}.Pos}\right)
							- \left(\vv{N_a} \cdot \left(\vv{V'_i.Pos}-\vv{V'_{i-1}.Pos}\right)\right) \cdot \vv{N_a}
							\right|
							} \\
		%\wedge \ m<k \wedge n>k	\\
		%\wedge \ V_m.Markierung = wahr \wedge \neg\exists q: \left( V_q.Markierung = wahr \wedge q>m \wedge q \leq k \right) \\
		%\wedge \ V_n.Markierung = wahr \wedge \neg\exists q: \left( V_q.Markierung = wahr \wedge q<n \wedge q > k \right) \\		
	\end{aligned}
\end{equation}

\ \\
Die Normalen der $V'_i$ werden �quivalent zu Formel \ref{F_NormalenADS} berechnet.
Nach der Erzeugung dieser finalen Andockstellenvertices
erfolgt die Generierung des Dreiecksnetzes des Adapters entsprechend Algorithmus \ref{A_Gangadapter}.
Die Triangulation nutzt, wie schon im urspr�nglichen Ansatz und in Abbildung \ref{B_AdapterPrinzip} aufgezeigt,
die Winkel�hnlichkeit der Vertices zur Zuordnung.
Im Fall, dass der aktuelle Triangulationsschritt zwei Dreiecke generiert, wird die Diagonalachse so gew�hlt,
dass sie die beiden Endpunkte mit dem kleineren Winkelunterschied verbindet.

\ \\
Das Gesamtverfahren hat eine Laufzeitkomplexit�t von $\mathcal O(n+m)$,
mit $n$ als Anzahl der urspr�nglichen Andockstellenvertices und $m$ als Anzahl der Gangsegmentvertices.
Einige Beispiele f�r fertige Adapter sind in Abbildung \ref{B_AdapterShots} zu sehen.

\begin{algorithm}
\caption{Triangulierung eines Gangadapters}
\label{A_Gangadapter}
\algorithmicrequire{Andockstellenvertices $V'_0$ ... $V'_{n-1}$, korrespondierende Gangsegmentvertices}\\
\algorithmicensure{Dreiecksnetz des Gangadapters}
\begin{algorithmic}[1]
	
	\State $K \gets V'_{0}.Korrespondenz$

	\For{$i = 1 \to n-1$}
		\If{$V'_{i}.Korrespondenz = K$} \Comment {Korr. gleichbleibend: ein Dreieck einf�gen}
			 \State $\textbf{create}\ Dreieck(V'_i,V'_{i-1},K)$			 
		\Else \Comment {Korr.-Wechsel: zwei Dreiecke einf�gen}
				\State $\beta_1 \gets \Delta_{\varangle}\left( V'_i.\alpha , K.\alpha \right) $ \Comment {Diagonalachse bestimmen}
				\State $\beta_2 \gets \Delta_{\varangle}\left( V'_{i-1}.\alpha , V'_{i}.Korrespondenz.\alpha \right) $
				\If{$\beta_1 < \beta_2$}
					\State $\textbf{create}\ Dreieck(V'_i,K,V'_{i-1})$	
					\State $\textbf{create}\ Dreieck(V'_i,K,V'_{i}.Korrespondenz)$
				\Else
					\State $\textbf{create}\ Dreieck(V'_{i-1},V'_{i}.Korrespondenz,V'_i)$
					\State $\textbf{create}\ Dreieck(V'_{i-1},V'_{i}.Korrespondenz,K)$	
				\EndIf
				\State $K \gets V'_{i}.Korrespondenz$				
		\EndIf
	\EndFor
	
	\If{$V'_{0}.Korrespondenz = K$}	 \Comment {Abarbeitung des verbleibenden Vertex}
			 \State $\textbf{create}\ Dreieck(V'_0,V'_{n-1},K)$			 
		\Else
				\State $\beta_1 \gets \Delta_{\varangle}\left( V'_0.\alpha , K.\alpha \right) $ \Comment {Diagonalachse bestimmen}
				\State $\beta_2 \gets \Delta_{\varangle}\left( V'_{n-1}.\alpha , V'_{0}.Korrespondenz.\alpha \right) $
				\If{$\beta_1 < \beta_2$}
					\State $\textbf{create}\ Dreieck(V'_0,K,V'_{n-1})$	
					\State $\textbf{create}\ Dreieck(V'_0,K,V'_{0}.Korrespondenz)$
				\Else
					\State $\textbf{create}\ Dreieck(V'_{n-1},V'_{0}.Korrespondenz,V'_0)$
					\State $\textbf{create}\ Dreieck(V'_{n-1},V'_{0}.Korrespondenz,K)$	
				\EndIf
		\EndIf

\end{algorithmic}
\end{algorithm}

\begin{figure}[htbp]
  \centering  
  \hfill
  \subfloat{\includegraphics[width=7.4cm]{Bilder/Screenshot_Adapter_SternVoll}}
  \hfill
  \subfloat{\includegraphics[width=7.4cm]{Bilder/Screenshot_Adapter_SternGitter}}  
  \hfill \\
  
  \hfill
  \subfloat{\includegraphics[width=7.4cm]{Bilder/Screenshot_Adapter_DreieckVoll}}
  \hfill
  \subfloat{\includegraphics[width=7.4cm]{Bilder/Screenshot_Adapter_DreieckGitter}}  
  \hfill \\
  
  \hfill
  \subfloat{\includegraphics[width=7.4cm]{Bilder/Screenshot_Adapter_DreieckInnen}}
  \hfill
  \subfloat{\includegraphics[width=7.4cm]{Bilder/Screenshot_Adapter_ViereckGitter}}  
  \hfill
    
	\caption[Adapter aus verschiedenen Andockstellen]{\emph{Adapter aus verschiedenen Andockstellen}: v.l.n.r. \& v.o.n.u.
	(a)+(b) sternf�rmige Andockstelle, (c)+(d) Dreieck als Andockstelle,
	(e) Blick in einen Gang mit einem Adapter aus einer dreieckigen Andockstelle,
	(f) Viereck als Andockstelle}
	\label{B_AdapterShots}
\end{figure}

%WICHTIG: Winkel der Punkte werden erst in Adapterroutine berechnet -> damit Ausrichtung korrekt
%
%- Verbinden Andockstelle mit Spline
%
%- Erkl�rung des Verfahrens:
%  - Erstellung der Ebene
%  - Sortierung aufsteigend nach Winkel
%  - Zuordnung der Punkte nach Winkel
%[!HERLEITUNG: Formel f�r Winkel: Splineende]
%[!HERLEITUNG: Formel f�r Winkel: Andockstelle - siehe Hefter Do 22.09.2011
%wichtig: Winkel vom immer Links-Vektor aus gesehen
%	- links-Vektor liegt grunds�tzlich in der XZ-Ebene und befindet sich in linker Richtung bzgl. der Normalen
%	(dies trifft auf auch auf Splineende zu)
%
%-> Projektion Vektor auf Ebene per Hessescher Normalform, dann Berechnung des Winkels]
% 
%Eingehen auf erste Idee:
%- ohne Einf�gen zus�tzlicher Punkte
%- TexturkoordinateX wird interpoliert �ber Winkelverh�ltnis
%- Problem: Sonderf�lle (z.B. Viereck mit 1,0,1,0 TextureX
%-> und um 45� gedrehte Andockstelle -> ergibt 0.5, 0.5, 0.5, 0.5 als TextureX -> keine sinnvolle Texturierung)
%
%Behebung des Problems:
%- Einf�gen zus�tzlicher St�tzpunkte an der ADS ohne Silhouette dieser zu ver�ndern
%- St�tzpunkte werden an den entsprechenden Winkeln des Splineendes eingef�gt
%[!HERLEITUNG: Formel f�r St�tzpunkt-Koordinaten]
%- die St�tzpunkte �bernehmen die TextureX-Koordianten von den entsprechenden Punkten des Splineendes
%
%Berechnung der fehlenden TextureX-Koordinaten der anderen Punkte:
%- Interpolation entsprechend der L�nge auf dem Kurvenzug zwischen 2 St�tzpunkten
%-> TextureX wird gleichm��ig �ber den Kurvenzug verteilt
%Verbesserung: Kurvenzug wird auf Adapterebene projeziert -> dort wird die L�nge gemessen
%
%Berechnung der TextureY-Koordinaten durch Abstand Mittelpkt. ADS <-> Splineende
%Verbesserung: Abstand Adapter-Vertex <-> (eventuell tempor�rer) Spline-Vertex wird zur Berechnung verwendet
%[!SKIZZE: Warum Verbesserung notwendig: wenn eine Seite der ADS deutlich n�her am Spline als andere Seite]
%
%Erstellung der Dreiecke:
%- Zuordung der entsprechenden Vertices von ADS und Splineende per �hnlichem Winkel
%[!SKIZZE, n�here Beschreibung des Verfahrens]

\subsubsection{Material�berblendungen}

Gang und Adapter bilden separate Geometrien.
Ein zweites Paar an Texturkoordinaten f�r den Adapter w�re zum Textur-Blending zwischen beiden am Adapter angeschlossenen Geometrien sinnvoll.
Die \emph{Andockstellenseitigen-Vertices} erhalten als zweite Y-Texturkoordinate $1,0$ , die \emph{Gangsegmentseitigen-Vertices} $0,0$.
Dadurch ist ein Textur-Blending per zunehmendem Alphawert in Y-Richtung m�glich.

Allerdings unterst�tzt das f�r den Geometrieexport verwendete Wavefront OBJ-Format, wie auch viele andere 3D-Mesh-Formate, nur ein Paar Texturkoordinaten pro Vertex.
Deswegen m�ssen die zweiten Texturkoordinaten beim Reimport der Daten neu erzeugt werden.
Hierf�r ist es erforderlich, die Vertices in Andockstellenseitige"=Vertices und Gangsegmentseitige"=Vertices unterteilen zu k�nnen.

Die Andockstellenseitigen-Vertices unterscheiden sich in ihrer bestehenden Y"=Texturkoordinate vom ersten und letzten Vertex des Gangs.
Die Gangsegmentseitigen-Vertices des Adapters tun dies dagegen nicht, da sie exakte Kopien der ersten bzw. letzten Gangvertices darstellen.
Somit ist eine exakte Zuordnung m�glich.

%================================================================================================
\subsection{Objekte in G�ngen}

In Dungeons gibt es oft eine Vielzahl von Detailobjekten.
Dies k�nnen verzierende Objekte wie z.B. Fackeln oder Statuen sein,
aber auch interaktive Objekte wie Monster und Fallen.
Gesucht ist eine Methode, Objekte entlang von G�ngen zu platzieren.

F�r das Platzieren von Objekten wird die Gangkurve $\vv{P(t)}$ �quivalent zu Algorithmus \ref{A_Gang} abgelaufen,
nur ohne die Generierung von Vertices und Dreiecken.
Stattdessen werden die Objekte postiert.
Als Abstand zwischen zwei Gangsegmenten $w$ wird exakt der gleiche Abstand wie bei der Ganggenerierung verwendet,
um genau den gleichen Kurvenverlauf zu erhalten.
In bestimmten, bei Bedarf zuf�lligen Intervallen werden nun die Objekte des gew�nschten Typs platziert.
Die Intervalle werden dabei in Vielfachen von Gangsegmenten gez�hlt.

\ \\
Gegeben sei ein Ausgangsobjekt mit einer 2D-Position bez�glich des Gangquerschnitts von $\vv{Q}$,
einer Ausgangsrotation von $\left(\beta_x,\beta_y,\beta_z \right)$ als Eulerwinkel und einer Skalierung von $\vv{S}$.
Gesucht sind die resultierende Skalierung, die Position und die Rotation in Eulerwinkeln eines jeden auf dem Ausgangsobjekt basierenden Objekts $O$ im Gang.

Die Skalierung bleibt unver�ndert und entspricht somit $\vv{S}$.
Bei der Platzierung entlang des Gangs berechnen sich die Position $\vv{O.Pos}$
und die Rotation in Eulerwinkeln $\vv{O.Rot}$ wie in Formel \ref{F_DetailobjektPosRot} angegeben.
$\varangle_M$ sei eine gegebene Umwandlungsfunktion f�r Rotationsmatrizen in Eulerwinkel $\left(\alpha_x,\alpha_y,\alpha_z\right)$, ihre Umkehrfunktion sei $\varangle_M^{-1}$.
$R_{Kurve}$ sei die Rotationsmatrix, die einen Vektor $(0,0,1)$ nach $\vv{P'(t)}$ rotiere.

\begin{equation}
 \begin{aligned}
  \label{F_DetailobjektPosRot}
 		\vv{O.Pos} &= \vv{P(t)} - \vv{Q}.x \cdot \vv{Links} + \vv{Q}.y \cdot \vv{Oben} \\
 		\vv{O.Rot} &= \varangle_M \left( R_{Kurve} \cdot R_{Objekt} \right) \\
 							 &= \varangle_M \left( R_{Kurve} \cdot \varangle_M^{-1} \left( \beta_x,\beta_y,\beta_z \right) \right) 								 
	\end{aligned}
\end{equation}

\ \\
Bei der Einbindung von Detailobjekten %werden als Meshes in die Szene eingebunden.
ist die Verwendung von Token f�r die Repr�sentation komplexerer Objekte, wie z.B. detailreiche animierte Meshes oder Partikelsysteme, sinnvoll.
Die Token stellen sehr einfache Meshes mit wenigen Polygonen dar und dienen als Platzhalter.
Nach dem Einladen in die jeweilige Anwendung kann das Token entsprechend ersetzt werden.
Die Identifizierung des Tokens erfolgt �ber den Objektnamen und die Objekt-ID, die das Objekt im Szenegraphen besitzt (siehe Anhang \ref{KA_Export}).

Die Vorteile der Token liegen sowohl in der schnelleren Darstellung im Generatorprogramm
als auch in der Erzeugung eines geringeren Overheads beim Speichern der Szenendatei,
da ein solches Token weniger Parameter erfordert als z.B. das entsprechende Partikelsystem.
Im Dungeongenerator-Programm werden daher Detailobjekte grunds�tzlich durch nichtanimierte Meshes repr�sentiert,
da sich diese optimal als Token eignen.
Beispiele f�r Objekt-Token entlang von G�ngen werden in Abbildung \ref{B_Token} dargestellt.


%Beispiele f�r Token:
%- statischer Mesh: Fackel
%  Erweiterungen: Partikelsystem f�r Darstellung von Feuer
%                 Lichtquelle
%
%- statischer Mesh: einfacher Mesh (z.B. Tetraeder) "`Monster"'
%  Ersetzung: dynamischer Mesh eines Monsters (M�glichkeit f�r zuf�lliges Ausw�hlen eines Monsters) 
%- .3ds sehr gebr�uchliches 3D-Format, von Irrlicht importierbar
%- statischer 3D-Mesh: Standard-Fall f�r Detailobjekte

\begin{figure}[htbp]
  \centering  
  \hfill
  \subfloat{\includegraphics[width=7.4cm]{Bilder/Screenshot_Token_Gang}}
  \hfill
  \subfloat{\includegraphics[width=7.4cm]{Bilder/Screenshot_Token_Einzeln}}  
  \hfill
	\caption[Objekte in G�ngen]{\emph{Objekte in G�ngen}: v.l.n.r. (a) Objekt-Token entlang eines Gangs, (b) einzelnes Token in Nahansicht}
	\label{B_Token}
\end{figure}


%================================================================================================
%================================================================================================

\section{R�ume als Subszenen}

R�ume sollen die Integration eigener Umgebung erlauben und als fertig ausgestattete \emph{Sub\-szenen} in den Dungeon eingef�gt werden k�nnen.
Jede Subszene dient dabei als Bauvorlage f�r die im Dungeon platzierten R�ume.

Diese Subszenen sind vorgefertigte Irrlicht-Szenen im .irr-Dateiformat.
Sie sind fertig gestaltet, inklusive verzierenden und interaktiven Objekten.
Es k�nnen so beliebige Szenen eingebunden werden.
Prinzipiell kann auch eine Gangkreuzung oder ein extern modellierter H�hlenabschnitt als Raum definiert und verwendet werden.

Jeder Raum hat vier definierte \emph{Andockstellen} f�r den Anschluss von G�ngen, in Richtung Nord, Ost, S�d und West.
Weiterhin f�r gibt es f�r jede dieser Himmelsrichtungen zwei Geometrien:
Die \emph{Verschlussgeometrie} wird verwendet, wenn kein Gang angedockt wird. Ein Beispiel ist eine geschlossene Wand.
Die \emph{Andockgeometrie} wird verwendet, wenn ein Gang angedockt wird. Ein Beispiel hierf�r ist ein Torbogen mit einer Falle am Boden.
 
%  	- interaktive Objekte per ID's definierbar wie in \cite{T_BachelorHoenig}
%  	- keine Kollision mit ID's von anderen R�umen, wenn im Szenegraph immer vom Knoten des Raums aus gesucht wird (Subknoten)
%  - hierdurch Einbindung eigener Dungeonteile m�glich
%- in R�umen werden keine weiteren Objekte platziert

\ \\
Zur Generierung der R�ume wird zuerst die Raumvorlage, also die Subszene, eingeladen und die vier Andockstellen werden ermittelt.
Bei der konkreten Generierung des Dungeons werden Kopien der eingeladenen Subszenen erzeugt und mit gew�nschter Transformation in den Szenegraphen eingef�gt.
Die Transformation wird ebenfalls auf die Andockstellen angewendet.
Je nachdem in welchen Himmelsrichtungen G�nge angeschlossen werden, wird dort entweder die Verschlussgeometrie oder die Andockgeometrie entfernt.
Diese Kopien stellen die finalen R�ume dar.
Ein Beispiel f�r einen Raum mit an die entsprechenden Andockstellen angeschlossenen Adaptern ist in Abbildung \ref{B_RaumAdapter} zu sehen.

\subsection{Ermittlung der Andockstellen}

Eine Andockstelle besteht, wie in Kapitel \ref{KK_Adapter} erl�utert,
aus einer Folge von Vertices $V_i$, einem Mittelpunkt $\vv{M_a}$ und der normalisierten Normale $\vv{N_a}$.  
%- zus�tzlich wird die Ausdehnung $H_a$ ermittelt, die angibt wie weit der wei
Andockstellen werden aus Meshes mit einer genau definierten Randkurve generiert.
Diese Meshes sind in der Subszene durch die Benennung "`Nord\_Andockstelle"', "`Ost\_Andockstelle"', "`Sued\_Andockstelle"' bzw. "`West\_Andockstelle"' definiert.

Die Randkurve besteht aus Vertices, die durch Kanten verbunden sind, die nur zu einem einzigen Dreieck geh�ren, d.h. Randkanten sind.\footnote{Kanten
geh�ren ansonsten i.d.R. zu genau zwei Dreiecken, die links und rechts der Kante liegen.}
$\vv{M_a}$ sei definiert als der Nullpunkt des Meshes.
Die Normale verlaufe in Richtung der Z-Achse des Meshes, wobei sowohl positive als auch negative Richtung erlaubt sind.
Die finale Normalenrichtung wird so gew�hlt, dass die Normale von der Mitte der Subszene weg zeigt.
  
\ \\
F�r die Ermittlung der Andockstellenvertices sei der Mesh durch seine Vertexliste und Dreiecksliste gegeben, wie in Kapitel \ref{KK_Dreiecksnetze} erl�utert.
Zun�chst wird die Adjazenzliste f�r alle Vertices aufgebaut.
Diese gibt f�r jeden Vertex an, welche Vertices mit diesem �ber eine Kante verbunden sind.
F�r den Aufbau der Adjazenzliste wird die Dreiecksliste durchlaufen und es werden jeweils f�r alle drei Kanten eines Dreiecks die entsprechenden Eintr�ge generiert.
Falls ein Eintrag mehrfach vorkommt, weil eine Kante zu mehreren Dreiecken geh�rt, so wird dieser Eintrag als \emph{mehrfach} markiert.

Nach dem Aufbau der Adjazenzliste wird ein Vertex $V_i$ bestimmt,
der mindestens einen Eintrag besitzt, welcher nicht als mehrfach markiert ist.
Dieser Vertex ist ein Andockstellenvertex, also wird er zur Vertexmenge der Andockstelle hinzugef�gt.
Nun wird ein adjazenter Vertex zu $V_i$ bestimmt, der zu diesem mit einer nicht-mehrfachen Kante verbunden ist und damit fortgefahren.
So wird entlang der Vertices der Randkurve gehangelt, bis der Startvertex wieder erreicht ist.
Um zu vermeiden, dass zwischen zwei Vertices hin- und hergesprungen wird, darf keine Kante mehrmals verwendet werden.
Ergebnis ist die Folge der Andockstellenvertices.

Diese Andockstellenvertices sind nun entweder gegen oder im Uhrzeigersinn angeordnet.
Falls sie gegen den Uhrzeigersinn angeordnet sind, wird die Reihenfolge umgekehrt.
Zum Schluss werden die Andockstellenvertices, der Mittelpunkt und die Normale mittels der Transformation,
mit welcher der Andockstellenmesh in die Subszene eingebunden wurde, transformiert.
Die Andockstellenmeshes werden nach Ermittlung der Andockstellen aus der Subszene gel�scht.

\ \\
Ausgehend von einem maximalen Knotengrad je Vertex von $g$, der Anzahl der Dreiecke des Andockstellenmeshs $m$
und der Anzahl seiner Vertices $n$,
liegt die Laufzeitkomplexit�t des implementierten Algorithmus in $\mathcal O(g^2\cdot m+n)$.
Der Aufbau der Adjazenzliste liegt dabei in $\mathcal O(g^2\cdot m)$,
wobei der Faktor $g^2$ aus der Pr�fung auf mehrfache Eintr�ge resultiert.
Es wird trivial jeder neue Eintrag mit den vorhandenen verglichen.
Hier ist Verbesserungspotential vorhanden.
Die Ermittlung der Andockstellenvertices liegt in $\mathcal O(m+n)$
und die Transformation der Andockstelle in $\mathcal O(n)$.

Der Algorithmus ist im Gesamtkontext des Generierungsprogramms nicht sehr zeitintensiv,
da er erstens relativ selten aufgerufen wird (nur beim erstmaligen Einlesen von Subszenen)
und zweitens die Adaptermeshes i.d.R. vergleichsweise wenige Vertices und Dreiecke besitzen.

%einer kleinen konstanten Anzahl von Nachbarvertices und anliegenden Dreiecken je Vertex,
%hat der Aufbau der , mit $n$ als Anzahl der Vertices des Andockstellenmeshes.
%Die 
%und somit auch der Gesamtalgorithmus.

%- bestimme alle Vertices auf Rand des Meshes:
%  alle Vertices die zu Kante geh�ren, die nur ein Dreieck als Nachbarn hat
%  - gehe �ber alle Dreiecke
%  - bilde f�r jeden Vertex Adjazenzliste
%  - markiere alle Vertices mit doppelten Eintr�gen in der Adjazenzliste f�r ung�ltig
%   (da es hier mindenstens 2 Dreiecke mit dieser Kante gibt -> keine Randkante)
%  - alle g�ltigen Vertices bilden Rand
%- Randkantenzug bilden:
%  - filtere aus Adjazenzliste der g�ltigen Vertices alle ung�ltigen heraus
%  - nimm g�ltigen Vertex als Start
%  - hangle durch adjazente Vertices bis wieder am Startvertex
%- wiederhole Verfahren, bis alle Andockstellen ermittelt    

%\subsection{Transformation der Andockstellen}
%wie werden Andockstellen bei Raumplatzierung transformiert?

\begin{figure}[htbp]
  \centering  
  \hfill
  \subfloat{\includegraphics[width=7.4cm]{Bilder/Screenshot_Raum_AdapterGitter}}
  \hfill
  \subfloat{\includegraphics[width=7.4cm]{Bilder/Screenshot_Raum_AdapterVoll}}  
  \hfill \\
  
  \hfill
  \subfloat{\includegraphics[width=7.4cm]{Bilder/Screenshot_Raum_GangVonRaum}}
  \hfill
  \subfloat{\includegraphics[width=7.4cm]{Bilder/Screenshot_Raum_RaumVonGang}}  
  \hfill \\
  
  \hfill
  \subfloat{\includegraphics[width=7.4cm]{Bilder/Screenshot_Raum}}
  \hfill
  \subfloat{\includegraphics[width=7.4cm]{Bilder/Screenshot_Raum_GangVonRaum2}}  
  \hfill
    
	\caption[Raum mit Adaptern]{\emph{Raum mit Adaptern}: v.l.n.r. \& v.o.n.u. (a) Drahtgitter des Adapters, (b) Adapter komplett,
	(c) Gang aus Raumsicht, (d) Raum aus Gangsicht, (e) Raum von innen, (f) gleiche Andockstelle wie bei (a)-(d) um $180^\circ$ gedreht verwendet}
	\label{B_RaumAdapter}
\end{figure}									% Struktur und Prinzip der G�nge
\chapter{Hinzuf�gen von R�umen und G�ngen zum Dungeon}
\label{K_Architekt}

%Architekt:
%- Platziert R�ume 
%- verbindet mittels G�ngen:
%  - Raum mit H�hle
%  - Raum mit Raum 
Aus der erzeugten H�hle soll die Konstruktion des Dungeons mittels G�ngen und R�umen erfolgen.
Diese sollen um die H�hle herum angeordnet sein und mit dieser verbunden werden.
%Gegeben: H�hle im Voxelraum
%Gesucht: G�nge und R�ume um die H�hle herumplatziert und mit ihr verbunden
Das Verfahren dazu besteht aus vier Schritten.

\ \\
\emph{Erster Schritt}: Der Voxelraum wird in Sektoren unterteilt.
Dazu ist die Erweiterung des bisherigen Voxelraums, in dem die H�hle gezeichnet werden kann, n�tig.
Dieser bisher betrachtete Voxelraum sei als \emph{Normalvoxelraum} bezeichnet.
Der Normalvoxelraum wird in jeder Richtung unendlich weit erweitert.
Alle Voxel, die nicht im Normalvoxelraum liegen, seien grunds�tzlich mit $0$ belegt.
%- der erweiterte Voxelraum sei als \emph{Gesamtvoxelraum} bezeichnet

Die Aufteilung des Voxelraums erfolgt in $S_x \times S_y \times S_z$ Sektoren.
Gefordert ist ein Zentralsektor, also eine ungerade Anzahl an Sektoren entlang jeder Dimension $X$, $Y$ und $Z$.\footnote{Diese
Anforderung beruht auf der Annahme, dass so m�gliche zentrale Freir�ume der H�hle besser erkannt und genutzt werden k�nnen.}
Die Mitte des Zentralsektors soll dabei in der Mitte des Normalvoxelraums liegen.
Es sollen mindenstens zwei Sektoren Rand um den Normalvoxelraum bzgl. der X- und Z-Koordinate bestehen,
damit um die H�hle herum gen�gend Platz f�r die R�ume des Dungeons ist. %\\
%
Mit einer Sektorausdehung $A$ entlang jeder Dimension und einer Ausdehnung des Normalvoxelraums von $V_x \times V_y \times V_z$ 
ergibt sich dadurch als Anzahl der Sektoren entlang einer Dimension (hier $X$, $S_z$ �quivalent, bei $S_y$ ohne Aufaddierung des Randes):
\begin{equation}
 \begin{aligned}
		S_x = \begin{cases} 
							\left\lceil \frac{V_x}{A}\right\rceil + 4 & falls\ \left\lceil \frac{V_x}{A}\right\rceil\ ungerade \\ 
							\left\lceil \frac{V_x}{A}\right\rceil + 5 & sonst
					\end{cases}			
	\end{aligned}
\end{equation}

\ \\
%Die Sektorkoordinaten liegen im Bereich $[0,S_x-1]\times [0,S_y-1] \times [0,S_z-1]$.
Ausgehend vom Zentralsektor $\left(\left\lfloor \frac{S_x}{2}\right\rfloor,\left\lfloor \frac{S_y}{2}\right\rfloor,\left\lfloor \frac{S_z}{2}\right\rfloor \right)$
berechnen sich die Grenzen eines Sektors $(i,j,k)$ dadurch wie folgt (hier Dimension $X/i$, $Y/j$ und $Z/k$ �quivalent):
\begin{equation}
 \begin{aligned}
		Min_x &= \left\lfloor \frac{V_x}{2}\right\rfloor + \left(i-\left\lfloor \frac{S_x}{2}\right\rfloor\right) \cdot A - \left\lfloor \frac{A}{2}\right\rfloor \\
		Max_x &= Min_x + A - 1
	\end{aligned}
\end{equation}

\ \\
Zwei Sektoren sind benachbart, wenn eine gemeinsame Grenzfl�che entlang der X- oder Z-Achse existiert.
Also hat jeder Sektor bis zu vier Nachbarsektoren: in positiver und negativer X- sowie positiver und negativer Z-Richtung.
Nachbarsektoren in Y-Richtung existieren im Kontext der Dungeonkonstruktion nicht, da diese Achse die H�henachse darstellt,
�ber die keine R�ume direkt benachbart sein sollen. \\
Nach der Aufteilung in Sektoren werden diese wie folgt charakterisiert: \\
\begin{compactitem}
	\item	\emph{H�hle}: Der Sektor enth�lt mindestens einen Voxel mit Belegung $1$.
	\item \emph{H�hlennachbar}: Der Sektor ist kein \emph{H�hle}-Sektor, hat aber einen Nachbarsektor der mit \emph{H�hle} charakterisiert ist.
	\item \emph{H�hlennachbar2}: Der Sektor ist weder \emph{H�hle}- noch \emph{H�hlennachbar}-Sektor,
														aber ein Sektor in zwei Sektoren Abstand entlang der X- oder Z-Achse ist ein \emph{H�hle}-Sektor.
	\item \emph{Frei}: Alle anderen Sektoren erhalten diese Charakterisierung.
\end{compactitem}

\ \\
\emph{Zweiter Schritt}: R�ume in H�hlenn�he werden platziert und mit dieser verbunden.
Dazu laufe �ber alle Sektoren.
Wenn ein Sektor \emph{H�hlennachbar} ist, so bestimme per Zufall eine beliebige Subszene.
Die M�glichkeit zur Wahl keiner Subszene (die leere Subzene) existiert dabei ebenfalls.
Falls eine nicht-leere Subszene bestimmt wurde, so platziere eine Instanz der Subszene als Raum im Sektor.
Nun versuche von diesem Raum eine Verbindung per Gang zur H�hle herzustellen,  wobei die Himmelsrichtungen in zuf�lliger Reihenfolge abgearbeitet werden.
Falls die Verbindung nicht m�glich ist, l�sche den Raum wieder und markiere den Sektor mit \emph{Frei}.
Wenn die Verbindung m�glich ist, markiere den aktuellen Sektor mit \emph{Raum} und erstelle den Gang.
Falls dieser Gang durch einen \emph{Frei}-Sektor verl�uft, so markiere diesen Sektor mit \emph{Gang}.
Nun erstelle, falls m�glich, mit Wahrscheinlichkeit $p_{gh}$ weitere G�nge vom Raum zur H�hle, wobei die �brigen Himmelsrichtungen verwendet werden.

Dieser Vorgang wird wiederholt, nur dieses Mal werden \emph{H�hlennachbar2}-Sektoren verwendet.
Durch diese Aufteilung in zwei Durchl�ufe werden Sektoren mit direkter Nachbarschaft zur H�hle priorisiert.
Resultat dieses Schritts sind platzierte R�ume, die direkt mit der H�hle durch G�nge verbunden sind.

%teste ob eine Verbindung zur H�hle m�glich ist (siehe Kapitel \ref{KK_HoehleAndock}).
%Wenn nein, dann markiere Sektor mit \emph{Frei} und fahre mit n�chstem Sektor fort.
%Wenn ja, so 
%Platziere eine Instanz der Subszene als Raum in diesem Sektor und markiere den Sektor mit \emph{Raum}.
%Erstelle einen Gang vom Raum zur H�hle,

\ \\
\emph{Dritter Schritt}: R�ume neben anderen R�umen werden platziert und mit diesen verbunden.
Dazu laufe �ber alle Sektoren und markiere alle diejenigen \emph{Frei}-Sektoren, die Nachbarsektoren eines \emph{Raum}-Sektors sind, mit \emph{Raumnachbar}.
Nun laufe erneut �ber alle Sektoren.
Wenn ein Sektor mit \emph{Raumnachbar} markiert ist, so w�hle eine zuf�llige Subszene aus.
Falls eine nicht-leere Subszene bestimmt wurde, so platziere eine Instanz der Subszene als Raum in diesem Sektor und markiere den Sektor mit \emph{Raum}.
Verbinde diesen Raum per Gang zu einem beliebigen, zuf�llig ausgew�hlten Raum in einem Nachbarsektor.

Die Markierung per \emph{Raumnachbar} dient dazu, zu verhindern, dass gerade erstellte R�ume gleich als potentielle Nachbarr�ume dienen k�nnen.
Prinzipiell l�sst sich dieser Schritt mehrfach durchf�hren, um ein sehr weit verzweigtes Netz aus G�ngen und R�umen zu erhalten.
In der Implementierung des Dungeongenerators wird nur ein Durchlauf verwendet.
Resultat des Schritts sind platzierte R�ume, die nur mit anderen R�umen, aber nicht direkt mit der H�hle verbunden sind.
Dennoch sind alle R�ume zumindest indirekt mit der H�hle verbunden.

\ \\
\emph{Vierter Schritt}: Weitere G�nge zwischen R�umen werden eingef�gt.
Laufe �ber alle Sektoren und pr�fe, ob ein Sektor mit \emph{Raum} markiert ist und der Sektor in positiver X-Richtung ebenfalls ein \emph{Raum}-Sektor ist.
Weiterhin �berpr�fe, ob noch kein Gang zwischen den R�umen der beiden Sektoren existiert.
Bei Erf�llung der Bedingungen f�ge zwischen beiden Sektoren mit der Wahrscheinlichkeit $p_{gr}$ einen Gang ein.
F�hre den gleichen Test in positiver Z-Richtung durch, auch hier wird entsprechend ein Gang eingef�gt.

Durch diesen Schritt werden alle potentiellen Gangverbindungen zwischen den bestehenden R�umen des Dungeons getestet und per Zufall neue G�nge eingef�gt.
Resultat sind zus�tzliche Gangverbindungen zwischen den R�umen des Dungeons.
Dieser besitzt nun alle seine G�nge und R�ume.

%Platzierungsroutine beschreiben incl. Aufbau des Sektorenraums
%Scan und Konstruktion des Dungeons
%Wie werden R�ume aus SubSzenen generiert?
%Wie werden R�ume untereinander verbunden?

%================================================================================================

\section{Platzierung von R�umen und Verbindung durch G�nge}
\label{KK_PlatzVerbRaum}

Nach Auswahl einer Subszene erfolgt die Platzierung des Raums im entsprechenden Sektor.
Die Ausdehnung des Sektors sei $[Min_x,Max_x]\times[Min_y,Max_y]\times[Min_z,Max_z]$ mit einer L�nge von $A$ in jeder Dimension.

Zur Berechnung der Position des Raums wird von der Mitte des Sektors ausgegangen, damit in jeder Richtung m�glichst viel Platz f�r die G�nge bleibt.
Die H�henkoordinate $y$ wird per Zufall modifiziert, welches zu Formel \ref{F_RaumPos} f�hrt.
Die Funktion Zufallswert liefert einen Zufallswert im angegebenen Intervall.
\begin{equation}
 \begin{aligned}
		\vv{Pos}.x &= 0,5 \cdot (Min_x + Max_x) \\
		\vv{Pos}.y &= Zufallswert[Min_y+ 0,25 \cdot A, Max_y - 0,25 \cdot A] \\
		\vv{Pos}.z &= 0,5 \cdot (Min_z + Max_z)		
	\end{aligned}
	\label{F_RaumPos}
\end{equation}

\ \\
Die Bestimmung der Rotation erfolgt ebenfalls per Zufallselement.
Es findet nur eine Rotation um die H�henachse $Y$ statt, damit die R�ume stets waagerecht bleiben.
\begin{equation}
 \begin{aligned}
		\vv{Rot}.y &= Zufallswert(0^\circ,360^\circ]
	\end{aligned}
\end{equation}

\ \\
Beim Verbinden per Gang wird dieser an derjenigen Andockstelle angedockt, die in Richtung des anderen Raums bzw. des entsprechenden H�hlenbereichs zeigt.
Die Rotation des Raums erfordert hierzu eine Neubelegung der Himmelsrichtungen der Andockstellen.
Das Anpassen dieser Himmelsrichtungen $H$ wird folgenderma�en vorgenommen, mit
$Nord \equiv 0 \equiv X\ positiv, Ost \equiv 1 \equiv Z\ negativ, Sued \equiv 2 \equiv X\ negativ, West \equiv 3 \equiv Z\ positiv$:
\begin{equation}
 \begin{aligned}
		H_{neu} = \begin{cases} 
								(H_{alt} + 1)\ \text{mod}\ 4 & falls\  45^\circ < \vv{Rot}.y \leq 135^\circ\\ 
								(H_{alt} + 2)\ \text{mod}\ 4 & falls\  135^\circ < \vv{Rot}.y \leq 225^\circ\\ 
								(H_{alt} + 3)\ \text{mod}\ 4 & falls\  225^\circ < \vv{Rot}.y \leq 315^\circ\\ 
								H_{alt} & sonst
						\end{cases}			
	\end{aligned}
\end{equation}

\subsubsection{G�nge aus Andockstellen}
Die Erstellung eines Gangs erfolgt von einer Andockstelle $1$ zu einer anderen Andockstelle $2$.
Die Normale von Andockstelle $1$ sei $\vv{N_a}$, der Mittelpunkt $\vv{M_a}$ und der maximale Abstand eines Vertex der Andockstelle zur Ebene,
mit $\vv{M_a}$ als Punkt der Ebene und $\vv{N_a}$ als Normale der Ebene, sei $W_a$.
Die Breite des Gangs sei $B$.
Ein Faktor $g$ f�r die St�rke der Gangrichtung ist als Parameter im Generatorprogramm festlegbar.
F�r die Erstellung des Gangs werden die Parameter $\vv{P_1}$ und $\vv{P'_1}$ wie folgt bestimmt:

\begin{equation}
 \begin{aligned}
		\vv{P_1} &= \vv{M_a} + \vv{N_a} \cdot \left( W_a+B \right) \\
		\vv{P'_1} &= g \cdot A \cdot \vv{N_a}
	\end{aligned}
\end{equation}

\ \\
Die Einbeziehung von $W_a$ in die Positionsberechnung verhindert ein m�gliches Ineinanderschieben von Adapter und Gang.
Durch Aufaddieren von $B$ wird gen�gend Platz f�r die �berblendung der Querschnitte des Adapters erzeugt.
%\footnote{Dieser Wert erwies sich experimentell als am optimalsten.}
Die Einbeziehung der Sektorausdehnung $A$ in die Ableitung des Gangsplines bewirkt eine Art Skalierung des Gangs mit dieser.
$\vv{P_2}$ und $-\vv{P'_2}$ werden auf die gleiche Art und Weise von Andockstelle $2$ bestimmt.

Da die erstellten Andockstellen initial in Richtung Gang zeigen,
wird f�r die Erstellung des Adapters an $\vv{P_2}$ des Gangs
die Normale von Andockstelle $2$ mit $-1$ multipliziert und
die Reihenfolge der Andockstellenvertices invertiert,
was einer Spiegelung der Andockstelle entspricht.

%Da Andockstellen initial zum Gang zeigen, wird f�r die Erstellung der Adapter des Gangs
%die Normale von Andockstelle $2$ mit $-1$ multipliziert  und
%die Reihenfolge der Andockstellenvertices invertiert werden, was einer Spiegelung der Andockstelle entspricht.

%- Platzieren von R�umen
%  - Bestimmung von Rotation, Position und Skalierung
%  - Anpassen der Himmelsrichtungen
%- Verbindung Raum mit H�hle (aus Raumsicht, aus H�hlensicht s.u.)

%================================================================================================
%================================================================================================

\section{Test f�r Andocken an H�hle}
%\label{KK_HoehleAndockTest}

Der Platzierungsalgorithmus f�r R�ume und G�nge testet in Schritt $2$, ob Verbindungen zur H�hle m�glich sind.
Als Bedingungen f�r das Andocken eines Gangs an die H�hle muss die lokale Struktur der H�hle den Anschluss eines Gangs erm�glichen.
Auch darf der Gang nicht mit der H�hle, R�umen oder anderen G�ngen kollidieren.
%was muss beim Andocken beachtet und deswegen getestet werden?

Als Hilfsmittel f�r die notwendigen Tests und den eigentlichen Andockprozess werden \emph{Scankarten} verwendet.
Diese erm�glichen den schnellen Zugriff auf die relevanten Informationen der umliegenden H�hlenteile.
Das Prinzip einer Scankarte ist in Abbildung \ref{B_Scankarte}(a) dargestellt.
Alle \emph{H�hlennachbar}- und \emph{H�hlennachbar2}-Sektoren scannen in allen vier Himmelsrichtungen den Voxelraum.
Der Scan erfolgt zwei Sektoren tief, um in jedem Fall Voxel mit Belegung $1$ zu finden.
Der Bereich des eigenen Sektors wird ebenfalls in den Scan einbezogen,
da die Scankarte auf Grund diverser Tests etwas breiter und h�her als der eigentliche Sektor ist.
Somit werden die entsprechenden Teile der Nachbarsektoren mitgescannt.

Der Scan bildet ein Tiefenprofil der Voxelh�hle bezogen auf die jeweilige Richtung.
Jeder Eintrag der Scankarte berechnet sich wie in Formel \ref{F_Scaneintrag} (Beispiel f�r Scan in positiver X-Richtung, andere Richtungen �quivalent) angegeben.
$Offset_y$ und $Offset_z$ sind Offsets, die sich aus den Sektorgrenzen berechnen und die Abbildung der Scankartenkoordinaten auf Voxelkoordinaten erm�glichen.
%- $Min_x$ und $Max_x$ sind die Sektorgrenzen, Sektorausdehung ist $A$
Falls kein entsprechendes $k$ existiert, wird der Eintrag auf $\infty$ (bzw. $-\infty$, wenn Scan in negativer Richtung) gesetzt.

\begin{equation}
 \begin{aligned}
 \label{F_Scaneintrag}
		Scan_{x+}(i,j) = \text{min}(k) : Voxel(k,j+ Offset_y,i + Offset_z)=1 \\
		%\wedge \ k \geq Min_x \wedge k\leq Max_x + 2 \cdot A
		\wedge \ k \in [Min_x,Max_x + 2 \cdot A]
	\end{aligned}
\end{equation}

\ \\
Jeder Eintrag einer Scankarte besitzt ein Flag \emph{valide}, welches angibt,
ob der Eintrag noch g�ltig ist und verwendet werden kann. % oder nicht.
Anfangs sind alle Eintr�ge valide.

F�r den Andockversuch in Himmelsrichtung $H$ wird zun�chst ein zuf�lliger Startindex f�r die Scankarte von $H$ bestimmt.
Dann wird getestet, ob ein Andocken an dieser Position m�glich ist.
Wenn ja, wird angedockt, wenn nein, wird der Index weitergez�hlt.
Dies wird solange fortgef�hrt, bis alle Indizes gepr�ft wurden oder eine g�ltige Andockposition gefunden wurde.

\begin{figure}[hbtp]
  \centering  
	\includegraphics[width=14.5cm]{Bilder/Scankarte}
	\caption[Generierung und Revalidierung von Scankarten]{\emph{Generierung und Revalidierung von Scankarten}: v.l.n.r.\\	
	(a) Generierung: Gescannt wird von H�hlennachbar- und H�hlennachbar2-Sektoren in Richtung H�hlen-Sektoren. Scannender Sektor ist orange markiert.\\
	(b) Revalidierung: Durch den Gang (braun) vom roten Sektor muss
	eine Scankarte (blau) vom orangenen Sektor revalidiert werden, da der Gang Teile der H�hle verdeckt.}
	\label{B_Scankarte}
\end{figure}

\subsection{Filterung der Scankarten}
\label{KK_Scankartenfilterung}

Der Bereich der H�hle, an dem der Gang andockt, sollte keine zu starken Tiefenunterschiede besitzen, also m�glichst eben sein.
Gesucht ist somit ein Bereich auf der Scankarte mit m�glichst kleinen Abweichungen der Eintr�ge.
Weiterhin m�ssen $1$-Voxel, d.h. Teile der H�hle, in gesamten Andockbereich vorhanden sein.
Nicht zul�ssig sind folglich Eintr�ge mit dem Wert $-\infty$ oder $\infty$.

Beim Test der Andockposition auf der Scankarte wird ein Bereich auf der Scankarte der Breite und H�he $B+2$ untersucht,
also der Gangbreite mit einem Voxel Rand in jeder Richtung.
Der Rand wird verwendet, um potentielle Kollisionen zwischen Adapter und H�hle zu vermeiden, %den Adapter des Gangs h�hlenseitig zu umschlie�en,
der m�gliche Gang w�rde den Bereich  $[i+1,i+B]\times[j+1,j+B]$ belegen.
Als maximale Tiefendifferenz f�r das Andocken wird die Gangbreite $B$ verwendet, wodurch die Toleranz f�r den Test mit $B$ skaliert.
Jeder Eintrag der Scankarte im Bereich muss valide sein.
Als G�ltigkeitstest f�r eine Andockposition $(i,j)$ ergibt sich somit:

\begin{equation}
 \begin{aligned}
		G\ddot{u}ltig(i,j) = \begin{cases} 
													ja & falls\  Max_{scan} - Min_{scan} \leq B \\
													   & \wedge\  Min_{scan} > -\infty \wedge  Max_{scan} < \infty \\
													   & \wedge \neg\exists\ Scan(a,b)\equiv invalide\\ 
													nein & sonst
												\end{cases} \\
		mit\  Min_{scan} = \text{min}(Scan(a,b)) \wedge 	Max_{scan} = \text{max}(Scan(a,b)): \\		
	 % i \leq a \leq i+B+2 \wedge j \leq b \leq j+B+2
	 a \in [i,i+B+1] \wedge b \in [j,j+B+1]
	\end{aligned}
\end{equation}

\ \\
Im Allgemeinen sind G�nge in Dungeons in Bodenn�he mit der H�hle verbunden, anderenfalls kann ein Begehen des Gangs unm�glich sein.
Dies f�hrt zu der Idee, den Scantest um den Test der Bodenn�he zu erweitern
Der Boden wird durch eine gro�e Tiefe auf der Scankarte im Vergleich zum Bereich dar�ber gekennzeichnet.
Auf der Scankarte unterhalb des Andockbereiches muss also eine gr��ere Tiefe vorliegen als im Andockbereich.

%- getestet wird Schicht f�r Schicht
Der maximaler Abstand zum Boden sei $R$,
die gew�nschte Minimaltiefe f�r den Test sei $F$.
Die Andockposition $(i,j)$ sei g�ltig bzgl. des Abstands zum Boden wenn Formel \ref{F_Bodenabstand} wahr ergibt.
Getestet wird hierbei der Bereich direkt unter dem Gang.
$Min_{scan}$ und $Max_{scan}$ seien die durch den obigen Test bereits ermittelten Werte.
Das Gesamtverfahren ist in Abbildung \ref{B_AndockTestBasis} illustriert.

\begin{equation}
 \begin{aligned}
 \label{F_Bodenabstand}
	 & Test\ bei\ positiver\ Achsenichtung: \\
	 & \exists\ b \in [j-R,j-1] \left( \forall a \in [i+1,i+B] \left( Scan(a,b)-Max_{scan} \geq F \right)\right) \\
	 & Test\ bei\ negativer\ Achsenichtung: \\
	 & \exists\ b \in [j-R,j-1] \left( \forall a \in [i+1,i+B] \left( Min_{scan}-Scan(a,b) \geq F \right)\right) \\
 \end{aligned}
\end{equation}

\begin{figure}[hbtp]
  \centering  
	\includegraphics[width=10cm]{Bilder/BodenDecke}
	\caption[Bestimmung einer g�ltigen Andockposition]{\emph{Bestimmung einer g�ltigen Andockposition}: gezeigt wird die Seitenansicht eines H�hlenteils:
	Im Bereich von $B+2$ wird getestet ob $Max_{scan}-Min_{scan}$ klein genug ist.
	$R$ und $F$ dienen dem Test, ob die Andockposition in Bodenn�he ist.}
	\label{B_AndockTestBasis}
\end{figure}

\subsubsection{Revalidierung von Scankarten-Eintr�gen}
%- Erweiterung f�r Andocken an unteren H�hlenbereichen
Neu erzeugte G�nge und R�ume k�nnen H�hlenteile f�r andere Sektoren verdecken.
Dies ist in Abbildung \ref{B_Scankarte}(b) dargestellt.
Also m�ssen nach der Erzeugung neuer G�nge und R�ume betroffene Scankarten �berpr�ft und ggf. Teile mit \emph{invalide} markiert werden.

F�r den Test wird die Bounding Box des Gangs plus dem Sektor des erstellten Raums als Konfliktgebiet
$[K_{minX},K_{maxX}]\times[K_{minY},K_{maxY}]\times[K_{minZ},K_{maxZ}]$ genutzt.

\ \\
Die Erkl�rung erfolgt am Beispiel f�r Scankarten in positiver X-Richtung.
Ein Sektor ist potentiell betroffen, wenn folgende Bedingungen zutreffen: \\
\begin{compactenum}
	\item �berschneidung der Scankarte mit $[K_{minY},K_{maxY}]\times[K_{minZ},K_{maxZ}]$
	\item Grenze des Sektors $Min_x \leq K_{maxX}$, also �berschneidung bzgl. der X-Koordinate
\end{compactenum}

\ \\
F�r diese Sektoren werden die Scaneintr�ge �berpr�ft.
Ein Eintrag ist invalide, wenn er im oder hinter dem Konfliktgebiet liegt,
d.h. durch das Konfliktgebiet verdeckt wird:
\begin{equation}
 \begin{aligned}
		Scan_{x+}(i,j) \equiv \begin{cases} 
										invalide & falls\  K_{minX} \leq Scan_{x+}(i,j) \\
										valide & sonst
										\end{cases} \\
	\end{aligned}
\end{equation}

\ \\
Durch die Revalidierung der Scankarten werden potentielle Kollisionen der G�nge mit R�umen und anderen G�ngen vermieden.


%================================================================================================

\subsection{Test auf potentielle Kollision von Gang und H�hle}

Da G�nge und H�hle potentiell kollidieren k�nnen, ist ein Kollisionstest H�hle $\Leftrightarrow$ Gang 
vor dem Erstellen neuer G�nge erforderlich.
Falls der Test positiv verl�uft, soll kein Gang erstellt werden, da eine potentielle Kollision vorliegt.
Ziel des Tests ist eine \emph{false negative}-Rate von $0$, also soll jede Kollision entdeckt werden.
F�r einen schnellen Test sei eine \emph{false positive}-Rate $>0$ erlaubt.
Der Adapter kann durch seine Konstruktion nach Kapitel \ref{KK_AndockenHoehle} nicht mit der H�hle kollidieren
und muss daher nicht in den Test einbezogen werden.

Ein Gang soll in Himmelsrichtung $H$ an die H�hle andocken.
Im ersten Ansatz wird die Bounding Box des neuen Gangs genutzt und getestet, ob Elemente der $H$-Scankarte
in die Bounding Box hineinreichen oder sogar davor liegen.
In diesem Fall k�nnte der Gang in die H�hle ragen.
Problematisch ist, dass die Bounding Box keine besonders exakte Approximation des Gangs darstellt.
Potential f�r neue G�nge wird hier verschenkt, da die \emph{false positive}-Rate zu hoch ist.
Eine bessere Approximation w�re w�nschenswert.

\ \\
Wie in \cite[S.210, 425]{salomon2006curves} beschrieben wird, l�sst sich ein kubischer Hermite-Spline
durch eine kubische B�zierkurve darstellen.
Eine B�zierkurve ist nach \cite[S.178 ff.]{salomon2006curves} durch Formel \ref{F_Bezier} gegeben.
Die $B_{n,i}(t)$ werden Bernsteinpolynome genannt.
Die Punkte $\vv{Q_i}$ werden als Kontrollpunkte der Kurve bezeichnet.

\begin{equation}
 \begin{aligned}
 \label{F_Bezier}
		\vv{Q(t)} = \sum\limits_{i=0}^{n} \left( B_{n,i}(t) \cdot \vv{Q_i}\right) \ \ mit\ B_{n,i}(t) =  {n \choose i} t^i (1-t)^{n-i}\ ,\ t\in[0,1]
	\end{aligned}
\end{equation}

\ \\
Die B�zierkurve f�r vier Punkte ergibt sich dadurch als:

\begin{equation}
 \begin{aligned}
		\vv{Q(t)} = (1-t)^3 \vv{Q_0} + 3t(1-t)^2 \vv{Q_1} + 3t^2(1-t) \vv{Q_2} + t^3 \vv{Q_3}
	\end{aligned}
\end{equation}

\ \\
Setzt man die Parameter $\vv{P_1},\vv{P_2},\vv{P'_1},\vv{P'_2}$ des Hermite-Splines auf:

\begin{equation}
 \begin{aligned}
 \label{F_BezierHermiteParameter}
	\vv{P_1} = \vv{Q_0},\ 
	\vv{P_2} = \vv{Q_3},\ 
	\vv{P'_1} = 3(\vv{Q_1}-\vv{Q_0}),\ 
	\vv{P'_2} = 3(\vv{Q_3}-\vv{Q_2}) 
	\end{aligned}
\end{equation}

\ \\
So erh�lt man eingesetzt in Formel \ref{F_HermitePunkt}:

\begin{equation}
 \begin{aligned}
		\vv{P(t)} = &(2t^3-3t^2+1)\vv{Q_0} + (-2t^3+3t^2)\vv{Q_3} \\
		     	 &+ (t^3-2t^2+t)3(\vv{Q_1}-\vv{Q_0})+(t^3-t^2)3(\vv{Q_3}-\vv{Q_2})\\
		     = &(1-t)^3\vv{Q_0} + 3t(1-t)^2\vv{Q_1} + 3t^2(1-t)\vv{Q_2} + t^3\vv{Q_3} \\
		     = &\vv{Q(t)}
	\end{aligned}
\end{equation}

\ \\
Somit ist die Gleichheit von Hermite-Spline und B�zierkurve gegeben.
Stellt man Formel \ref{F_BezierHermiteParameter} um,
so erh�lt man als Punkte $\vv{Q_0},\vv{Q_1},\vv{Q_2},\vv{Q_3}$ der Bezierkurve:

\begin{equation}
 \begin{aligned}
	\vv{Q_0} = \vv{P_1},\ 
	\vv{Q_1} = \vv{P_1} + \frac{\vv{P'_1}}{3},\ 
	\vv{Q_2} = \vv{P_2} - \frac{\vv{P'_2}}{3},\ 
	\vv{Q_3} = \vv{P_2}
	\end{aligned}	
\end{equation}

\ \\
Als Eigenschaften der Bernsteinpolynome nach \cite[S.134]{B_GraDV1} gelten:
Bernsteinpolynome im Intervall $[0,1]$ sind eine Partition der $1$ und stets positiv.

\begin{equation}
 \begin{aligned}
		\forall t \in [0,1] &: \sum\limits_{i=0}^{n}  B_{n,i}(t) = 1 \\
		\forall t \in [0,1] &: B_{n,i}(t) \geq 0		
	\end{aligned}	
\end{equation}

\ \\
Nach \cite[S.71-72]{B_RourkeCompGeo} wird eine konvexe H�lle wie in Definition \ref{D_konvexeHuelle} definiert.
Die Eigenschaften der $k_i$ entsprechen genau denen der Bernsteinpolynome.
Somit verl�uft eine B�zierkurve bzw. der entsprechende Hermite-Spline innerhalb der konvexen H�lle
der Punkte $\vv{Q_0},\vv{Q_1},\vv{Q_2},\vv{Q_3}$.
Diese Eigenschaft kann f�r den Kollisionstest genutzt werden.
Wenn die H�hle nicht mit einer den Gang umschlie�enden H�lle kollidiert, dann auch nicht mit dem Gang selbst.

\begin{TheoremDefinition}[Konvexe H�lle]
\ \\
Eine Konvexkombination von Punkten $\vv{Q_0}, \vv{Q_1}, \vv{Q_2}, ..., \vv{Q_n}$ ist definiert als:
\begin{equation}
 \begin{aligned}
		\sum\limits_{i=0}^{n} \left(k_i\cdot \vv{Q_i}\right) \ \ mit\  \sum\limits_{i=0}^{n}k_i = 1 \wedge \forall i: k_i \geq 0
	\end{aligned}	
\end{equation}
Die konvexe H�lle dieser Punkte ist die Menge aller ihrer Konvexkombinationen.
  \label{D_konvexeHuelle}
\end{TheoremDefinition}

%================================================================================================

\subsubsection{Schneller Kollisionstest}
Ziel ist das schnelle Testen aller Scankarteneintr�ge auf die Kollision mit dem Gang.
Die konvexe H�lle von $\vv{Q_0},\vv{Q_1},\vv{Q_2},\vv{Q_3}$ umschlie�t nur den Gangspline, nicht den Gangs selbst.
Also ist ein Hinzurechnen eines Abstands $G$ zur konvexen H�lle notwendig.
Dieser berechnet sich aus dem maximalen Gangradius plus Sicherheitsabstand nach Formel \ref{F_KollisionG}.
Die Gangbreite sei $B$, der Sicherheitsabstand sei $\epsilon$.
Parameter $\epsilon$ sollte auf Grund der Koordinatenverwacklung der H�hlenvertices mindestens $0,5$ betragen.
In der Implementierung wurde $\epsilon = 1,0$ gew�hlt.
%!!!BILDSKIZZE Wurzel 2 Faktor notwendig, siehe Hefter 21.04.2012

\begin{equation}
 \begin{aligned}
		G = \sqrt{2} \cdot 0,5 \cdot B + \epsilon =  \frac{1}{\sqrt{2}} \cdot B + \epsilon
	\end{aligned}
	\label{F_KollisionG}
\end{equation}

\ \\
Im Folgenden sei das Beispiel f�r das Andocken in positiver X-Richtung aus Raumsicht gegeben,
d.h. der Gang verl�uft von der H�hle aus in negativer X-Richtung.
Das Prinzip des Tests wird in Abbildung \ref{B_Pyramide} aufgezeigt.
Es sei vorausgesetzt, dass $\vv{Q_3}.x < \vv{Q_0}.x$ und $\vv{Q_2}.x < \vv{Q_0}.x$ gilt,
anderenfalls ist die Andockposition ung�ltig.
$\vv{Q_1}.x < \vv{Q_0}.x$  gilt nach Konstruktion (siehe Abschnitt \ref{KK_AndockenHoehle}) immer und braucht nicht getestet zu werden.
%$Q_1.x < Q_0.x$ gilt nach Konstruktion (siehe Andocken, Abschnitt \ref{KK_AndockenHoehle}) immer und braucht nicht getestet zu werden.
F�r den Test wird eine Pyramide konstruiert, um die konvexe H�lle von $\vv{Q_0},\vv{Q_1},\vv{Q_2},\vv{Q_3}$ zu umschlie�en.
Die Anstiege $m_{z+},m_{z-},m_{y+},m_{y-}$ aller vier Seitenfl�chen der Pyramide werden berechnet.
Formel \ref{F_KollisionAnstieg} gibt das Beispiel f�r $m_{y+}$ an, die anderen Ansteige werden �quivalent ermittelt.
F�r das Berechnen der Anstiege werden die Richtungsvektoren $\vv{Q_2}-\vv{Q_0}$ und $\vv{Q_3}-\vv{Q_0}$ verwendet.
Der Richtungsvektor $\vv{Q_1}-\vv{Q_0}$ zeigt per Konstruktion in Richtung $(-1,0,0)$ %(siehe Abschnitt \ref{KK_AndockenHoehle})
und entspricht somit dem Anstieg $-\infty$.
%!!!BILDSKIZZE (siehe Hefter Fr 13.01.2012)

\begin{equation}
 \begin{aligned}
 		\vv{r_1} &= \vv{Q_2}-\vv{Q_0},\ \vv{r_2} = \vv{Q_3}-\vv{Q_0} \\
 		m_{y+} &= \begin{cases} 
								\text{max}\left(\vv{r_1}.x/\vv{r_1}.y,\vv{r_2}.x/\vv{r_2}.y\right) & falls\ \vv{r_1}.y > 0 \wedge \vv{r_2}.y > 0\\
								\vv{r_1}.x/\vv{r_1}.y & falls\ \vv{r_1}.y > 0 \wedge \vv{r_2}.y \leq 0\\
								\vv{r_2}.x/\vv{r_2}.y & falls\ \vv{r_1}.y \leq 0 \wedge \vv{r_2}.y > 0\\
								 -\infty & sonst
							\end{cases}
	\end{aligned}
	\label{F_KollisionAnstieg}
\end{equation}

\ \\
Jeder Eintrag $Scan_{x+}(i,j)$ der Scankarte wird getestet,
ob er innerhalb der auf die Scankarte projizierten Bounding Box $[K_{minY},K_{maxY}]\times[K_{minZ},K_{maxZ}]$ der konvexen H�lle,
in jeder Dimension um $G$ erweitert, liegt.
Alle Scankarteneintr�ge, die darin liegen, %und nicht im Andockbereich f�r den Gang 
werden nun in einem zweiten Test auf ihre Position bzgl. der Pyramide getestet.
Je nach dem, ob diese Eintr�ge links, rechts, oben oder unten vom Gangandockbereich liegen,
wird die dortige Tiefe der Pyramide unter Verwendung der jeweiligen Anstiege berechnet.

Formel \ref{F_KollisionPy} gibt ein Beispiel f�r einen Punkt $(i,j)$ auf der Scankarte an,
der in positiver Z- und positiver Y-Richtung vom Andockbereich liegt.
Parameter $i+Offset_z$ ergibt die Voxelkoordinate $z$ des Scaneintrags,
$j+Offset_y$ die Voxelkoordinate $y$.
Die Position der Pyramidenspitze wird durch $\vv{Q_0}$ definiert.
$P_h$ ist die resultierende Tiefe der Pyramide.
%!!!BILDSKIZZE: schneidende Pyramidenseiten -> resultierende Tiefe
Durch die Einbeziehung von $G$ resultiert ein Pyramidenstumpf als getestete H�lle.

\begin{equation}
 \begin{aligned}
 	P_{hy} &= m_{y+} \cdot (j+Offset_y -\vv{Q_0}.y -G) \\
 	P_{hz} &= m_{z+} \cdot (i+Offset_z -\vv{Q_0}.z -G) \\
  P_h &= \vv{Q_0}.x + G + \text{min}(P_{hy}, P_{hz}, 0)
	\end{aligned}
	\label{F_KollisionPy}
\end{equation}

\ \\
Diese Tiefe wird mit dem Wert auf der Scankarte verglichen.
Ist der Scankartenwert kleiner, so liegt m�glicherweise ein Teil der H�hle innerhalb des Pyramidenstumpfes.
Dadurch ergibt sich eine potentielle Kollision und der Gang wird nicht erstellt.

%Die Verwendung einer Pyramdie bzw. eines Pyramidenstumpfes erm�glicht ein schnelleres Testen
%aller potentiell betroffenen Eintr�ge der Scankarte,
%als wenn die konvexe H�lle direkt verwendet werden w�rde.

\begin{equation}
 \begin{aligned}
		Kollision(i,j) &= \begin{cases} 
									ja & Scan_{x+}(i,j) < P_h\\ 
									nein & sonst
									\end{cases} \\
	\end{aligned}	
\end{equation}

\begin{figure}[hbtp]
  \centering  
	\includegraphics[width=13cm]{Bilder/Pyramide}
	\caption[Test der konvexen H�lle per Pyramide]{\emph{Test der konvexen H�lle per Pyramide}: gezeigt wird die Seitenansicht eines H�hlenteils:
	Um die konvexe H�lle der Kontrollpunkte $\vv{Q_0},\vv{Q_1},\vv{Q_2},\vv{Q_3}$ des braunen Gangs wird eine Pyramide gelegt. Die roten Voxel liegen im Kollisionsgebiet.\\
	Anmerkung: $\vv{Q_0}$ liegt ein St�ck von der H�hle entfernt, da sich zwischen Gang und H�hle noch der Adapter befindet.
	%Dieser kann durch sein Konstruktionsprinzip nicht mit der H�hle kollidieren.
	}
	\label{B_Pyramide}
\end{figure}

\ \\
Als Verbesserungsm�glichkeit f�r einen exakteren Test k�nnte man, statt der Verwendung einer Pyramide,
um die konvexe H�lle zu umschlie�en, die konvexe H�lle direkt f�r den Kollisionstest verwenden.
Die verwendete konvexe H�lle der Kontrollpunkte l�sst sich weiter verfeinern:
Durch den de-Casteljau-Algorithmus \cite[S.136]{B_GraDV1} k�nnen weitere Kontrollpunkte errechnet werden.
Dieser l�sst sich prinzipiell beliebig oft anwenden.
Mit jeder Anwendungsstufe erh�lt man mehr Kontrollpunkte als in der vorherigen.
Zudem liegen diese Kontrollpunkte immer n�her am eigentlichen Splineverlauf.
Somit approximiert die konvexe H�lle dieser neu errechneten Punkte die konvexe H�lle des Gangsplines genauer.
Diese Verbesserungen sind allerdings mit h�herem Rechenaufwand verbunden.

%Die konvexe H�lle l�sst sich weiter verfeinern, in dem mehr und exaktere Eckpunkte dieser berechnet werden.
%Dazu kann der verwendet werden ,
%welcher sich beliebig oft anwenden l�sst.
%Mit jeder Anwendungsstufe kommen weitere Kontrollpunkte hinzu, die gleichzeitig Eckpunkte der konvexen H�lle darstellen.


%================================================================================================
\section{Andocken an H�hle}
\label{KK_AndockenHoehle}

Beim Andocken des Gangs an die H�hle muss diese angepasst werden, um eine �ffnung f�r den Gang zu bieten.
Des Weiteren ist eine Andockstelle f�r den Gang zu erstellen.
%Aus den Andockstellenparametern erfolgt die Berechnung der Gangparameter $\vv{P_1}$, $\vv{P'_1}$ bzw. $\vv{P_2}$, $-\vv{P'_2}$
%mittels der in Kapitel \ref{KK_PlatzVerbRaum} beschriebenen Methode.
Das Andockverfahren beschreibt Algorithmus \ref{A_AndockenGangHoehle} und ist in Abbildung \ref{B_Andocken} skizziert.
Angegeben ist das Beispiel f�r das Andocken in positiver X-Richtung aus Raumsicht, die anderen Richtungen werden �quivalent behandelt.

\begin{figure}[hbtp]
  \centering  
	\includegraphics[width=15cm]{Bilder/Andocken}
	\caption[Andocken an H�hle]{\emph{Andocken an H�hle}: gezeigt wird die Seitenansicht eines H�hlenteils:
	v.l.n.r (a) H�hle vor dem Andocken mit $Min_{scan}$ (gelb) und $Max_{scan}$ (orange),
	(b) H�hle nach dem Andocken mit Markierung $10$ (rot), ausgefr�stes St�ck (blau) und Andockstellenvertices (gr�n)	
	}
	\label{B_Andocken}
\end{figure}


\begin{algorithm}
\caption{Andocken Gang an H�hle, in positiver X-Richtung}
\label{A_AndockenGangHoehle}
\algorithmicrequire{Voxelraum, Scankarte, $Min_{scan},Max_{scan}$,\\ Andockbereich $[i,i+B+1]\times [j,j+B+1]$ mit Gangbreite $B$} \\
\algorithmicensure{ge�ndertes Voxelgebiet, Andockstelle f�r Gang}
\begin{algorithmic}[1]

	\State $X_{min} \gets Min_{scan}-1$ \Comment {Grenzen des zu �ndernden Voxelgebietes}
	\State $X_{max} \gets Max_{scan}-1$
	
	\State $Y_{min} \gets j+Offset_y + 1$ \Comment {Bereich in Gangbreite}
	\State $Y_{max} \gets Y_{min} + B -1$
	\State $Z_{min} \gets i+Offset_z + 1$
	\State $Z_{max} \gets Z_{min} + B -1$
	
	\State 
	
	\State $T \gets \infty$ \Comment {Suche nach nahestem Gangrandst�ck}
	\ForAll{$(a,b) \in [i,i+B+1]\times [j,j+B+1]$}  
	\If{$(a = i) \vee (b = j) \vee (a = i+B+1) \vee (b=j+B+1)$}
	  \If{$Scan_{x+}(a,b) < T$}
	    	\State $T \gets Scan_{x+}(a,b)$
	  \EndIf
	\EndIf
	\EndFor
	
	\State \Comment {Voxelh�hle f�r Andocken anpassen}
	
	\ForAll{$(x,y,z) \in [X_{min},X_{max}]\times [Y_{min},Y_{max}]\times [Z_{min},Z_{max}]$} 
		 \If{$x < T$}
			   \State $Voxel(x,y,z) \gets 10$	\Comment {Bereich markieren}
		  \ElsIf {$x < Scan_{x+}(z-Offset_z, y-Offset_y)$}
			    \State $Voxel(x,y,z) \gets 1$  \Comment {Bereich ausfr�sen}
		  \EndIf
	\EndFor
	
	\State 
	
	\State $\vv{N_a} \gets (-1,0,0)$	\Comment{Andockstelle erstellen}
	\State $\vv{M_a} \gets (T ,(Y_{min}+Y_{max}+1)/2,(Z_{min}+Z_{max}+1)/2)$ \Comment{$+1$ durch Voxel zu Vertex}
	\For{$z = (Z_{max}+1) \to Z_{min}$}
			\State $\textbf{createAndAdd}\ Vertex(T,Y_{min},z)$	
	\EndFor
	\For{$y = (Y_{min}+1) \to Y_{max}$}
			\State $\textbf{createAndAdd}\ Vertex(T,y,Z_{min})$	
	\EndFor
	\For{$z = Z_{min} \to (Z_{max}+1)$}
			\State $\textbf{createAndAdd}\ Vertex(T,Y_{max}+1,z)$	
	\EndFor
	\For{$y = Y_{max} \to (Y_{min}+1)$}
			\State $\textbf{createAndAdd}\ Vertex(T,y,Z_{max}+1)$	
	\EndFor
	
\end{algorithmic}
\end{algorithm}

\ \\
Zuerst werden die Grenzen des zu �ndernden Bereich des Voxelraums bestimmt.
Diese liegen im Gangbereich und $Min_{scan}$, $Max_{scan}$, um einen Voxel verschoben. %Minscan, Maxscan(berechnet wie in Kapitel \ref{KK_Scankartenfilterung})

Dann wird der Randbereich des Andockgebietes untersucht, welcher den zuk�nftigen Gang h�hlenseitig umschlie�t.
Die minimale X-Koordinate dieses Bereiches wird bestimmt.
Diese Koordinate stellt eine Grenze $T$ dar.\footnote{Dieser
Wert $T$ wird in der eigentlichen Implementierung schon mit $Min_{scan}$ und $Max_{scan}$ zusammen ermittelt.}
Alles das, was im Gangbereich weiter weg ist, wird "`freigeschaufelt"' bzw. ausgefr�st.
Dieses Ausfr�sen findet bis zur H�hlengrenze statt, welche durch die Scankarte bekannt ist.
Es wird hier kein Bereich innerhalb der bestehenden H�hle ver�ndert, um dortige Strukturen nicht zu entfernen.

Der n�here Bereich bzgl. der Grenze wird als Andockstelle im Voxelraum markiert, mit einer Belegung von $10$.\footnote{Prinzipiell
sind andere Werte ungleich $0$, $1$ oder $2$ m�glich.
Der Wert $10$ ist aus diversen berechnungstechnischen Gr�nden vorteilhafter als niedrigere Werte.}
Dadurch sind im Gangbereich keine $0\Leftrightarrow 1$-�berg�nge mehr vorhanden, wodurch dort keine Dreiecke gebildet werden (vgl. Kapitel \ref{KK_Umwandlung}).
Es entsteht an dieser Stelle also eine �ffnung der H�hle.
%!!!BILDSKIZZE: Ausfr�sbereich s.Hefter 24.11.2011

Dann erfolgt die Bildung der Andockstelle.
Deren Mittelpunkt liegt in der Mitte des Gangbereiches und auf H�he der Ausfr�sgrenze $T$.
Die Normale zeigt in die negative X-Richtung, also in Richtung des Sektors, aus dem der Gang kommt.
Die Vertices der Andockstelle werden im Uhrzeigersinn bez�glich der Normalenrichtung angeordnet und umschliessen genau die entstandene �ffnung.
Ihre Koordinaten werden nach derselben deterministischen Methode berechnet wie die Vertexkoordinaten der H�hle (vgl. Kapitel \ref{K_VertexkoordinatenBerechnen}).
Somit werden Gang und H�hle l�ckenlos miteinander verbunden.
%!!!BILDSKIZZE: (Andockstelle: Voxel<->Vertex Beziehung n Voxel -> n+1 Vertices)
%!!!BILDSKIZZE: Andockstelle: Winkel durch Clamping auch bei Verwacklung zirkul�r aufsteigend s.Hefter 24.11.2011
Ein an die H�hle angedockter Gang ist in Abbildung \ref{B_HoehleAngedockt} zu sehen.

\begin{figure}[htbp]
  \centering  
  \hfill
  \subfloat{\includegraphics[width=7.4cm]{Bilder/Screenshot_GangHoehle_Hoehle}}
  \hfill
  \subfloat{\includegraphics[width=7.4cm]{Bilder/Screenshot_GangHoehle_Gang}}  
  \hfill
	\caption[An H�hle angedockter Gang]{\emph{An H�hle angedockter Gang}: v.l.n.r. (a) aus H�hlensicht, (b) aus Gangsicht}
	\label{B_HoehleAngedockt}
\end{figure}

\ \\
Die Laufzeitkomplexit�t des gesamten Andockprozesses inklusive Andocktest wird im Wesentlichen von den Tests bestimmt.
Im Worst Case m�ssen alle $n$ Positionen auf der Scankarte �berpr�ft werden,
wobei der Test jeder Position durch den Kollisionstest in $\mathcal O(n)$ liegt.
Daraus resultiert eine Gesamtkomplexit�t f�r die Laufzeit von $\mathcal O(n^2)$ bzw.
$\mathcal O\left((A+B+R)^2\cdot(A+B)^2\right)$, da die Scankarte eine Ausdehnung von $[A+B+1] \times [A+B+1+R]$ besitzt.

% evtl. erkl�ren warum Ausfr�sen und nicht beliebige Rang -> Verhindern von �berlappenden Dreiecken
% bei jetziger Methode nicht m�glich -> Beweis
%\ \\
%Anmerkung: durch genauere Betrachtung lassen sich die Schleifengrenzen minimal weiter verfeinern.
%Der �bersichtlichkeit sind diese so wie angegeben dargestellt.
%Zugriff auf Voxel au�erhalb des Voxelraums stellen nicht-definierte Sonderf�lle dar


%\subsection{Verbinden von Raum mit H�hle}
%
%[alt:]
%Passende Andockstelle an H�hle finden:
%- Variante 1:
%	- beim Zeichnen der Voxelh�hle Kandidaten f�r Andockpunkte mitspeichern:
%		- unteres Gangende (Strichmitte - up*radius) mit anderer Farbe zeichnen
%	- diese Andockpunkte sp�ter als unterstes Ende des Ganges verwendbar
%- Probleme:
%	-	durch �berlagerung der Strich kann ein Andockpunkt mitten in der H�hle liegen und nicht am Boden dieser
%	
%- Variante 2:
%	- sende Strahl mit Gangdicke von Voxelraum-/Gebietsrand in Richtung Voxelh�hle
%	- bei Treffer:
%		- verschiebe soweit nach unten wie m�glich, um Boden zu treffen
%		- ermittle passende Gangrichtung: abh�ngig von der Lage der getroffenen Voxel (diese als eine Ebene mit entsprechender Normale betrachten)
%		
%- nach Finden der Andockstelle an der H�hle: eventuell Entfernung von Voxeln im Andockgebiet (wenn erforderlich)		
%
%------
%\subsection{Konkrete Platzierung von R�umen und G�ngen}
%
%[alt:]
%	- Unterteilung des Voxelraums in Sektoren (Octree der Umwandlung Voxel->Dreiecksnetz wiederverwendbar)
% 		- Sektor ist belegt oder frei
% 	Variante 1: 	
% 	- f�r jeden Raum eine Platzierungswahrscheinlichketi angebbar
% 	  - es existiert auch eine Wahrscheinlichkeit f�r keinen Raum
% 	  - alle freien Sektoren werden durchlaufen:
% 	    - Platzierung eines Raums oder eines Freiraums per Zufall mit entsprechender Wahrscheinlichkeit
% 	- maximale Anzahl R�ume angebbar (-> so oft wird gew�rfelt)
%  Variante 2:
%  - festgelegte Anzahl f�r jeden Raum 
%
%---
\section{Grenzen des Verfahrens}

Die Kollision von G�ngen bei Raum-zu-Raum-Verbindungen ist i.d.R. unproblematisch, da der Gang nur zwei benachbarte Sektoren verbindet.
Allerdings sollten nicht zu gro�e Ableitungsfaktoren $g$ verwendet werden, damit der Gang nicht �ber die Sektorgrenzen hinausragt bzw. in R�ume hinein.
Das Gleiche gilt f�r die Ableitungsfaktoren $g$ bei der Raum-zu-H�hle-Verbindung.
Da nur der Bereich der Scankarte getestet wird, also der Sektorausdehnung plus Rand,
k�nnen bei zu ausladenden G�ngen eventuell Kollisionen nicht erkannt werden.

Des Weiteren sollte das Verh�ltnis von Sektorausdehnung und Raumgr��e passend gew�hlt werden.
Die Sektorausdehnung muss stets gr��er als die maximale Raumausdehnung sein,
da sonst Kollisionen zwischen R�umen auftreten k�nnen.

	
%- sinnvoll: Sonderf�lle herausfiltern:
%	- bei denen die begrenzende Seitenfl�chen nicht alle in die gleiche Richtung bez�glich der Zentralachse laufen
%	-> Test einfacher (Sicherheitsabstand immer in Richtung Zentralposition dazurechnen, und nicht in die andere Richtung)
\begin{figure}[htbp]
  \centering  
  \subfloat{\includegraphics[width=14cm]{Bilder/Screenshot_Dungeon_komplex}}\\    
  \subfloat{\includegraphics[width=14cm]{Bilder/Screenshot_Dungeon_einfach}}  
	\caption[Beispieldungeons von au�en]{\emph{Beispieldungeons von au�en}:
	Gezeigt sind durch das beschriebene Verfahren erstellte Dungeons.
	In Spielen w�ren solche Dungeons nur von innen erkundbar.
	}
	\label{B_BspDungeons}
\end{figure} 					% Hinzuf�gen von R�umen und G�ngen
%\include{K_ZusatzObjekte}					% Platzierung von Detailobjekten und Interaktionsobjekten
\chapter{Methoden zur Reduzierung des Renderaufwandes}
\label{K_LOD}

Ein genereller Konflikt in der Computergrafik ist der Widerspruch zwischen m�glichst hoher Detailf�lle
und m�glichst schneller Darstellung.
Diese Darstellung muss beispielsweise in Computerspielen in Echtzeit erfolgen, also mit mindestens 30 Bildern pro Sekunde.
F�r das Erreichen dieses Ziels sind bestimmte Techniken notwendig.
So sollen nicht sichtbare Objekte nach M�glichkeit von vornherein vom Zeichnen ausgeschlossen werden.
Sichtbare Objekte sollten in desto niedrigerer Detailstufen dargestellt werden, je weiter weg sie sich von der Kamera befinden.
Die Verwendung verschiedener Detailstufen stellt das Prinzip des \emph{Level of Detail} dar.
Das Weglassen von nicht sichtbaren Objekten ist das Prinzip des \emph{Culling}.
Arten des Culling sind beispielweise \emph{Frustum Culling} (Ausschlu� von Objekten au�erhalb des Sichtbereiches der Kamera)
und \emph{Occlusion Culling} (Ausschlu� von verdeckten Objekten).

%- M�glichkeiten f�r Level-of-Detail
%  - Zerlegen des Gebietes in Sektoren
%  - Erstellung verschiedener Detailstufen f�r jeden Sektor
%  - im Spiel: abh�ngig von der Entfernung des Sektors, diesen in passender Detailstufe darstellen
%  	- je weiter weg, desto weniger Details notwendig

\section{Sichtbarkeitsgraph}

F�r die Entfernungs- und Sichtbarkeitstest ist es von Vorteil, einen \emph{Sichtbarkeitsgraphen} zu haben.
%In seiner oft verwendeten Form stellt dieser einen Graphen dar,
%der angibt, welche Ecken von Polygonen andere Polygonecken sehen k�nnen \cite[S.322 ff.]{B_CompGeo} \cite[S.315]{B_AlgGeo}.
%Das Konzept wird nun auf Szenegraphen erweitert.
Dieser enth�lt die Informationen, welches Objekt mit welchem anderen Objekt, r�umlich betrachtet, direkt verbunden ist
und von welchem Objekt zu welchem anderen gesehen werden kann.
%!!!Valve Visleaf Methode + Quelle
Die Objekte im Sichtbarkeitsgraphen des Dungeons sind: jedes Subnetz der H�hle, alle G�nge und alle R�ume.

\ \\
Es wird bestimmt, wie die Objekte miteinander verbunden sind und welche Sichtbarkeitseigenschaften sie besitzen:

F�r jeden \emph{Raum} speichere, in welcher Himmelsrichtung welcher Gang angeschlossen ist.\footnote{Es
werden die urspr�nglichen Himmelsrichtungen der Subszene verwendet.}
Der Aufbau des Raums, z.B. von welcher Andockstelle zu welcher eine Sichtverbindung existiert, ist extern zu ermitteln,
da die zugrundeliegenden Subszenen importiert werden.

F�r jeden \emph{Gang} speichere, an welchem Ende welcher Raum bzw. die H�hle angeschlossen ist.
Weiterhin speichere $\vv{P_1}$, $\vv{P_2}$, $\vv{P'_1}$ und $-\vv{P'_2}$, um zu wissen, wo der Gang endet und in welche Richtung er dort zeigt.
Zus�tzlich wird getestet, ob man von einem Ende des Gangs zum anderen Ende sehen kann.

F�r jedes \emph{H�hlensubnetz} speichere die Grenzen $X_{min}$, $X_{max}$, $Y_{min}$, $Y_{max}$, $Z_{min}$ und $Z_{max}$
des zugrundeliegenden Voxelgebietes.
Zus�tzlich wird �berpr�ft, in welcher der sechs Koordinatenachsenrichtungen (X+,X-,Y+,Y-,Z+,Z-) man von diesem Subnetz andere Subnetze sehen kann,
also ob dort ein anderes Subnetz mit direkter Verbindung existiert.
  
\ \\
Importiert man diese Informationen in die gew�nschte Anwendung, so l�sst sich ein Sichtbarkeitsgraph
beispielsweise wie folgt erzeugen:
F�r die Knoten des Graphen werden die Verbindungsstellen zwischen direkt miteinander verbundenen Objekten verwendet.
Zum Beispiel w�rde eine Andockstelle Gang zu Raum einen Knoten repr�sentieren.
Die Kanten werden durch die Objekte oder deren Teile gebildet.
Beispielsweise w�rde ein Gang eine Kante darstellen, ein Raum entspricht bis zu sechs Kanten
(zwischen jedem Paar von Andockstellen eine Kante).
Als Kantengewicht wird der euklidische Abstand zwischen beiden Knoten genutzt.
Falls man von einer Verbindungsstelle nicht zur anderen sehen kann, dann wird keine Kante eingef�gt.
Wenn man also von einem Ende eines Gangs nicht zur anderen Ende sehen kann,
dann wird der Gang nicht als Kante verwendet.

Der Sichtbarkeitsgraph kann in einem Spiel oder einer anderer Anwendung per Breitensuche oder Dijkstra-Algorithmus
(dieser wird beispielsweise in \cite[S.670 ff.]{B_AlgoEinf} beschrieben) traversiert werden,
um zu testen, wie weit ein Objekt vom Objekt, in dem sich die Kamera befindet, effektiv entfernt ist
bzw. ob es �berhaupt sichtbar ist (dies ist nur der Fall, wenn es im Graphen erreichbar ist).

%Sichtbarkeits- und Entfernungsberechnungen f�r die einzelnen Teile des Dungeons:
%- Welcher Teil ist wie weit von welchem entfernt?
%- Sind zwei Abschnitte des Dungeons so verbunden, dass man von einem Abschnitt in den anderen sehen kann?
%
%- !Erkl�rung: Dijkstra-Algorithmus f�r Wegfindung
%- Gebietsnachbarschaften per Bitmatrix speichern (true = Gebiete benachbart)
%- Dijkstra-Algorithmus f�r Abst�nde zwischen 2 Gebieten
%-> finale Abstandsmatrix (wichtig f�r LOD)
%
%Problem: Abstandsmatrix teuer bei vielen Levelbestandteilen
%-> Modifikation: Graph wird gespeichert, Spiel f�hrt Dijkstra selbst durch
%
%jedes H�hlenteil: bis zu 6 Nachbarknoten Nord,Ost,S�d,West,Oben,Unten  + Nachbarknoten f�r G�nge
% - Test ob 2 H�hlenteile im Graph benachbart:
%   - sind die H�hlenteile benachbart im Voxelraum ? (Test 1)
%   - sind Voxel im Grenzbereich zwischen diesen beiden Teilen gesetzt? (Test 2)
%	 -> beide Tests positiv: Nachbarn
%
%jeder Raum: bis zu 4 Nachbarknoten: Nord,Ost,S�d,West
%jeder Gang: 2 Nachbarknoten
%
%jeder Knoten im Graph kennt seinen Typus: H�hlenteil, Gang (nicht-blickdicht), Gang (blickdicht), Raum (hier weiss der Raumersteller, ob blickdicht oder nicht)
%
%Strategie f�r Nutzung des Graphen in einem Spiel:
%- bestimme Knoten in dem sich der Spieler momentan befindet
%- wenn sich Knoten ge�ndert hat:
%  - berechne sichtbare Kinder (per Dijkstra)
%  - lege passende Detailstufe f�r jedes sichtbare Kind fest (abh�ngig von der Entfernung)
%- stelle alle sichtbaren Knoten mit der jeweils festgelegten Detailstufe dar

\subsection{Sichtbarkeitsinformationen der H�hlensubnetze}

Das Verfahren f�r die Sichtbarkeitsinformationen der H�hlensubnetze ist vergleichsweise einfach umsetzbar.
Es wird f�r jede Achsenrichtung getestet, ob an irgendeiner Position 
auf beiden Seiten der Grenzfl�che ein Voxel mit Belegung $1$, also ein Freiraumvoxel, existiert.
Falls ja, w�rde die H�hle �ber die Grenzfl�che hinweg verlaufen
und man k�nnte entlang der Achse von dem einen Subnetz in das andere schauen.

Formel \ref{F_VisHoehle} beschreibt diesen Test.
Der Parameter $Vis$ gibt f�r jede Achsenrichtung an, ob andere H�hlensubnetze sichtbar sind.
Wenn der Wert $wahr$ betr�gt, ist eine solche Sichtm�glichkeit vorhanden.

\begin{equation}
 \begin{aligned}
 \label{F_VisHoehle}
	Vis_{X+} =& \exists (y,z) \in [Y_{min},Y_{max}]\times [Z_{min},Z_{max}]:\\
	          & Voxel(X_{max},y,z) = 1 \wedge Voxel(X_{max}+1,y,z) = 1 \\
	Vis_{X-} =& \exists (y,z) \in [Y_{min},Y_{max}]\times [Z_{min},Z_{max}]:\\
						& Voxel(X_{min},y,z) = 1 \wedge Voxel(X_{min}-1,y,z) = 1 \\
	Vis_{Y+} =& \exists (x,z) \in [X_{min},X_{max}]\times [Z_{min},Z_{max}]:\\
						& Voxel(x,Y_{max},z) = 1 \wedge Voxel(x,Y_{max}+1,z) = 1 \\
	Vis_{Y-} =& \exists (x,z) \in [X_{min},X_{max}]\times [Z_{min},Z_{max}]:\\
						& Voxel(x,Y_{min},z) = 1 \wedge Voxel(x,Y_{min}-1,z) = 1 \\
	Vis_{Z+} =& \exists (x,y) \in [X_{min},X_{max}]\times [Y_{min},Y_{max}]:\\
						& Voxel(x,y,Z_{max}) = 1 \wedge Voxel(x,y,Z_{max}+1) = 1 \\	
	Vis_{Z-} =& \exists (x,y) \in [X_{min},X_{max}]\times [Y_{min},Y_{max}]:\\
						& Voxel(x,y,Z_{min}) = 1 \wedge Voxel(x,y,Z_{min}-1) = 1 	
	\end{aligned}
\end{equation}

%Anmerkung: sehr allgemeiner Test

\subsection{Sichtbarkeitsinformationen der G�nge}

Es soll gepr�ft werden, ob man von einem Ende des Gangs zum anderen Ende sehen kann.
Dazu wird erst einmal ein beliebiges St�ck eines polygonalen Schlauchs betrachtet.
Gegeben seien drei Segmente $R_1$, $R_2$ und $R_3$ des Schlauchs.

Die �berlegung ist folgende:
Um $R_1$ und $R_3$ wird eine Bounding Box gelegt, welche im Folgenden als \emph{Viewing Volume} bezeichnet wird.
Wenn die Querschnittsfl�che von $R_2$ dieses Viewing Volume nicht schneidet, dann kann kein Lichtstrahl von $R_1$ nach $R_3$ gelangen.
Somit k�nnte man nicht von $R_1$ nach $R_3$ sehen und
dieser Abschnitt des polygonalen Schlauchs w�re blickdicht.
Das Prinzip ist in Abbildung \ref{B_Vistest} dargestellt.

\begin{figure}[hbtp]
  \centering  
	\includegraphics[width=14cm]{Bilder/Vistest}
	\caption[Blickdichtigkeitstest f�r G�nge]{\emph{Blickdichtigkeitstest f�r G�nge}:
	 Zu sehen sind die drei Segmente des Schlauchs $R_1$, $R_2$, $R_3$ (gr�n) und das Viewing Volume (blau).
	 Der eigentliche Kollisionstest kann mit den auf eine Ebene projizierten Elementen durchgef�hrt werden.}
	 \label{B_Vistest}
\end{figure}

\ \\
Diesen Test kann man nun f�r alle Segmente des Gangs plus die beiden Adapter durchf�hren.
Problematisch ist die Laufzeitkomplexit�t bei insgesamt $n$ Segmenten von $\mathcal O(n^3)$
(durch $\mathcal O(n^2)$ Viewing Volumes mit je $\mathcal O(n)$ Tests).\footnote{Ausgegangen
wird von einer konstanten Laufzeit f�r den eigentlichen Schnitttest.}
F�r alle Segmente ist das Verfahren dadurch in der Praxis nicht anwendbar.

Dies f�hrt zu der Idee, die Tests auf "`interessante"' Segmente, im folgenden \emph{Region of Interest} genannt, zu beschr�nken.
Als Regions of Interest sollen beide Andockstellen und die Segmente des Gangs f�r $\vv{P(0)}$ und $\vv{P(1)}$ gelten.

Hinzu kommen Segmente, deren Abstand zur Achse $\overline{P(0)P(1)}$ m�glichst gro� ist,
da dort die Wahrscheinlichkeit f�r eine Blickverdeckung am h�chsten ist.
Gesucht werden somit die Segmente an $\vv{P(t)}$ f�r die $D_{Achse}$ nach Formel \ref{F_LODAbstandAchse} ein lokales Maximum darstellt.
F�r eine einfachere Berechnung wird der quadratischen Abstand $\left(D_{Achse}\right)^2$ zur Bestimmung der Maxima verwendet.
Die dadurch ermittelten Maxima sind identisch mit denen des einfachen Abstands.

\begin{equation}
 \begin{aligned}
 \label{F_LODAbstandAchse}
		D_{Achse} = \left|\vv{P(t)}-\vv{P(0)}-  \left( \vv{N_{Achse}} \cdot (\vv{P(t)}-\vv{P(0)}) \right) \cdot  \vv{N_{Achse}}\right|  \\
		mit\ \vv{N_{Achse}} = \frac{\vv{P(1)}-\vv{P(0)}}{\left|\vv{P(1)}-\vv{P(0)}\right|}
	\end{aligned}
\end{equation}

\ \\
Die Anzahl der lokalen Maxima ist maximal $2$.
%- Erl�uterung (!siehe Notizbuch Do 15.12.2011):
\emph{Beweis}: Der quadratischer Abstand zur Achse ist ein Polynom 6. Ordnung in $t$.
Zur Bestimmung der Extrema wird erste Ableitung gleich $0$ gesetzt, welche ein Polynom 5. Ordnung darstellt.
Gesucht werden also die Nullstellen eines Polynoms 5. Ordnung, wovon es maximal $5$ geben kann.
Es gibt somit maximal $5$ Extrema.
Davon sind mindestens $2$ lokale Minima (quadr. Abstand = 0): der Anfangspunkt $\vv{P(0)}$ und der Endpunkt $\vv{P(1)}$.

Zwischen $2$ bestehenden Maxima muss stets ein Minimum liegen, da die Kurve stetig ist.
Wenn $2$ lokale Maxima vorliegen, dann muss folglich dazwischen ein lokales Minimum existieren, d.h. $5$ lokale Extrema w�rden vorliegen.
Mehr Maxima kann es nicht geben, da dann mehr als $5$ Extrema existieren w�rden.
Folgerung: Es gibt maximal $2$ lokale Maxima. \emph{q.e.d.}

%Es gibt drei F�lle.
%Wenn kein Maximum existiert verl�uft die Kurve auf der Achse.
%Wenn $1$ lokales Maximum existiert, dann gibt es keine weiteren lokalen Minima, d.h. $3$ lokale Extrema liegen vor.
%  - 1.Ableitung = 0 -> Nullstellen eines Polynoms 5. Ordnung gesucht -> max. 5 Nullstellen -> max. 5 lok. Extrema

\ \\
Da die exakten Nullstellen eines Polynoms 5. Ordnung schwer zu bestimmen sind,
soll die Ermittlung der lokalen Maxima mit einer numerischen L�sung erfolgen.
Beim der Erstellung des polygonalen Schlauchs wird w�hrend des Abtastens der Kurve der Abstand von $\vv{P(t)}$ zur Achse �berwacht.
Liegt der Abstand im letzten und folgenden Schritt niedriger als im aktuellen, wurde ein lokales Maximum gefunden.

Insgesamt gibt es also maximal sechs Regions of Interest, was eine konstante Maximalanzahl n�tiger Tests bedeutet.
Die Laufzeitkomplexit�t des Tests f�r $n$ Segmente liegt nun bei $\mathcal O(1)$.

\subsubsection{Durchf�hrung des Tests mittels projizierten Bounding Boxes}

Die Regions of Interest $R_1$, $R_2$, ... , $R_n$ seien nach ihrer Reihenfolge entlang des Gangs angeordnet.
Zuerst erfolgt die  Konstruktion der Viewing Volumes aus allen Paaren $(R_i,R_j)$ mit $i+1<j$, also mit mindestens einer Region of Interest dazwischen.
Dann werden f�r jedes Viewing Volume aus $(R_i,R_j)$ alle Regions of Interest $R_k$ mit $i<k<j$ getestet.
Der eigentliche Schnitttest wird mit projizierten Bounding Boxes durchgef�hrt.
%Gegeben seien die drei Region of Interests $R_1$, $R_2$ und $R_3$.

Eine Region of Interest sei gegeben durch die Position $\vv{P_R}$,
die Vektoren des lokales Koordinatensystems $\vv{Links_R}$ und $\vv{Oben_R}$
und die Bounding Box $[RMin_x,RMax_x]\times[RMin_y, RMax_y]$ des Segmentquerschnitts.
Diese Bounding Box ist auf das 2D-Koordinaten"-system mit $-\vv{Links_R}$ als X-Achse, $\vv{Oben_R}$ als Y-Achse
und $\vv{P_R}$ als Nullpunkt bezogen.
  
Ein Viewing Volume wird durch seinen Bezugspunkt $\vv{P_V}$,
die Vektoren des lokalen Koordinatensystems $\vv{Links_V}$ und $\vv{Oben_V}$
und seine Bounding Box $[VMin_x,VMax_x]\times[VMin_y,VMax_y]$ gekennzeichnet.
Diese Bounding Box ist auf das 2D-Koordinaten"-system mit $-\vv{Links_V}$ als X-Achse, $\vv{Oben_V}$ als Y-Achse
und $\vv{P_V}$ als Nullpunkt bezogen.

\ \\
Zur Konstruktion des Viewing Volumes aus $R_i$ und $R_j$ wird zuerst das lokale Koordinatensystem nach Formel \ref {F_ViewVolKoordsys} ermittelt.
Der Normalenvektor der Viewing Volume-Ebene verl�uft entlang der Achse $\overline{P_{Ri}P_{Rj}}$.
Bezugspunkt sei $\vv{P_V} = \vv{P_{Ri}}$.

\begin{equation}
 \begin{aligned}
		\vv{Links_V} = \frac{\vv{vorn} \times (0,1,0)}{\left|\vv{vorn} \times (0,1,0)\right|},\ 
		\vv{Oben_V} = \vv{Links_V} \times \vv{vorn}\\
		\ mit\ \vv{vorn} = \frac{\vv{P_{Rj}}-\vv{P_{Ri}}}{\left|\vv{P_{Rj}}-\vv{P_{Ri}}\right|}
 \end{aligned}
 \label{F_ViewVolKoordsys}
\end{equation}

\ \\
Nun werden die beiden Bounding Boxes der Region of Interests $R_i$ und $R_j$ in die Ebene des Viewing Volumes projiziert.
Die generelle Projektion einer Bounding Box einer Region of Interest $R$ erfolgt durch die Projektion
aller vier Bounding Box-Eckpunkte auf die Ebene und das Bilden einer neuen Bounding Box um diese Projektionen.
Formel \ref{F_BBProj} gibt das Beispiel f�r den unteren linken Eckpunkt an.
Dessen Projektion sei als $\vv{Q_{proj}}$ bezeichnet.
Die anderen Eckpunkte werden �quivalent behandelt.

\begin{equation}
 \begin{aligned}
 		\vv{temp} =& \vv{P_R} - \vv{P_V} -RMin_x \cdot \vv{Links_R} + RMin_y \cdot \vv{Oben_R}\\
 		\vv{Q_{proj}} =& \left(-\vv{temp} \cdot \vv{Links_V}, \vv{temp} \cdot \vv{Oben_V}\right)
 		\label{F_BBProj}
 \end{aligned}
\end{equation}

\ \\
Bei den Bounding Boxes von $R_i$ und $R_j$ kann dabei der Term $\vv{P_R} - \vv{P_V}$ aus Zeile $1$
weggelassen werden, da dieser in Normalenrichtung der Ebene verl�uft und somit bei
der Projektion entf�llt.
Die Viewing Volume-Bounding Box $[VMin_x,VMax_x]\times[VMin_y,VMax_y]$ wird durch eine Bounding Box
um alle acht projizierten Eckpunkte gebildet.

\ \\
F�r den Test wird nun die Region of Interest $R_k$ auf die Ebene des Viewing Volumes projiziert
und der Schnitttest nach Formel \ref{F_VisG} durchgef�hrt.
Die projizierte Bounding Box der Region of Interest
sei $[R_kMin_x^{pr},R_kMax_x^{pr}]\times[R_kMin_y^{pr},R_kMax_y^{pr}]$.

Wenn sich die Bounding Boxes �berschneiden ($Vis_G = wahr$), ist der Gang m�glicherweise nicht blickdicht.
Ansonsten liegt \emph{definitive Blickdichtigkeit} vor,
d.h. man kann in keinem Fall von einem Ende des Gangs durch den Gangs zum anderen Ende sehen.

\begin{equation}
 \begin{aligned}
		Vis_G = & \left(R_kMin_x^{pr} < VMax_x\right) \wedge \left(VMin_x < R_kMax_x^{pr}\right) \\
		        &\wedge \left(R_kMin_y^{pr} < VMax_y\right) \wedge \left(VMin_y < R_kMax_y^{pr}\right)
 \end{aligned}
 \label{F_VisG}
\end{equation}

\ \\
Der Vorteil der Projektion liegt in einem schnelleren Test.
Wenn ein Viewing Volume einmal erstellt ist,
k�nnen alle gew�nschten Regions of Interest per zweidimensionalen �berschneidungstest �berpr�ft werden,
anstatt dreidimensional, wie es ohne Projektion der Fall w�re.
In Abbildung \ref{B_VisGangShot} ist ein Vergleich zwischen einem blickdichten
und einem nicht blickdichten Gang aufgezeigt.
F�r das Bestimmen dieser Informationen wurde der obige Test verwendet.

\begin{figure}[htbp]
  \centering  
  \hfill
  \subfloat{\includegraphics[width=7.4cm]{Bilder/Screenshot_Gang_Blickdicht_false}}
  \hfill
  \subfloat{\includegraphics[width=7.4cm]{Bilder/Screenshot_Gang_Blickdicht_true}}  
  \hfill
	\caption[Blickdichtigkeit von G�ngen]{\emph{Blickdichtigkeit von G�ngen}: v.l.n.r. (a) nicht blickdicht, (b) blickdicht}
	\label{B_VisGangShot}
\end{figure}


%============================================================================================
\section{Reduktion von Dreiecksnetzen}

%Die meisten Objekte des Dungeons bestehen aus Dreiecksnetzen.
%Aus Dreiecksnetzen aufgebaute Objekte werden f�r die Generierung weiterer Detailstufen reduziert.
%Die Reduktion von Dreiecksnetzen dient der Erstellung zus�tzlicher Detailstufen f�r die jeweiligen Objekte.
F�r die Generierung verschiedener Detailstufen von aus Dreiecksnetzen aufgebauten Objekten werden deren Netze reduziert.
Wichtig bei dieser Reduktion ist zum einen, dass das reduziertes Objekt m�glichst weit dem Original entspricht,
und zum anderen, dass es m�glichst einfach gestaltet ist, d.h. m�glichst wenig Dreiecke und Vertices besitzt.

\ \\
F�r die \emph{R�ume} bleibt die Reduktion dem Raumgestalter �berlassen.
Es erfolgt f�r diese keine Reduktion im Generator.

F�r die \emph{G�nge} erm�glicht die schlauchartige Struktur eine effiziente Reduktion.
F�r Reduktionsstufe $i$ wird der Gang erneut erstellt, mit gleichem Gangsegmentabstand $w$ wie im Original,
aber nur jedes $2^i$-te Segment wird wirklich eingef�gt.
Zus�tzlich werden alle lokalen Maxima eingef�gt, damit die Blickdichtigkeitsberechnung f�r alle reduzierten Versionen ebenso gilt.
Resultat ist eine geringere Anzahl ben�tigter Vertices und Dreiecke bei prinzipiell gleichem Gangverlauf.
Diese Anzahl wird bei jeder zus�tzlichen Reduktionsstufe halbiert.

F�r die \emph{H�hle} muss ein anderer Algorithmus eingesetzt werden.
Ziel ist die Generierung mehrerer Detailgrade f�r jedes Subnetz der H�hle.
Wichtig ist, dass die Meshes aller Detailstufen (inkl. Original) an den R�ndern zusammenpassen.
D.h. Position und Normalen m�ssen an diesen Stellen exakt gleich sein,
damit beliebige Detailstufen in Kombination verwendet werden k�nnen.
F�r jede Detailstufe und alle m�glichen Randbeziehungen je eine extra Variante zu generieren
w�re deutlich zu speicheraufw�ndig, da bis zu 26 Nachbarmeshes m�glich sind
(wobei 6 davon eine Grenzfl�che teilen, also eine gro�e Zahl gemeinsamer Grenzvertices besitzen k�nnen).

\subsection{Reduktion der H�hlensubnetze}

% - Verwendung von CGAL? (Algorithmus: Triangulated Surface Mesh Simplification)
% http://www.cgal.org/Manual/latest/doc_html/cgal_manual/Surface_mesh_simplification/Chapter_main.html
% - eventuelles Problem: �berg�nge Gang/Raum <-> H�hle:
%	- meist unregelm��ig, kein harter Schnitt
%- �bersicht �ber die Verfahren
%- umgesetztes Verfahren
%  - Laufzeitbetrachtung
Die Topologie eines Meshs kann nach \cite[S.14-17]{B_LoD3D} �ber den Genus und die Mannigfaltigkeit beschrieben werden.
Der \emph{Genus} ist die Anzahl von L�chern im Mesh. Eine Kugel hat einen Genus von $0$, ein Torus einen von $1$.
In \emph{mannigfaltigen} Meshes hat jede Kante exakt zwei benachbarte Dreiecke und jedes Dreieck teilt seine Kanten mit exakt drei benachbarten Dreiecken. \\
\emph{Topolgogie-erhaltende} Algorithmen bewahren die Mannigfaltigkeit und den Genus in jedem Schritt.
Manche Algorithmen sind \emph{topologie-tolerant}, sie ignorieren nicht-mannigfaltige Topologie und verhalten sich ansonsten topologie-erhaltend. \\
\emph{Topologie-modifizierende} Algorithmen k�nnen die Mannigfaltigkeit und den Genus von Meshes ver�ndern,
was �blicherweise schlechtere visuelle Ergebnisse zur Folge hat.
 
W�nschenswert w�re ein topologie-erhaltender bzw. topologie-toleranter Algorithmus, damit die H�hlenstruktur besser erhalten bleibt.
Da die �ffnungen f�r G�nge und die Grenzen zu anderen Subnetzen im Gegensatz zum Rest des Meshes nicht-mannigfaltig sind
(die Grenzkanten haben nur ein benachbartes Dreieck),
muss bei Verwendung eines topologie-erhaltenden Algorithmus dieser zu einem topologie-toleranten Algorithmus modifiziert werden.

%Fidelity-based simplification
%Budget-based simplification
\ \\
Viele Reduktionsalgorithmen verwenden folgendes Prinzip:
Der Reduktionsprozess setzt sich aus der Durchf�hrung einer meist gr��eren Anzahl von Anwendungen lokaler Reduktionsoperatoren zusammen.
Die Auswahl der Region, an der der Operator angewendet wird, erfolgt i.d.R. anhand einer Fehlermetrik.
Die Region mit dem kleinsten zu erwartenden Fehler wird als n�chstes verwendet.
\cite[S.21 ff.]{B_LoD3D} beschreibt folgende lokalen Reduktionsoperatoren:

\begin{figure}[hbtp]
  \centering  
	\includegraphics[width=13cm]{Bilder/EdgeCollapse}
	\caption[Edge Collapse]{Edge Collapse \cite{garland1997surface}}
	\label{B_EdgeCollapse}
\end{figure}

\ \\
\emph{Edge Collapse}: Als Reduktionsmethode wird die Kontraktion einer Kante verwendet (siehe Abbildung \ref{B_EdgeCollapse}).
Die zwei Endpunkte einer Kante werden zu einem Vertex verschmolzen.
Die Position des resultierenden Vertex  kann nach verschiedenen Methoden berechnet werden, z.B. als Mitte zwischen beiden Ausgangsvertices.
Beide Endpunkte der Kante d�rfen nicht mehr als zwei gemeinsame Nachbarn haben, da sonst nicht-mannigfaltige Strukturen entstehen
(siehe Abbildung \ref{B_EdgeCollapseMannig}).\\
\emph{Vertex-Pair Collapse}: Die Reduktionsmethode funktioniert �quivalent zu Edge Collapse,
doch hier brauchen die verschmelzenden Vertices nicht zwangsl�ufig durch eine Kante verbunden zu sein.\footnote{Diese
Beschreibung folgt \cite{garland1997surface}.
\cite{B_LoD3D} nennt diesen Operator explizit f�r zwei nicht-verbundene Vertices.}\\
\emph{Triangle Collapse}:
Alle drei Vertices eines Dreiecks werden zu einem einzigen Vertex verschmolzen.
Die Methode ist �quivalent zu zwei hintereinander ausgef�hrten Edge Collapses. \\
\emph{Cell Collapse}:
Der Raum wird mittels eines gleichm��igen Gitters in Zellen unterteilt.
Alle Vertices, die zusammen in einer Zelle sind, werden zu einem Vertex verschmolzen. \\
\emph{Vertex Removal}:
Ein Vertex wird entfernt.
Das entstehende Loch des Meshes wird retrianguliert und dadurch wieder geschlossen. \\
\emph{Polygon Merging}:
Benachbarte, m�glichst weit koplanare Polygone werden zu einem gr��eren Polygon verschmolzen.
Dieses wird sp�ter trianguliert. \\
\emph{General Geometric Replacement}:
Eine Menge von adjazenten Dreiecken wird durch eine andere Menge von adjazenten Dreiecken ersetzt,
so dass die Begrenzung die gleiche bleibt.
Die neue Menge von Dreiecken besitzt dabei weniger Dreiecke als die alte.

\begin{figure}[hbtp]
  \centering  
	\includegraphics[width=14cm]{Bilder/Mannigfaltigkeit}
	\caption[Verlust der Mannigfaltigkeit bei mehr als zwei gemeinsamen Nachbarn]{\emph{Verlust der Mannigfaltigkeit bei mehr als zwei gemeinsamen Nachbarn}:
	 betrachten wird ein Ausschnitt eines Dreiecksnetzes: v.l.n.r. \\
	 (a) Gezeigt ist die Ausgangstriangulation. Die zu kontrahierende Kante ist rot dargestellt. Beide Endpunkte der Kante (gr�n) haben drei gemeinsame Nachbarpunkte (gelb). \\
	 (b) Durch den Edge Collapse verschmelzen das rote und blaue Dreieck zum orangenen Dreieck.
	 Die Mannigfaltigkeit ist verloren gegangen, da nun zwei Kanten dieses Dreiecks zu keinen weiteren Dreiecken geh�ren.}
	 \label{B_EdgeCollapseMannig}
\end{figure}

\ \\
Im Vergleich lassen sich die Collapse-Operatoren am einfachsten implementieren.
Cell Collapse stellt das einfachste Verfahren dar, ist aber, wie auch Vertex-Pair Collapse, ein topologie-modifizierender Reduktionsschritt.
Edge Collapse und Triangle Collapse sind topologie-erhaltend.
%Zudem ist er nicht invariant gegen�ber Rotation und Translation, welches aber im Kontext der H�hlendreicksnetze irrelevant w�re.
Edge Collapse und Vertex-Pair Collapse lassen sich besser steuern als Triangle Collapse, da sie pro Schritt kleinere �nderungen vornehmen.

%Tabelle mit Eigenschaften (diese auch erkl�ren)
%- Laufzeit
%- Speicherbedarf
%- Topologieerhaltung
%- erfordert Retriangulierung
%- ...
%\ \\
%Error Metrics

\ \\
M�glichkeiten f�r Fehlermetriken werden in \cite[S.47 ff.]{B_LoD3D} genannt.
Beispiele daf�r sind: die Abst�nde zwischen vom Reduktionsoperator betroffenen Vertices, die Abst�nde
vom betroffenen Vertex und der durch seine Umgebung gelegte Fl�che, die Normalen�nderungen
im betroffenen Gebiet.

Die lokalen Reduktionsoperatoren lassen sich ggf. in Verbindung mit entsprechenden Fehlermetriken zu konkreten
Algorithmen, wie sie in \cite[S.121 ff.]{B_LoD3D} beschrieben werden, kombinieren.
\emph{Vertex Clustering}, \emph{Vertex Decimation} und \emph{Quadric Error Metrics} stellen dabei drei der wichtigsten Verfahren dar.

% (mit Vorteilen und Nachteilen)
\ \\
\emph{Vertex Clustering}, wie es von Rossignac und Borrel in \cite{rossignac1992multi} beschrieben wurde,
verwendet die Methode Cell Collapse zur Reduktion.
Zuerst wird jedem Vertex eine Wichtigkeit bzw. Wichtung zugeordnet.
Vertices gro�er Fl�chen oder an ausgepr�gten Kr�mmungen werden st�rker gewichtet. 
Dann wird ein 3D-Gitter �ber das zu reduzierende Objekt gelegt und dieses so in Zellen aufgeteilt.
Alle Vertices in einer Zelle werden zum am h�chsten gewichteten Vertex kollabiert,
d.h. der wichtigste Vertex repr�sentiert nun alle Vertices der Zelle.
Degenerierte Dreiecke werden entfernt.
Der Algorithmus ist einer der schnellsten bekannten Reduktionsalgorithmen.
%und l�uft in $\mathcal O(n)$ (mit $n$ als Anzahl der Vertices), falls jeder Vertex eine 
Nachteil ist vor allem, dass der Algorithmus stark topologie-modifizierend ist.

\ \\
\emph{Vertex Decimation} wird in  \cite{schroeder1992decimation} beschrieben und nutzt den Operator Vertex Removal.
Der Algorithmus verwendet mehrere Durchl�ufe.
In jedem Durchlauf werden Vertices nach bestimmten Kriterien entfernt und danach die entstandenen L�cher retrianguliert.
Als Hauptkriterium wird die Distanz des zu entfernenden Vertex zur Ebene, die durch die benachbarten Vertices gelegt wird, genommen.
Ist diese unter einem Grenzwert $\epsilon$, so wird der Vertex entfernt.
Das $\epsilon$ wird in jedem Durchlauf erh�ht.
Der Algorithmus terminiert beim Erreichen des gew�nschten Reduktionslevels.
%Die Laufzeit des Algorithmus betr�gt $\mathcal O(m\cdot n)$ mit $m$ als Anzahl der Durchl�ufe und $n$ als Anzahl der Vertices
Er erweist sich in der Praxis als recht schnell.
Nachteil ist die schlechte Steuerbarkeit �ber den $\epsilon$-Wert.
Ist dieser zu gro�, wird das Modell zu stark vereinfacht. Bei zu kleinen Werten werden zu viele Durchl�ufe ben�tigt, wodurch der Algorithmus langsam wird.
Im Kontext der H�hlenmeshes ist anzumerken, dass durch ihre zerkl�ftete Struktur ein �hnlicher Distanz-Wert f�r die meisten Vertices vorliegt.
Weiterhin w�rde der Algorithmus lokal stark ausgepr�gte Voxelkanten hervorheben und eher weichere Kanten verschwinden lassen.
Dadurch ist diese Methode hierf�r ungeeignet.
  
\ \\
\emph{Quadric Error Metrics} wurde von Garland  und Heckbert 1997 in \cite{garland1997surface} vorgestellt.
Der Algorithmus basiert auf Vertex-Pair Collapses.
F�r die Kontraktion zugelassene Vertex-Paare sind alle Paare, die �ber eine Kante verbunden sind,
sowie alle Paare, deren Abstand zueinander einen spezifierten Grenzwert unterschreitet.
Als Fehlermetrik wird die Quadric Error Metric benutzt. Jedem Vertex ist eine Quadrik $Q$ zugeordnet,
die kumulativ speichert, welche Verzerrungen durch diesen Vertex und seine bisherigen Collapses verursacht werden.
Die Quadriken werden am Anfang f�r jeden Vertex berechnet.
Bei jedem Collapse von Vertex $a$ und Vertex $b$ zu Vertex $c$ ergibt sich $Q_c = Q_a + Q_b$.
Die Vertex-Paare $(a,b)$ werden in einer Priorit�tswarteschlange,
sortiert nach steigendem Quadric Error $Q_a + Q_b$, gespeichert und abgearbeitet.
Der Algorithmus terminiert, wenn die gew�nschte Anzahl an Dreiecken erreicht wurde.
Das Verfahren erzeugt qualitativ sehr hochwertige Ergebnisse und die erzeugte Dreiecksanzahl l�sst sich exakt steuern.
Allerdings hat die Abarbeitung der Vertex-Paare durch die Verwendung einer Priorit�tswarteschlange
eine Laufzeitkomplexit�t von $\mathcal O(n \log n)$ ($n$ als Anzahl der Vertices), was bei einer hohen Vertexanzahl schnell zu vergleichsweise hohen Laufzeiten f�hrt.

\subsubsection{Reduktionsalgorithmus}

%Idee: Triangle Collapse
%- Dreiecke werden zuammengemergt
%  - jedes Dreieck, bei dem alle 3 Vertices noch original sind (also nicht durch Collapse entstanden)
%- ansonsten keine weitere Metrik
%  - Vermutung: mit Metrik w�rde man entweder auf Collapses verzichten m�ssen, damit Nachbarcollapses m�glich sind
%							oder viel Rechenaufwand betreiben m�ssen, um die optimale Kontraktion zu finden (Vermutung: NP-schwer)
%- damit keine 1-Voxel-dicken S�ulen zu arg deformiert werden:
%	- bei Umwandliung in Dreiecksnetz darauf achten, dass Vorder- und R�ckseite versetzt trianguliert (Projektionen der Diagonalen schneiden sich)							
%-> durch die Voxelstruktur �hnliches Prinzip wie Vertex Clustering
%- lineare Laufzeit, Speicherplatz-sparend
%
%- die verwendete Fehlermetrik in Verbindung mit der regelm��igen Struktur des Netzes f�hrt zu einer Art Vertex Clustering,
%welches aber im Gegensatz zu echtem Vertex Clustering Topologie erhaltend ist
%
%- falls Vertex-Beleuchtung verwendet wird, k�nnte man eine modifizierte Fehlermetrik verwenden:
%  je �hnlicher die Vertex-Normalen desto eher wird der Triangle-Collapse durchgef�hrt
Der implementierte Ansatz verwendet Edge Collapse als Reduktionsoperator.
Als Fehlermetrik wird in jeder Reduktionstufe jeder Vertex h�chtens einmal kollabiert.
Dies bedeutet: Wenn von einer Kante mindestens ein Vertex schon kollabiert ist, dann erfolgt keine Kontraktion.
Diese Metrik sei als \emph{Single Contraction} bezeichnet.
Vertices an Meshgrenzen werden nicht kollabiert, damit die �ffnungen f�r G�nge und Grenzen zu anderen H�hlenmeshes erhalten bleiben.

Durch die auf Grund der Voxelbasiertheit sehr regelm��ige Struktur des Dreiecksnetzes,
wirkt diese Methode �hnlich dem Cell Collapse bzw. der Vertex Clustering-Methode,
ist aber topologie-erhaltend.
Zus�tzlich wird eine Gl�ttung des Meshs erreicht, statt Kanten der Voxel weiter hervorzuheben.

\begin{algorithm}
\caption{Reduktion eines Dreicksnetzes, per Edge Collapse und Single Contraction}
\label{A_LODHoehle}
\algorithmicrequire{Dreicksnetz aus Dreiecken $(i,j,k)$ und Vertices $V_i$, Adjazenzliste aller Vertices.\\ $Adr$, $Flag$ und $Flag'$ sind mit $-1$ initialisiert.} \\
\algorithmicensure{reduziertes Dreicksnetz aus Dreiecken $(i,j,k)$ und Vertices $V'_i$}
\begin{algorithmic}[1]
		
	\ForAll{$(i,j) \in Kanten\ aller\ Originaldreiecke$}
		\If{$Adr(i)+Adr(k)+\vv{V_i.T}.x+\vv{V_j.T}.x+\vv{V_i.T}.y+\vv{V_j.T}.y < -5.5$}
			
			\ForAll{$k \in Adjazenzliste(i)$}
				\If{$Adr(k)<0$}			
						\State $Flag(k) \gets i$ \Comment {Originalvertex markieren}
				\Else
						\State $r \gets Adr(k)$	 \Comment {kontrahierten Vertex markieren}
						\State $Flag'(r) \gets i$			
				\EndIf
			\EndFor
			
			\State $doppelt \gets 0$
			
			\ForAll{$k \in Adjazenzliste(j)$}
				\If{$Adr(k)<0$}	
						\State $\textbf{if}\ Flag(k)=i\ \textbf{then}\ doppelt \gets doppelt+1$	
				\Else
						\State $r \gets Adr(k)$
						\State $\textbf{if}\ Flag'(r)=i\ \textbf{then}\ doppelt \gets doppelt+1$		
				\EndIf
			\EndFor
			
			\If{$doppelt<3$}	\Comment {Bedingung f�r Erhaltung der Mannigfaltigkeit}
				\State $\textbf{create}\ Vertex\ V'_{r}$ \Comment {Durchf�hrung des Edge Collapse}
				\State $\vv{V'_r.Pos} \gets (\vv{V_i.Pos} + \vv{V_j.Pos}) / 2$
				\State $\vv{V'_r.T} \gets (-1,-1)$		 \Comment {kein Gang�ffnungs- oder Grenzvertex}
				\State $Adr(i) \gets r$
				\State $Adr(j) \gets r$
			\EndIf		
			
		\EndIf
	\EndFor
	
	\ForAll{$V_i \in Originalvertices$} \Comment {Nicht kontrahierte Vertices �bernehmen}
		\If{$Adr(i)<0$}
			\State $\textbf{create}\ Vertex\ V'_{r}$
			\State $\vv{V'_r.Pos} \gets \vv{V_i.Pos}$
			\State $\vv{V'_r.T} \gets \vv{V_i.T}$
			\State $\vv{V'_r.N} \gets \vv{V_i.N}$		
			\State $Adr(i) \gets r$			
		\EndIf
	\EndFor
	
	\ForAll{$(i,j,k) \in Originaldreiecke$} \Comment {Nicht degenerierte Dreiecke �bernehmen}
			\If{$ Adr(i)\neq Adr(j)\wedge Adr(i)\neq Adr(k)\wedge Adr(j)\neq Adr(k)$}
				\State $\textbf{create}\ Dreieck'(Adr(i),Adr(j),Adr(k))$			
			\EndIf	
	\EndFor
	
\end{algorithmic}
\end{algorithm}

\begin{figure}[htbp]
  \centering  
  \subfloat{\includegraphics[width=14cm]{Bilder/Screenshot_LOD_0}}\\    
  \subfloat{\includegraphics[width=14cm]{Bilder/Screenshot_LOD_1}}  
	\caption[Reduktion der H�hlennetze]{\emph{Reduktion der H�hlennetze}:
	 v.o.n.u. (a) Originaldetailstufe, (b) eine Reduktionsstufe}
	 \label{B_LODH1}
\end{figure} 

\begin{figure}[htbp]
  \centering  
  \subfloat{\includegraphics[width=14cm]{Bilder/Screenshot_LOD_2}}\\    
  \subfloat{\includegraphics[width=14cm]{Bilder/Screenshot_LOD_4}}  
	\caption[Reduktion der H�hlennetze ff.]{\emph{Reduktion der H�hlennetze ff.}:
	 v.o.n.u. (a) zwei Reduktionsstufen, (b) vier Reduktionsstufen}
	\label{B_LODH2}
\end{figure}

\ \\
Das Verfahren ist in Algorithmus \ref{A_LODHoehle} aufgezeigt und wird f�r jedes Subnetz der H�hle separat angewendet.
Das Subnetz sei gegeben durch seine Vertex- und Dreiecksmenge wie in Kapitel \ref{KK_Dreiecksnetze} erl�utert.
Die Grenzvertices haben eine X-Texturkoordinate von $1$ (siehe Kapitel \ref{KK_VoxelNormalen}), andere von $-1$.
�quivalent sollen Vertices an Gang�ffnungen eine Y-Texturkoordinate von $1$ haben, andere von $-1$ \footnote{Dies
l�sst sich bei der Umwandlung der Voxel in das Dreiecksnetz einfach bewerkstelligen:
Vertices mit solchen Nachbarvoxeln, die Belegung $10$ haben, sind Gang�ffungsvertices (vgl. Kapitel \ref{KK_AndockenHoehle}).}.

Vorm Start des Algorithmus erfolgt der Aufbau der Adjazenzliste.
Dazu laufe �ber alle Dreiecke und f�ge f�r die drei Eckpunkte die Adjazenzen hinzu.
Es wird nur eine Umlaufrichtung betrachtet, damit keine doppelten Adjazenzen auftreten.
Aus Dreieck $\vv{P_1},\vv{P_2},\vv{P_3}$ ergeben sich als Adjazenzen:
f�r $\vv{P_1}$ ist $\vv{P_2}$ adjazent, f�r $\vv{P_2}$ ist $\vv{P_3}$ adjazent und f�r $\vv{P_3}$ ist $\vv{P_1}$ adjazent.

Der Algorithmus unterscheidet zwischen Originalvertices $V_i$ und Vertices $V'_i$, die dem reduzierten Netz angeh�ren.
Jeder $V_i$ erh�lt einen Verweis $Adr$ auf einen Vertex aus dem reduzierten Netz und eine Markierung $Flag$, die der Markierung von Nachbarschaften dient.
Jeder $V'_i$ erh�lt eine Markierung $Flag'$, ebenfalls f�r die Markierung von Nachbarschaften.
Alle $Adr$-, $Flag$- und $Flag'$-Werte sind anfangs mit $-1$ initialisiert.
Alle Vertices besitzen Indices $\geq0$, d.h. solange f�r einen Vertex $V_i$ $Adr(i)<0$ gilt, ist dieser Vertex noch nicht kontrahiert,
anderenfalls tr�gt $Adr$ die neue Adresse des kontrahierten Vertex.
$Flag$ und $Flag'$ sind durch diese Markierung anfangs keiner Nachbarschaft zugeordnet.

\ \\ 
Der Algorithmus muss topologie-tolerant sein, d.h. Grenzvertices und Gang�ffnungsvertices m�ssen erhalten bleiben.
Weiterhin d�rfen nur Originalvertices (also noch nicht verschmolzene Vertices) kontrahiert werden.
Somit l�sst sich folgender Test durchf�hren:
Addiere die X-Texturkoordinaten und Y-Texturkoordinaten sowie die $Adr$-Werte beider Vertices einer Kante.
Wenn die Summe $-6$ ergibt, sind alle diese Bedingungen erf�llt, bei $-5$ oder mehr nicht.
Weiterhin muss getestet werden, ob maximal zwei gemeinsame Nachbarn vorliegen, damit die Mannigfaltigkeit des Meshs erhalten bleibt.
Dazu werden erst f�r einen Vertex der Kante alle Nachbarn markiert (per $Flag$ bzw. $Flag'$).
Dann werden vom anderen Vertex aus alle entsprechenden Markierungen in der Nachbarschaft gez�hlt.
Diese Zahl ist die Anzahl der gemeinsamen Nachbarn.
Die Laufzeitkomplexit�t dieser Tests ist linear in der Anzahl der gepr�ften Vertices,
bei einem trivialen Test aller Nachbarschaftspaare w�re diese quadratisch.

Wenn auch dieser Test positiv verl�uft, so wird die Kante kontrahiert.
Als Position des neuen Vertex wird die Mitte der Kante verwendet.
Die Texturkoordinaten bleiben identisch zu denen der Ausgangsvertices und entsprechen somit $(-1,-1)$.
Nachdem alle m�glichen Kontraktionen durchgef�hrt wurden, werden alle nicht kontrahierten Vertices f�r den neuen Mesh �bernommen.
Danach werden alle diejenigen Dreiecke, die nicht zu Linien oder Punkten degeneriert wurden, ebenfalls �bernommen.

\ \\
Nach dem Durchlauf des Algorithmus m�ssen die Normalen berechnet werden.
Dazu wird das Verfahren aus \ref{KK_VoxelNormalen} verwendet.
Allerdings bleiben die Normalen der Grenzvertices unver�ndert, da diese f�r die Meshes aller Detailstufen gleich bleiben m�ssen.

Die Anzahl der Dreiecke der H�hlenmeshes liegen konstruktionsbedingt in $\mathcal O(n)$, mit $n$ als Anzahl der Vertices des Meshs.
Der Algorithmus besitzt inklusive des Aufbaus der Adjazenzliste
und Berechnung der Normalen somit eine Laufzeitkomplexit�t von $\mathcal O(g\cdot n)$, mit $g$ als maximalem Knotengrad je Vertex.
Der Faktor $g$ resultiert aus der Pr�fung auf mehrfache Nachbarn.
%
In Abbildung \ref{B_LODH1} und \ref{B_LODH2} sind Ergebnisse des Verfahrens gezeigt.
%In jeder Reduktionsstufe wird die Anzahl der Dreiecke und Vertices nahezu halbiert.
In der ersten Reduktionsstufe wird die Anzahl der Dreiecke und Vertices nahezu halbiert.
Der Reduktionsfaktor sinkt in den folgenden Reduktionsstufen langsam ab,
da die nicht kontrahierten Grenz- und Gang�ffnungsvertices einen immer h�heren Anteil an der Gesamtmenge der Vertices einnehmen.

\ \\
Zu den Problemen des Verfahrens geh�rt das verst�rkte Hervortreten der Meshgrenzen mit steigenden Reduktionsstufen,
da diese keiner Reduktion unterliegen.
Weiterhin kann es, falls keine Filterung schwebender Fragmente durchgef�hrt wurde, dazu kommen,
dass kleinere solcher Fragmente zu Tetraedern kollabieren.
Diese werden im n�chsten Reduktionsschritt zu zwei aufeinanderliegenden Dreiecken kontrahiert
und verschwinden im Schritt darauf vollst�ndig.
F�r diesen Sonderfall ist die Topologieerhaltung lokal nicht gew�hrleistet.
%
Ebenfalls kann es bei sehr vielen Reduktionsstufen zu einer starken Verformung des H�hlenmeshs kommen.
R�ume und G�nge, die an vielen Seiten von der H�hle umschlossen sind, k�nnen so m�glicherweise mit dieser kollidieren.
Im Allgemeinen wird die Anzahl der Reduktionsstufen jedoch nicht so hoch sein, dass dieser Fall auftritt.
Das gr��te Problem stellen Mesh Foldover dar.

\subsubsection{Mesh Foldover}

Mesh Foldover werden in \cite[S.22]{B_LoD3D} beschrieben und stellen das Ineinanderschieben benachbarter Dreiecke dar.
Dies kann bei den Reduktionsoperatoren Edge Collapse, Vertex-Pair Collapse und Triangle Collapse auftreten.
\cite{B_LoD3D} beschreibt die folgende Methode zur Verhinderung:
Es wird getestet, ob es Normalen mit mehr als $90^\circ$ �nderung durch den Collapse gibt.
Wenn ja, dann liegt ein potentieller Foldover vor und der Collapse wird nicht durchgef�hrt.
Allerdings ist dies kein generelles Kriterium.
\cite{A_SimplificationTokyo} berichtet von einem g�nstigen Grenzwert, der f�r diese Normalen�nderung gefunden werden muss.
Dort wird $\frac{\pi}{3}$ verwendet.
%!!!BILDSKIZZE Hefter  Do 19.04.2012 bei jedem Grenzwert m�glich

Bei relativ gleichm��igen Meshes mit eher lokal flachen Oberfl�chen funktioniert diese Methode des Foldover-Verhinderns recht gut.
Ein Foldover ist hier i.d.R. mit hohen Normalen�nderungen verbunden, da das entsprechende Dreieck "`umklappt"'.
Bei stark unebenen, zerkl�fteten Meshes kann aber eine kleine Normalen�nderung schon einen Foldover bedeuten.
Dies ist in Abbildung \ref{B_Foldover} dargestellt.

%Hier auch das Problem des H�hlenmeshes: das Verwackeln ist mittels Clamping gerade so abgeschw�cht,
%das hier gerade so keine Konflikte auftreten -> kleine �nderungen k�nnen schon Foldovers bedeuten.
%!!!Bild Meshfoldover prinzipiell immer m�glich, siehe Hefter 13.11.2011 ???

\ \\
Das Problem ist: Wenn der Grenzwert zu niedrig ist, werden kaum Reduktionsoperationen durchgef�hrt.
Wenn er zu hoch liegt, entstehen Mesh Foldover.
Im schlimmsten Fall k�nnen mit dieser Methodik keine Reduktionsoperationen durchgef�hrt werden, ohne dass es zu Mesh Foldovern kommt,
und zwar in dem Fall, wenn die Operation mit der niedrigsten Normalen�nderung schon einen Mesh Foldover bewirkt.

Eine sichere Methode diese Foldover zu verhindern, besteht darin, alle durch die Kontraktion betroffenen Dreiecke zu testen, ob sie danach andere Dreiecke
schneiden w�rden. Dies h�tte aber betr�chtlichen Rechenaufwand zur Folge.

\begin{figure}[hbtp]
  \centering  
	\includegraphics[width=14cm]{Bilder/MeshFoldover}
	\caption[Mesh Foldover]{\emph{Mesh Foldover}: v.l.n.r. \\
	 (a) Gezeigt ist die Ausgangstriangulation. Die zu kontrahierende Kante ist rot dargestellt.\\
	 (b) Nach dem Edge Collapse gibt es einen Konflikt zwischen dem roten und blauen Dreieck.
	 Die Normale des blauen Dreiecks wurde durch den Collapse nicht ver�ndert, die des roten Dreiecks deutlich unter $90^\circ$.
	 Durch horizontale Kontraktion der Skizze lassen sich F�lle mit Mesh Foldover bei beliebig kleiner Normalen�nderung erzeugen.}
	 \label{B_Foldover}
\end{figure}

%Mesh-Foldovers k�nnen auch bei Vertex Decimation-Algorithmen auftreten. Siehe mesh\_simplification\_slides.pdf
%(da Vertex Removal wie Edge Collapse betrachtbar, bei bestimmten Triangulationen)
\ \\
Durch solche Foldover kommt es oft zu starken Beleuchtungsunterschieden.
Bei der Verwendung von zweidimensionalen Texturen kommt es ebenfalls zu Texturunterschieden, bei dreidimensionalen Texturen nicht.
Foldover k�nnen bei bestimmten H�hlenfarben und -texturierungen (speziell: Eis) vom Betrachter als weitere kristalline Strukturen wahrgenommen werden.

Mesh Foldover entstehen insbesondere bei sehr stark zerkl�fteten Meshes, d.h. bei Verwendung von Erosion,
und nach mehreren Reduktionsstufen, da sich der Fehler aufsummiert.
Da die stark reduzierten Reduktionsstufen i.d.R. in gr��erer Entfernung angezeigt werden, ist das Problem dadurch etwas abgemindert.
In Abbildung \ref{B_LODH2}(b) sind solche Foldover an den lokalen Beleuchtungsunterschieden gut erkennbar.

%Dennoch w�re es w�nschenswert diese Foldover zu verhindern.
%M�glichkeit zur Verringerung der Foldovers bei erw�nschter Erosion: hohe Detailstufen mit Erosion, niedrige ohne Erosion erstellen.
%!!!"'Foldover"' mit nicht adjazenzten Dreiecken (siehe Hefter 13.11.2011)

% Weitere Probleme
%\subsubsection{Weitere Probleme}
%Enstehung weiterer degenerierter Dreiecke wie Linien, trotzdem mit unterscheidlichen Eckpunkten
%Einzelstehende Voxel k�nnen zu Tetraedern mutieren, diese dann zu zwei aufeinanderdliegenden Dreiecken durch den Kollaps einer
%Bei zunehmenden Reduktionsstufen treten die R�nder der Meshbuffer hervor, da deren Kanten nicht kollabiert werden.

%Als Sonderfall k�nnen durch einzelne bzgl. 18er-Nachbarschaft freistehende 0-Voxel Strukturen entstehen, die nach und nach zu Tetraedern kollabieren.
%Diese w�rden zu zwei aufeinanderliegenden Dreiecken mutieren, falls eine weitere Kante kontrahiert wird.
%Um dies zu verhindern, wird zus�tzlich ein Test durchgef�hrt, ob beide gemeinsamen Nachbarvertices der Kanteneckpunkte zueinander benachbart sind.

						% Level of Detail

\chapter{Programmstruktur}

Abbildung \ref{B_UML} zeigt den prinzipiellen Programmaufbau per Klassendiagramm nach UML 2.4.1 \cite{W_UMLInfra}\cite{W_UMLSuper}.
Abgebildet sind die wichtigsten Klassen mit ihren grunds�tzlichen Beziehungen zueinander.
Auf Angabe der Attribute und Operationen wurde aus Platzgr�nden verzichtet.
Der Aufbau des Programms gliedert sich grunds�tzlich in zwei Teile:
zum einen die Klassen f�r die Ansteuerung, Eingabe und Ausgabe,
zum anderen die Generatorklassen plus generierte Objekte.

\emph{CEventReceiver} empf�ngt Inputs und leitet diese an die anderen Objekte weiter.
\emph{CGUI} definiert die Anordnung der GUI-Elemente.
\emph{COptionen} steuert die einzelnen Klassen an, zus�tzlich werden hier sowohl Einstellung als auch Ergebnisse geladen und gespeichert.
In dieser Klasse ist zugleich der Pipeline-Ablauf definiert.
\emph{CSzene} stellt sowohl die GUI als auch die Dungeonszene dar.
Hierzu sind Methoden vorhanden, um die Dungeonszene aus den einzelnen Teilen zusammenzusetzen.
Die \emph{main}-Klasse enth�lt neben Initialisierungen eine Programmschleife zum Betrieb der Irrlicht-Engine.

\emph{CFraktalGenerator} erstellt L-Systeme und deren Ableitungen, welche die Turtle-Grafik-Zeichenanweisungen darstellen.
\emph{CVoxelRaum} verwaltet den Voxelraum, zeichnet die Turtle-Grafik-Zeichenanweisungen und f�hrt Erosion sowie Filterung schwebender Fragmente durch.
\emph{CDreiecksMesh} verwaltet die H�hlensubnetze, f�hrt die Umwandlung Voxel zu Dreiecksnetz durch, berechnet Sichtbarkeitsinformationen 
der H�hle und f�hrt die Reduktion der H�hlensubnetze durch.
\emph{CArchitekt} verwaltet Vorlagen f�r R�ume, Detailobjekte sowie die Grundkonfiguration der G�nge
und konstruiert aus einer gegebenen Voxelh�hle durch das Platzieren von G�ngen und R�umen den Dungeon.
Bei der Erstellung von Andockstellen an der H�hle wird auf Vertexkoordinaten-Berechnungsroutinen von CDreiecksMesh zur�ckgegriffen.
Die R�ume \emph{SDungeonRaum} tragen einen Verweis auf die zugrundeliegende Subzene \emph{CSubSzene}, welche die Geometrie des Raums enth�lt.
Jeder Gang ist eine Instanz von \emph{CSpline}.
F�r deren Erstellung werden von CArchitekt aus Gangparameter,
ermittelte Andockstellen sowie Informationen aus Detailobjektvorlagen \emph{SSplineDetailobjektVorlage} an CSpline �bergeben.
Hier erfolgt die Generierung der Ganggeometrie und Adaptergeometrien sowie die Platzierung der Detailobjekte als \emph{SSplineDetailobjekt}.
Sichtbarkeitsinformationen der G�nge und Gangreduktionsstufen werden ebenfalls in CSpline berechnet.

Es existieren zwei separate Zufallsgeneratoren vom Typ \emph{CZufallsGenerator}:
einer ist f�r CDreiecksMesh zust�ndig und wird dort f�r das Verwackeln der Eckpunkte eingesetzt,
der andere liefert Zufallswerte f�r die restlichen Klassen.


%welche Verfahren wurden implementiert?
%-> besonders wichtig, falls Alternativen genannt wurden, wie z.B. beim Zeichnen von Strichen im Voxerkaum
%Bibliotheken
%- Irrlicht
%Coding-Konventionen:
%- eigene
%- Irrlicht: http://irrlicht.sourceforge.net/faq.html
%UML, Pap
%VoxelRaum:
%- direkter Zugriff auf das Voxel-Array statt LeseVoxel() 
%- MinVoxelRandOffset == 3:
%-> keine Spezialf�llen an R�ndern
%--> deutliche Beschleunigung (ca. Faktor 4-10, je nach Prozedur)
%- Eigene Irrlichklassen erl�utern
%OBJ-Format erkl�ren:
%- Irrlicht-Exporter
%- eigener Exporter
%Wichtig: Beim Durchlauf durch Voxelraum: (bezieht sich auf Visual Studio 2010, vermutlich bei den meisten anderen Compiler gleiche Ablage im Speicher)
%- Schleifen-Reihenfolge von aussen nach innen m�glichst immer X,Y,Z -> effektiverer Speicherzugriff (nacheinander folgende Elemente stehen auch nacheinander im Speicher) -> merkbar schneller


\begin{figure}[hbtp]
  \centering   % trim: links unten rechts oben
	\includegraphics[width=15cm,trim = 10mm 6mm 6mm 10mm,clip]{Bilder/UML-DunGen}
	\caption[Programmaufbau - UML Klassendiagramm]{\emph{Programmaufbau - UML Klassendiagramm}:
	Dargestellt ist der prinzipielle Aufbau des Dungeongenerator-Programms.
	Auf die Darstellung von Klassen der Irrlicht-Engine wurde verzichtet.
	Der Inputteil ist \emph{rot} markiert, die eigentlichen Generatoren \emph{blau} und der Outputteil \emph{gr�n}.
	}
	\label{B_UML}
\end{figure}


\section{Anmerkung zum Rendering}

Die generierten Dreiecksnetze der H�hle sowie die der G�nge mit ihren Adaptern sind
bez�glich der Positionen ihrer Vertices bereits korrekt zueinander ausgerichtet.
Das bedeutet, dass sie untransformiert in den Szenegraphen eingebunden werden k�nnen.
Beim Rendering per Shader entfallen somit die Transformationen zur Berechnung der resultierenden Vertexpositionen und Normalen.
Es braucht nur die Projektionstransformation von 3D-Raumkoordinaten zu 2D-Bildschirmkoordinaten durchgef�hrt zu werden.
Durch das Weglassen der �brigen Transformationen ist eine h�here Performanz erzielbar.

%H�hlenmesh besitzt im allgemeinen viele Vertices und Dreiecke
%-> Falls im Shader Vertex-Positionen oder Vertex-Normalen benutzt werden:
%  - Mesh m�glichst untransformiert einbinden oder vorher transformieren (wenn statische Transformation)
%	-> keine Transformation im Shader n�tig -> spart Berechnungsaufwand
%- F�r Beleuchtung:
%  - im Normalfall reicht es aus die Lichtquelle zu transformieren, diese muss relativ zur H�hle korrekt positioniert und ausrichtet befinden
%  
%trifft auch f�r G�nge zu:
%alle vom Generator prozedural generierten Objekte brauchen keine Transformationen, sondern sind relativ zueinander korrekt positioniert und ausgerichtet
%
					% Programmstruktur
%\include{Kapitel9}									% Ergebnisse

%\include{Kapitel9}									%% Verwendung im Spiel
\chapter{Auswertung}

Dieses Kapitel besch�ftigt sich mit konkreten Laufzeittests,
Anmerkungen zu Parametern und weiteren Problemen.
Danach wird eine Zusammenfassung get�tigt und ein Ausblick
f�r fortsetzende Arbeiten gegeben.

\section{Untersuchungen zu Parametern und Laufzeiten}

Tabelle \ref{T_MessDungeon} zeigt einen �berblick �ber die durchgef�hrten Laufzeittests.
Die Ergebnisse wurden f�r jeden Dungeon �ber $500$ Messungen gemittelt.
Die Abbildungen der verwendeten Dungeons finden sich in Anhang \ref{A_Beispiele},
die zugeh�rigen XML-Konfigurationsdateien sind auf dem beiliegenden Datentr�ger mitgeliefert, siehe Anhang \ref{A_DVD}.

Die aufgef�hrten Dungeons ben�tigen jeweils weniger als eine halbe Minute zur Erstellung.
Tendenziell der gr��te Posten f�r die Laufzeit ist das Zeichnen der Voxelh�hle.
Die Laufzeit hierf�r ist stark abh�ngig von der Anzahl der Zeichenanweisungen, insbesondere Symbol $'F'$, sowie dem Strichradius.
Beispiel \ref{B_DungeonBsp1} hat sehr viele Zeichenanweisungen bei einem Strichradius von $14$,
und Beispiel \ref{B_DungeonBsp5} hat sehr wenige Zeichenanweisungen, aber einen Strichradius von $32$.
%
Die Erosion ist algorithmisch vergleichsweise invariant gegen�ber der Struktur des Voxelgebildes
und ben�tigt so unabh�ngig vom Dungeon eine ann�hernd konstante Laufzeit.
%
Filterung (schwebender Fragmente) 2 braucht �hnliche Zeiten f�r alle Dungeons.
Es sind hier i.A. nur wenige, kleine schwebende Fragmente zu entfernen.
Die vorhandenen Laufzeitunterschiede sind daher prim�r auf die unterschiedliche Ausdehnung der �u�eren H�lle zur�ckzuf�hren.
Bei Filterung 1 schwankt die Laufzeit st�rker.
Beispiel \ref{B_DungeonBsp5} besitzt ein sehr gro�es schwebendes Fragment im Inneren, daher wird signifikant mehr Rechenzeit ben�tigt.

Der Aufbau der Sektor- und Scankarten h�ngt insbesondere von der Anzahl der zu erstellenden Scankarten ab.
Beispiel \ref{B_DungeonBsp4} stellt einen sehr flachen Dungeon dar, wodurch weniger Scankarten erstellt werden m�ssen.
Dadurch ist die ben�tigte Laufzeit hier geringer.
%
Die Laufzeit der Generierung von G�ngen und R�umen ist nicht prinzipiell prim�r von der Anzahl der generierten G�nge abh�ngig.
Falls starke Restriktionen bez�glich der Andockpositionen f�r die G�nge an der H�hle
vorliegen, k�nnen die Andocktests den Prim�ranteil der Laufzeit bilden.
Beispiel \ref{B_DungeonBsp5} ist ein solcher Fall.
Hier wurde per Bodenabstandstest festgelegt, dass die G�nge nur im unteren Bereich an der H�hle andocken sollen.
Dadurch werden die meisten Andockpositionen verworfen.
%
Wie sich zeigt, ist der Blickdichtigkeitstest f�r die G�nge recht effektiv, im Schnitt wurden �ber 84 Prozent aller G�nge
als definitiv blickdicht erkannt.

Die Laufzeiten der Umwandlung und Reduktion der H�hlennetze skalieren in etwa linear der Anzahl der Vertices und Dreiecke.
Die restlichen Berechnungen nehmen einen nahezu konstanten Anteil an der Gesamtlaufzeit ein.

%- Bisektionsverfahren, welches auch bei den Detailobjekten

%Zufallsregel: F->(FF+-)* mit zuf�lligem Winkel != 90�,180�,270� (m�glichst kein Vielfaches von 10 -> unregelm��iger 
%(z.B. Primzahlen?, kein Teiler von 360, modulo 90 kein Teiler von 90)) sorgt i.d.R. f�r ebene H�hlen
%- zus�tzlich Drehungen um die anderen Achsen -> nicht-ebene H�hlen
%--> Tests: Regeln fest, dann Winkel variieren (z.B. mit Drachenkurven-Regeln aus Ausgangsregeln)
%
%Falls im Spiel keine Flugm�glichkeit/Leitern:
%- Abstand Gang zu Boden festlegen (nicht zu gross)
%- eher flache H�hlen konstruieren (Zeichen \$ hilft -> wieder horizontal ausrichten)



%!!!DATEINAMEN DER TESTS ANGEBEN, feste Seeds!
%- Einfluss von Parametern auf konkrete Ergebnisse und Laufzeiten
%!!!hier nur Parameter und Ergebnisse erkl�ren, gro�e Screenshots der H�hlen in Anhang
%- LOD H�hle austesten (welche Reduktion, Anzahl Vertices, Anzahl Dreiecke, Zeit)
%
%!!!H�hlen todo:
%- Zerghive-H�hle (Spirale)
%
%Laufzeit: 3 Beispiele (im Anhang zu sehen)
%Zeichnen,
%Erosion,
%Fragmente filtern (mit Erosion),
%Umwandlung (mit und ohne Erosion)
%Normalen berechnen
%Level of Detail
%
%- Einfluss von Parametern auf konkrete Ergebnisse und Laufzeiten
%!!!hier nur Parameter und Ergebnisse erkl�ren, gro�e Screenshots der Dungeons in Anhang
%
%- einzelne Perfomanz f�r verschiedene Programmteile testen
%  - eigentliche Gangerstellung
%  - Vistest Gang
%  - Kollisiontest Gang<->H�hle


\begin{table}
\centering
\begin{tabular}[htbp]{|l|r|r|r|r|r|}
  \hline
  Dungeon & \ref{B_DungeonBsp1} & \ref{B_DungeonBsp2} & \ref{B_DungeonBsp3} & \ref{B_DungeonBsp4} & \ref{B_DungeonBsp5} \\
  \hline\hline 
  \rowcolor{FarbeZeit}
  Erstellung Zeichenanweisungen																										& 39 & 32 & 27 & 26 & 28 \\
  Anzahl Zeichenanweisungen				& 826992 & 181989 & 39061 & 21845 & 1367 \\
  \rowcolor{FarbeZeit}
  Zeichnen																																				& 17219 & 1778 & 4239 & 3146 & 7164 \\
  \rowcolor{FarbeZeit}
  Erosion																																					& 1586 & 1585 & 1579 & 1576 & keine \\
  \rowcolor{FarbeZeit}
  Filterung 1																																			& 590 & 625 & 555 & 590 & 3300 \\
  \hline
  \rowcolor{FarbeZeit}
  Sektoren- und Scankartenaufbau																									& 1337 & 1317 & 1766 & 820 & 2465\\
  \rowcolor{FarbeZeit}
  Generierung R�ume und G�nge																											& 1799 & 1830 & 2994 & 1156 & 2285\\	
  \rowcolor{FarbeZeit}
  \ \ \ Generierung G�nge (inkl. LOD)																										& 908 & 1026 & 1795 & 739 & 293\\
  \rowcolor{FarbeZeit}
  \ \ \ Platzierung Detailobjekte																												& 23 & 24 & 42 & 18 & 5\\	  
  \rowcolor{FarbeZeit}
  \ \ \ Rest (R�ume, Andocken, u.a.)																										& 868 & 780 & 1157 & 399 & 1987\\	  
  Anzahl R�ume 												& 35 & 42 & 66 & 27 & 16\\
  Anzahl G�nge 												& 49 & 56 & 97 & 40 & 16\\
  \ \ \ davon definitiv blickdicht  	& 44 & 44 & 84 & 39 & 11\\
  \hline
  \rowcolor{FarbeZeit}
  Filterung 2 																																		& 588 & 625 & 556 & 589 & 733 \\
  \hline
  \rowcolor{FarbeZeit}
  Umwandlung Voxelh�hle																														& 1890 & 2285 & 1911 & 2151 & 5002 \\
  \rowcolor{FarbeZeit}
  \ \ \ davon Normalenberechnung		 																										& 286 & 464 & 316 & 434 & 824 \\
  Anzahl Vertices											& 400724& 683354 & 441594 & 621232 & 1315602\\
  Anzahl Dreiecke											& 822748& 1390848 & 892120 & 1274544 & 2596588\\
  Anzahl H�hlensubnetze 							& 13 & 36 & 36 & 28 & 56\\
  \hline
  \rowcolor{FarbeZeit}
 	Reduktion H�hlennetze Stufe 1																										& 458 & 813 & 586 & 732 & 1399 \\
  Anzahl Vertices											& 226202& 387130 & 253981 & 353852 & 730974\\
  Anzahl Dreiecke											& 473704& 798400 & 516894 & 739784 & 1427332\\
  \rowcolor{FarbeZeit}
 	Reduktion H�hlennetze Stufe 2																										& 314 & 568 & 434 & 515 & 901 \\
  Anzahl Vertices											& 129322& 221471 & 149665 & 206312 & 400090\\
  Anzahl Dreiecke											& 279944& 467082 & 308262 & 444704 & 765564\\
  \rowcolor{FarbeZeit}
  Reduktion H�hlennetze Stufe 3																										& 229 & 429 & 346 & 392 & 613 \\
  Anzahl Vertices											& 79594& 136054 & 95790 & 130910 & 225497\\
  Anzahl Dreiecke											& 180488& 296248 & 200512 & 293900 & 416378\\
  \rowcolor{FarbeZeit}
  Reduktion H�hlennetze Stufe 4																										& 189 & 358 & 302 & 331 & 462 \\
  Anzahl Vertices											& 54817& 93064 & 68611 & 93689 & 134969\\
  Anzahl Dreiecke											& 130934& 210268 & 146154 & 219458 & 235322\\
  \hline
  \hline
  \rowcolor{FarbeZeit}
  Restliche Berechnungen																													& 517 & 528 & 519 & 524 & 552 \\
   \rowcolor{FarbeZeit}
  Gesamt																																					& 26755 & 12773 & 15814 & 12548 & 24904 \\
  \hline
\end{tabular}
\caption{Laufzeiten und weitere Daten f�r die Generierung von Dungeons}
 {Laufzeiten sind gr�n unterlegt und jeweils in Millisekunden angegeben. \\
 Gemessen auf: \emph{Prozessor}: AMD Phenom II X4 945, 3 Ghz , \\ \emph{RAM}: 3,25 GiB Dual Channel
 DDR2-800 , \\ \emph{Betriebssystem}: Microsoft Windows 7 Professional 64 Bit mit Service Pack 1}
 \label{T_MessDungeon}
\end{table}

\ \\
Beim Entwurf von Dungeons ist es sinnvoll, zuerst die H�hle zu erstellen und zu verfeinern.
Erst dann, wenn die H�hle die gew�nschte Struktur aufweist,
sollten die Parameter f�r die Generierung der G�nge und R�ume entsprechend anpasst werden.

Bei der Konstruktion von H�hlen zeigt sich, dass L-Systeme der Form
$G=\langle \{\cdots\},F,$ $\{F \rightarrow \left(F,+,-,o,u,g,z\right)^*\} \rangle$
relativ oft zu realistisch anmutenden H�hlen f�hren (siehe hierzu Beispiel \ref{B_DungeonBsp3} und \ref{B_DungeonBsp4}).
Die Winkel sollten dabei m�glichst unregelm��ig gew�hlt werden, d.h. keine Vielfachen von $90^\circ$.
Bei Verwendung von nur $'+','-'$ als Symbole f�r Drehungen entstehen ebene H�hlen.
Bei zus�tzlicher Verwendung von $'o','u','g','z'$ ist das H�hlenskelett in alle Dimensionen ausgedehnt.
Symbol $'\$'$ kann dabei helfen, eine sonst sehr unebene H�hle ebener zu gestalten (vgl. Beispiel \ref{B_DungeonBsp1}).
Der Einsatz von Erosion ist nicht zwingend erforderlich, sondern von der Art der gew�nschten H�hle abh�ngig.
Gleiches gilt f�r die Filterung schwebender Fragmente.

Bei der Konstruktion des restlichen Dungeons sollte zun�chst die Wahl der Sektorgr��e abh�ngig von der H�hlenbeschaffenheit durchgef�hrt werden.
Falls z.B. mittlere Freir�ume ausgenutzt werden sollen (vgl. Beispiel \ref{B_DungeonBsp2}), muss die Sektorgr��e entsprechend klein gew�hlt werden.
Die Skalierung der Subszenen ist dabei entsprechend anzupassen, damit diese gut in die Sektoren passen.
Als Faktor $g$ f�r die St�rke der Gangsplineableitungen zeigen Werte zwischen $0,5$ und $1,5$ gute Resultate.
Parameter $R$ und $F$ f�r den Abstandstest zum Boden sollten nach Abh�ngigkeit von der H�hlenstruktur angepasst werden,
bei gr��eren Zeichenradien oder dem Anwenden von Erosion sind u.U. h�here Werte sinnvoll.

\section{Weitere Probleme}

\begin{figure}[hbtp]
  \centering  
	\includegraphics[width=14cm]{Bilder/Screenshot_Bug_Luecken}
	\caption[L�cken bei Raumadaptern]{\emph{L�cken bei Raumadaptern}: Hintergrundfarbe ist Wei�.
	Zu sehen ist ein schmaler wei�er Spalt zwischen Adapter und umschlie�ender Raumgeometrie.
	}
	\label{B_BugLuecken}
\end{figure}

%Dieses Problem betrifft die aktuell verwendeten 3D-Modelle der Subszenen:
\noindent An Adaptern mit Anschluss an R�umen kann es zu feinen L�cken zwischen Raum und Adapter kommen, siehe Abbildung \ref{B_BugLuecken}.
Diese L�cken sind bei kontrastreichem Hintergrund zu erkennen.
Blendet man die Raumgeometrie aus und l�scht den dem Adapter zugrundeliegenden Andockstellenmesh nicht,
so sind keine L�cken zwischen Adapter und Andockstellenmesh sichtbar.\footnote{Bei ausgeschaltetem Antialiasing.
Bei aktiviertem Antialiasing k�nnen dennoch Artefakte auftreten. Ursache:
Die f�r diese Arbeit erstellten Subszenen haben in ihren Materialeigenschaften das Flag "`Antialiasing"' nicht aktiviert,
da die verwendete IrrEdit-Version dies nicht unterst�tzt.
Daher wird das Antialiasing nicht auf sie angewendet, sondern nur auf den Adapter.
Dies f�hrt zu Artefakten an den gemeinsamen Kanten.}
Die L�cken bestehen also schon zwischen Andockstellenmesh und der umschlie�enden Subszenengeometrie.
Sie sind somit bereits in der importierten .irr-Szene enthalten.
Ungenauigkeiten beim Zusammenbau der jeweiligen Szenen im genutzten Szeneneditor Ambiera IrrEdit konnten durch Betrachtung der exportierten Daten ausgeschlossen werden.
Ursache sind demnach Ungenauigkeiten beim Modelling und Export der 3D-Daten aus dem verwendeten 3D-Modelling-Progamm Autodesk 3DS Max 2010.\footnote{Die
Andockstellenmeshes und die umschlie�ende Subszenengeometrie wurden gr��tenteils aus den gleichen Ausgangsobjekten mittels
unterschiedlich verkn�pften CSG-Operationen erstellt.}
%oder Ungenauigkeiten beim Zusammenbau und Export der Szene aus dem verwendeten Szeneneditor IrrEdit.
Eine m�gliche L�sung w�re die Schaffung eines Plugins
zur Modellierung sowie Export von Andockstellen und umgebender Geometrie f�r das jeweilig benutzte 3D-Modelling-Programm.
%die Nutzung eines eigenen Szeneneditors sowie 

%- die Ursache ist wahrscheinlich numerischer Natur
%- wenn man die Andockstellenmeshes nicht l�scht sind trotzdem feine L�cken sichtbar siehe Abbildung !!!(b)
%- eine Fehlerursache ist somit vermutlich der zu ungenaue Export der Meshes aus dem verwendeten 3D-Modelling-Programm 3DS Max
%  - Adapter und umschliessende Andockgeometrie (Wand) werden hier als separate Meshes exportiert  
%- dieser Fehler durch Transformationen weiter verst�rkt
%- insbesondere durch die Tatsache, dass die Transformation der Andockstelle schon auf der CPU berechnet wird,
%die Transformation des Raums aber ausschlie�lich auf der GPU beim Zeichnen

%Bild: v.l.n.r.
%(a) ScreenshotBugLuecken
%(b) ScreenshotBugLueckenADSeingeblended
%
%Fakten:
%
%- L�cken sind nicht sichtbar, wenn statt dem Raum die ADS-Meshes eingeblendet sind
%- leichtes Durchflimmern des Hintergrunds bei gleicher Variante und anschaltetem Antialiasing
% -> Ursache: Antialiasing wird nicht auf Meshes der Subszenen angewandt, diese werden per Irrlicht Standardmaterial ohne Effekte gerendert
% ** bei den hier erstellten Test-Subszenen!!
%
%- L�cken ansonsten durch ungenauen Export aus 3DS Max,
%-> Divergenz zwischen Andockstellenmesh und ungebender Wand
%** eventuell auch zu ungenaues Modelling?

%Code zum Test (Szene.cpp):
%ueberfluessigeGeoNode = SceneManager->getSceneNodeFromName("Nord_offen",tempSceneNode);
%if (ueberfluessigeGeoNode!=NULL) ueberfluessigeGeoNode->remove();
%ueberfluessigeGeoNode = SceneManager->getSceneNodeFromName("Nord_geschlossen",tempSceneNode);
%if (ueberfluessigeGeoNode!=NULL) ueberfluessigeGeoNode->remove();
%ueberfluessigeGeoNode = SceneManager->getSceneNodeFromName("Ost_offen",tempSceneNode);
%if (ueberfluessigeGeoNode!=NULL) ueberfluessigeGeoNode->remove();
%ueberfluessigeGeoNode = SceneManager->getSceneNodeFromName("Ost_geschlossen",tempSceneNode);
%if (ueberfluessigeGeoNode!=NULL) ueberfluessigeGeoNode->remove();
%ueberfluessigeGeoNode = SceneManager->getSceneNodeFromName("Sued_offen",tempSceneNode);
%if (ueberfluessigeGeoNode!=NULL) ueberfluessigeGeoNode->remove();
%ueberfluessigeGeoNode = SceneManager->getSceneNodeFromName("Sued_geschlossen",tempSceneNode);
%if (ueberfluessigeGeoNode!=NULL) ueberfluessigeGeoNode->remove();
%ueberfluessigeGeoNode = SceneManager->getSceneNodeFromName("West_offen",tempSceneNode);
%if (ueberfluessigeGeoNode!=NULL) ueberfluessigeGeoNode->remove();
%ueberfluessigeGeoNode = SceneManager->getSceneNodeFromName("West_geschlossen",tempSceneNode);
%if (ueberfluessigeGeoNode!=NULL) ueberfluessigeGeoNode->remove();
%ueberfluessigeGeoNode = SceneManager->getSceneNodeFromName("Basisraum",tempSceneNode);
%if (ueberfluessigeGeoNode!=NULL) ueberfluessigeGeoNode->remove();
%
%zus�tzlich: ADS-Mesh l�schen auskommentieren (SubSzene.cpp)

\section{Zusammenfassung und Bewertung}

% auch kritische Bewertung:
%Gitterstruktur weiteraufbrechen

Das entwickelte Verfahren zur Erstellung von Dungeons folgt dem nat�rlichen Entstehungsprozess:
Zuerst findet die Generierung der H�hle statt, danach erfolgt das Hinzuf�gen von G�ngen sowie R�umen.

Das Vorgehen f�r die H�hlenerstellung basiert auf Lindenmayer-Systemen.
Als Erstes erfolgt die Generierung von Turtle-Grafik-Zeichenanweisungen auf der Basis eines L-Systems.
Die L-Systeme zur Vorgabe der H�hlengrundstrukturen erm�glichen hierbei realistisch anmutende H�hlensysteme.
Dann findet das Zeichnen der Struktur im Voxelraum statt.
%Dieser erm�glicht es, sich selbst �berlagernde Strukturen problemlos umzusetzen.
Danach wird eine Nachbearbeitung per Erosion und dem Herausfiltern schwebender Fragmente durchgef�hrt,
um die entstandenen Strukturen realit�tsnaher zu gestalten.

Unter Ausnutzung der Voxel-Informationen werden R�ume um die H�hle herum platziert
und G�nge zwischen R�umen und der H�hle verlegt.
Die R�ume werden als Kopien vorgefertigter Szenen angelegt und erm�glichen so die Einbindung eigener Dungeonteile.
Die G�nge sind polygonale Schl�uche, die dem Verlauf kubischer Hermite-Splines folgen.
Hierdurch sind gewundene Gangstrukturen m�glich.

Die Voxelh�hle wird danach in ein Dreiecksnetz umgewandelt, welches aus mehren Subnetzen besteht.
Bei der Umwandlung werden die Koordinaten der Eckpunkte verwackelt.
Der Verwacklungsprozess kann zur Gl�ttung der entstehenden Netze eingesetzt werden.
Die entstehenden Strukturen ergeben zerkl�ftete H�hlen.
Diese lie�en sich in Verbindung mit prozeduralen Texturen realistischer gestalten.

F�r die Verwendung von Techniken des Culling und des Level of Detail werden Information zum Aufbau eines Sichtbarkeitsgraphen generiert.
Diese geben an, welche Dungeonbestandteile mit welchen anderen wie verkn�pft sind
und welche Sichtbarkeitseigenschaften die Dungeonteile besitzen.
Die Dreiecksnetze der G�nge und der H�hle werden f�r die Erstellung niedrigerer Detailstufen reduziert.
Die Reduktion der Gangnetze ist qualitativ hochwertig, die der H�hlennetze noch nicht.
Hier verhindern insbesondere Mesh Foldover bei zunehmenden Reduktionsstufen hohe Darstellungsqualit�t.

Die generierten Dungeons sind visuell ansprechend und lassen sich in Computerspielen verwenden.
Die H�hlen sind realen H�hlen �hnlich, wie ein Vergleich der Lavah�hle aus Abbildung \ref{B_Hoehlenbsp}(b)
und den erodierten H�hlen aus Abbildung \ref{B_DungeonInnenBsp1} zeigt.
Andere H�hlenarten lassen sich durch Modifizierung von L-Systemen, Erosion und Texturierung bilden.
Die erstellten G�nge sind geschwungen und f�r Dungeons sehr gut geeignet.
Die entworfenen und implementierten Algorithmen sind effizient in Laufzeit und Speicherbedarf.
Dungeons lassen sich i.d.R. in weniger als einer Minute komplett erstellen.

Die fertig erstellten Dungeons lassen sich exportieren.
Allerdings sind die exportierten Wavefront OBJ-Meshes sehr speicherplatzintensiv, da sie als Text und nicht bin�r gespeichert werden.
Damit ben�tigt auch der Gesamtdungeon viel Speicherplatz zur Ablage.
Als L�sung ergibt sich hier zum einen die Verwendung eines anderen Exporters (bin�res Format statt Text, z.B. 3DS\footnote{Allerdings besitzt Irrlicht keinen 3DS-Exporter.}).
Auch k�nnte der Dungeongenerator als Programmbibliothek verwendet werden.
Der gesamte Dungeon lie�e sich dann bei Bedarf aus einer XML-Konfigurationsdatei deterministisch generieren,
wodurch der ben�tigte Speicher zur Ablage der Dungeons extrem reduziert w�re.\footnote{I.d.R. werden statt �ber 100 MiB beim OBJ-Export nur noch $\sim$15 KiB ben�tigt.}
Die Generierungsgeschwindigkeit der Algorithmen reicht aus, um diese Methode praktikabel zu verwenden.

\section{Ausblick}

Die Parameter der L-Systeme, insbesondere die Winkel, k�nnte man beim Zeichnen zuf�llig variieren,
um weitere M�glichkeiten zur Vorgabe von H�hlenstrukturen schaffen.
Die Erosion im Voxelraum l�sst sich besser steuerbar und realistischer gestalten,
indem die Erosionswahrscheinlichkeit f�r jede Richtung (oben, unten, ...) einzeln festlegbar ist.

Eine bessere Aufteilung bei der Zerlegung der H�hle in Subnetze ist w�nschenswert.
Statt eines Octrees w�re ein Verfahren, welches die Struktur der H�hle besser einbezieht, geeigneter.
So w�ren weniger Grenzvertices n�tig und eine bessere Aufteilung f�r den Sichtbarkeitsgraphen m�glich.
Die Sektoraufteilung f�r die Raumplatzierung l�sst sich ebenfalls adaptiv gestalten.
Die Aufteilung in Sektoren w�rde nicht per Raumgitter erfolgen, sondern auch hier die Struktur der H�hle einbeziehen.

Die Heuristik zum Finden geeigneter Andockpositionen kann weiter verbessert werden.
Beispielsweise lie�e sich zus�tzlich der Bereich links und rechts der Andockposition betrachten, statt nur der Bereich in Richtung Boden.
Beim Andockvorgang von G�ngen an die H�hle l�sst sich der Querschnitt exakter ausfr�sen.
Statt der Verwendung eines quadratischen Andockbereiches k�nnte man hier den Gangquerschnitt per Voxeln abbilden.
Weiterhin sinnvoll ist das Zulassen beliebiger Andockrichtungen, nicht nur entlang der X- bzw. Z-Achse.

Prozedurale R�ume und prozedurale Detailobjekte sind ebenfalls denkbar.
Die R�ume k�nnten �hnlich dem Verfahren zur Generierung prozeduraler Geb�ude aus \cite{T_MasterJanusch} ebenfalls mit L-Systemen erstellt werden.
Auch eine "`beliebige"' Anzahl von Andockstellen f�r R�ume w�re m�glich.
Die Platzierung von Detailobjekten in H�hlen, beispielsweise Statuen an Gangm�ndungen oder Monster und Fallen auf relativ ebenen Fl�chen,
ist ein weiterer Fortf�hrungspunkt.
Daf�r lie�en sich die Informationen aus dem Voxelraum bzw. f�r Gangm�ndungen zus�tzlich die direkten Splineparameter nutzen.
Denkbar w�ren Parameterangaben wie: Objekt mit Ausdehnung $X\times Y \times Z$ soll $N$ Voxel �ber dem Boden platziert werden, mit einer Toleranz von $M$ Voxeln.

Statt nur die H�hlenvertices zu verwackeln, k�nnte auch der Querschnitt des Gangs optional verwackelt werden.
Auch das Ausprobieren und Vergleichen anderer Umwandlungstechniken von Voxeln in Dreicksnetze,
wie beispielsweise des Marching Cubes-Algorithmus \cite{lorensen1987marching}, ist denkbar.
Andere Zufallsgeneratoren, separat f�r jeden Teil des Dungeongenerators, k�nnten getestet werden.

Die Texturierung der H�hle durch prozedurale 3D-Texturen sollte weiter betrachtet werden.
Die Feinstruktur der H�hle l�sst sich beispielsweise durch prozedurale 3D-Displacement-Texturen verbessern.
Zus�tzliche Untersuchungen bez�glich Reduktionstechniken f�r Dreiecksnetze sind notwendig,
besonders um mit wenig Rechenaufwand Mesh Foldover verhindern zu k�nnen.

%- Streckung in XY-Ebene (bzw. Stauchung entlang Z-Achse): Hauptbewegungsebene der Spielfigur
%- danach Reduzierung des Netzes in gew�nschte Aufl�sung ? -> bei einstellbarer Abtastung nicht unbedingt n�tig (unperformanter), allerdings Ausprobieren sinnvoll	
%	- hierbei darauf achten, dass eventuelle Objekte immer noch korrekt platziert sind (nicht in der Wand verschwindend oder zu weit davor)
%- freiere Platzierung von R�umen (bessere Ausnutzung von L�cken der H�hle)
	

%Weitere Untersuchungen zu H�hlenstrukturen?

							% Auswertung 

\begin{appendix}
\chapter{Export}
\label{KA_Export}

Der Export des Dungeons erfolgt als (a) .irr-Datei f�r die komplette Dungeon-Szene,
(b) Wavefront OBJ-Dateien f�r die erstellten H�hlen- und Ganggeometrien sowie
(c) Zusatzinformationen zum Aufbau des Sichtbarkeitsgraphen als XML-Datei.
Zur weiteren Verwendung der exportierten Daten m�ssen sich alle im Folgenden genannten Dateien
sowie die in den genutzten Subszenen enthaltenen Texturen und Meshes im gleichen Ordner befinden
und auch so weitergegeben werden.

\ \\
Die Benennung der .irr-Szenendatei erfolgt als: \emph{"`[Dungeonname].irr"'}. \\
Der Aufbau des Szenegraphen der exportierten Dungeon-Szene wird im Folgenden dargestellt.
Die Benennung und ID's zum Wiederauffinden jedes Dungeonbestandteils sind
dabei geschrieben als \emph{(Benennung, ID)}.
Eine ID von -1 bedeutet, dass keine explizite ID vergeben wird, eine ID-Angabe von
0++ bedeutet aufsteigende ID's, beginnend mit 0. \\

\begin{compactitem}
	\item \emph{(Dungeon,-1)} $\Rightarrow$ kompletter Dungeon
	\begin{compactitem}
			\item \emph{(Raeume,-1)} $\Rightarrow$ alle R�ume
				\begin{compactitem}
						\item \emph{([gew�hlte Benennung],0++)} $\Rightarrow$ nacheinander alle R�ume mit ihrer Bezeichnung und ihren ID's
				\end{compactitem}
				
			\item \emph{(GangDetailstufe,0++)} $\Rightarrow$ alle vorhandenen Detailstufen der G�nge
				\begin{compactitem}
						\item \emph{(Gang,0++)} $\Rightarrow$ nacheinander alle G�nge mit ihren ID's
				\end{compactitem}
			\item \emph{(GangAdapter,-1)} $\Rightarrow$  alle vorhandenen Gangadapter
				\begin{compactitem}
						\item \emph{(Adapter0,0++)} $\Rightarrow$ nacheinander alle Adapter an Ende $\vv{P(0)}$ des Gangs mit der zum Gang geh�rigen ID
						\item \emph{(Adapter1,0++)} $\Rightarrow$ nacheinander alle Adapter an Ende $\vv{P(1)}$ des Gangs mit der zum Gang geh�rigen ID
				\end{compactitem}
				
			\item \emph{(Detailobjekte,-1)} $\Rightarrow$  alle vorhandenen Detailobjekte
				\begin{compactitem}
						\item \emph{(DetailobjektAnGang,0++)} $\Rightarrow$ nacheinander alle Detailobjekte des Gangs mit der zum Gang geh�rigen ID
						\begin{compactitem}
								\item \emph{([gew�hlte Benennung],0++)} $\Rightarrow$ nacheinander alle Detailobjekte mit ihrer Benennung, die ID's werden von $\vv{P(0)}$ bis $\vv{P(1)}$ aufsteigend gez�hlt
						\end{compactitem}
				\end{compactitem}
				
			\item \emph{(HoehlenDetailstufe,0++)} $\Rightarrow$ alle vorhandenen Detailstufen der H�hle
				\begin{compactitem}
						\item \emph{(Hoehlenteil,0++)} $\Rightarrow$ nacheinander alle Subnetze der H�hle mit ihren ID's
				\end{compactitem}
	\end{compactitem}	 
\end{compactitem}	 

\ \\
Die Benennung der exportierten H�hlengeometrien erfolgt als: \\
\emph{"`[Dungeonname]\_HoehlenMesh\_[SubnetzID]\_Detailstufe\_[DetailstufenNummer].obj"'}. \\
Die Benennung der exportierten Gangmeshes erfolgt als: \\
\emph{"`[Dungeonname]\_Gang\_[GangID]\_Detailstufe\_[DetailstufenNummer].obj"'}. \\
Die Benennung der exportierten Adaptermeshes erfolgt als:\\
\emph{"`[Dungeonname]\_Gang\_[GangID]\_Adapter\_[0 bzw. 1].obj"'}. \\
Zu den OBJ-Meshes zugeh�rig ist die ebenfalls erzeugte Datei \emph{"`DungeonBasisMaterial.mtl"'}, die ein Standardmaterial enth�lt.

\ \\
Die Benennung der Datei f�r die Zusatzinformationen erfolgt als: \\
\emph{"`[Dungeonname]\_Zusatzinformationen.xml"'}. \\
Der Aufbau der Zusatzinformationen ist Folgender:

\lstset{tabsize = 2, basicstyle=\footnotesize}%,  escapeinside={<!--}{-->)} } %,language=XML}
\begin{lstlisting}
<DungeonInformationen>
	<Raueme>
	  <!--Aufz�hlung aller R�ume:-->
		<Raum ID="[ID des Raums]">
			<GangIDs Gang0Nord="[ID des angeschlossenen Gangs, bzw. -1 f�r keinen]"
				Gang1Ost="[...]" Gang2Sued="[...]" Gang3West="[...]" />
		</Raum>
	</Raueme>
		
	<Gaenge GangBreite="[Breite der G�nge in Voxeln]">
		<!--Aufz�hlung aller G�nge:-->
		<Gang ID="[ID des Gangs]"
			DefinitivBlickdicht="[0 oder 1]">
			<RaumIDs RaumAndockstelle0="[ID des angeschlossenen
				Raums, bzw. -1 f�r H�hle]"	 RaumAndockstelle1="[...]" />
			<Position0 X="[...]" Y="[...]" Z="[...] />
			<Position1 X="[...]" Y="[...]" Z="[...]" />
			<Ableitung0 X="[...]" Y="[...]" Z="[...]" />
			<Ableitung1 X="[...]" Y="[...]" Z="[...]" />
		</Gang>
	</Gaenge>
	
	<Hoehle>
		<!--Aufz�hlung aller H�hlensubnetze:-->
		<HoehlenTeil ID="[ID des Subnetzes]">
			<VoxelMin X="[Min X des zugeh�rigen Voxelgebietes]"
				Y="[...]" Z="[...]" />
			<VoxelMax X="[...]" Y="[...]" Z="[...]" />
			<BlickdichtNegativeRichtung X="[kann nicht nach X-negativ
				gesehen werden? 0 oder 1]" Y="[...]" Z="[...]" />
			<BlickdichtPositiveRichtung X="[...]" Y="[...]" Z="[...]" />
		</HoehlenTeil>
	</Hoehle>	
</DungeonInformationen>	
\end{lstlisting}

%\subsubsection{Formate}
%- Export: Dateien als Archiv exportieren
%* 3D-Modelle (Levelst�cke)
%* eventuell Texturen (wenn keine Standardtexturen)
%* zus�tzliche Daten (Abstandsmatrix etc.)
%* bei Zusatzobjekte (Detail, Interaktion) speichern in welchem Gebiet
%* Szenendatei (Irrlicht)
%- Detailobjekte per ID-Nummer und Name im Szenegraph auffindbar -> 
%Verwendung als Interaktionsobjekte wie in \cite{T_BachelorHoenig}
%- XML-Beschreibung (DTD)

\subsubsection{Weitere Exportm�glichkeiten}

Die H�hle l�sst sich auch einzeln in einer OBJ-Datei mit frei w�hlbarem Dateinamen exportieren.
Alle H�hlensubmeshes werden dabei in diese Datei geschrieben.
Die Datei \emph{"`DungeonBasisMaterial.mtl"'} wird hier ebenso erzeugt und ist zur OBJ-Datei zugeh�rig.

\ \\
Die aus L-Systemen erzeugten Turtle-Grafik-Zeichenanweisungen lassen sich ebenfalls exportieren.
Als Dateiformat wird XML verwendet.
Die zum Zeichnen relevanten Parameter werden ebenfalls gespeichert.
Der Dateiaufbau gestaltet sich wie folgt:

\begin{lstlisting}
<LSystemAbleitungen>
	<LSystemParameter>
		<Winkel WinkelGier="[...]" WinkelNick="[...]" WinkelRoll="[...]" />
		<Radius StartRadius="[...]" RadiusFaktor="[...]"
			RadiusDekrementor="[...]" />
	</LSystemParameter>
	
	<!--alle generierten Ableitungen, beginnend mit Iteration 0:-->
	<Iteration Nummer="[Iterationsstufe]"
		String="[Generierter String]" />		
</LSystemAbleitungen>		
\end{lstlisting}									% Export
\chapter{Weitere Beispiele f�r Dungeons und H�hlen}
\label{A_Beispiele}

Im Folgenden finden sich einige Beispiele f�r Dungeons, die mittels des Dungeongenerator-Programms generiert wurden.
Links ist jeweils die H�hle abgebildet, rechts der fertige Dungeon.
Zu jeder H�hle ist das zugrundeliegende L-System angegeben.
Anschlie�end finden sich einige Screenshots aus dem Inneren von Dungeons.

\begin{figure}[htbp]
  \centering  
  \hfill
  \subfloat{\includegraphics[width=7.4cm]{Bilder/Screenshot_Bsp_Aufsteigend_Hoehle}}
  \hfill
  \subfloat{\includegraphics[width=7.4cm]{Bilder/Screenshot_Bsp_Aufsteigend_Dungeon}}  
  \hfill
	\caption[Dungeon Beispiel 1]{\emph{Dungeon Beispiel 1}:
	$G=\langle \{F,X,Y,+,-,o,\$\},YYFYF,$ \\ $
	\{
	F\rightarrow F-YX-X---,
	X\rightarrow F\$F++F-X,
	Y\rightarrow oYX--XX++
	\} \rangle$ \\
	mit $\alpha_{Gier} = 250^\circ$, $\alpha_{Nick} = 1^\circ$, Startradius $=14$, 8. Iteration
	}
	\label{B_DungeonBsp1}
\end{figure}

\begin{figure}[htbp]
  \centering  
  \hfill
  \subfloat{\includegraphics[width=7.4cm]{Bilder/Screenshot_Bsp_Ring_Hoehle}}
  \hfill
  \subfloat{\includegraphics[width=7.4cm]{Bilder/Screenshot_Bsp_Ring_Dungeon}}  
  \hfill
  \caption[Dungeon Beispiel 2]{\emph{Dungeon Beispiel 2}:
	$G=\langle \{F,+,o,u,g,z,[,],!\},F+[!+uFoF]gFz+[!+uFozzzzF]ggFzz+[!+uF]gFz+F+[!+oFuF]zF,$ \\
	$\{F\rightarrow F-F+FF+F-F\} \rangle$ \\
	mit $\alpha_{Gier} = 60^\circ$, $\alpha_{Nick} = 45^\circ$, $\alpha_{Roll} = 20^\circ$, $r_f = 1$, $r_d = 4$, \\ Startradius $=16$, 5. Iteration
	}
	\label{B_DungeonBsp2}
\end{figure}


\begin{figure}[htbp]
  \centering  
  \hfill
  \subfloat{\includegraphics[width=7.2cm]{Bilder/Screenshot_Bsp_Tief_Hoehle}}
  \hfill
  \subfloat{\includegraphics[width=7.2cm]{Bilder/Screenshot_Bsp_Tief_Dungeon}}
  \hfill
	\caption[Dungeon Beispiel 3]{\emph{Dungeon Beispiel 3}:
	$G=\langle \{F,+,-,o,u,g,z\},F,\{	F\rightarrow F+FoFg-FuzF\} \rangle$
	mit $\alpha_{Gier} = 285^\circ$, $\alpha_{Nick} = 255^\circ$, $\alpha_{Roll} = 255^\circ$, Startradius $=13$, 6. Iteration
	}
	\label{B_DungeonBsp3}
\end{figure}


\begin{figure}[htbp]
  \centering  
  \hfill
  \subfloat{\includegraphics[width=7.2cm]{Bilder/Screenshot_Bsp_Weit_Hoehle}}
  \hfill
  \subfloat{\includegraphics[width=7.2cm]{Bilder/Screenshot_Bsp_Weit_Dungeon}}  
  \hfill
	\caption[Dungeon Beispiel 4]{\emph{Dungeon Beispiel 4}:
	$G=\langle \{F,+\},F,\{F\rightarrow F+FFF\} \rangle$ \\
	mit $\alpha_{Gier} = 68^\circ$, Startradius $=16$, 7. Iteration
	}
	\label{B_DungeonBsp4}
\end{figure}

\begin{figure}[htbp]
  \centering  
  \hfill
  \subfloat{\includegraphics[width=7.2cm]{Bilder/Screenshot_Bsp_Doppelkugel_Hoehle}}
  \hfill
  \subfloat{\includegraphics[width=7.2cm]{Bilder/Screenshot_Bsp_Doppelkugel_Dungeon}}  
  \hfill
	\caption[Dungeon Beispiel 5]{\emph{Dungeon Beispiel 5}:
	$G=\langle \{F,W,X,Y,+,o,[,]\},W,$ \\ $
	\{
	W\rightarrow W+[X],
	X\rightarrow YYYYYYYYY,
	Y\rightarrow FoFoFoFo
	\} \rangle$ \\
	mit $\alpha_{Gier} = 10^\circ$, $\alpha_{Nick} = 10^\circ$, Startradius $=32$, 20. Iteration
	}
	\label{B_DungeonBsp5}
\end{figure}

\begin{figure}[htbp]
  \centering  
  \subfloat{\includegraphics[width=14cm]{Bilder/Screenshot_BspShot_GaengeErosion}}\\    
  \subfloat{\includegraphics[width=14cm]{Bilder/Screenshot_BspShot_GaengeErosion2b}}  
	\caption[Dungeons von innen]{\emph{Dungeons von innen}: beide H�hlen unter Einfluss von Erosion, Umwandlung per Verwackeln und Gl�tten}
	 \label{B_DungeonInnenBsp1}
\end{figure}


\begin{figure}[htbp]
  \centering  
  \subfloat{\includegraphics[width=14cm]{Bilder/Screenshot_BspShot_GaengeErosion3}}\\    
  \subfloat{\includegraphics[width=14cm]{Bilder/Screenshot_BspShot_DungeonBizzar}}  
	\caption[Dungeons von innen ff.]{\emph{Dungeons von innen ff.} v.o.n.u. (a) H�hle unter Einfluss von Erosion, (b) H�hle ohne Einfluss von Erosion,
	beide: Verwackeln und Gl�tten}
\label{B_DungeonInnenBsp2}
\end{figure} 

								% Beispiele f�r Dungeons und H�hlen
%\chapter{Weitere Screenshots}

\begin{figure}[htbp]
  \centering  
  \subfloat{\includegraphics[width=14cm]{Bilder/Screenshot_BspShot_GaengeErosion}}\\    
  \subfloat{\includegraphics[width=14cm]{Bilder/Screenshot_BspShot_GaengeErosion2}}  
	\caption[Screenshot Dungeon 1]{Dungeons von innen}
	 \label{B_VerwackelnGlaettenShots}
\end{figure}


\begin{figure}[htbp]
  \centering  
  \subfloat{\includegraphics[width=14cm]{Bilder/Screenshot_BspShot_GaengeErosion3}}\\    
  \subfloat{\includegraphics[width=14cm]{Bilder/Screenshot_BspShot_DungeonBizzar}}  
	\caption[Screenshot Dungeon 1]{Dungeons von innen ff.}
	 \label{B_VerwackelnGlaettenShots}
\end{figure} 

							% weitere Screenshots
%\chapter{Bedienungsanleitung des Dungeongenerators}

- Aufl�sung und Antialiasing per Startparameter f�r Dungeongenerator.exe festlegbar

Startoptionen:
- Aufl�sung des Dungeongeneratorfensters: Breite H�he
- Antialiasfaktor, der von der Grafikkarte gefordert wird, 0 .. aus
\footnote{Antialiasing wird nicht auf alle Objekte der Szene angewendet, da die verwendeten Materialien dies teilweise nicht zulassen.}

Syntax:
"`Dungeongenerator.exe Breite H�he Antialiasfaktor"'
Vorlage: Datei "`DunGen.bat"' kann modifiziert werden

%================================================================================================
\subsubsection{Allgemeine Hinweise}

Alle Zahlen werden mit Punkt statt Komma angegeben

Objekt-Benennung: Symbole f�r Fraktale, Dateinamen etc.: nach M�glichkeit keine Sonderzeichen, da Daten als XML
gespeichert werden

Angabe von Dateinamen:
- relative Pfade verwenden: in Dateiauswahl-Dialogfenster den absoluten Teil l�schen, bevor OK geklickt wird
- Gro�/Kleinschreibung: in Dateiauswahl-Dialogfenster Schreibweise �ndern, bevor OK geklickt wird

%================================================================================================
\subsubsection{Tastenbelegung}

F�r die Verwendung der folgenden Tastaturbefehle darf kein GUI-Fenster fokussiert sein.
Fokussierte Fenster sind an der blauen Titelleiste erkennbar.
Falls ein Tastendruck nicht funktionieren sollte, in die Szene klicken um die GUI zu defokussieren.

F1 ... H�hlenmeshes: h�here Detailstufe
F2 ... H�hlenmeshes: niedrigere Detailstufe
F3 ... Gangmeshes: h�here Detailstufe
F4 ... Gangmeshes: niedrigere Detailstufe

F5 ... Kamera: Freiflugmodus
hier Steuerung der Kamera mittels Pfeiltasten und Maus
F6 ... Kamera: feste Position und Ausrichtung
F7 ... Kamera: schnellere Geschwindigkeit f�r Freiflugmodus
F7 ... Kamera: langsamere Geschwindigkeit f�r Freiflugmodus

F9 ... Kamera: alle Werte auf Original
F11 ... Kamera: Position und Ausrichtung merken
F12 ... Kamera: Position und Ausrichtung wiederherstellen

1 ... Drahtgittermodell f�r H�hlenmeshes an/aus
2 ... Drahtgittermodell f�r Gangmeshes an/aus
3 ... Drahtgittermodell f�r Adaptermeshes an/aus
4 ... Ein-/Ausblenden von Detailobjekten in G�ngen
5 ... Ein-/Ausblenden von R�umen

6 ... schalte definitiv blickdichte G�nge und zugeh�rige Adapter auf Drahtgittermodus, andere auf normale Anzeige
7 ... schalte definitiv blickdichte G�nge und zugeh�rige Adapter auf normale Anzeige, andere auf Drahtgittermodus

8 ... fertige Screenshot mit GUI an, im Verzeichnis des Programms unter "`Screenshot\_JJJJ-MM-DD\_HH-MM-SS.bmp"' zu finden
9 ... fertige Screenshot ohne GUI an, gleiches Verzeichnis und gleiche Benennung wie oben

0 ... Wechsle H�hlenmaterial: Innen-/Aussenansicht wird vertauscht


%================================================================================================
\subsubsection{Men�: Dateizugriff - Einstellungen}

- laden und speichern der einzelnen einzelnen Einstellungen als .xml-Dateien
- �berall dort wo ein H�kchen gesetzt wird, diese Einstellungen werden geladen bzw. gespeichert
- beim Laden werden nur die Einstellungen wirklich geladen, die auch in der .xml-Datei enthalten sind

- die Datei "`Einstellungen\_Standard.xml"' wird automatisch beim Start des Programms geladen

Fraktalparameter:
- Untermen� "`FraktalEditor"'
- alle Einstellungen aus denen sich ein Lindenmayer-System generieren l�sst
- Winkel, Startradius, Radiusfaktor, Radiusdekrementor, Startfolge, Ersetzungsregeln, Iterationsstufe
- f�r ein gegebenes Fraktal reicht es also aus diese Einstellungen zu speichern bzw. zu laden

Weitere FraktalEditor-Einstellungen
- Untermen� "`FraktalEditor"'
- alle �brigen Einstellungen dieses Untermen�s, wie z.B. die gew�nschten Minimum und Maximumwerte und Anzahl der Nachkommastellen

VoxelRaum:
- Untermen� "`VoxelRaum"'
- alle Einstellungen

G�nge:
- Untermen� "`G�nge"'
- alle Einstellungen der linken Seite, also alles was nur die G�nge und nicht die Detailobjekte betrifft

Detailobjekte:
- Untermen� "`G�nge"'
- alle Einstellungen der rechten Seite, also alles was nur die Detailobjekte betrifft

Architekt:
- Untermen� "`Architekt"'
- alle Einstellungen der linken Seite, also alles was nur den Architekt und nicht die Subszenen betrifft

Subszenen:
- Untermen� "`Architekt"'
- alle Einstellungen der rechten Seite, also alles was nur die Subszenen betrifft

Erstellung:
- Untermen� "`Erstellung"'
- alle Einstellungen

Weiteres:
- Untermen� "`Weiteres"'
- alle Einstellungen

%================================================================================================
\subsubsection{Men�: Dateizugriff - Ergebnisse}

- laden und speichern der Ergebnisse

Fraktal-Iterationen Zeichenanweisungen:
- Speichern der generierten Turtlegrafik-Zeichenanweisungen (s. Anhang \ref{KA_Export})

Dichtekarten:
- speichern der generierten Dichtekarten
- Format wird durch Dateiendung festgelegt: .jpg, .png, .tga, .bmp, .pcx, .ppm

Voxelh�hle Dreiecksmesh:
- speichert den Dreickesmesh der gesamten H�hle im .obj-Format
%- jedes Subnetz wird in einem extra Meshbuffer abgelegt

Dungeon:
- speichert die Szene, Zusatzinformationen und Meshes des Dungeons
(s. Anhang \ref{KA_Export})

Test auf korrekten Export:
- l�dt .irr-Szene ein
- wenn es sich um einen exportierten Dungeon handelt, so wird der Szenegraph neu traversiert
  - so lassen sich Detailstufen und Drahtgitteransichten umschalten, Deatailobjekte und R�ume ein- und ausblenden
  - die Zusatzinformationen werden nicht eingelesen, so lassen sich blickdichte G�nge nicht markieren

%================================================================================================
\subsubsection{Men�: FraktalEditor}

Parameter :
- Wert ... konrekter aktueller Wert
- Zuf�lliges Bestimmen:: Minimum, Maximum, Nachkommastellen werden beim Ausw�rfeln per Button "`Zufallsparameter"' verwendet
%\footnote{!!!Erkl�rung Berechnungsprinzip. Es kann hierbei zu Rundungsfehlern kommen.
%Beispiel f�r Visual Studio 2010: mit "`int x = 18467;"' ergibt "`(double)x * 0.01"' den Wert "`184.669998"'}
% FIXED mit Service Pack 1!
als Nachkommastellen sind auch negative Werte zul�ssig, diese reduzieren dann die Zahl der signifikanten Stellen vor dem Komma

- Startstring: mit welchem Symbolen soll begonnen werden
- Zuf�lliges Bestimmen �ber "`Zuf�lliger Startstring"'
  - Erlaubte Symbole werden verwendet, je �fter ein Symbol erhalten ist, desto wahrscheinlicher ist sein Vorkommen
  - L�nge des Strings liegt zwischen Minimum und Maximum\footnote{Ausser bei Folgen von nur '\$' oder nur '\|'}
  - sinnfreie Kombinationen im String wie '+-', '[]' oder '\$\$' werden automatisch verhindert    
  
- Aktuelle Regeln zeigen Ersetzungsregeln an
- Ersetzungsinfo gibt die Ersetzung f�r aktuell ausgew�hlte Regel komplett wieder
  Dies ist notwendig, falls die Regel zu lang f�r die gro�e Anzeigebox ist.

- neue Regel generieren:
  - Symbol und Ersetzung angeben und auf "`Regel hinzuf�gen"' klicken
  - Zuf�llige Ersetzung f�r angegebenes Symbol �ber "`Zuf�llige Regel hinzuf�gen"': s.o. bei zuf�lliger Startstring

"`Regeln l�schen"' l�scht alle Regeln

Gew�nschte Iterationstufe:
- bis zu welcher Iterationstufe (0...Startstring) soll maximal generiert werden
- die Generierung bricht ab, wenn mehr als 1.000.000 Zeichenanweisungen vorliegen w�rden
- dann ist die erreichte Iterationsstufe u.U. niedriger, der Eintrag wird enstprechend aktualisiert

Generierung der Turtle-Grafik-Zeichenanweisungen nach Parametervorgabe per Button "`Generiere FraktalIterationen"'
Erzeugung zuf�lliger Parameter nach Vorgaben plus Generierung der Turtle-Grafik-Zeichenanweisungen  per Button "`Generiere Zufallsparameter und Zeichenanweisungen"'
- hierbei werden die erlaubten Symbole f�r Startstring und Ersetzung nach Zeichen durchsucht, die keine Zeichenanweisung nach Tabelle !!!??? darstellen.
- f�r jedes dieser Zeichen und 'F' wird eine Erstezungsregel erstellt


%Anmerkung: Die Turtle hat das Intialkoordinatensystem $\vv{Vorn}=(1,0,0)$, $\vv{Links}=(0,0,1)$, $\vv{Oben}=(0,1,0)$


%================================================================================================
\subsubsection{Men�: VoxelRaum}

linke Seite: Zeichnen und Nachbearbeiten im Voxelraum
rechte Seite: Umwandlung Voxel zu Dreiecksmesh

%---

Voxelrand: gibt an welcher Rand beim Zeichnen und Erodieren mit Belegung 0 bleiben soll
empfohlene Einstellung: Minimaler VoxelRand (= 3) + Anzahl Erosionschritte + zus�tzliche Ausfr�stiefe (s. Men�: Architekt)
-> so k�nnen die Erosion und die Ausfr�stiefe voll greifen, und brauchen nicht per Algorithmus beschr�nkt zu werden

Minimaler Strichradius: Minimalradius beim Zeichnen der Turtle-Grafik

Erosionszufallsseed: Seed f�r Erosion, beim gleichen Seed und gleichem Ausgangs-Voxelgebilde gleiche Ergebnisse
Erosion: Vorgebenen Seed nutzen: an: obiger Seed wird verwenden, aus: zuf�lliger Seed wird verwendet

Erosionswahrscheinlichkeit: Wahrscheinlichkeit $p_e$ mit der erodiert wird

Anzeigebox Aktuelle Erosionschritte: zeigt die geplanten Erosionschritte in ihrer Reihenfolge an
"`F�ge Erosionschritt hinzu"' f�gt neuen Schritt mit angegebener Erosionswahrscheinlichkeit hinzu
"`Erosionschritte l�schen"' l�scht alle Erosionschritte

Schwebende Fragmente entfernen: wenn H�kchen gesetzt wird nach der Erosion die Filterung schwebender Fragmente durchgef�hrt

"`Berechne Voxel"' zeichnet die Turtle-Grafik im Voxelraum und f�hrt danach die Erosion und ggf. Filterung schwebender Fragmente durch

"`Dichte berechnen"' berechnet die Dichte von Voxeln ungleich $0$ bzgl. der jeweiligen Ebene und zeigt diese als Bild an
die Dichte wird dabei normalisiert: weiss .. h�chste Dichte, schwarz ... Dichte ist $0$ (keine entsprechenden Voxel in diesem Bereich)

%---

Verwackle Vertices: gibt an, ob die Vertices beim Umwandeln verwackelt werden sollen
Verwacklungsst�rke: Verwacklungst�rke $K$ f�r das Verwackeln

Zufallsgeratorparameter: a,c,m, sowie der Seed

Gl�tte Mesh beim Verwackeln: gibt an ob Gl�ttungsmethode angewendet werden soll

Normalen-Berechung Wichtungsmethode: gibt an, nach welcher Methode die Vertexnormalen berechnet werden sollen

Anzahl zus�tzlicher Detailstufen: wieviele reduzierte Versionen sollen erstellt werden (0 ... nur Original)

"`Berechne Dreiecksnetz aus Voxeln"' f�hrt die Umwandlung mit den eingestellten Parametern durch
"`Zeige berechnetes Dreiecksnetz"' zeigt das Dreiecksnetz der H�hle an


%================================================================================================
\subsubsection{Men�: G�nge}

linke Seite: Einstellungen f�r G�nge
rechte Seite: Detailobjekte f�r G�nge

G�nge, Detailobjekte: 
- lokale 2D-Position in [-1,1]x[-1,1]  angeben
Resultierende Position wird durch Multiplizierung mit Gangbreite/2 bestimmt
Gangbreite $B$ wird durch das entsprechende Feld festlegt

%---

Gangquerschnitt:
- Punkte m�ssen bez�glich ihres Polarwinkels zirkul�r aufsteigend im Uhrzeigersinn angeordnet sein

"`Punkt hinzuf�gen"' f�gt neuen Punkt mit aktuellen Werten (lokale Position, Tetxurekoord. X) zum Querschnitt hinzu (an letzter Position)
"`Punkt l�schen"' l�scht markierten Punkt
"`�nderung �bernehmen"' �bernimmt die aktuellen Werte f�r den angew�hlten Punkt
%TIPP: durch Markieren eines Punktes und Klicken von "`Punkt hinzuf�gen"' wird eine Kopie des Punktes eingef�gt

Gangsegmentabstand: ist der Gangsegmentabstand $w$
Texturkoordinate Y Erh�hung pro Abstand 1: gibt das $TY_{IncPro1}$ an

Anzahl zus�tzlicher Detailstufen: wieviele reduzierte Versionen des Gangs sollen erstellt werden (0 ... nur Original)

Gangtextur: gibt die Textur f�r Gang und Adapter an, kann direkt eingetragen werden, oder per "`Texturdatei ausw�hlen"' angew�hlt werden

%---

Anmerkung: Die Reihenfolge der Detailobjekte hat keine Relevanz
"`Detailobjekt hinzuf�gen"' f�gt ein Detailobjekt mit den aktuellen Einstellungen hinzu,
 dazu muss eine Mesh-Datei ausgew�hlt werden (in einem von Irrlicht unterst�tzten Format, siehe \cite{W_Irrlicht}, z.B. .3ds-Format)
"`Detailobjekt l�schen"' l�scht das aktuell markierte Detailobjekt
"`�nderungen �bernehmen"' �bernimmt die aktuellen Werte f�r das angew�hlte Detailobjekt

Benennung: Benennung des Detailobjekts im Szenegraph
Skalierung: gew�nschte Skalierung
Rotation: Initialrotation, wenn der Gang in Richtung $(0,0,1)$ verlaufen w�rde

Platzierung von Detailojekten:
	
- Abtastung: �quivalent Gangsegmentabstand des Gangs $w$ aus !!!Kapitel
- beide Werte sollten nach M�glichkeit identisch gew�hlt werden, f�r exakte Platzierung
- Faktor Num Abstand: gibt an, in welchen Vielfachen der Abstastung Objekte Platziert werden sollen
- dabei Faktor f�r Num Abstand = $2^n$ f�r exakte Platzierung f�r $n$ zus�tzliche Detailstufen des Gangs

Num Abstand Min/ Num Abstand Max:
Platzierung eins Objektes alle $n$ Segmente mit
$n = (Faktor Num Abstand) * Zufallswert[Num Abstand Min, Num Abstand Max]$
wobei $n$ nach Platzierung eines Objektes neu bestimmt wird

f�r das erste Element gibt es separate Werte, um Objekte direkt am Gangeingang platzieren zu k�nnen
Objekt am Ende des Gangs definitiv platzieren: ein Objekt wird bei $P(1)$ platziert, egal ob der Abstand passen w�rde oder nicht

Pfadinfo: gibt den Dateipfad des Meshes des aktuell ausgew�hlten Detailobjektes wieder

%================================================================================================
\subsubsection{Men�: Architekt}

links: allgemeine Architekt-Einstellungen
rechts: Einstellungen f�r Subszenen

%---

Konstruktions-Zufallsseed: Seed f�r Konstruktion des Dungeons, also Platzieren der R�ume und Verlegen der G�nge
 bei gleichem Seed, gleicher H�hle und gleichen Einstellungen gleiches Ergebnis
Konstruktion: Vorgebenen Seed nutzen: an: obiger Seed wird verwenden, aus: zuf�lliger Seed wird verwendet

Raumsektorausdehnung: gibt die Ausdehnung $A$ der Sektoren in Voxeln an
Minimalabstand Raum <-> H�hle: welchen Abstand soll die Andockposition vom andockenden Sektor mindestens haben
- sinnvolle Werte liegen zwischen $0$ und $A$, z.B $A/2$

Andockabstand zum Boden: gibt maximaler Abstand zum Boden $R$ an (0 ... kein Test diesbez�glich)
Tiefe f�r Abstandstest: gibt gew�nschte Minimaltiefe f�r den Test $T$ an

Zus�tzlicher Sicherheitsabstand: wird als Faktor angegeben $1.0$ entsprechen 100 Prozent,
 gibt an als mit welcher Ausdehnung der Gang in den Kollisionstest zwecks Andocken einbezogen werden soll
 - Nutzen: damit Detailobjekte, bei denen Teile aus dem Gang herausragen nicht mit der H�hle kollisieren
 - wenn alles im Bereich [-1,1]x[-1,1] keine Zusatzausdehung ben�tigt
 - kann im Normalfall auf $1.0$ gelassen werden

Wahrscheinlichkeit f�r zus�tzlichen Gang:
gibt die Wahrscheinlichkeiten $p_{gh}$ und $p_{gr}$ an

Platzierung von Subszenen:
- jede Subszene (auch die leere Subszene ... nichts wird platziert) hat einen Platzierungswahrscheinlichkeitsfaktor $c$
- getrennt f�r: H�hle (Zweiter Schritt der Dungeonkonstruktion nach !!!Kapitel), an Raum (Dritter Schritt)
- die Wahrscheinlichkeit f�r die Platzierung einer Subszene ist $c/C$ mit $C$ als Summe aller Faktoren

%---
Anmerkung: Die Reihenfolge der Subszenen hat keine Relevanz
"`Subszene hinzuf�gen"' f�gt eine Subszene mit den aktuellen Einstellungen hinzu,
 dazu muss eine .irr-Datei ausgew�hlt werden (Anforderungen an diese, s.u.)
"`Subszene l�schen"' l�scht die aktuell markierte Subszene
"`�nderungen �bernehmen"' �bernimmt die aktuellen Werte f�r die angew�hlte Subszene

Benennung: Benennung der Subszene im Szenegraphen
Skalierung: Skalierung der Subszene
St�rke der Gangsplineableitungen: Faktoren $g$ f�r die Ableitungen der Gangsplines
Pfadinfo: Dateipfad der markierte Subszene


%- Minimal-Abstand zwischen Raum und H�hle festlegbar (empfohlen: halber Sektor f�r st�rungsfreien Gangverlauf, min. 0, max. Sektorausdehnung)
%  -> auch Sektoren mit Voxelh�hle in 2 Sektoren Abstand m�ssen als FREI\_MIT\_VOXELNACHBAR katalogisiert werden
%  (damit auch zu R�ndern verbindbar)  
%  (dennoch nur 2 Sektoren tiefer Scan)
% -> 2 Sektoren Rand zu Voxelraum n�tig 

%================================================================================================
\subsubsection{Men�: Erstellung}

In diesem Men� lassen sich die einzelnen Generatorschritte ausw�hlen
"`Dungeon generieren"' f�hrt diese Generatorschritte in der angegebenen Reihenfolge durch
"`Dungeon anzeigen"' zeigt die aktuell erstellte 3D-Szene an, also den generierten Dungeon


%================================================================================================
\subsubsection{Men�: Weiteres}

Zufallsgenerator-Parameter f�r FraktalGenerator, Erosion und Architekt einstellbar
Farbe der H�hle festlegbar (die Werte sollten im Normalfall zwischen 0.0 und 1.0 liegen, k�nnen aber auch dar�ber oder darunter liegen)
Hintergrundfarbe des Renderfensters: RGB-Farbe im Bereich $[0,255]\times[0,255]\times[0,255]$

"`�nderungen �bernehmen"' �bernimmt die jeweiligen Farb�nderungen

"`Meshcache \& Texturebuffer l�schen"' l�scht alle gespeicherten Meshes und alle Texturen,
die Dungeonszene und alle Subszenen werden ebenso gel�scht, da sie sonst auf invalide Daten zugreifen w�rden

Im Normalfall nicht n�tig, nur wenn viele verschiedene Subszenen und/oder Texturen eingeladen und wieder gel�scht werden,
da die Meshes und Texturen im Cache bzw. Buffer bleiben.

%================================================================================================
\subsubsection{Anforderungen an Subszenen}

Subszenen sind als .irr-Szenendateien abgelegt

Folgende Geometrien m�ssen vorhanden sein und bilden die Adaptermeshes:
"`Nord\_Andockstelle"'
"`Ost\_Andockstelle"'
"`Sued\_Andockstelle"'
"`West\_Andockstelle"'
Der Mesh muss den jeweiligen Namen tragen.

- Andockstellenmeshes m�ssen so ausgerichtet sein, dass ihre Normalen m�ssen ungleich der Y-Achse sind
(sonst lokales Koordinatensystem uneindeutig)
- ihre Punkte m�ssen bez�glich ihres Polarwinkels zirkul�r aufsteigend im oder gegen den Uhrzeigersinn angeordnet sein

Folgende Geometrien k�nnen vorhanden sein und stellen die Verschlussgeometrie bzw. Andockgeometrie dar:
Andockgeometrien:
"`Nord\_offen"'
"`Ost\_offen"'
"`Sued\_offen"'
"`West\_offen"'
Verschlussgeometrien:
"`Nord\_geschlossen"'
"`Ost\_geschlossen"'
"`Sued\_geschlossen"'
"`West\_geschlossen"'
Jeder beliebige Knoten im Szenegraph darf diese Namen tragen.

 				% Anleitung f�r den Dungeongenerator
%\chapter{Anleitung zur Kompilierung des Dungeongenerators}

Wie wird das Programm mittels Visual Studio 2010 SP 1 kompiliert?

- Pfade zu include setzen
- Pfade zu lib setzen
- Arbeitverzeichnis anpassen	% Anleitung f�r die Kompilierung des Programms
%\include{A_Anmerkungen}						% Bugs und Anmerkungen
\chapter{CD-Inhalt}
\label{A_DVD}

Inhalt des beiliegenden Datentr�gers: % !!!mit Dateinamen und Verzeichnissen

\begin{itemize}
	\item Ausarbeitung als digitale Version: \\
				Datei "`\textbackslash MaxHoenig\_Masterarbeit.pdf"'
	\item Anleitung f�r das beiliegende Programm: \\
				Datei "`\textbackslash Anleitung\_Dungeongenerator.pdf"'
	\item Anleitung zum Kompilieren des Programms: \\
				Datei "`\textbackslash Anleitung\_Kompilierung.pdf"'
	\item �bersicht �ber die Namenskonventionen, die innerhalb des Programms verwendet wurden: \\
				Datei "`\textbackslash Namenskonventionen.pdf"'
	\item Programmcode in C++ als Projekt f�r Microsoft Visual Studio 2010 mit Service Pack 1: \\
				Ordner "`\textbackslash Programmcode\textbackslash"'
	\item F�r Windows x86 kompiliertes Programm (ben�tigt OpenGL): \\
				Ordner "`\textbackslash Dungeongenerator\textbackslash"'
	\item Irrlicht, Version 1.7.3 (verwendete Version) \cite{W_Irrlicht}:\\
				Ordner "`\textbackslash Extern\textbackslash irrlicht-1.7.3\textbackslash"'
	\item Szeneneditor IrrEdit (erm�glicht die  Erstellung von .irr-Szenen), Version 1.5 \cite{W_IrrEdit}:\footnote{Bugs in Irredit 1.5:
1.) Das Einlesen von Meshes, die Meshbuffer mit mehr als $\sim$ 32k Vertices erzeugen, resultiert in inkorrekten Meshes.
Vermutete Ursache: Verwendung von \emph{signed short} statt \emph{unsigned short} als Datentyp f�r die Vertex-Indices.
%Au�erdem wird nur der erste Meshbuffer dargestellt.
%Im 3D-Modelling-Programm 3ds Max 2010 funktioniert das Einlesen problemlos.
2.) Das �ndern von Materialeigenschaften kann die Texturkoordinaten eines Objekts auf falsche Werte setzen.
}\\
				Ordner "`\textbackslash Extern\textbackslash irrEdit-1.5\textbackslash"'
	\item Modellings der f�r die Subszenen verwendeten Basisgeometrien, Andockgeometrien, Verschlussgeometrien
				und Andockstellen als Autodesk 3DS Max 2010 .max-Dateien: \\
				Ordner "`\textbackslash Modellings\textbackslash"'
\end{itemize}


\subsubsection{Beispieldateien}

Beispieldateien f�r H�hlen, Dungeons und H�hlenfarben befinden als XML-Dateien sich im Verzeichnis 
"`\textbackslash Dungeongenerator\textbackslash"'.

\ \\
Beispiele f�r h�hlenartige L-Systeme und bekannte L-Systeme:
\begin{compactitem}
	\item Abbildung \ref{B_LSys_Hoehlenbasis}: "`LSystem\_Hoehle\_Flach1.xml"'
	\item Abbildung \ref{B_BspDungeons} unten: "`LSystem\_Hoehle\_Flach2.xml"'
	\item Abbildung \ref{B_BspDungeons} oben: "`LSystem\_Hoehle\_Tief1.xml"'
	\item viele der Abbildungen, die eine H�hle von innen zeigen: "`LSystem\_Baum.xml"' \cite[S.25]{B_LindSys}
	(stellt von au�en betrachtet eine Baumstruktur dar, eignet sich aber von innen betrachtet auch als H�hle)	
	\item Drachenkurve: "`LSystem\_Drachenkurve.xml"' \cite[S.11]{B_LindSys}
	\item Waben: "`LSystem\_Waben.xml"' \cite{W_LSys_HexaKolam} %http://jsxgraph.uni-bayreuth.de/wiki/index.php/Hexagonal_kolam
	\item Koch-Kurve: "`LSystem\_Koch-Kurve.xml"'
	\item Kochsche Schneeflocke: "`LSystem\_Kochsche-Schneeflocke.xml"'
\end{compactitem}

%\ \\
%Parameter zur Erstellung von h�hlenartigen, zuf�lligen L-Systemen:
%-> als Voreinstellung
%in Anleitung und dieser Arbeit andere Symbole erw�hnt

\ \\
Beispiele f�r Dungeons:
\begin{compactitem}
	\item Abbildung \ref{B_DungeonBsp1}: Datei "`Dungeon\_Aufsteigend.xml"'
	\item Abbildung \ref{B_DungeonBsp2}: Datei "`Dungeon\_Ring.xml"'
	\item Abbildung \ref{B_DungeonBsp3}: Datei "`Dungeon\_Tief.xml"'
	\item Abbildung \ref{B_DungeonBsp4}: Datei "`Dungeon\_Weit.xml"'
	\item Abbildung \ref{B_DungeonBsp5}: Datei "`Dungeon\_Doppelkugel.xml"'
	\item Dungeon aus Abbildung \ref{B_BspDungeons} oben, etwas modifiziert: "`Dungeon\_Verdreht.xml"'
	\item Abbildung \ref{B_DungeonInnenBsp2}(b): "`Dungeon\_Bizarr.xml"'
\end{compactitem}

\ \\
Farben f�r H�hlen:
\begin{compactitem}
	\item Bernstein : "`Farbe\_Bernstein.xml"'
	\item Eis : "`Farbe\_Eis.xml"'
	\item Rubin : "`Farbe\_Rubin.xml"'
	\item Smaragd : "`Farbe\_Smaragd.xml"'
	\item Amethyst : "`Farbe\_Amethyst.xml"'
\end{compactitem}

%\ \\
%Die Einstellungen der Datei "`Einstellungen\_Standard.xml"' werden bei Programmstart automatisch geladen.


%Gute Generierungsvorgaben f�r L-System-H�hlen
%------
%
%- L-Systeme inkl. Parametern
% - bekannte L-Systeme
% - f�r H�hlen aus K3, 8
% 
%- Dungeons aus Anhang 2
%
%- Materialien
%
%- Generierungsparamerter f�r h�hlenartige L-Systeme
 
								% Inhalt der beliegenden DVD

\end{appendix}

%\appendix
%\include{AnhangA}
%\begin{thebibliography}{----------} 
\bibitem[UrbanHoen] {pbuh} 
     Maik Urbannek, Max H�nig. 
     {\em Einstieg in die interaktive immersive Manipulation}. 
     Burg Giebichenstein Kunsthochschule Halle, Sommersemester 2010
\bibitem[Eidinger] {pbeid} 
     Jens Eidinger. 
     {\em Radioreparatur}. 
     Burg Giebichenstein Kunsthochschule Halle, Oktober 2009     
\end{thebibliography}

\listoffigures
%\nocite{*}													% alle Quellen anzeigen (auch nicht benutzte)
\bibliography{quellen}							% quellen.bib  



\end{document}
